\documentclass{amsart}
\usepackage[utf8]{inputenc}
\usepackage{lmodern}
\usepackage{amssymb,mathtools,array,longtable,tabularx,booktabs}
\usepackage{mathrsfs}
\usepackage{enumitem,graphicx}
\usepackage[hidelinks]{hyperref}
\newcommand{\GoldbachTitlePhrase}{1 + 1.8}
\hypersetup{pdftitle={Theorem (\GoldbachTitlePhrase) on the Goldbach conjecture},pdfauthor={Ruiyi Zhang}}
\allowdisplaybreaks
\emergencystretch=2em
\setlength{\LTpre}{1em}
\setlength{\LTpost}{1em}
\newcommand{\PP}{\mathbb{P}}
\newcommand{\UC}{U_C}
\DeclareUnicodeCharacter{2212}{\ensuremath{-}}
\DeclareUnicodeCharacter{2264}{\ensuremath{\le}}
\DeclareUnicodeCharacter{2265}{\ensuremath{\ge}}
\DeclareUnicodeCharacter{2260}{\ensuremath{\ne}}
\DeclareUnicodeCharacter{221E}{\ensuremath{\infty}}
\DeclareUnicodeCharacter{2208}{\ensuremath{\in}}
\DeclareUnicodeCharacter{2286}{\ensuremath{\subseteq}}
\DeclareUnicodeCharacter{00D7}{\ensuremath{\times}}
\DeclareUnicodeCharacter{2211}{\ensuremath{\sum}}
\DeclareUnicodeCharacter{220F}{\ensuremath{\prod}}
\DeclareUnicodeCharacter{220E}{\ensuremath{\square}}
\DeclareUnicodeCharacter{2032}{\ensuremath{\prime}}
\DeclareUnicodeCharacter{2011}{-}
\DeclareUnicodeCharacter{2013}{--}
\DeclareUnicodeCharacter{2014}{---}
\newtheorem{theorem}{Theorem}
\newtheorem{lemma}{Lemma}[section]
\newtheorem{proposition}[lemma]{Proposition}
\newtheorem{corollary}[lemma]{Corollary}
\numberwithin{equation}{section}
\numberwithin{table}{section}
\theoremstyle{definition}
\newtheorem{definition}[lemma]{Definition}
\title{Theorem (\GoldbachTitlePhrase) on the Goldbach conjecture}
\author{Ruiyi Zhang}
\address{Independent Researcher, Columbia, Maryland, USA}
\email{ryzfsu@gmail.com}
\subjclass[2020]{Primary 11P32; Secondary 11N36, 11N13}
\keywords{Goldbach problem, almost-primes, sieve methods, well-factorable weights, shifted primes}
\begin{document}
\begin{abstract}
Every sufficiently large even
integer $N$ has a representation $N=p+rq$.
Here $p,q$ are primes and $r$ is $1$ or a prime,
with $r\le q^{a-1}=q^{4/5}=q^{0.8}$. This is Proposition $(1+a)$ of Li--Liu with
$a=1.8=9/5$. We use a weighted sieve on shifted primes.
Its level of distribution exceeds one half. The proof combines
thirteen row estimates with a switched upper bound. For a fixed source
cutoff, the certified D-unit endpoint is
$D_{\rm lower}=1.83399665352036682372$. Since the source mass is
$4(1-\varepsilon_A)U_C$, division by four gives
$(1-\varepsilon_A)D_{\rm lower}=1.833813253855014787037628-o(1)$.
The balanced-semiprime cost is at most
$1.788209644611262112809348242339+o(1)$ in those units. The remaining
coefficient is
$0.04560360924375267422827-o(1)$.
\end{abstract}
\maketitle
\section{Introduction}
Chen announced his theorem in 1966 \cite{Chen1966}.
His proof appeared in 1973 \cite{Chen1973}.
By this theorem, every sufficiently large even integer is a prime plus an integer
with at most two prime factors.
Li and Liu \cite{LiLiu2026} formulate Proposition $(1+a)$:
\[
 N=p+rq,\qquad p,q\in\mathbb P,\quad r\in\{1\}\cup\mathbb P,
 \quad r\le q^{a-1}.
\]
Li and Liu take $a=1.9$. Here $a=1.8=9/5$, so
$r\le q^{a-1}=q^{4/5}=q^{0.8}$.
Theorem~\ref{thm:goldbach} follows from the strict comparison
\[
 \begin{aligned}
 D_{\rm lower}&\ge1.83399665352036682372,\\
 T&=8\log\frac{50009}{39991}<1.78838848345960807362<D_{\rm lower}.
 \end{aligned}
\]
It holds at the fixed source cutoff below.
We use the Chen-constant normalization
\[
 C(N)=\prod_{\ell>2}\left(1-\frac1{(\ell-1)^2}\right)
       \prod_{\substack{\ell\mid N\\\ell>2}}\frac{\ell-1}{\ell-2},
 \qquad U_C=\frac{C(N)N}{(\log N)^2};
\]
There is no leading factor of two. With $\varepsilon_A=10^{-4}$,
our selected endpoint and the pointwise division by four give
\[
 P_{2,\varepsilon_A}(N)\ge
 [1.833813253855014787037628-o(1)]C(N)N/(\log N)^2,
\]
where $P_{2,\varepsilon_A}(N)=\#\{p<(1-\varepsilon_A)N:
p\in\mathbb P,\ \Omega(N-p)\le2\}$.
Our normalization uses $C(N)$ without a leading factor $2$, as do Wu's lower
bounds $D_{1,2}(N)\ge0.836C(N)N/(\log N)^2$
\cite[Theorem~2]{Wu2004} and
$D_{1,2}(N)\ge0.899C(N)N/(\log N)^2$ \cite{Wu2008}.
Here $D_{\rm lower}$ is in D units. The source mass is
$4(1-\varepsilon_A)U_C$, so pointwise division by four contributes the
factor $1-\varepsilon_A$ to $D_{\rm lower}$.
This calculation uses the specified direct and switched distribution inputs,
with domain-dependent levels up to $153/256$.
The separate balanced-semiprime subtraction costs
$8(1-10^{-4})\log(50009/39991)$, leaving coefficient
$0.0456036092437526742\ldots-o(1)$ for the restricted $(1+1.8)$
representations.
Li--Liu's $(1+1.9)$ argument uses the same singular-series convention and
the corresponding limiting exclusion threshold
$8\log(10/9)=0.8428841252626104\ldots\approx0.84288$
\cite[(5.42)--(5.43)]{LiLiu2026}.

Lichtman \cite{Lichtman} develops distribution estimates beyond one half.
We use the signed estimates of \cite[Theorems~1.1--1.2]{Z1}
and the local estimate \cite[Theorem~3.2, (3.13)]{Z1}.
In outline, the proof runs
\[
 \text{thirteen signed rows}\ \Longrightarrow\ \text{switching}
 \ \Longrightarrow\ D_{\rm lower}>T
 \ \Longrightarrow\ N=p+rq.
\]

\subsection*{Input from \textup{[Z1]}}
\phantomsection\label{sec:z1-input}
We apply \cite[Theorem~1.2]{Z1} to the complete coefficient family below.
Fix $A,B>0$, $0<\varepsilon\le10^{-4}$, and
\[
 \begin{gathered}
 0<u\le1/500,\quad 1/2<\theta<\Theta(u)=(5-7/32-3u)/8,\\
 \eta_\Theta=\Theta(u)-\theta,\quad
 0<\sigma\le\min\{10^{-10},\eta_\Theta/200,\sqrt{u/(2\theta)}\},\\
 \tau=\sigma^9,\quad D=N^\theta,\quad z=N^u,\quad
 Y=D^{\sigma^2},\quad K_{\rm scale}=B+1,\\
 C_{\rm box}\in\mathbb Z,\qquad C_{\rm box}\ge\lceil2A+32\rceil.
 \end{gathered}
\]
Thus $C_{\rm box}\ge2(A+16)$ and $C_{\rm box}\ge A+18$,
as required in \cite[(3.24), (B.18)--(B.19)]{Z1}.
For large even $N$ and $\nu\in\{+,-\}$, put
\begin{equation*}\tag*{(W)}
\begin{split}
 w_N^\nu(q)={}&
 \sum_{\substack{d<D^\sigma\\d\mid P_N(Y)}}
       \phi_d^\nu\mathbf1_{q=d}\\
 &+\sum_{\mathbf D\in\mathcal D^\nu(D)}
 \sum_{\substack{d<D^\sigma\\d\mid P_N(Y)}}
       \lambda_d^\nu(\mathbf D)
 \sum_{\substack{p_i\in\mathcal P_N\\
       D_i\le p_i<\min(D_i^{1+\tau},z)}}
       K_{\mathbf D}(\mathbf p)\mathbf1_{q=dp_1\cdots p_k}.
\end{split}
\end{equation*}
Here $\mathcal P_N=\{p\text{ prime}:p\nmid2N\}$ and
$P_N(t)=\prod_{p<t,\ p\in\mathcal P_N}p$.
We use Iwaniec's grid, vectors and coefficients \cite[Theorem~4]{ref5},
with depth $R=1$, mesh $\sigma,\tau$, and $D_i\ge Y$.
All smooth divisors and the smooth-only term are kept.
The coefficient bounds and distinctness mask are
\[
 |\phi_d^\nu|,|\lambda_d^\nu(\mathbf D)|\le1,\qquad
 K_{\mathbf D}(\mathbf p)=
 \prod_{\substack{i<j\\D_i=D_j}}(1-\mathbf1_{p_i=p_j}).
\]
An integer shift $a$ is admissible if
\[
 N\le a\le N\log N/3,\qquad
 (a,q)=1\quad(q\in\operatorname{supp}w_N^\nu).
\]
The shift $a$ is distinct from the Li--Liu parameter $a=1.8=9/5$ below.
Put $I_N=(N/(\log N)^B,(1-\varepsilon)N)$.
Let $\pi_{I_N}$ count the primes in this interval.
Let $\pi_{I_N}(q,a)$ count those in the residue class $a$.
For fixed admissible $\sigma$,
\[
 \begin{gathered}
 \operatorname{supp}w_N^\nu\subseteq
 \{q<D^{1+\sigma+\tau}:(q,2N)=1\},\quad
 D^{1+\sigma+\tau}<N^{\Theta(u)},\\
 |w_N^\nu(q)|\ll_{\sigma,u,\theta}1,\qquad
 Y\le N^{u/2}<z,\quad \log D/\log z=\theta/u,
 \end{gathered}
\]
\[
 \left|\sum_q w_N^\nu(q)
 \left\{\pi_{I_N}(q,a)-\frac{\pi_{I_N}}{\varphi(q)}\right\}\right|
 \ll_{\sigma,A,u,\theta,B,\varepsilon,C_{\rm box}} N(\log N)^{-A}.
\]
The sum runs over the full support, uniformly in admissible shifts.
The implied constant and the threshold depend on the displayed fixed parameters.
Through Siegel--Walfisz they may be ineffective.
For $u=1/500$, this is \cite[Theorem~1.1]{Z1}, with
\[
 \Theta(1/500)=19101/32000,\qquad
 \sqrt{u/(2\theta)}=\sqrt{1/(1000\theta)}.
\]
By \cite[Corollary~1.3]{Z1}, the same estimate holds with the comparator
$\operatorname{Li}((1-\varepsilon)N)-\operatorname{Li}(N/(\log N)^B)$.

For the row S1-0 we take
\[
 \begin{gathered}
 u=1/500,\quad\theta=19101/32000-10^{-4},\quad
 B=30,\quad\varepsilon=10^{-4},\quad K_{\rm scale}=31,\\
 a=N,\quad\eta_\Theta=10^{-4},\quad \theta/u=298.403125.
 \end{gathered}
\]
Since $(q,2N)=1$, this shift is admissible.
The choice $0<\sigma\le10^{-10}$ made later meets the remaining restrictions.
We keep $C_{\rm box}\ge\lceil2A+32\rceil$.
Each prefixed or product row is checked separately below.
The Heath--Brown decomposition \cite{ref6} and packing produce the boxes.
BFI and Siegel--Walfisz give the first dispersion.
The spectral estimate gives the second square.
Type-I estimates and the full-family sum finish the bound.
Details are in \hyperref[p2:dispersion-proof]{the local proof below} and
\cite[\S\S3--5 and Appendix~B]{Z1}.

\subsection{Statement and sources}

\begin{theorem}[Goldbach representation $(\GoldbachTitlePhrase)$]\label{thm:goldbach}
In the Proposition $(1+a)$ convention, $a=1.8=9/5$, so $a-1=0.8=4/5$.
Every sufficiently large even integer $N$ has a representation
\[
 N=p+rq,\qquad p,q\in\mathbb P,\qquad r\in\{1\}\cup\mathbb P,
 \qquad r\le q^{a-1}=q^{4/5}=q^{0.8}.                    \tag{1.1}
\]
\end{theorem}

\subsection{Normalization and fixed source}

Put $L=\log N$. The singular series has no leading factor of two:
\[
 C(N)=C_2\prod_{\substack{\ell\mid N\\\ell>2}}
          \frac{\ell-1}{\ell-2},\qquad
 C_2=\prod_{\ell>2}\left(1-\frac1{(\ell-1)^2}\right)>0,
 \qquad U_C(N)=\frac{C(N)N}{L^2}.                            \tag{1.2}
\]
Here $\ell,p,q$ denote primes.
Uniformly over even $N$, $C(N)\ge C_2$ and $U_C(N)\to\infty$.
Fix
\[
 \varepsilon_A=\eta_L=\frac1{10000},\qquad a=1.8=\frac95,
 \qquad b=\frac{a-1}{a}-\varepsilon_A
          =\frac{39991}{90000}>\frac13.                    \tag{1.3}
\]
The cutoff $b$ separates the balanced semiprimes, and $\eta_L$ is the row buffer.
Here $T$ is the fixed exclusion threshold and $\tau$ is the mesh width.
In Sections~2--4, $D$ denotes a sieve level; $D_{\rm lower}$ is the
normalized endpoint. The local cutoff $H$, switched integral $H_L$,
source majorant $g(e)$, and Li--Liu threshold $g_3$ are distinct quantities.
The source is
\[
 \mathcal A_{\varepsilon_A}
   =\{N-p:p\in\mathbb P,\ p<(1-\varepsilon_A)N\}.
                                                               \tag{1.4}
\]
Thus $\varepsilon_A N<n<N$ on the source.
Exponents are measured against $N$, and a local shell has scale $X$.

The weight O1 uses the rational parameters
\[
 \alpha=\frac{100}{1149},\quad \beta=\frac{50}{309},\quad
 \mu=\frac{25}{128},\quad\lambda=\frac{93}{400},\quad
 \tau_S=\frac{57}{224},\quad u=\frac1{500},\quad s_{\rm O1}=\frac{11}{5}.
                                                               \tag{1.5}
\]
For $r\in[\beta,\lambda)$, put
\[
 B(r)=\begin{cases}
 \beta,&\beta\le r<\mu,\\
 \min\{\beta,\max\{\alpha,C_0-5r/2\}\},&\mu\le r<\lambda,
 \end{cases}
 \qquad C_0=\frac{57}{64}-\frac1{500}-s_{\rm O1}\alpha.
                                                               \tag{1.6}
\]
Since $2(C_0-\alpha)/5>\lambda$, the lower truncation is inactive.
The S10 proof writes $B_0$ for this same clipped function on $[\mu,\lambda)$.
The other breakpoint is
\[
 v_\beta=2(C_0-\beta)/5=506848823/2366940000\in(\mu,\lambda).
\]
Hence $B$ is continuous, non-increasing, and piecewise affine.

\subsection{The weight and its elementary counting inequality}

For squarefree $d$ with $(d,N)=1$, define the cofactor count
\[
 S(d;z)=\#\{n=dm\in\mathcal A_{\varepsilon_A}:
             P^-(m)\ge z,\ (m,Nd)=1\}.                    \tag{1.7}
\]
Set $P^-(1)=+\infty$ and $t_i=\log p_i/L$.
Each row $S_i$ of the table is a sum of (1.7),
with each admissible divisor counted separately.
For $S_1,S_2$, take $d=1$.

\begin{center}
\small\renewcommand{\arraystretch}{1.15}\setlength{\tabcolsep}{3pt}
\begin{tabular}{@{}lc p{.51\textwidth}l@{}}
\toprule
Row & Coefficient $c_i$ & Divisor and domain & Cofactor cutoff \\
\midrule
$S_1$ & $3$ & $d=1$ & $N^\alpha$ \\
$S_2$ & $1$ & $d=1$ & $N^\beta$ \\
$S_3$ & $-2$ & $d=p_1,\ \alpha\le t_1<\mu$ & $N^\alpha$ \\
$S_4$ & $-1$ & $d=p_1,\ \mu\le t_1<1/4$ & $N^\alpha$ \\
$S_5$ & $-1$ & $d=p_1,\ 1/4\le t_1<\tau_S$ & $N^\alpha$ \\
$S_6$ & $-1$ & $d=p_1,\ \tau_S\le t_1<1/3$ & $N^\alpha$ \\
$S_7$ & $-1$ & $d=p_1,\ \mu\le t_1<\lambda$ & $N^\alpha$ \\
$S_8$ & $1$ & $d=p_1p_2,\ \alpha\le t_2<t_1<\beta$ & $N^\alpha$ \\
$S_9$ & $1$ & $d=p_1p_2,\ \alpha\le t_2<\beta\le t_1<\mu$ & $N^\alpha$ \\
$S_{10}$ & $1$ & $d=p_1p_2,\ \mu\le t_1<\lambda,\ \alpha\le t_2<B(t_1)$ & $N^\alpha$ \\
$S_{11}$ & $-2$ & $d=p_1p_2,\ \lambda\le t_1<1/3,\ t_1<t_2<(1-t_1)/2$ & $p_2$ \\
$S_{12}$ & $-1$ & $d=p_1p_2,\ \alpha\le t_1<1/3\le t_2<(1-t_1)/2$ & $p_2$ \\
$S_{13}$ & $-1$ & $d=p_1p_2,\ \beta\le t_1<\lambda\le t_2<(1-t_1)/2$ & $(N/(p_1p_2))^{1/2}$ \\
\bottomrule
\end{tabular}
\end{center}

Let $R_i(n)$ count the contributions of $n$ to row $i$. Set
\[
 W_{13}(n)=\sum_{i=1}^{13}c_iR_i(n),\qquad
 W(n)=W_{13}(n)-\mathbf1_{\Omega(n)\ge3}\max\{W_{13}(n),0\}.
                                                               \tag{1.8}
\]
All these weights vanish unless $P^-(n)\ge N^\alpha$. The switched excess is
\[
 S_{\rm sw}(N)=\sum_{\substack{n\in\mathcal A_{\varepsilon_A}\\
                  P^-(n)\ge N^\alpha,\ \Omega(n)\ge3}}
                         \max\{W_{13}(n),0\},
 \qquad \Sigma_W=\sum_{i=1}^{13}c_iS_i-S_{\rm sw}.           \tag{1.9}
\]
The row dictionary consists of $S_1,\ldots,S_{13}$; $S_{14},S_{15}$ are absent.
The next lemma turns the weighted sum into a lower bound for the prime count.

\begin{lemma}[Pointwise count]\label{lem:pointwise-count}
For all sufficiently large $N$,
\[
 \#\{p<(1-\varepsilon_A)N:p\in\mathbb P,\ \Omega(N-p)\le2\}
 \ge \frac{\Sigma_W}{4}-O(N^{1-\alpha}L).                 \tag{1.10}
\]
\end{lemma}
\begin{proof}
First take squarefree $n$ with $(n,N)=1$.
If $\Omega(n)\ge3$, then $W(n)\le0$ by (1.8).
For $\Omega(n)\le2$, the bound $n>\varepsilon_A N$ excludes $n=1$.
A prime $n$ exceeds $N^{1/3}$ for large $N$,
so only $S_1,S_2$ contribute and $W(n)=4$.

Let $n=rq$ with primes $r<q$.
Then $q>(\varepsilon_A N)^{1/2}>N^{1/3}$.
Put $t=\log r/L$. The first ten rows give
\[
 \begin{array}{c|cccc}
 t&[\alpha,\beta)&[\beta,\lambda)&[\lambda,1/3)&[1/3,1/2)\\ \hline
 \text{weight}&1&2&3&4.
 \end{array}                                                \tag{1.11}
\]
The positive pair rows cannot use $q$,
and the single-prime rows supply the subtractions in this display.
The last three rows contribute zero when $n=rq$, hence $W(n)\le4$.

Since $12\alpha>1$, every retained $n<N$ has $\Omega(n)\le11$.
So there are at most $11$ single primes and $\binom{11}{2}=55$ pairs, and
\[
 |W_{13}(n)|\le3+1+2\cdot11+4\cdot11
                     +3\cdot55+2\cdot55+55+55=455,
 \qquad |W(n)|\le455.
\]
The terms group as $S_1,S_2$; $S_3$; $S_4$--$S_7$;
$S_8$--$S_{10}$; and $S_{11},S_{12},S_{13}$, respectively.
The non-squarefree exceptions number at most
\[
 \sum_{j\ge N^\alpha}N/j^2=O(N^{1-\alpha}).
\]
Integers with a prime factor in common with $N$ number at most
$\omega(N)N^{1-\alpha}=O(N^{1-\alpha}L)$,
and their weighted sum has the same order.
On the remaining integers,
$W(n)\le4\mathbf1_{\Omega(n)\le2}$.
Summing and dividing by four proves the lemma.
\end{proof}

The switched main term is $8(1-\varepsilon_A)I_{\rm HL}U_C$, where
\[
 I_{\rm HL}=\sum_{k=5}^{11}
   \int_{\substack{\alpha\le t_1<\cdots<t_k\\t_1+\cdots+t_k=1}}
   \frac{\max\{W_{13}(\mathbf t),0\}}{t_1\cdots t_k}
                         \,dt_1\cdots dt_{k-1}.             \tag{1.12}
\]
Here $t_k=1-\sum_{i<k}t_i$, and the integral uses the coordinate measure shown.
For the prime factors of an actual source integer $n$,
$\sum t_i=\log n/L$, and this radial coordinate need not equal $1$.
Section~5.5 proves the cell decomposition, the restriction $k\ge5$
and the prime-product estimate.
For $k=3,4$, it uses a one-sided comparison at $\sum t_i=\log n/L<1$,
where the slice dictionary need not be an identity.
After normalization, (1.12) contributes $-2I_{\rm HL}$ to $D$.

\subsection{Balanced semiprimes}
This subsection removes the semiprimes whose smaller factor is too large.
Let $P_{2,\varepsilon_A}(N)$ denote the count in (1.10).
Let $G_a(N)$ count primes $p<(1-\varepsilon_A)N$ admitting a representation as in (1.1), with $r\le q^{a-1}$.
For a prime remainder, take $r=1$.
For $n=rq$, $r\le q$, the bound $r<N^b$ gives
\[
 q>\varepsilon_A N^{1-b},\qquad
 q^{a-1}>\varepsilon_A^{a-1}N^{b+a\varepsilon_A}>N^b>r       \tag{1.13}
\]
provided
\[
 N>\varepsilon_A^{-(a-1)/(a\varepsilon_A)}.                \tag{1.14}
\]
Here $(1-b)(a-1)-b=a\varepsilon_A$.
For (1.3), $\log N\ge44445$ suffices, since
$(a-1)/(a\varepsilon_A)=40000/9$ and $\log10000<10$.

For the remaining semiprimes put
\[
 g(e)=\#\{(r,q)\in\mathbb P^2:
       N^b\le r\le N^{1/2},\quad
       \varepsilon_A N<rq<N,\quad rq=e\}.                  \tag{1.15}
\]
The majorant allows $q<r$ and counts both choices of $r$.
Thus $0\le g(e)\le2$, and every remaining bad semiprime is counted.
Consequently
\[
 G_a(N)\ge P_{2,\varepsilon_A}(N)
                 -\sum_e g(e)\mathbf1_{\mathbb P}(N-e).     \tag{1.16}
\]
Prime squares remain in (1.15);
discarding them would cost $O(N^{1/2})$.
For $a<2$, they fail $r=q\le q^{a-1}$.

Section~5.6 proves the following untruncated estimate for the majorant.
For each fixed $A>0$ there is a fixed $B$ such that, with $Q=N^{1/2}L^{-B}$,
\[
 \sum_{d\le Q}\max_{(h,d)=1}
 \left|\sum_{e\equiv h\ (d)}g(e)
       -\frac1{\varphi(d)}\sum_{(e,d)=1}g(e)\right|
       \ll_{A,\varepsilon_A} NL^{-A}.                       \tag{1.17}
\]
Lemma~\ref{lem:complete-support-bv} gives (5.24).
Section~5.6 applies it at the scales and precisions specified after (5.47).
The partition (5.45)--(5.46) and the error (5.47) then give (5.48).
Take $b_{\rm G4}=b=39991/90000>1/3$ and $\varepsilon=\varepsilon_A$.
The multiplicity in (5.44) includes the endpoint $e=\varepsilon_AN$,
where $|g-g_N|\le2$.
So each discrepancy changes by at most $2+2/\varphi(d)$,
and the total change is $O(Q)=O(N^{1/2})$.

Put $X_g=\sum_e g(e)$. Section~5.6 also proves
\[
 \begin{aligned}
 X_g
 &=\sum_{\substack{N^b\le r\le N^{1/2}\\r\in\mathbb P}}
       \{\pi(N/r)-\pi(\varepsilon_A N/r)\}\\
 &=\frac{(1-\varepsilon_A)N}{L}
       \int_b^{1/2}\frac{dt}{t(1-t)}+O_{\varepsilon_A}(N/L^2)\\
 &=\frac{(1-\varepsilon_A)N}{L}
       \log\frac{1-b}{b}+O_{\varepsilon_A}(N/L^2).
 \end{aligned}                                               \tag{1.18}
\]
Uniform PNT gives the inner count with relative error $O_{\varepsilon_A}(L^{-1})$,
and partial summation in $r$ gives the integral.
Strict prime endpoints are absorbed in the error term.
Let $P_N(z)=\prod_{2<\ell<z,\ \ell\nmid N}\ell$ and $z=Q^{1/2}$.
The sequence takes each $N-e$ with multiplicity $g(e)$; call it $\mathcal B$.
For squarefree $d\mid P_N(z)$,
\[
 |\mathcal B_d|=\frac{X_g}{\varphi(d)}+r_d,\qquad
 r_d=\Delta_g(d,N)-\frac1{\varphi(d)}
                                      \sum_{(e,d)>1}g(e).    \tag{1.19}
\]
Eventually every prime factor of $e$ is at least
$\min(N^b,\varepsilon_A N^{1/2})>N^{1/4}>z$,
so the correction vanishes.
Section~5.6 applies the upper sieve with
\[
 \begin{gathered}
 V_N(z)=\frac{2e^{-\gamma}C(N)}{\log z}(1+o(1)),\quad s=2,\quad F(2)=e^\gamma,\\
 V_N(z)F(2)=\frac{8C(N)}L(1+o(1)).
 \end{gathered}
\]
The bound includes the small-prime term $2z$ and the sieve error
$O((\log Q)^{-1/3})$. Hence
\[
 \sum_e g(e)\mathbf1_{\mathbb P}(N-e)
 \le\left[8(1-\varepsilon_A)\log\frac{1-b}{b}
                   +r_3(N)\right]U_C,\qquad r_3(N)\to0.     \tag{1.20}
\]
This is the constant-eight estimate.
The same upper-sieve normalization bounds the switched excess in Section~5.5.
With singular series $2C(N)$ the multiplier becomes four,
and the absolute bound is unchanged.

\subsection{The coefficient criterion}
An auxiliary choice is \emph{admissible} if it is independent of $N$,
satisfies (1.27)--(1.29), and meets the row hypotheses in Sections~2--5.
The rational parameters are (1.3)--(1.6),
and the separate S2 argument is $s_0=19037/64$.
The next criterion reduces the theorem to one numerical inequality.

\begin{theorem}[Coefficient reduction criterion]
Fix an admissible choice. Suppose
\[
 \Sigma_W\ge[4(1-\varepsilon_A)D_{\rm cert}-r_\Sigma(N)]U_C,
 \qquad r_\Sigma(N)\to0,                                   \tag{1.21}
\]
where $D_{\rm cert}$ is fixed, and suppose (1.20).
If
\[
 D_{\rm cert}>T,\qquad
 T=8\log\frac{1-b}{b}=8\log\frac{50009}{39991},              \tag{1.22}
\]
then (1.1) holds for every sufficiently large even $N$.
\end{theorem}
\begin{proof}
By Lemma~\ref{lem:pointwise-count}, (1.16), and (1.20)--(1.21),
\[
 G_{9/5}(N)\ge
   [(1-\varepsilon_A)(D_{\rm cert}-T)-R(N)]U_C,
 \qquad R(N)\to0,                                         \tag{1.23}
\]
where $R$ is a sum of absolute remainders.
These include $|r_\Sigma|/4$, $|r_3|$, and the rough non-coprime error
$O(N^{-\alpha}L^3)$ from (5.40). Exceptional direct factors and switched
source tags contribute separately
$O_\sigma(L^3/N+N^{-\alpha/2+o(1)})$ after normalization, as in the
error account. The square exceptions also contribute $o(1)$.
Put $M=(1-\varepsilon_A)(D_{\rm cert}-T)>0$.
Choose $N_0$ so that $R(N)\le M/2$ for all even $N\ge N_0$.
Enlarge $N_0$ so that (1.14) holds and
\[
 (M/2)C_2N/L^2\ge1.
\]
Then $G_{9/5}(N)\ge1$, and this is an integer count of primes $p$.
A prime remainder gives $r=1$;
otherwise (1.13) bounds the smaller prime factor.
Thus the positive normalized weight yields an actual representation.
\end{proof}

The threshold has the directed enclosure
\begingroup\footnotesize
\[
 \begin{gathered}
 1.78838848345960807361670991332947582426562068913553428811089076453<T\\[-2pt]
 T<1.78838848345960807361670991332947582426562068913553428811089076454.
 \end{gathered}                                             \tag{1.24}
\]
\endgroup
For $x=5009/45000=(50009-39991)/(50009+39991)$,
\[
 T=16\sum_{j\ge0}\frac{x^{2j+1}}{2j+1}.
\]
Truncating at $j=39$, bounding the tail by
$16x^{81}/[81(1-x^2)]$ and comparing rationals gives (1.24).
The raw coefficient
$(1-\varepsilon_A)T=1.7882096446112621128\ldots$
is compared with $(1-\varepsilon_A)D_{\rm cert}$.
The limiting threshold is
$8\log(5/4)=1.7851484105136780461\ldots$.
A value such as $D=1.786$ exceeds that limit but fails (1.22).
Thus the rounded $1.785148$ cannot replace the exact threshold.

The threshold algebra extends to $3/2<a<2$ and
$b=(a-1)/a-\varepsilon_A>1/3$.
Writing $t=\exp(-D_{\rm cert}/8)$, it gives
\[
 a>\frac{1+t}{1-\varepsilon_A(1+t)},                        \tag{1.25}
\]
when the denominator is positive.
Smaller fixed cutoffs would require a uniform coefficient bound.
The construction here fixes $\varepsilon_A=10^{-4}$ and proves $a=9/5$ only.

\subsection{Fixed parameters}\label{p2:sec:fixed-parameters}
This subsection fixes the auxiliary parameters in the order the error terms require.
The constants and the final threshold may be ineffective.
Theorem~\ref{thm:sieve-mass} in Appendix~B.2 proves
\[
 W_{\rm AA}\le
 \frac{4000\cdot585\cdot10240090}{1-1/10000}<10^{14}=W_* .    \tag{1.26}
\]
The bound is uniform in the fixed mesh parameter and uses \cite[\S~2.4, (2.28)--(2.30)]{Z1}.
The row applications are in Section~6.2.
First fix (1.3)--(1.6), all row domains and the classical input constants,
including the finite AA37 constant $C_E$.
Put $g_{\rm AA}=0.00002853328$.
Choose a fixed $\sigma>0$ within the source, depth, HB and packing limits, with
\[
 \sigma\le10^{-10},\qquad
 \sigma\le\frac{4g_{\rm AA}}{C_EW_*}\quad\text{if }C_E>0.  \tag{1.27}
\]
The bounds $\sigma\le\eta_L/2500$ and $\sigma<3^{-R}$ are automatic:
$\eta_L/2500=4\cdot10^{-8}$ and the AA37 depth is $R=1$
\cite[Theorem~2.6]{Z1}. If $C_E=0$, only $\sigma\le10^{-10}$ is needed.
Set $\tau=\tau_{\rm mesh}=\sigma^9$, $\rho=1+\tau$, and $\delta=30\sigma$.
The cost $C_E\sigma$ enters $g_{\rm AA}$ once;
the AA37 term $o_\sigma(1)$ vanishes separately.

Next fix the HB degrees and moment exponents.
Take
\[
 \begin{gathered}
 c_{\rm bdry}=163,\quad c_{\rm LOG}=164,\quad A_{\rm raw}=3,\\
 C_b=334\ge\max(A_{\rm raw}+c_{\rm LOG}+2,6)=169.
 \end{gathered}
\]
With $C_{\rm box}=C_b=334$, the generic aggregate input introduced above
is restricted to $0<A\le151$ by $C_{\rm box}\ge\lceil2A+32\rceil$.
The external count and local savings are
\begingroup\small
\[
 K_0=24C_b+26=8042,\qquad
 A'=A_{\rm raw}+24C_b+28=8047,\qquad
 A_{\rm work}=A'+4=8051.                                    \tag{1.28}
\]
\endgroup
Each local theorem includes all three slot-modulus dyadic splits.
The $O((\log X)^3)$ slot-shell triples use $A_{\rm work}$,
and their sum gives $A'$ with one spare logarithm.
The count $K_0$ includes two redundant slot-shell factors.
An external slot-1 convention would give
$K_{\rm II}=24C_b+27=8043$ instead.
Only the $K_0$ convention is used, and $K_0+2-A'=-3$.
Here $A_{\rm raw}=3$ specifies the normalized target saving.
Summing the arbitrary-precision local estimate of
\cite[Theorem~3.2, (3.13)]{Z1}, as in the first inequality of
\cite[(3.92)]{Z1}, gives a raw logarithmic error bounded by
\[
 N L^{K_0-A'}=N L^{8042-8047}=N L^{-5}.
\]
Division by $W_{\rm norm}=4(1-\varepsilon_A)U_C\gg N/L^2$
gives $L^{-3}$ in (1.30).
The local precision $A'=8047$ gives this bound.
The aggregate theorem in \cite{Z1} has its own calibration
$A'=A+24C_{\rm box}+31$.
Fix the actual Type-I count $K_{\rm I}$, then take $A_1,A_2\ge K_{\rm I}+4$.
Let $c_{\rm pre}$ be the coefficient loss preceding the exceptional cutoff.
Choose, in order,
\[
 \begin{gathered}
 A_E\ge2A_{\rm work}+2C_b+c_{\rm pre}+2,\quad C_1\ge A_E+c_1+3,\\
 C_M\ge\max\{5C_1+A_E+C_b+3,
                   3C_1+C_b+A_E+c_{\rm diag}+1\},\\
 A_3\ge A_E+C_b+2,
 \end{gathered}                                              \tag{1.29}
\]
Here $c_1$ is the moment exponent in \cite[(3.20)]{Z1}.
The symbol $c_{\rm diag}$ denotes a fixed exponent that dominates
the preliminary logarithmic loss in \cite[(3.41)]{Z1}.
It is chosen before $A_E,C_1,C_M$, as in \cite[\S3.10]{Z1}.
One may use $c_1=2^{125}+2^{126}+2$.
Then fix the BFI savings, the logarithmic powers of the moduli,
the SW exponents and the derivative bounds.
The later GR loss $C_{\rm GR}$ may depend on $C_1$;
it differs from $c_{\rm pre}$ in (1.29).

At these fixed parameters, Section~6 proves
the Class-II estimate below, where $C_{\rm pow}$ is the fixed logarithmic
exponent after the conversion $X\ge N/L^{31}$.
\[
 \begin{gathered}
 E_{\rm II}^D(N)\ll
 L^{K_0+2-A'}
 +N^{-c_\sigma}L^{K_0+C_{\rm pow}+2},\\
 c_\sigma=\frac{\ell_0\sigma^2}{4},\qquad
 \ell_0=\frac{264897}{2060000}>0.
 \end{gathered}                                             \tag{1.30}
\]
The first term is $L^{-3}$.
The second allows $X\ge N/L^{31}$.
The resulting fixed factor $L^{31c_\sigma}$ is absorbed in $C_{\rm pow}$.
Both terms tend to zero.
Section~6.5 gives the fixed-parameter closure and all eventual thresholds.
The endpoint enclosures of Appendix~A and the persistent losses of Section~6.4 hold uniformly for
$0<\sigma\le10^{-10}$, $0\le\tau\le10^{-90}$.
The effect of $\tau$ on the residual profile after (A.10a) is at most $5\cdot10^{-86}$.
The selected residual slack absorbs it.
The separate G1-3 PR peel is counted in (4.57) and (6.6).
Thus the numerical $D_{\rm cert}$ is fixed before $\sigma$,
and the choice (1.27) uses only $C_E<\infty$.

Every fixed cost is normalized by $4(1-\varepsilon_A)U_C$.
A raw error $(1-\varepsilon_A)eU_C$ contributes $e/4$ to $D$,
and each normalized error enters once.
The separate balanced-semiprime bound gives (1.22).

\setcounter{equation}{30}
\subsection{Distribution inputs}

This subsection states the distribution estimates used by the rows.
Row S1-0 uses \hyperref[sec:z1-input]{the complete coefficient \textup{(W)}} with
\[
\begin{gathered}
 u=1/500,\quad \Theta(u)=19101/32000,\quad\theta=19101/32000-10^{-4},\\
 B=30,\quad\varepsilon=10^{-4},\quad a=N,\quad K_{\rm scale}=31.
\end{gathered}
\]
Here $0<\sigma\le10^{-10}$ satisfies both restrictions, and
$\tau=\sigma^9$, $Y=N^{\theta\sigma^2}<N^{u/2}$.
We take $C_{\rm box}\ge\lceil2A+32\rceil$ and keep the grid, distinctness masks and smooth-divisor term.
Since $(q,2N)=1$, the shift $a=N$ is admissible.
Then \cite[Theorems~1.1--1.2, Corollary~1.3]{Z1} applies to the complete sum,
and the corollary gives the same Li comparator.
The remaining rows use the following local input.

\begingroup
\renewcommand{\thetheorem}{A}
\begin{theorem}[Local dispersion estimate]\label{thm:distribution-input}
Let $0<\sigma\le10^{-10}$ and let $N$ be large and even.
Put
\[
\begin{gathered}
 \tau=\sigma^9,\quad\delta=30\sigma,\quad\vartheta_{\rm sp}=7/32,
 \quad\eta_L=10^{-4},\\
 \ell_0={53\over412}-{\eta_L\over2}={264897\over2060000},
 \quad c_\sigma={\ell_0\sigma^2\over4},\\
 N/(\log N)^{31}\le X\le N,\quad L=\log X,\quad
 a=N,\quad z_0=X^{1/(\log\log X)^3},\quad u_0=N^{\ell_0\sigma^2}.
\end{gathered}
\]
For an interior Class-II Heath--Brown box $B$, define
\[
\begin{gathered}
 \mathcal S_t(q;a)=1_{t\equiv a\ (q)}-{1_{(t,q)=1}\over\varphi(q)},\quad
 \Delta_B(q;a)=\sum_{\mathbf v\in B}w(\mathbf v)\mathcal S_{\prod v_i}(q;a),\\
 S_B^{\rm PSDT}=\sum_qC_b(q)\Delta_B(q;N).
\end{gathered}
\]
Assume the local conditions below.
Then \cite[Theorem~3.2, (3.13)]{Z1} gives, for fixed $A'>0$,
\begin{equation}\label{eq:distribution-input}
 |S_B^{\rm PSDT}|\ll_{\sigma,31,C_{\rm box},A',\mathcal C_{\rm src}}
 XL^{-A'}+X^{1-c_\sigma}L^C.
\end{equation}
Here $C=C(\sigma,C_{\rm box},A',31,\mathcal C_{\rm src})$ is fixed,
and the constants in $\mathcal C_{\rm src}$ include $C_{\rm ctr},C_{\rm slot}$.
The bound is uniform in boxes, even $N$ and the shift $a=N$.
The constant and threshold may be ineffective.
First $\sigma>0$ is fixed, and then $N\to\infty$.
\end{theorem}
\endgroup

\paragraph{Application checklist.}
Each application uses literal prime sets and the row's own source.
The correspondence with \cite[\S3]{Z1} is
\begin{center}
\begin{tabular}{@{}ll@{}}
\toprule
Condition & Data checked below \\
\midrule
H1 & C1: exact unrestricted coefficient, repeats, $(q,2a)=1$ \\
H2 & C2: labelled factors, last two $u_0$-rough \\
H3 & C2--C3: packing, HB coefficients, majorant, doubled support, LEN \\
H4 & C4: positive shift and its uniform size \\
H5 & C5: selected window \\
H6 & Full-source logarithmic and boundary estimates \\
H7 & Full-family outer coefficient sum \\
\bottomrule
\end{tabular}
\end{center}
\textbf{C1.} For fixed row, prefix cell, AA37 sign and tail vector $b$ we require
\[
\begin{gathered}
 C_b(q)=\sum_{d,\mathbf p}\lambda_d1_{q=d\prod_{j\in J}p_j},\quad |\lambda_d|\le1,\quad p_j\ge u_L\ge u_0,\\
 u_L=L_{\rm rel}^{\sigma^2},\quad q\le X^{0.6},\quad(q,2N)=1.
\end{gathered}
\]
Here $L_{\rm rel}=L_c,L_b,D_c$ for DIRECT, S2, G1.
The labelled primes vary independently over fixed literal sets.
The divisor $d\mid P(u_L;N)$ is squarefree, with the AA37 cap, where
$P(y;N)=\prod_{p<y,\,p\nmid2N}p$.
Selectors are constant on blocks.
The equal-tail, prefix-diagonal, NC, cut and selector-transfer corrections are separated exactly,
and each is estimated separately.

\textbf{C2.} Let $J=J_1\sqcup J_2\sqcup J_3$ be a partition, with $d$ attached to slot 1.
We require, including repeated-prime multiplicities,
\[
 C_b=C_{\rm slot}(\gamma_1*\gamma_2*\gamma_3),\quad
 C_{\rm slot}=\prod_{i=1}^3|J_i|!=O_\sigma(1),\quad |\gamma_i|\le1.
\]
Slots 2 and 3 are $u_0$-rough.
An empty rough slot is $1_{m=1}$.
If $J_1=\varnothing$, put $\gamma_1(m)=\sum_d\lambda_d1_{m=d}$.
The dyadic scales $(Q,R,S)$ belong to slots $(1,3,2)$.
For the test scale $T=T_{\rm II}$ we require
\[
\begin{gathered}
 Q<TX^{-29\sigma},\quad QRS<M_I X^{-89\sigma},\quad
 N_I\ge QX^{\sigma/3},\quad M_I\ge X^{1/2},\\
 Q<TX^{-\sigma},\quad T^2R^2S<X^{1-8\sigma},\quad A_0=X^{1+\sigma/20},\\
 A_0^{\vartheta_{\rm sp}}TQR^2S^5<X^{2-15\sigma},\quad
 (TR/S)^{2\vartheta_{\rm sp}}TQR^2S^5<X^{2-15\sigma}.
\end{gathered}
\]
The labelled-tuple map proves the convolution identity.

\textbf{C3.} The source is either
\[
 1_{\mathbb P}(p)1_{N/(\log N)^{30}<p<(1-\varepsilon_A)N},\quad
 \varepsilon_A=10^{-4},
\]
or the plain product measure
\[
 A_G^{\rm cut}(e)=\sum_{r=3}^4\sum_{\mathbf p\in\mathcal D_r}
 w_r(\mathbf p)1_{e=p_1\cdots p_r}
 1_{\varepsilon_A N<e\le N\{1-(\log N)^{1-C_{\rm box}}\}}.
\]
The ordered $\mathcal D_r$ and bounded $w_r$ are fixed.
Each coordinate is at least $N^{100/1149}$.
All representation multiplicities are retained.
The plain measure is used with its full source predicates.
We apply Heath--Brown's finite identity \cite{ref6} with disjoint boxes of relative width $L^{-C_{\rm box}}$
and bound boundary boxes separately.
The selected $n$ is the eligible variable with largest lower-endpoint exponent.
Ties are broken by variable index, and $m$ is the complement.
For $\zeta\ge z_0$ we require
\[
\begin{gathered}
 \alpha_n\in\left\{\mu(n),1,{\log n\over L}\right\}
 1_{P^-(n)\ge\zeta}1_{I_N}(n),\quad |\alpha_n|\le1,
 \quad|\beta_m|\ll\tau_{23}(m),\\
 I_N\subset(T,2T],\quad N_I\asymp TL^{-C_{\rm box}},\\
 \psi_I(m)=\phi_0((m-m_0)/M_I),\quad M_I\asymp(X/T)L^{-C_{\rm box}},\\
 |m_0|\le C_{\rm ctr}X/T,\qquad 1\le|mn-N|\le X^3.
\end{gathered}
\]
Fixed HB coefficients and tags may be exterior scalars.
The function $\phi_0\in C_c^\infty([-1,1])$ is fixed independently of all parameters,
so its derivative seminorms are fixed and finite.
We require $0\le\phi_0\le1$ and $\phi_0=1$ on $[-1/2,1/2]$.
The complementary interval is majorized by $\psi_I$ on its doubled support.
The shift gap holds throughout the doubled support.
The same integer interval is used, with
\[
 \textbf{LEN:}\qquad N_I\ge(2T)^{1-2^{-126}}.
\]
Tagged comparisons and removed faces have separate estimates.

\textbf{C4--C5.} Every original modulus is coprime to $a=N$, and
\[
 0<a\le X(\log N)^{31}\le A_0,\qquad
 N^{60\sigma}\le T\le N^{1/3+14\sigma},\qquad
 X^{60\sigma}\le T\le X^{1/3+15\sigma}.
\]
For G1, C1--C5 and LEN are checked on the plain $A_G^{\rm cut}$.
This is \textbf{G1-LOCAL}.
Source faces are removed before interior boxes are selected.
The check covers exact HB coefficients, ambient intervals and the fixed majorant.
Lemma~\ref{lem:s2-preparation} gives the direct-source verification,
and Section~4 gives the product-source verification.

\paragraph{Local mechanism.}\label{p2:dispersion-proof}
The argument is \cite[(3.26)--(3.90)]{Z1}.
A first dispersion separates the diagonal and off-diagonal terms.
A second square and Poisson summation give moving-frequency Kloosterman sums.
The spectral estimate and C2 give \eqref{eq:distribution-input}.
The Kloosterman and spectral inputs follow the Deshouillers--Iwaniec
framework and Drappeau's dispersion estimate \cite{ref1,ref3};
the spectral-gap input is Kim--Sarnak \cite{ref7}.
The BFI and Siegel--Walfisz inputs in [Z1] are cited at their sources
\cite{ref2,Siegel1935,Walfisz1936,ref11}.
For S2 the local Mellin data are kept:
\[
\begin{gathered}
 q\asymp Q,\quad r_i\asymp R,\quad s_i\asymp S,\quad n_1,n_2\asymp T,\\
 f=(n_1-n_2)/q\ne0,\quad\eta=\sigma/1000,\quad H_*=40QR^2S^2X^\eta/M_I.
\end{gathered}
\]
On $H_0<|h|\le\min(2H_0,H_*)$, put
$\varsigma_h=\operatorname{sgn}(h)$ and $v=|h|S/(H_0s_1)$.
Let $w$ be smooth, supported in $(1/4,4)$ and equal to 1 on $[1/2,2]$.
Then
\begin{equation}\label{p2:eq:dispersion-mellin}
\begin{gathered}
 W_Y(v)=w(v)\widehat\phi_0(Y_2v)e(Y_1v),\quad
 Y_1={\varsigma_hH_0(a/n_1-m_0)\over Sqr_1r_2s_2},\quad
 Y_2={\varsigma_hH_0M_I\over Sqr_1r_2s_2},\\
 |Y_1|+|Y_2|\ll L^{C_{\rm box}+31}X^\eta=:L_{\rm osc},\qquad
 \int_{\mathbb R}\sup_Y|\widetilde W_Y(it)|\,dt\ll L_{\rm osc}.
\end{gathered}
\end{equation}
Both nonmodular phases remain in $W_Y$.
Put $m(t)=\sup_Y|\widetilde W_Y(it)|$.
We normalize by $m(t)>0$ and restore it before integration;
for $m(t)=0$ the integrand vanishes.
The details are in \cite[(3.56)--(3.65)]{Z1} and \S\ref{sec:s2-3}.

Type I uses \cite[Lemmas~3.3--3.5, \S3.11]{Z1} only after preparation.
Lemma~\ref{lem:s2-preparation} verifies its source and scale hypotheses.
The product source has only Class-II interior boxes (Section~4).
The respective choices are
\[
 D_f=X^{1/20},\quad V\ge qX^{1/10};\qquad D_f=X^{7.5\sigma/2}.
\]
At these prepared scales \cite[(3.102)--(3.115)]{Z1} applies.

\paragraph{Global conditions.}
The two global conditions bound the outer coefficients and the logarithmic and boundary errors.
For H7 we sum every outer coefficient:
{\scriptsize
\[
 \begin{gathered}
 \sum_j|c_j|\ll_{\sigma,C_{\rm box},\mathcal C_{\rm src}}
 (\log N)^{K_0},\\
 K_0\le24C_{\rm box}+26,\qquad A'=A+24C_{\rm box}+31.
 \end{gathered}
\]
}
Then \cite[(3.92)]{Z1} gives
\[
 E_{\rm II}\ll NL^{K_0-A'}+N^{1-c_\sigma}L^{K_0+C}.
\]
Since $K_0-A'\le-A-5$, the first term is $\ll NL^{-A-5}$.
For fixed $\sigma>0$, the second is also $O(NL^{-A-5})$.
The numerical application uses the per-box precision in (1.28).
For H6 we sum all labels, shells, boxes and representation multiplicities.
We require fixed $c_{\log},c_R$ with
\[
\begin{gathered}
 E_{\rm log}\ll N(\log N)^{c_{\log}-C_{\rm box}},\quad
 E_{\rm bdry}\ll N(\log N)^{c_R-C_{\rm box}/2},\\
 C_{\rm box}\ge\max\{A+c_{\log},2(A+c_R)\}.
\end{gathered}
\]
For the prime source, $(c_{\log},c_R)=(18,16)$
\cite[\S3 and Appendix~B]{Z1}.
Section~4 verifies the product-source values separately.
Together these give the aggregate $O(N(\log N)^{-A})$ error.

The next lemma checks packing on these rows.

\begin{lemma}[Packing on the actual row domains]\label{thm:packing-input}
Consider a row in the table below, including the three empty prefixes.
Let $c$ be its retained cell, at $\Theta_c=\inf_c\widetilde H$,
with its literal source domain and block selector.
For a nonempty prefix, let $t_1\ge\cdots\ge t_r$, $1\le r\le3$.
For either AA37 sign $\nu$, let
$\mathbf D\in\mathcal D^\nu(L_c)$ be any finite tail of length $k\ge0$.
Put
\[
\begin{gathered}
 \ell_c=\Theta_c-\eta_L-\sum_{i\le r}t_i\ge\ell_0,\quad L_c=N^{\ell_c},
 \quad z\le L_c\ (\text{upper}),\quad z^2\le L_c\ (\text{lower}),\\
 \rho=1+\sigma^9,\quad \delta=30\sigma,\quad\vartheta_{\rm sp}=7/32,
 \quad n\in[60\sigma,1/3+14\sigma],\\
 T=N_{\rm II}=N^n,\quad h=\rho(\Theta_c-\eta_L)+2\delta,
 \quad g=n-\delta,\quad A_{\max}=X^{1+\sigma/20}.
\end{gathered}
\]
Sort the endpoint exponents as $e_1\ge\cdots\ge e_m$, $m=r+k$.
Repeats are kept as distinct labelled atoms, and the list is padded with zeros.
Assume the DIRECT or plain G1 intervals satisfy
\[
 N_I\gg T(\log X)^{-C_{\rm box}},\qquad
 \log_NM_I\ge2/3-14.02\sigma,\qquad M_I\ge X^{1/2}.
\tag{P6}
\]
These bounds are uniform in row labels and boxes.
Then a deterministic labelled partition satisfies C2 at the same $\Theta_c$.
It attaches $d_{AA}$ to slot 1 and allows unit slots.
After dyadic and base transfer, the four spectral surpluses are at least
\[
 \left({7271\over256},\ 52-{1\over100},\
 {95221\over1280}-{1\over50},\ {22407\over256}-{1\over50}\right)\sigma
 >(28,51,74,87)\sigma.
\]
All four lower-scale conditions hold, with $\Delta D_{\rm PACK}=0$.
For G1-3 the cutoff is $z=N^{t_3}$ after the PR peel.
\end{lemma}

\begin{proof}
Put
\[
\begin{gathered}
 u=1/500,\quad\alpha=100/1149,\quad\beta=50/309,\quad\mu=25/128,\quad\lambda=93/400,\\
 t_*=7109/28000,\quad t_{\rm dd}=25/[32(1+3\rho)],\quad\overline H=153/256,\quad T_0=509/1000-\eta_L.
\end{gathered}
\]
Here $v,w,H,\widetilde H$ are as in \eqref{p2:eq:row-level-function}.
Each interval $[N^{t_i},N^{\rho t_i})$ is clipped to its row domain and keeps its label.
Missing coordinates lie in $[0,w_z]$, where $z=N^{w_z}$.
For $f_1\ge\cdots\ge f_m\ge0$, write
\[
 \mathrm{WF}(y):\quad\sum_{i<j}f_i+2f_j\le y\quad(1\le j\le m).
\tag{P1}
\]
The proof follows \cite[Theorem~4.1, (4.1)--(4.5)]{Z1},
with auxiliary identities \cite[Lemmas~4.2--4.4]{Z1}
and two-bin step \cite[Lemma~4.5]{Z1}.
The two-bin rule retains labels, repeats and zero padding.
Put $d_*=h-2\delta$ and $\Sigma_s=\sum_{i\in J_s}e_i$.
We verify the hypotheses of that theorem on these actual rows.

By \cite[Lemma~4.2]{Z1}, $v$ is concave and $\widetilde H\le\overline H$.
For $x_1\in[1/10,\rho t_*]$,
\[
 v(x_1)\ge\min\{v(1/10),v(127/500)\}>509/1000.
\]
For $x_1<1/10$, the active $w$ branch gives
\[
 \min\{(132-68x_3)/214,(153-96x_3)/256,1-2x_2\}
 \ge717/1280\ge11/20\ge(1+x_1)/2.
\]
Hence
\[
 509/1000\le\Theta_c\le\overline H,\qquad
 \Theta_c\ge m(t_1):=\min\{(1+t_1)/2,57/64-3\rho t_1/2\}.
\tag{P5}
\]
The first branch is smaller for $t_1\le t_{\rm dd}$.
At rank 1, $2t_1<2t_*=7109/14000<T_0$.
The remaining prefix checks are
\[
 \sum_{i<j}t_i+2t_j<\Theta_c-\eta_L\quad(1\le j\le r).
\tag{P4}
\]
For S1, S2, G1 the cutoff is $z=N^u$ at $r\le2$.
At $r=3$ it is $z=N^{t_3}$, after PR for G1.
The other rows use $z=N^\alpha$.

\begin{longtable}{@{}p{0.29\textwidth}p{0.66\textwidth}@{}}
\toprule
Rows and literal domains & Prefix checks and selectors \\
\midrule\endhead
S1-1, S1-2, S1-3; $[u,\alpha)$
& $(j+1)t_1\le4\alpha=400/1149<T_0$. \\[3pt]
S2-1, S2-2, S2-3; $[u,\beta)$
& Rank 2: $t_1+2t_2\le3\beta=150/309<T_0$.
Rank 3: $t_1<\beta<t_{\rm dd}$ and the whole-$K$ selector
$t_1+2t_2+(4+2\tau)t_3\le249/250+2\tau\beta$.
Thus $t_1+2t_2+4t_3<1-2\eta_L$;
$1-2\eta_L-249/250=19/5000>2\tau\beta$. \\[3pt]
G1-1, G1-2, G1-3; $[u,\mu)$
& Rank 2: $t_1+2t_2\le3t_1$,
$\mu<(1-2\eta_L)/5$, and
$3/256-\eta_L-(3/2)\tau\mu>0$.
Rank 3 requires $\Theta_c-\eta_L-t_1-t_2-(3+2\tau)t_3\ge0$.
Then $t_1+t_2+2t_3\le\Theta_c-\eta_L-(1+2\tau)t_3$,
$t_3\ge u/\rho>0$. \\[3pt]
S3-A: $[\alpha,\mu)$;
S4: $[\mu,1/4)$;
S5-low: $[1/4,t_*)$;
S7: $[\mu,\lambda)$
& One prefix each; $2t_1<2t_*<T_0$. \\[3pt]
S8-A: $\alpha\le t_2<t_1<\beta$
& $t_1+2t_2\le3\beta=150/309<T_0$. \\[3pt]
S9-A: $\beta\le t_1<\mu$, $\alpha\le t_2<\beta$
& $t_1<1-4\beta-2\eta_L$ for the first branch;
$3886931/49440000-(3/2)\tau\mu>0$ for the second. \\[3pt]
S10-Ret (the retained S10 row): $\mu\le t_1<\lambda$, $\alpha\le t_2<\beta$
& $t_2\le\operatorname{sel}(t_1)$,
$C_0=57/64-u-(11/5)\alpha$,
$\operatorname{sel}(t)\le B_\tau(t)\le\max\{\alpha,C_0-(5/2+3\tau/2)t\}$.
For $\Theta_c-u-t_1-t_2\ge(11/5)\alpha$,
the second branch in (P5) is active, since $t_1\ge\mu\ge t_{\rm dd}$.
Its affine alternative gives equality;
the $\alpha$ alternative gives $265391/9192000-(3/2)\tau\lambda>0$.
Thus $t_1+2t_2\le\Theta_c-u-(11/5)\alpha+t_2<\Theta_c-\eta_L$,
with binding prefix-mask gap
$u+(11/5)\alpha-\eta_L-\beta=12449531/394490000>0.03155$. \\[3pt]
S1-0, S2-0, G1-0; empty prefix
& $\Theta=H(u,u,u)=19101/32000$; tail endpoints $\le u$.
The active $w$ branch gives
$\widetilde H(\bar e_1,\bar e_2,\bar e_3)\ge
\min\{(132-68u)/214,(153-96u)/256,1-2u\}=19101/32000$. \\
\bottomrule
\end{longtable}

A prefix endpoint lies in $[t_i,\rho t_i]$.
A tail endpoint satisfies $\epsilon'_j\le\min\{\rho s_j,w_z\}$,
from $[N^{s_j},\min\{N^{\rho s_j},z\})$.
Every prefix prime is at least $z$.
The AA37 masks give
\[
 \sum_{j\le2l}s_j+3s_{2l+1}<\ell_c\ (\nu=+),\quad
 \sum_{j\le2l-1}s_j+3s_{2l}<\ell_c\ (\nu=-),\quad 2s_1<\ell_c.
\]
We drop one copy at cubic ranks and use $s_{j-1}\ge s_j$ between them.
Then $\sum_{i<j}s_i+2s_j<\ell_c$ at every tail rank.
The prefix sum is then added at these ranks.
If $s_1,\ldots,s_q$ precede $t_i$, then $t_i\le s_q$ and
\[
 \sum_{a<i}t_a+\sum_{b\le q}s_b+2t_i
 \le\sum_{a\le i}t_a+\left(\sum_{b<q}s_b+2s_q\right)
 <\Theta_c-\eta_L.
\]
Endpoint inflation by $\rho$ gives $\mathrm{WF}(d_*)$.
The first $r$ sorted endpoints lie in $[t_j,\rho t_j]$.
The rest are tails or zero pads in $[0,w_z]$.
Their first three entries therefore lie in the closed parent box.
Hence
\[
\begin{gathered}
 \Theta_c\le\widetilde H(\bar e_1,\bar e_2,\bar e_3),\quad
 \Theta_c-h=\eta_L-\tau(\Theta_c-\eta_L)-60\sigma>0.97\eta_L,\\
 1/2<h\le\overline H,\quad 0<\delta\le3\cdot10^{-9},\quad
 g=n-\delta\in[\delta,1/3-\delta/2].
\end{gathered}
\]
These verify \cite[Theorem~4.1, (4.2)]{Z1} for every finite tail.
The mixed-prefix checks above verify \cite[Theorem~4.1]{Z1} directly.
The pure-tail result \cite[Lemma~4.7]{Z1} is not used.
For the empty prefixes, the table gives the same hypotheses.

We attach $d_{AA}<L_c^\sigma=N^{\sigma\ell_c}$ to slot 1.
Then
\[
\begin{gathered}
 \log_NQ\le\Sigma_1+\sigma\overline H,\quad\log_NS\le\Sigma_2,\quad\log_NR\le\Sigma_3,\\
 n=g+\delta,\quad\log_NA_{\max}\le1+\sigma/20.
\end{gathered}
\]
By \cite[(4.3)]{Z1}, the four excess exponents are
\[
\begin{aligned}
 -(30-\overline H)\sigma&=-(1+7271/256)\sigma,\qquad -2\delta=-60\sigma,\\
 -3\delta+(\vartheta_{\rm sp}/20+\overline H)\sigma
 &=-(15+95221/1280)\sigma,\\
 -(3+2\vartheta_{\rm sp})\delta+\overline H\sigma
 &=-(15+22407/256)\sigma.
\end{aligned}
\]
Dyadic representatives within a factor two cost $24\log2/\log N<\sigma/1000$.
For $k_X=\log_XN=1/(1-\omega)$, $\omega=\log_N(N/X)$,
$0\le k_X-1<0.01\sigma$.
The first exponent improves.
The second costs at most $0.01\sigma$, and the last two cost at most $0.02\sigma$.
These give the stated spectral margins.

The factorization follows \cite[Lemma~4.8, (4.12)--(4.13)]{Z1}.
The small-prime divisor is unique in slot 1.
Each labelled slot has at most $|J_s|!$ multiset orderings.
Thus its finite-sum identity holds for these literal sets, including repeats.
Normalization gives $|\gamma_s|\le1$; C2 covers empty slots.
Lemma~4.8 bounds tail-collision support only; source corrections remain separate.
Also,
\[
 L_c^{\sigma^2}\ge u_0,\quad p_{\rm prefix}\ge N^u,\quad
 s_j\ge\sigma^2\ell_c,\quad\sum_js_j<\ell_c,\quad
 k<\sigma^{-2},\quad C_{\rm slot}=O_\sigma(1).
\]
Hence slots 2 and 3 are $u_0$-rough.

Condition (P6) is verified separately for DIRECT and plain G1.
The first margin gives $Q<TX^{-29\sigma}$, so
\[
 {N_I\over QX^{\sigma/3}}\gg
 (\log X)^{-C_{\rm box}}X^{(29-1/3)\sigma}\longrightarrow\infty.
\]
The last WF mask gives $\sum_s\Sigma_s\le d_*<\overline H$.
Thus
\[
\begin{gathered}
 \log_N(QRS)\le153/256+\sigma,\\
 \log_NM_I-\log_N(QRS)\ge
 2/3-153/256-15.02\sigma=53/768-15.02\sigma>89\sigma.
\end{gathered}
\]
Since $\sigma\le10^{-10}$, $QRS<M_I X^{-89\sigma}$.
Condition (P6) also gives $M_I\ge X^{1/2}$.
The same rule covers zero pads and unit slots.
For G1-3, the comparison $N^{\rho t_3}\to N^{t_3}$ is proved separately.
Packing uses the latter cutoff at the original $\Theta_c$,
so its level loss is zero.
\end{proof}

Sections~2--3 verify direct-row source conditions; Section~4 verifies plain G1.
Each source comparison and removed boundary term has its separate estimate.

\section{The sieve rows S1--S10 on the actual source}
\label{p2:sec:rows}
This section treats the direct sieve rows on the actual source.
Put $\rho=1+\tau$, $\mu=25/128$. The rows are grouped as
\[
\begin{array}{c|l}
\mathrm{DIRECT}&\mathrm{S1\hbox{-}0\hbox{--}S1\hbox{-}3,S3\hbox{-}A,S4,S5\hbox{-}low,S7,S8\hbox{-}A,S9\hbox{-}A,S10\hbox{-}Ret}\\
\mathrm{S2}&\mathrm{S2\hbox{-}0\hbox{--}S2\hbox{-}3}\quad\text{(whole-$K$ at depth three)}\\
\mathrm{BV}&\mathrm{S5\hbox{-}top,S6}.
\end{array}
\]
Sections~3 and~5 treat S2 and BV.
Conditions C1--C5 are those of Theorem~\ref{thm:distribution-input}.
The internal Lemma~\ref{thm:packing-input} gives the slots.

\subsection{Source replacement and the exact signed coefficient}
This subsection replaces the literal source by the plain one.
\begin{lemma}[Source replacement and coefficients]
\label{lem:source-replacement}
Let $\mathcal P_N=\{p>2:p\nmid N\}$ and
\[
 A_D(v)=\mathbf1_{\mathbb P}(v)
 \mathbf1_{N/\log^{30}N<v<(1-\epsilon_A)N}.
\]
For a squarefree prefix $b$ from $\mathcal P_N$, set
\[
\begin{aligned}
\mathcal A_b&=\{v\in\mathbb P:v<(1-\epsilon_A)N,\ v\nmid N,\ b\mid N-v,
                         \ p^2\nmid N-v\ (p\mid b)\},\\
\mathcal A_b^{\rm pl}&=\{v:A_D(v)=1,\ b\mid N-v\}.
\end{aligned}
\]
For $A\subseteq\mathcal A_b\cup\mathcal A_b^{\rm pl}$, the notation
$S(A,\mathcal P_N,z)$ means
\[
 \#\left\{v\in A:\left(\frac{N-v}{b},P_N(z)\right)=1\right\}.
\]
Thus this sieve tests the cofactor, not the prime $v$ or the integer $N-v$.
If every prefix prime is $\ge z$, the literal sieve is
$S(\mathcal A_b,\mathcal P_N,z)$. Its source comparison is
\[
 S(\mathcal A_b,\mathcal P_N,z)-S(\mathcal A_b^{\rm pl},\mathcal P_N,z)
 =S(\mathcal A_b\setminus\mathcal A_b^{\rm pl},\mathcal P_N,z)
  -S(\mathcal A_b^{\rm pl}\setminus\mathcal A_b,\mathcal P_N,z).
 \tag{2.17a}\label{p2:eq:direct-source-comparison}
\]
The normalized unsifted bounds, in Section~6 units, are
\[
\begin{array}{c|c}
v\le N/\log^{30}N&O_u(L_N^{-28})\\
v\mid N&O_u(L_N^3/N)\\
p^2\mid N-v\text{ for a prefix prime}&O_u(N^{-u/(1+\tau)}L_N^3).
\end{array}
\]
Fix an AA37 sign $\nu$, a row cell and tail labels.
All literal clipped sets and the smooth divisor are kept. Put
\[
 C_c(q)=\sum_{d,(p_j)}\lambda_d\mathbf1_{q=d\prod_jp_j},\qquad |\lambda_d|\le1.
\]
Let $K_T$ be the equal-label tail mask,
and let $K_P$ impose complete equal-label prefix distinctness.
For a complete permutation-invariant block, let $M_c$ be its prefix orbit factorial.
Then
\[
 O_c(q)=\frac{C_c(q)}{M_c}-\Xi_{U,c}(q),\qquad
 \Xi_{U,c}=\frac1{M_c}\sum_{\rm full}\lambda_d(1-K_TK_P)
                  \mathbf1_{q=d\prod_jp_j}.
 \tag{2.17b}\label{p2:eq:direct-c1}
\]
With $\gamma_c=C_c/C_{\rm slot}$, the full remainder is
\[
 R_{{\rm AA},c}=\frac{C_{\rm slot}}{M_c}
       \sum_q\gamma_c(q)\Delta_P(q)
       -\sum_q\Xi_{U,c}(q)\Delta_P(q),
 \qquad \gamma_c=C_c/C_{\rm slot}.
 \tag{2.17c}\label{p2:eq:direct-remainder}
\]
The external multiplier $c\nu$ acts on both terms.
This includes the four S2 children and the whole-$K$ last child.
The same identity holds for nonnegative occurrence weights.
On the common occurrence space, $a_r,b_r$ are the original and plain weights.
The prefix and sifting data are fixed, and the weights are extended by zero. Put
\[
 h_r=\min(a_r,b_r),\qquad
 e_r^+=(a_r-b_r)_+,\qquad e_r^-=(b_r-a_r)_+.
\]
Then $a_r=h_r+e_r^+$, $b_r=h_r+e_r^-$ and $e_r^+e_r^-=0$.
Linearity gives $S_z(a_r)=S_z(b_r)+S_z(e_r^+)-S_z(e_r^-)$,
including the G1 representation multiplicities.
\end{lemma}
\begin{proof}
For $m=(N-v)/b$, reduction modulo a common prime gives
\[
 (m,N)=1\iff v\nmid N,\qquad
 (m,b)=1\iff p^2\nmid N-v\quad(p\mid b).
\]
Here $v\nmid b$, since otherwise $v\mid N$.
These conditions concern the ambient source; cofactor sieving is separate.
Additivity proves \eqref{p2:eq:direct-source-comparison}.
The two differences have the three stated exceptional conditions.
Each point is assigned to the first true condition,
and the resulting unsifted masses are bounded.
These source terms are removed before HB.
A common source point removed by sieving creates no source error.
All later Buchstab and distribution calculations use $A_D$.

For (2.17b) we use $1=K_TK_P+K_T(1-K_P)+(1-K_T)$.
Overlapping repetitions are assigned to the first forbidden-tail condition.
Each modulus of $\Xi_{U,c}$ repeats a prime $\ge L_c^{\sigma^2}$.
Theorem~\ref{thm:repeated-primes} bounds the complete nineteen-row union,
including prefix multiplicities, once.
The smooth divisor lies wholly in slot 1.
Then the finite slot convolution is $C_c/C_{\rm slot}$ on every modulus, repeats included.
Empty rough slots are unit atoms, and each realization and factorial occurs once.

Put $C_A(q)=\sum_{v\equiv N\,(q)}A_D(v)$,
$X_A=\sum_vA_D(v)$, $Y_A(q)=\sum_{(v,q)=1}A_D(v)$ and
\[
 \Delta_0=C_A-X_A/\varphi(q),\qquad \Delta_P=C_A-Y_A/\varphi(q).
\]
Prefixes are $\ge z$ and tails are $<z$.
Totient multiplicativity converts the relative AA37 comparator into $\Delta_0$.
Theorem~\ref{thm:repeated-primes} gives $q<N^{3/5}$ on the full labelled support.
Eventually $v>N/L_N^{30}>N^{3/5}$.
Thus $Y_A=X_A$, so $\Delta_0=\Delta_P$ before HB.
This proves \eqref{p2:eq:direct-remainder}.
\end{proof}

For the actual levels put
\[
\begin{aligned}
v(t)&=\min\{(1+t)/2,(32-34t)/39,57/64-3t/2\},\\
\psi(x,y)&=x\mathbf1_{x\ge y},\\
w(a,b,c)&=\psi\!\left(\min\{(132-68c)/214,(153-96c)/256,1-2b\},(1+a)/2\right),\\
H(a,b,c)&=\max\{v(a),v(b),v(a+b),v(a+b+c),\\
 &w(a,b,c),w(a,c,b),\\
 &\psi(v(a+c),a+2b+c),\\
 &\psi(v(b+c),2a+b+c)\}.
\end{aligned}
 \tag{2.17d}\label{p2:eq:row-level-function}
\]
Let $\widetilde H=H$ when the largest coordinate is $\le\mu$,
and above $\mu$ let $\widetilde H(a,b,c)=v(a)$.
A parent label $t_i$ retains $[N^{t_i},N^{\rho t_i})$, clipped by its row.
At depth $r<3$, complete the ordered triple in $[0,w_z]$, where $z=N^{w_z}$.
Let $\Theta_c$ be the infimum on this closed ordered parent domain.
At depth three the infimum is taken on the three parent intervals.
This infimum is taken before any branch is selected, including crossings at $\mu$.
At depth zero, with $z=N^u$, $\Theta_c=H(u,u,u)=19101/32000$.
Then
\[
 L_c=N^{\ell_c},\qquad \ell_c=\Theta_c-\eta_L-\sum_{i\le r}t_i,
 \qquad D_c=N^{\Theta_c-\eta_L}.
 \tag{2.17e}\label{p2:eq:row-relative-level}
\]
For nonempty DIRECT rows, $t_1<t_*=7109/28000$ and
$509/1000\le\Theta_c\le153/256$. Also
\[
 \Theta_c\ge m(t_1):=\min\{(1+t_1)/2,57/64-3\rho t_1/2\}.
 \tag{2.17f}\label{p2:eq:row-level-floor}
\]
The first branch is active through $t_{\rm dd}=25/[32(1+3\rho)]$.
The required prefix masks are
\[
 t_1+\cdots+t_{j-1}+2t_j<\Theta_c-\eta_L
 \qquad(1\le j\le r).
 \tag{2.17g}\label{p2:eq:direct-prefix-mask}
\]
Their row checks are
\[
\begin{array}{c|l}
\text{rank one}&2t_1<2t_*=7109/14000<509/1000-\eta_L\\
\mathrm{S1\hbox{-}2,S1\hbox{-}3}&t_i<\alpha,\quad4\alpha<509/1000-\eta_L\\
\mathrm{S2\hbox{-}2,S2\hbox{-}3}\ (j=2),\ \mathrm{S8\hbox{-}A}
 &3\beta<509/1000-\eta_L\\
\mathrm{S2\hbox{-}3}\ (j=3)&t_1<\beta<t_{\rm dd}\\
 &t_1+2t_2+(4+2\tau)t_3\le249/250+2\tau\beta<1-2\eta_L\\
\mathrm{S9\hbox{-}A}&t_1<\mu,\ t_2<\beta:\quad t_1+2t_2<m(t_1)-\eta_L\\
\mathrm{S10\hbox{-}Ret}&\Theta_c-u-t_1-t_2\ge(11/5)\alpha\\
 &t_2<\beta<u+(11/5)\alpha-\eta_L.
\end{array}
\]
S2 at rank three uses the first branch of \eqref{p2:eq:row-level-floor}.
S9 substitutes in both branches.
For S10 we use \S\ref{p2:sec:2.10}; its binding prefix-mask gap is
\[
 u+(11/5)\alpha-\eta_L-\beta
 =\frac{12449531}{394490000}>0.03155.
\]
These checks include clipped cells and weak selection endpoints.
AA37 tail labels $D_j=N^{s_j}$ satisfy $\sum_{i<j}s_i+2s_j<\ell_c$.
Together with \eqref{p2:eq:direct-prefix-mask}, prefix/tail separation gives
$\sum_{i<j}\lambda_i+2\lambda_j<\Theta_c-\eta_L$ for merged sorted labels.
Sorted upper endpoints are at most $\rho$ times their labels,
and their mask is strictly below $\rho(\Theta_c-\eta_L)$.
This proves packing geometry, including empty prefixes.
The tested-source and boundary hypotheses are proved below.

For the Section~4.1 representation domains let $R_j$ count ordered representations.
In $R_j^{\rm tag}$ the coordinates are prime to $N$
and $e/(p_1p_2)$ is prime to $p_1p_2$. Put
\[
\begin{aligned}
 A_G&=(2R_{11}+R_{12}+R_{13})\mathbf1_{\epsilon_AN<e<N},\\
 w_{\rm tag}&=(2R_{11}^{\rm tag}+R_{12}^{\rm tag}+R_{13}^{\rm tag})
                    \mathbf1_{\epsilon_AN<e<N},\\
 w_{\rm pl}&=A_G^{\rm cut}=A_G\mathbf1_{e\le N_{\rm cut}},\quad
 N_{\rm cut}=N(1-(\log N)^{1-C_b}),\quad C_b=334,\\
 U_G&=w_{\rm tag}\mathbf1_{e>N_{\rm cut}},\quad
 E_G=(A_G-w_{\rm tag})\mathbf1_{e\le N_{\rm cut}}.
\end{aligned}
\]
Thus $U_G,E_G\ge0$ and $w_{\rm tag}-w_{\rm pl}=U_G-E_G$.
For $\chi_{b,z}(e)=\mathbf1_{b\mid N-e}\mathbf1_{(N-e,P_N(z))=1}$,
\[
 \sum_e w_{\rm tag}\chi_{b,z}=\sum_e w_{\rm pl}\chi_{b,z}+\Delta_G(b,z),
 \qquad |\Delta_G(b,z)|\le\sum_e(U_G+E_G).
\]
Here $U_G$ is the upper cut; $E_G$ counts coordinate divisibility and repetitions.
Section~4 bounds their unsifted masses and local hypotheses.
The source is replaced before Buchstab, and the signed parent correction is retained.
C1--C5 test the untagged discrepancy $\Delta[A_G^{\rm cut}]$ on interior boxes.
Fixing two coordinates at a prime $p$ leaves at most two positive coordinates.
The elementary divisor bound, over all pairs and four source forms, gives
\[
 O\!\left(\sum_{p\ge N^\alpha}\frac N{p^2}(1+\log N)^2\right)
 =O(N^{1-\alpha}(1+\log N)^2).
\]
Division by $\mathcal W\gg N/(\log N)^2$ gives $O(N^{-\alpha}(\log N)^4)$.
This is within the Section~6 source-restoration allowance.

We use the classical BFI hypothesis (A2), \cite{ref2}.
We fix $\mathscr G$ and $B_\gamma>0$.
For every $A>0$ it asks for a uniform $C_A$ such that
\[
 \left|\sum_{\substack{n\equiv\ell\ (k)\\(n,d)=1}}\gamma_n
 -\frac1{\varphi(k)}\sum_{(n,dk)=1}\gamma_n\right|
 \le C_A\|\gamma\|_2T^{1/2}\tau(d)^{B_\gamma}(\log(2T))^{-A}.
\]
Here $\gamma\in\mathscr G$, $\operatorname{supp}\gamma\subseteq[1,T]$, $T\ge2$,
$d,k\ge1$, $\ell\ne0$, $(\ell,k)=1$.
The constant is uniform in the sequence, subinterval and auxiliary coprimality divisor.
For Class II we use the selected elementary factor
and exchange factor names as in the classical theorem.
Theorem~\ref{thm:distribution-input} requires this $\ell^2$ normalization.
The complementary convolution has its separate norm bound.


\subsection{Common direct preparation}
This subsection verifies the local and global hypotheses for the direct families.
\begin{lemma}[Common preparation]
\label{lem:s2-preparation}\label{lem:direct-preparation}
Let $0<\sigma\le10^{-10}$, $\tau=\sigma^9$, $C_b=334$.
Then $C_b\ge\lceil2A+32\rceil$ for $0<A\le(334-32)/2$.
The per-box precision $A'$ is arbitrary; it is distinct from the
prime-source target $A$ in the condition $C_b\ge\lceil2A+32\rceil$.
Each DIRECT and S2 family satisfies H1--H5 and LEN of
\cite[Theorem~3.2]{Z1} on retained Class-II boxes.
For each family $\mathscr F$, separately, H6 and H7 hold with
\[
\begin{gathered}
 E_{\rm log,\mathscr F}\ll_\sigma NL_N^{18-C_b},\quad
 E_{\rm bdry,\mathscr F}\ll_\sigma NL_N^{16-C_b/2},\quad L_N=\log N,\\
 \sum_{j\in\mathcal J_{\mathscr F}}|c_j|\ll_{\sigma,C_b}L_N^{K_*},\qquad
 K_*=6(C_b+1)+2=2012<24C_b+26=8042.
\end{gathered}
\]
Every other nonboundary box is Class I-1 or I-2.
For fixed $A_I>0$, its signed discrepancy before external coefficients is
$O_{\sigma,A_I}(X(\log X)^{-A_I})$.
The entire S2 Class-I error is $o(W_N)$, with
\eqref{p2:eq:s2-normalization} defining $W_N$.
Section~3 treats the strict BV complement separately.
DIRECT and S2 are each counted once.
\end{lemma}
\begin{proof}
Put $L=\log N$, $H=N/L^{30}$, $\Gamma_k=C_k/M_k$,
$P=\sum_{i\le r}t_i$, $\ell=\Theta_c-\eta_L-P\ge\ell_0$.
The source is $A_D(v)=\mathbf1_{\mathbb P}(v)\mathbf1_{H<v<(1-\epsilon_A)N}$.
Source exceptions have been removed.
All prefix, sieve and whole-cell conditions are modulus-side.
The row masks and AA37 bound $\sum s_j<\ell$ give
\[
\begin{gathered}
 q<N^{\rho P+(\rho+\sigma)\ell}
 \le N^{(1+\tau+\sigma)(153/256-\eta_L)}<N^{e_D},\\
 e_D=153/256-\eta_L/2<153/256.
\end{gathered}
\]
For the dyadic $X$-shell covering $v>H$, one has $X>H/2$;
eventually $q\le X^{3/5}$ and $q<H$.
All factors avoid $2N$.
The tested prime exceeds $q$, so its comparators agree before HB.
Let $u_c=L_c^{\sigma^2}\ge N^{\ell_0\sigma^2}$.
Primes below $u_c$ determine the smooth divisor uniquely.
Labelled primes are $\ge u_c$, since $u_c<N^u$.
Thus $u_c^m\le q<N^{3/5}$ gives
\[
 |\Gamma_k(q)|\le m!/M_k\le C_\sigma,\qquad
 m\le\left\lceil\frac3{5\ell_0\sigma^2}\right\rceil.
\]
This includes repeated primes and each numerical realization once.
Three prefixes, $k<\sigma^{-2}$ tails and two signs give $O_\sigma(1)$ blocks.
Whole-cell conditions remove blocks, and smooth divisors remain inside coefficients.

We apply \cite[(B.7)--(B.13)]{Z1} with
$B=30$, $\varepsilon=\epsilon_A$, $a=N\ne v$.
For $j\le3$ we take $s=2j\le6$, $n\in[X,2X)$, $m_i\le2X^{1/3}$.
Factors are $z_0(X)$-rough, including the unit atom.
Put
\[
\begin{gathered}
 \delta_X=(\log X)^{-C_b},\quad J_X=\lceil\log2/\log(1+\delta_X)\rceil,
 \quad r_X=2^{1/J_X},\quad\Delta_X=r_X-1,\\
 \delta_X/3\le\delta_X(1-2\delta_X)\le\Delta_X\le\delta_X.
\end{gathered}
\]
The complete tested-side faces are
\[
 n=H,\quad n=(1-\epsilon_A)N,\quad n=X,2X;\qquad
 m_i=2X^{1/3}\quad(1\le i\le j).
\]
Crossing boxes are removed.
Roughness stays in the coefficients, and whole-$K$ stays modulus-side.
The bound $2X(1+\delta_X)^6<3X$ makes $n\le8X$ redundant.
The cited construction uses these data and the proved coefficient bound,
so it permits prefixed moduli.
Proper powers contribute $O_\sigma(N^{1/2+o(1)})$:
$O(N^{1/2}\log N)$ terms, each $N^{o(1)}$.
The signed HB coefficient is
\[
 c_j(\mathbf v)=(-1)^{j-1}\binom3j
       \prod_{i\le j}\mu_{z_0}(m_i)\log n_1,\qquad |c_j(\mathbf v)|\le3L.
\]

For H6 put $F(n)=d(N-n)+2(1+L)$.
Divisor counting and $\sum_{q\le N^{3/5}}1/\varphi(q)\ll1+L$ give
$|\sum_q\Gamma_k(q)\mathcal S_n(q;N)|\ll_\sigma F(n)$.
Each product fixes its shell by \cite[(B.11)]{Z1}.
We use $\sum_{n\le y}\tau_h(n)\le y(1+\log y)^{h-1}$,
$\tau_6^2\le\tau_{36}$ and $d^4\le\tau_{16}$.
Cauchy--Schwarz gives $(35+15)/2=25$, $(5+25)/2=15$, and
\[
 \sum_{n<N}\tau_6(n)\ll NL^5,\quad
 \sum_{n<N}\tau_6(n)F(n)^2\ll NL^{25},\quad
 \sum_{n<N}\tau_6(n)F(n)\ll NL^{15}.
\]
For $\lambda_B=\sum_i\log V_i$, we have
$\log n,\lambda_B\asymp L$, $0\le\log n-\lambda_B\le6\delta_X$,
$\log X=L+O(\log L)$.
Hence
\[
 |E_{\rm log,\mathscr F}|\ll_\sigma
 L^{-C_b-1}\sum_{n<N}\tau_6(n)F(n)\ll_\sigma NL^{14-C_b}.
\]
This is stronger than stated.
Each tuple has one shell and one box.
For faces we use \cite[Lemma~B.2, (B.14)--(B.16)]{Z1}, with
$2\le s\le6$, $Y_0=N^{1/7}$, $\eta_{\rm dist}\ll L^{-C_b}$.
Source faces are singletons, and the lattices are
$(\log2)\mathbb Z$ and $\log2+(\log2/3)\mathbb Z$.
Active product coordinates exceed $(H/2)^{1/6}$.
At caps, $m_i\asymp X^{1/3}$ and $X>H/2$.
Both exceed $Y_0$ eventually. Thus
\[
\begin{gathered}
 \#\mathcal U\ll NL^{5-C_b}+N^{6/7}L^5\ll NL^{5-C_b},\\
 |E_{\rm bdry,\mathscr F}|\ll_\sigma
 L(\#\mathcal U)^{1/2}(NL^{25})^{1/2}\ll_\sigma NL^{16-C_b/2}.
\end{gathered}
\]
Here $|c_j|/\lambda_B\ll L$.
Finite faces and unique shells introduce no box count.
Three patterns and $O_\sigma(1)$ blocks alter only constants.
This proves H6 separately for both families.

The classification uses \cite[(B.20)--(B.24), (B.27)--(B.28), (B.30)]{Z1}.
We choose $N_0(\sigma)$ so that
\[
 \begin{gathered}
 \omega=\log_N(N/X),\quad\kappa=(\log2+6\log r_X)/L,\\
 \omega,\kappa\le\sigma/202,\quad
 \sigma L\ge2^{130}(C_b+1)(\log L+1).
\end{gathered}
\]
The product and M\"obius caps are exactly (B.21).
The selected variable has the largest eligible lower exponent in $[60\sigma,1/3+14\sigma]$,
with ties broken by index.
The exclusions are $14\sigma>\kappa$, $360\sigma<1-\omega-\kappa$,
$1+42\sigma>1-\omega+\kappa$.
Every other box has one or two large smooth variables.
For the selected Class-II cell $[T,r_XT)$ put
\[
 I_T=(T,r_XT],\quad N_I=\lfloor r_XT\rfloor-\lfloor T\rfloor,\quad
 M_I=M(r_X^r-1),\quad m_0=M(1+r_X^r)/2.
\]
Here $M$ is the complementary lower product, $TM\asymp X$.
Its nonunit-cell count satisfies $1\le r\le5$,
since eventually $n>H>N^{1/3+15\sigma}$.
By \cite[Lemma~B.3, (B.22)--(B.24)]{Z1}, internal endpoints are irrational.
Rough coefficients vanish at powers of two.
Both integer counts differ from $T\Delta_X$ by less than one.
The threshold gives $T\Delta_X>2$ and $T^{2^{-126}}\ge12(\log X)^{C_b}$.
Consequently
\[
\begin{gathered}
 \frac{T}{6(\log X)^{C_b}}\le N_I\le2T(\log X)^{-C_b},\qquad
 M_I\asymp(X/T)(\log X)^{-C_b},\\
 N_I\ge(2T)^{1-2^{-126}},\quad X^{60\sigma}\le T\le X^{1/3+15\sigma},\\
 \log_NM_I\ge2/3-(14+1/202)\sigma-o(1)\ge2/3-14.02\sigma,\\
 M_I\gg X^{2/3-15\sigma}(\log X)^{-C_b}\ge X^{1/2}.
\end{gathered}
\]
The majorant in \cite[Appendix~B, immediately before (B.27)]{Z1}
equals one on $[M,Mr_X^r]$,
with support $[M-M_I/2,Mr_X^r+M_I/2]$.
For a source point $n_0$, enlarged products satisfy
\[
 n_0r_X[1+\tfrac32(r_X^r-1)]\le n_0(1+11\delta_X)<(1-\epsilon_A/2)N.
\]
Thus $N-mn\ge\epsilon_A N/2$, $|mn-N|\le X^3$, $|m_0|\ll X/T$.
Eventually $N\le XL^{31}\le X^{1+\sigma/20}$.
The row masks and Lemma~\ref{thm:packing-input} give exact repeated-prime slots,
including S1-0 and S2-0, and all spectral inequalities.
In particular,
\[
 Q<TX^{-29\sigma},\quad QRS<M_IX^{-89\sigma},\qquad
 \frac{N_I}{QX^{\sigma/3}}\ge
 \frac{X^{(29-1/3)\sigma}}{6(\log X)^{C_b}}\longrightarrow\infty.
\]
For rough M\"obius, unit and normalized-log factors
we use \cite[(3.25), (3.47)--(3.48)]{Z1}.
The complement is bounded by $\tau_5$.
This proves H1--H5 and LEN separately from H6/H7.

For Class I we apply \cite[Lemmas~3.3 and~3.5, \S3.11, (B.30)]{Z1}.
Here $k<\sigma^{-2}$, $|\Gamma_k|\le C_\sigma$, complement $\le\tau_5$.
The nonnegative weights $1,\log v/L$ have bounded variation and supremum.
Predicates are constant and factors are $z_0(X)$-rough.
The discrepancy remains $\mathcal S_{\prod v_i}(q;N)$.
Assign proper powers, LOG, boundary, II, I-1, I-2 in order.
The classification gives
\[
 e_*-e_D-(1-\omega)/10\ge0.9-e_D-0.9\omega-\kappa-300\sigma>0.3023,
\]
so $V\ge qX^{1/10}$.
For I-2 we order $N_1\ge N_2$. Then
\[
 \begin{gathered}
 M\le X^{(240\sigma+\kappa)/(1-\omega)}\le X^{2.401\cdot10^{-8}},\\
 N_2\ge X^{1/3},\quad N_1N_2M\le2^{30}X,\quad q<X^{0.5977}.
\end{gathered}
\]
Larger $N_0$ absorbs fixed dyadic factors into $\kappa$.
Thus \cite[(3.100)--(3.101)]{Z1} hold, and (3.105) and (3.111) apply.
For H7, at most six coordinates and three patterns have
$O(\log\log N)\le O(L)$ shells.
Each coordinate has $O(L)$ dyadic intervals and $J_X=O(L^{C_b})$ cells.
Labels, signs, binomial factors and slot factorials cost $O_\sigma(1)$.
The sole logarithmic scalar is $L/\lambda_B=O(1)$.
Hence
\[
 \sum_{j\in\mathcal J_{\mathscr F}}|c_j|
 \ll_{\sigma,C_b}L^{6(C_b+1)+1}\le C_{\sigma,C_b}L^{K_*}.
\]
The extra unit covers the I-2 small-product split.
Realizations and smooth divisors remain inside coefficients.
PF and the selected boundary region form one family.
With $A_I=K_*+6$, the Class-I total is $O_\sigma(NL^{-6})$.
Since $W_N\gg NL^{-2}$, its normalized bound is $O_\sigma(L^{-4})$.
Normalized LOG and boundary errors are
$O_\sigma(L^{20-C_b}+L^{18-C_b/2})=o(1)$.
\end{proof}


\subsection{Four-term Buchstab and the S1 comparison}
\label{p2:sec:2.1}\label{p2:sec:2.2}
This subsection proves the Buchstab chains and bounds the cutoff comparison.
\begin{lemma}[Four-term chains]
Set $S_{\rm pl}(b,z)=S(\mathcal A_b^{\rm pl},\mathcal P_N,z)$.
The plain-source identity is
\[
\begin{aligned}
3S_{\rm pl}(1,N^\alpha)
={}&3S_{\rm pl}(1,N^u)
 -3\sum_{u\le r_1<\alpha}S_{\rm pl}(p_1,N^u)\\
&+3\sum_{u\le r_2<r_1<\alpha}S_{\rm pl}(p_1p_2,N^u)
 -3\sum_{u\le r_3<r_2<r_1<\alpha}S_{\rm pl}(p_1p_2p_3,p_3).
\end{aligned}\tag{2.19}\label{p2:eq:2.19}
\]
Here $r_i=\log p_i/\log N$ and all $p_i\in\mathcal P_N$.
The S2 identity (3.2) replaces $\alpha$ by $\beta$ and divides by $3$.
Both use $L_c=N^{\Theta_c-\eta_L-P}$, $P=\sum_{i\le r}t_i$, $A_R=c/2$.
Their parameters are
\begin{center}
\begin{tabular}{@{}cllll@{}}\toprule
Row & $c$; direction & Cutoff & Floor for $\ell_c$ & $M_c$\\\midrule
S1-0 & $3$; lower $f$ & $N^u$ & $19101/32000-\eta_L$ & $1$\\
S1-1 & $-3$; upper $F$ & $N^u$ & $5243851/11490000$ & $1$\\
S1-2 & $3$; lower $f$ & $N^u$ & $1414617/3830000$ & $1,2$\\
S1-3 & $-3$; upper $F$ & $E_3=N^{t_3}$ & $3743851/11490000$ & $1,2,6$\\
S2-0 & $1$; lower $f$ & $N^u$ & $95489/160000$ & $1$\\
S2-1 & $-1$; upper $F$ & $N^u$ & $1294691/3090000$ & $1$\\
S2-2 & $1$; lower $f$ & $N^u$ & $264897/1030000$ & $1,2$\\
S2-3 & $-1$; upper $F$ & $E_3=N^{t_3}$ & whole-$K$: \S\ref{sec:s2-3} & $1,2,6$\\\bottomrule
\end{tabular}
\end{center}
The S1 functional has signs $+,-,+,-$ and densities
$f/u,-F/u,+f/u,-\Phi$. Appendix~A.2 gives
\[
 12.9020516<\mathcal S_1<12.9023390,\qquad\tfrac34\mathcal S_1>9.6765387.
\]
Its cutoff allowance is $62500\tau D$.
The directed prefix comparison costs zero and has multiplicity one.
The true one-prefix conversion error is zero.
The alternative $-f$ repair costs less than $10^{-390}D$.
\end{lemma}
\begin{proof}
For S1, \eqref{p2:eq:direct-source-comparison} gives
$3S_1=3S(\mathcal A_1^{\rm pl},\mathcal P_N,N^\alpha)+3\Delta_1^{\rm src}$.
For S2 we use Section~\ref{sec:s2-1}:
$S_{2,\rm raw}=S(\mathcal A_{2,1},\mathcal P_N,N^\beta)$.
Both parents have prefix $1$, $v<(1-\epsilon_A)N$, $v\nmid N$.
The corrections are low-prime mass minus exceptional $v\mid N$ mass.
Prefix-square restoration is empty, and its nonnegative allowance $O_u(N^{-u/(1+\tau)}L_N^3)$ is retained.
These unsifted bounds are included once.
Descendants use only $A_D(v)$, $b\mid N-v$ and cofactor sieving.
Buchstab's identity holds for every sequence, without squarefree or coprime
restrictions \cite{Buchstab1937}.
Three iterations give (2.19).
The coefficients are $(3,-3,3,-3)$ and $(1,-1,1,-1)$.
In subsequent sieve calculations we write $S=S_{\rm pl}$.
We apply \cite[Theorem~4]{ref5} in the order lower, upper, lower, upper.
For S1-0, $\ell-2u=94849/160000>0$.
For S1-1 and S1-2, respectively, $s>228.19$ and $s>184.67$.
Equal-label pairs retain $[N^u,N^\alpha)$ or $[N^u,N^\beta)$.
The full summand, including $L_c$, depends on unordered products and labels.
Thus \eqref{p2:eq:direct-c1} gives $M_c=2$ there, and $1$ otherwise.
Its diagonal is assigned to prefix squares.

We fix $E_3$ before the structured sieve.
Since $p_3\ge E_3$,
\[
 -3S(p_1p_2p_3,p_3)\ge-3S(p_1p_2p_3,E_3).
 \tag{2.20}\label{p2:eq:2.20}
\]
For S1-3, $\ell-t_3\ge2743851/11490000>0$; the upper sieve is legal.
Fixing $E_3$ makes the complete tail domain permutation-invariant.
The factorials are $1,2,6$.
Formula \eqref{p2:eq:direct-remainder} has scalar $-3\nu C_{\rm slot}/M_c$.
Using actual $p_3$ before this comparison would lose invariance.
All four terms have fixed modulus-side labels.
Lemma~\ref{lem:s2-preparation} supplies their boxes.
The equal-prefix terms have disjoint prefix diagonals.
Theorem~\ref{thm:repeated-primes} bounds their union;
Theorem~\ref{thm:sieve-mass} bounds the AA37 main error.

For the S1-3 comparison put
\[
\begin{aligned}
\Theta_3^{\rm saved}(a,b,c)=\max\{&v(a)-u,v(b)-u,v(a+b)-u,v(a+b+c)-u,\\
 &\psi(v(a+c)-u,a+2b+c),\psi(v(b+c)-u,2a+b+c)\}.
\end{aligned}
\]
\hypertarget{p2:def:six-branch-saved}{}
Here both $w$ branches are deleted.
Every retained $v$ is reduced by $u$, including the guarded first inputs.
Let $\Theta_3^{\rm saved}[\tau\text{-bin}(r)]$ be its infimum
on the full ordered comparison box $[r,\rho r]$.
This box differs from the inherited source cell. Put
\[
\begin{gathered}
 G_\ell(w)=F(\ell/w)/w,\quad
 \ell_p(r)=\Theta_3^{\rm saved}(r)-r_1-r_2-r_3,\\
 \Phi(r)=F\!\left(\frac{\Theta_3^{\rm saved}[\tau\text{-bin}(r)]-\sum r_i}
 {\rho r_3}\right)/r_3.
\end{gathered}
\]
On the full parent, $u\le a,b,c\le\rho\alpha$ and $2\rho\alpha<25/128$.
Thus $v(a+c),v(b+c)\ge1/2$.
Both thresholds are $\le4\rho\alpha<1/2-u$.
The original and saved guards are active throughout.
Adding $u$ restores six corresponding $H$ branches,
each with slopes of absolute value $\le3/2$.
Their maximum is $3/2$-Lipschitz in coordinate $\ell^1$ distance.
For $x\in[r,\rho r]$, $\sum_i|x_i-r_i|\le\tau\sum_i r_i\le3\alpha\tau$.
Therefore
\[
 \Theta_3^{\rm saved}+u\le H,\qquad
 \Theta_c\ge\Theta_3^{\rm saved}(r)+u-(9/2)\alpha\tau
              >\Theta_3^{\rm saved}(r)+\eta_L.
\]
Branchwise, $(v-u)(\sum_{i\in A}r_i)-\sum_i r_i$ is coordinatewise nonincreasing.
The inherited lower labels give the same residual comparison $\ell_p(r)\le\ell_c$.
The comparison-bin infimum is at most its point value.
Decreasing $F$ gives $F(\ell_p(r)/(\rho r_3))/r_3\le\Phi(r)$.
For $H_F(s)=sF(s)$, the delay equations give
$H_F'=0$ on $[1,3]$ and $H_F'=f(s-1)\in[0,1]$ above $3$.
Hence
\[
 G_{\ell_c}(t_3)\le G_{\ell_c}(r_3)+\tau/r_3
 \le(1+\tau)G_{\ell_c}(r_3)\le(1+\tau)\Phi(r).
 \tag{2.21}\label{p2:eq:2.21}
\]
For S1-3, $t_i<\alpha$.
On the full parent, $a+c\le2\rho\alpha<t_{\rm dd}<25/128$,
so $v(a+c)=(1+a+c)/2$.
Also $a+2b+c\le4\rho\alpha<1/2$, so its $\psi$ branch is active.
Since $a\ge t_1$ and $c\ge t_3$,
\[
 \Theta_c\ge\frac12+\frac{t_1+t_3}{2},\qquad
 \ell_c\ge\frac12-t_2-\frac{t_1+t_3}{2}-\eta_L
             \ge\frac12-2\alpha-\eta_L.
\]
Since $\eta_L<u$ and $r_3\le\rho\alpha$,
differentiation is legal:
$\ell_c/r_3>(1/2-u-2\alpha)/(\rho\alpha)=186101/(50000\rho)>1$.
The comparison-box guards also remain active.
Its numerator loses at most $(9/2)\alpha\tau$, and the argument stays $>3.7$.
Thus both monotonicities apply, using the actual bin infimum.
Integration over the same literal clipped partition gives
\[
 E_{\rm cut}\le\tfrac34(2/e^\gamma)\tau I^*<1821.2\tau\,D
 \le62500\tau\,D.
 \tag{2.22}\label{p2:eq:2.22}
\]
Here Appendix~A.2 proves $I^*\le I_0+B_J<2161.074150710003$,
uniformly for $0\le\tau\le10^{-90}$.
Reproduction requires the interval-arithmetic source, dependencies and directed output;
see Appendix~A's data-availability statement.
The allowance $62500\tau D$ includes (2.22),
and Section~6.3.5(c) includes the prefix comparison once.

The one-prefix profile is $-F$, already used by the comparison functional.
For the alternative $-f$ we write $-F=-f-(F-f)$.
The delay equations give $s(F-f)=2e^\gamma\varrho(s-1)$,
where $\varrho$ is Dickman's function.
Here $s\ge s_*=64441/309$, $v_*=s_*-1$.
We iterate $\varrho(v)\le\varrho(v-1)/v$ exactly $207$ times;
the factor is $4/u=2000$.
Since $\log(\alpha/u)\le4$, $\log(\beta/u)\le5$,
\[
 E_{f\to F}\le
 \frac{2000(\tfrac34\cdot4+\tfrac14\cdot5)}
 {s_*\prod_{j=0}^{206}(v_*-j)}<10^{-390.77}<10^{-390}.
 \tag{2.23}\label{p2:eq:2.23}
\]
The S2 part belongs to its atomic endpoint.
Section~6 subtracts the residual $10^{-390}D$ once.
These are the true functions $f,F$.
Equations (2.19)--(2.23) prove the source/profile comparison.
Section~6.3.5(b) uses the single joint S1 bound.

The S2 last child retains the invariant condition (3.3), $h\le K$.
This includes PF and all remaining selected blocks.
Here $E_3$ is fixed before C1, and the scalar is $-\nu C_{\rm slot}/M_c$.
The actual-$p_3$ comparison is kept separate.
The strict complement uses actual $p_3$ and BV.
A PF-only coefficient would omit selected source blocks.
Since $1-2\eta_L-249/250=19/5000$, $K<1-2\eta_L$.
Lemma~\ref{thm:packing-input} uses the literal whole-$K$ infimum.
Section~\ref{sec:s2} checks the local source hypotheses and Class I.
The $X^{1-c_\sigma}$ term of Theorem~\ref{thm:distribution-input} is kept.
The strict BV complement has no prepared Class-II family.
Sections~\ref{sec:s2-3}--\ref{sec:s2-7} prove $d_{T8}$.
They retain whole-$K$ norms, spectral parameters, BV coefficients and the eight-branch functional.
Retain $275009680\tau D$, one structured-sieve error share and tail allowance $D/20000$.
Its true one-prefix profile is $-F$.
Section~6.3.5(b) subtracts these errors once.
\end{proof}


\subsection{Signed fixed-cutoff profiles}
\label{p2:sec:2.3}\label{p2:sec:2.4}\label{p2:sec:2.5}\label{p2:sec:2.6}\label{p2:sec:2.7}\label{p2:sec:2.8}\label{p2:sec:2.9}
This subsection treats the fixed-cutoff rows.
\begin{lemma}[Fixed-cutoff rows]
The six fixed-cutoff rows and S10-Ret have the following data.
In all tables $P=\sum_{i\le r}t_i$, $L_c=N^{\Theta_c-\eta_L-P}$.
The source is $A_D$, with shift $N$.
Write $\mathcal E$ for the common errors listed below.
Let $U_j,L_j$ be the selected raw-profile upper and lower endpoints
of Appendix~A.3, respectively.
Their directed $D$ contributions appear in the last table.
The relative AA37 parameter is $A_R=c/2$.
Each full remainder is \eqref{p2:eq:direct-remainder}, with scalar $c\nu C_{\rm slot}/M_c$.
\begin{center}
\small
\begin{tabularx}{\textwidth}{@{}lXX@{}}\toprule
Row & Exact exponent domain & Parent/anatomy/selection; depth\\\midrule
S3-A & $\alpha\le r_1<\mu$ & Inherited A/B at $t_{\rm dd}$; $1$\\
S4 & $\mu\le r_1<1/4$ & Inherited A/B at $t_{\rm dd}$; $1$\\
S5-low & $1/4\le r_1<t_*$ & B, clipped at $N^{t_*}$; $1$\\
S7 & $\mu\le r_1<\lambda$ & Inherited A/B at $t_{\rm dd}$; $1$\\
S8-A & $\alpha\le r_2<r_1<\beta$ & Small-coordinate anatomy; $2$\\
S9-A & $\beta\le r_1<\mu$, $\alpha\le r_2<\beta$ & Inherited A/B; $2$\\
S10-Ret & $\mu\le r_1<\lambda$,\newline
$\alpha\le r_2<\beta$,\newline $r_2<B_0(r_1)$
 & Whole parent cells, $t_2\le\operatorname{sel}(t_1)$; $2$\\\bottomrule
\end{tabularx}
\end{center}
\begin{center}
\small
\begin{tabular}{@{}lccll@{}}\toprule
Row & $c$; $D$ multiplier & Direction & Cutoff & Relative level\\\midrule
S3-A & $-2$; $-1/2$ & upper $F$ & $N^\alpha$ & $L_c$\\
S4 & $-1$; $-1/4$ & upper $F$ & $N^\alpha$ & $L_c$\\
S5-low & $-1$; $-1/4$ & upper $F$ & $N^\alpha$ & $L_c$\\
S7 & $-1$; $-1/4$ & upper $F$ & $N^\alpha$ & $L_c$\\
S8-A & $1$; $1/4$ & lower $f$ & $N^\alpha$ & $L_c$\\
S9-A & $1$; $1/4$ & lower $f$ & $N^\alpha$ & $L_c$\\
S10-Ret & $1$; $1/4$ & lower $f$ & $N^\alpha$ & $L_c$\\\bottomrule
\end{tabular}
\end{center}
\begin{table}[htbp]
\caption{Direct floors, cutoff gaps and prefix orbits.}\label{tab:direct-parameters}
\small
\begin{tabularx}{\textwidth}{@{}lXXc@{}}\toprule
Row & Lower floor for $\ell_c$ & Cutoff gap; ratio $s$ & $M_c$\\\midrule
S3-A & $64359/160000$ (A), $>0.402243$ (B)
 & $\ell-\alpha>0.315211$; $s>4.6217$ & $1$\\
S4 & $64359/160000$ (A), $\ge0.265524$ (B)
 & $\ell-\alpha>0.178492$; $s>3.0508$ & $1$\\
S5-low & $>0.255792$
 & $\ell-\alpha>0.168760$; $s>2.9390$ & $1$\\
S7 & $64359/160000$ (A), $>0.309274$ (B)
 & $\ell-\alpha>0.222242$; $s>3.5535$ & $1$\\
S8-A & $264897/1030000$
 & $\ell-2\alpha\ge98366653/1183470000$; $s>2.9550$ & $1,2$\\
S9-A & $11886931/49440000$ (A), $>0.240430$ (B)
 & $\ell-2\alpha>0.066366$; $s>2.7625$ & $1$\\
S10-Ret & $11886931/49440000$ (A), $2221831/11490000$ (B)
 & $\ell-2\alpha\ge221831/11490000$; $s>2.2218$ & $1$\\\bottomrule
\end{tabularx}
\end{table}
\begin{center}
\begin{tabularx}{\textwidth}{@{}lXX@{}}\toprule
Row & Saved profile; directed comparison cost & Selected output; estimates outside that endpoint\\\midrule
S3-A & $-\mathcal I_{3,A}/2$; zero & $-U_{3,A}/2$; $\mathcal E$\\
S4 & $-\mathcal I_4/4$; zero & $-U_4/4$; $\mathcal E$\\
S5-low & $-\mathcal I_{5,\rm low}/4$; zero & $-U_{5,\rm low}/4$; $\mathcal E$\\
S7 & $-\mathcal I_7/4$; zero & $-U_7/4$; $\mathcal E$\\
S8-A & $\mathcal I_{8,A}/4$; zero & $L_{8,A}/4$; $\mathcal E$, prefix and tail squares\\
S9-A & $\mathcal I_{9,A}/4$; zero & $L_{9,A}/4$; $\mathcal E$, tail squares only\\
S10-Ret & $\mathcal I_{10}/4$; $7.160\tau$
 & $L_{10}/4$; $\mathcal E$, tail squares only.
 Selected raw interval:
 $\left[\begin{gathered}0.858847455777940,\\0.867486588163666\end{gathered}\right]$\\\bottomrule
\end{tabularx}
\end{center}
S8-A and S9-A have zero cell-selection cost.
S10 has zero profile and S2-tail error.
\end{lemma}
\begin{proof}
The proof applies \cite[Theorem~4]{ref5} and \cite[Theorem~2.6, (2.30)--(2.33b)]{Z1}.
Table~\ref{tab:direct-parameters} gives $\ell_c\ge\ell_0$,
$z\le L_c$ for negative rows and $z^2\le L_c$ for positive rows.
Prefixes are $\ge z$ and tails are $<z$.
Every predicate is modulus-label dependent.
Lemma~\ref{thm:packing-input} and Lemma~\ref{lem:s2-preparation}
give the slots and local/complementary bounds.
Section~6 includes each of the following errors once:
\[
\begin{array}{c}
\text{source restoration, low cut, repeated primes, logarithm freezing, boundary,}\\
\text{Class I, Class II, roughness, structured sieve, prime mass}.
\end{array}
\]
These form $\mathcal E$; its squareful terms have the stated row restrictions.

For one uniform comparison put
\[
\begin{gathered}
 g_3(r)=v(r)-u,\quad
 g_8(r_1,r_2)=\max\{v(r_1),v(r_2),v(r_1+r_2)\},\quad g_9(r)=v(r_1),\\
 g(r)=57/64-u-(5/2)r,\qquad u=1/500.
\end{gathered}
\]
We call the absolute S3-A/S8-A/S9-A profiles Option-A.
Every missing completion contains the corresponding unthresholded branches.
Lipschitz control gives
\[
 \ell_c\ge g_3(r)-r+19/10000-3\tau,\qquad
 \ell_c-[g_8(r)-u-r_1-r_2]\ge19/10000-3\tau>0.
\]
The S4/S7 A/B levels give the same buffer.
For S5-low we use $\Theta_c\ge57/64-3\rho t_1/2$ in B.
On S9-A, $4\beta+3\alpha>25/32$ gives $g_8=g_9$.
Decreasing $F$ and increasing $f$ give all six signed lower comparisons.
Inherited labels are kept through clipping and before branch selection.
For S3, S4, S7 we split at $t_{\rm dd}$.
For S9 we split the inherited anatomy.
Retain both literal two-prefix infima.
S4's clipping at $\mu$ retains both A/B bounds.
Section~6.3 sums its two remainder children once.
S4 and S7 overlap in the initial signed inequality;
their HB partitions remain separate and disjoint for each term.

Put $\mathcal W=4(1-\varepsilon_A)U_C$.
In the displays below, $D_j$ is the signed source contribution divided by
$\mathcal W$, and $\mathcal E_j$ is that row's allocated share of the
common error $\mathcal E$. Section~6 charges these shares once.
The signed sources and their normalized bounds are
\begin{gather}
 \mathcal W D_{3,A}:=-2\sum_{\alpha\le r_1<\mu}S(p_1,N^\alpha).
 \tag{2.24}\label{p2:eq:2.24}\\
 D_{3,A}\ge-\frac12\mathcal I_{3,A}-\mathcal E_{3,A}-o_\sigma(1),\qquad
 \mathcal I_{3,A}=\frac{2}{e^\gamma\alpha}
 \int_\alpha^\mu F\!\left(\frac{v(r)-u-r}{\alpha}\right)\frac{dr}{r}.
 \tag{2.25}\label{p2:eq:2.25}
\end{gather}
\begin{gather}
 \mathcal W D_4:=-\sum_{\mu\le r_1<1/4}S(p_1,N^\alpha).
 \tag{2.26}\label{p2:eq:2.26}\\
 D_4\ge-\tfrac14\mathcal I_4-\mathcal E_4-o_\sigma(1),\qquad
 \mathcal I_4=\frac{2}{e^\gamma\alpha}
 \int_\mu^{1/4}F\!\left(\frac{57/64-u-(5/2)r}{\alpha}\right)\frac{dr}{r}.
 \tag{2.27}\label{p2:eq:2.27}
\end{gather}
For S5 set $t_*=7109/28000$, $\tau_S=57/224$.
Before the structured sieve or HB we split
\[
 -\sum_{1/4\le r_1<\tau_S}S(p_1,N^\alpha)
 =-\sum_{1/4\le r_1<t_*}S(p_1,N^\alpha)
  -\sum_{t_*\le r_1<\tau_S}S(p_1,N^\alpha).
 \tag{2.28}\label{p2:eq:2.28}
\]
The first term is S5-low and the second is S5-top.
The endpoint $t_*$ belongs to BV.
Clipping at $N^{t_*}$ retains parent labels.
S5-low uses B and
\begin{gather}
 \mathcal W D_{5,\rm low}:=-\sum_{1/4\le r_1<t_*}S(p_1,N^\alpha),\notag\\
 D_{5,\rm low}\ge-\tfrac14\mathcal I_{5,\rm low}
                 -\mathcal E_{5,\rm low}-o_\sigma(1),\notag\\
 \mathcal I_{5,\rm low}=\frac{2}{e^\gamma\alpha}
 \int_{1/4}^{t_*}F\!\left(\frac{57/64-u-(5/2)r}{\alpha}\right)\frac{dr}{r}.
 \tag{2.29}\label{p2:eq:2.29}
\end{gather}
For S7,
\begin{gather}
 \mathcal W D_7:=-\sum_{\mu\le r_1<\lambda}S(p_1,N^\alpha).
 \tag{2.30}\label{p2:eq:2.30}\\
 D_7\ge-\tfrac14\mathcal I_7-\mathcal E_7-o_\sigma(1),\qquad
 \mathcal I_7=\frac{2}{e^\gamma\alpha}\int_\mu^\lambda
 F\!\left(\frac{57/64-u-(5/2)r}{\alpha}\right)\frac{dr}{r}.
 \tag{2.31}\label{p2:eq:2.31}
\end{gather}
Its profile loss is zero.
Appendix~A.3 proves its outward upper endpoint,
which Section~6.4 uses at the fixed positive parameter.
Equation (6.5b) uses the S3 comparison.

For the two-prefix rows,
\begin{gather}
 \mathcal W D_{8,A}:=+\sum_{\alpha\le r_2<r_1<\beta}S(p_1p_2,N^\alpha).
 \tag{2.32}\label{p2:eq:2.32}\\
 D_{8,A}\ge\tfrac14\mathcal I_{8,A}-\mathcal E_{8,A}-o_\sigma(1),\notag\\
 \begin{aligned}
 \mathcal I_{8,A}
   &=\frac{2}{e^\gamma\alpha}
     \iint_{\alpha\le r_2<r_1<\beta}
     f\!\left(\frac{g_8(r)-u-r_1-r_2}{\alpha}\right)\\
   &\qquad\qquad\times\frac{dr_1dr_2}{r_1r_2}.
 \end{aligned}
 \tag{2.33}\label{p2:eq:2.33}
\end{gather}
\begin{gather}
 \mathcal W D_9:=+\sum_{\beta\le r_1<\mu\atop\alpha\le r_2<\beta}
 S(p_1p_2,N^\alpha).
 \tag{2.34}\label{p2:eq:2.34}\\
 D_9\ge\tfrac14\mathcal I_{9,A}-\mathcal E_9-o_\sigma(1),\qquad
 \mathcal I_{9,A}=\frac{2}{e^\gamma\alpha}
 \int_\beta^\mu\int_\alpha^\beta
 f\!\left(\frac{v(r_1)-u-r_1-r_2}{\alpha}\right)
 \frac{dr_2dr_1}{r_1r_2}.
 \tag{2.35}\label{p2:eq:2.35}
\end{gather}
Only S8-A has equal-set prefix orbits.
Its equal-label sets are identically clipped to $[N^\alpha,N^\beta)$.
The full summand and level are symmetric before slotting.
Thus $M_c=2$ there and $1$ elsewhere.
Formula \eqref{p2:eq:direct-c1} removes strict triangular ordering.
Its prefix-square part is disjoint from the forbidden-tail union.
The range lies in the small-coordinate anatomy, and both squareful errors are included.
Its outward lower endpoint enters (6.4) once, with coefficient $+1/4$.
S9's literal prefix sets are disjoint, including $N^\beta$.
Hence $M_c=1$, the prefix diagonal is empty,
and both normalization differences vanish by \eqref{p2:eq:direct-remainder}.
Only its tail-squareful error occurs.
Its lower endpoint has coefficient $+1/4$ and multiplicity one;
see Sections~6.3.3 and~6.4.

S5-top and S6 use BV separately.
For S5-top the total level is $N^{1/2}L_N^{-B}$,
with relative level obtained by dividing by the actual prefix.
Here $\kappa=\kappa_{\rm BV}=0$, cutoff $N^\alpha$, profile $-F$.
Section~5 proves its full coefficient and error estimate.
Its prepared Class-II, packing and Class-I contributions are zero,
and its main term enters separately.
Section~6.3 includes S5-low errors once.

\end{proof}
\subsection*{The separate BV row S6}
This subsection bounds the separate BV row.
Its exact source is
\[
 \mathcal W D_6:=-\sum_{\tau_S\le r_1<1/3}S(p_1,N^\alpha),\qquad
 \tau_S=57/224,\quad c=-1,\quad A_R=-1/2.
 \tag{2.36}\label{p2:eq:2.36}
\]
Section~5.3 takes $\kappa_{\rm BV}=10^{-6}<61/766=1/6-\alpha$.
The margin $(1/6-10^{-6}-\alpha)/\alpha$ absorbs the logarithmic loss uniformly.
Put $Q_{\rm BV}=N^{1/2-\kappa_{\rm BV}}L_N^{-B}$,
$D_p=Q_{\rm BV}/p$, with cutoff $N^\alpha$.
Then
\[
 \frac{\log D_p}{\alpha\log N}
 =\frac{1/2-\kappa_{\rm BV}-r_1}{\alpha}
  -\frac{B\log\log N}{\alpha\log N}>1
\]
eventually, uniformly for $r_1<1/3$.
Thus the upper sieve $F$ is legal, and
\[
 \mathcal I_6(\kappa_{\rm BV})=
 \frac{2}{e^\gamma\alpha}\int_{\tau_S}^{1/3}
 F\!\left(\frac{1/2-\kappa_{\rm BV}-r}{\alpha}\right)\frac{dr}{r},
 \qquad D_6\ge-\mathcal I_6(\kappa_{\rm BV})/4-o(1).
 \tag{2.37}\label{p2:eq:2.37}
\]
The BV error includes comparison with the $1/2$-level profile.
Section~5 gives the full weighted BV remainder, prefix multiplicities and endpoint direction.
This row uses the classical proof, without structured slot normalization or C1--C5.
Its BV error is subtracted once, together with its main term.


\subsection{The retained curved S10 row}
\label{p2:sec:2.10}
This subsection bounds the retained curved row.
\begin{lemma}[Retained curved profile]
Put
\[
\begin{aligned}
 C_0&=57/64-u-(11/5)\alpha,\quad
 B_0(r)=\min\{\beta,\max\{\alpha,C_0-5r/2\}\},\\
 B_\tau(t)&=\min\{\beta,\max\{\alpha,C_0-(5/2+3\tau/2)t\}\},\\
 R(t)&=\min\{\rho t,\lambda\},\qquad
 \operatorname{sel}(t)=\min\{B_\tau(t),B_0(R(t))/\rho\}.
\end{aligned}
\]
The actual source is
\[
 \mathcal W D_{10}:=+\sum_{\mu\le r_1<\lambda\atop\alpha\le r_2<\beta,\ r_2<B_0(r_1)}
 S(p_1p_2,N^\alpha).
 \tag{2.38}\label{p2:eq:2.38}
\]
Ret consists of whole parent-label cells with $t_2\le\operatorname{sel}(t_1)$.
For the curved source-trimmed target
\[
 \mathcal I_{10}=\frac{2}{e^\gamma\alpha}
 \int_\mu^\lambda\int_\alpha^{B_0(r_1)}
 f\!\left(\frac{57/64-u-(5/2)r_1-r_2}{\alpha}\right)
 \frac{dr_2dr_1}{r_1r_2}.
 \tag{2.39}\label{p2:eq:2.39}
\]
we have
\[
 D_{10}\ge\mathcal I_{10}/4-7.160\tau-\mathcal E_{10}-o_\sigma(1).
 \tag{2.40}\label{p2:eq:2.40}
\]
The profile and S2-tail errors are zero.
\end{lemma}
\begin{proof}
On clipped support,
$r_2<\rho t_2\le B_0(R(t_1))\le B_0(r_1)$.
Thus $\mathrm{Src}=\mathrm{Ret}\sqcup\mathrm{Drop}$ at finite $N$.
The coefficient $+1$ permits discarding Drop before estimating Ret.
The literal prefixes $[N^\mu,N^\lambda)$, $[N^\alpha,N^\beta)$ are disjoint.
Hence $M_c=1$ and the prefix diagonal is empty.
Selection is constant on whole labels.
For anatomy A, $t_2\le\beta$ gives residual $13475869/473388000>0$.
Table~\ref{tab:direct-parameters} verifies the lower cutoff, including the crossing band.
The tested source stays $A_D$, with shift $N$.
On Ret,
\[
 \alpha(s_{\rm live}-s_{\rm saved})\ge
 (u-\eta_L)-\tau((5/2)\lambda+\beta)>0.
\]
Increasing $f$ gives the favorable direction.
For every discarded actual point,
\[
 t_2>\operatorname{sel}(t_1),\quad r_2\ge t_2,\quad
 r_2<B_0(r_1)\le B_0(t_1).
\]
The larger of the two selection deficits bounds the strip.
The projection $P$ onto $[\alpha,\beta]$ is $1$-Lipschitz, so $B_0$ is $5/2$-Lipschitz.
The two deficits satisfy
\[
\begin{aligned}
 B_0(t)-B_\tau(t)&\le(3/2)t\tau,\\
 B_0(t)-B_0(R(t))/\rho
 &\le(5/2)(R(t)-t)+\beta(1-1/\rho)\\
 &\le(5\lambda/2+\beta)\tau.
\end{aligned}
\]
The projection estimate covers the $\beta$ cap, affine middle and $\alpha$ floor.
The two displayed inequalities cover both branches of the minimum.
Also $R(t)-t\le\tau\lambda$, including the clipped case $R(t)=\lambda$,
and $1-1/(1+\tau)\le\tau$.
Therefore
\[
 \operatorname{width}\le(5\lambda/2+\beta)\tau=\frac{36737}{49440}\tau.
\]
This includes endpoint cells: $t_1\le r_1<\lambda$, $t_2\le r_2$.
The inherited labels at $\mu,\lambda$ are retained.
On the whole canonical source, $11/5\le s<11519839/3200000<4$.
Here $f(s)=2e^\gamma\log(s-1)/s$ increases on $[2,4]$.
The $D$-density is at most $\log(3)/(4\alpha r_1r_2)$.
Hence
\[
 \Delta D\le\frac{\log3}{4\alpha^2}(5\lambda/2+\beta)
                    \frac{\lambda-\mu}{\mu}\tau<5.137\tau.
\]
For an exact rational check, use $\log3<10987/10000$.
\[
\sum_{k=0}^{8}\frac{(10987/10000)^k}{k!}>3.
\]
Substitution of $\alpha,\beta,\mu,\lambda$ gives the strict $5.137\tau$ bound.
The independent Appendix~A.3 bound uses $f\le1$ on this same strip:
\[
 \Delta D\le\frac{e^{-\gamma}}{2\alpha^2}
 \left(\frac52\lambda+\beta\right)
 \log\frac{\lambda}{\mu}\,\tau
 <4.82\tau<5.137\tau.
\]
These are two bounds for the same loss. The error account retains and subtracts
only the single conservative allowance $7.160\tau$; the allocated common
error $\mathcal E_{10}$ is charged once in Section~6.
Section~6.3 assembles the endpoint and common errors.
Lemma~\ref{lem:direct-preparation} supplies preparation, with
$K_D=K_*=6(C_b+1)+2=2012\le24C_b+26$.
The coefficient identities and masks apply with the stated signed row coefficient.
\end{proof}

\section{The S2 group: a joint three-piece estimate}
\label{sec:s2}

This section bounds the four S2 rows jointly from below. These rows are
$\{\text{S2-0},\text{S2-1},\text{S2-2},\text{S2-3(K/PF)}\}.$
Their signed lower functional is \texttt{S2-T8/4}.

\subsection{Source partition and four rows}
\label{sec:s2-1}\label{sec:s2-2}

\begin{lemma}
The plain S2 source satisfies (3.2), with the disjoint partition (3.4).
The selected third row has level and cutoff (3.5).
\end{lemma}
\begin{proof}
Fix
\[
 u=\frac1{500},\quad \beta=\frac{50}{309},\quad
 \epsilon_A=\eta_L=\frac1{10000},\quad
 0<\sigma\le10^{-10},\quad \tau=\sigma^9,\quad \rho=1+\tau,
\]
where $\sigma$ is sufficiently small for the AA37 constant and is fixed before $N$.
Then $N$ tends to infinity through even integers. Put
\[
 \mathcal A_D(v)=\mathbf1_{\{v\ {\rm prime}\}}
 \mathbf1_{\{N/\log^{30}N<v<(1-\epsilon_A)N\}},
 \qquad \mathcal A_0=\{N-v:\mathcal A_D(v)=1\}.
\]
Sections~3--5 use this endpoint, with the normalization
\[
 C(N)=\prod_{p>2}\left(1-\frac1{(p-1)^2}\right)
       \prod_{\substack{p\mid N\\p>2}}\frac{p-1}{p-2},\quad
 U_C=C(N)\frac{N}{\log^2N},\quad W_N=4(1-\epsilon_A)U_C.
 \tag{3.1}\label{p2:eq:s2-normalization}
\]
The first product is positive, and the second is at least one.
Thus $W_N\gg N/\log^2N$ uniformly in even $N$.

The unsifted AA37 input is $\mathcal A^{\rm pl}_{2,1}=\mathcal A_0$.
By (1.7),
$S_{2,\rm raw}=S(\mathcal A_0,\mathcal P_N,N^\beta)+\Delta^{\rm src}_2$,
where $\Delta^{\rm src}_2$ contains low-tested-prime and exceptional $v\mid N$ cells.
Source-square restoration is empty at $b=1$.
For a Buchstab prefix $b$ over $\mathcal P_N$, put
\[
 S_{\rm pl}(b,z)=S(\mathcal A^{\rm pl}_{2,b},\mathcal P_N,z),
 \qquad \mathcal A^{\rm pl}_{2,b}=\{N-v:\mathcal A_D(v)=1,\ b\mid N-v\},
\]
where $N-v$ is sieved with the empty cofactor allowed. Then
\[
\begin{split}
 T_0&=S_{\rm pl}(1,N^u),\\
 T_1&=\sum_{N^u\le p_1<N^\beta}S_{\rm pl}(p_1,N^u),\\
 T_2&=\sum_{N^u\le p_2<p_1<N^\beta}S_{\rm pl}(p_1p_2,N^u),\\
 T_3&=\sum_{N^u\le p_3<p_2<p_1<N^\beta}S_{\rm pl}(p_1p_2p_3,p_3),\\
 S_{\rm pl}(1,N^\beta)&=T_0-T_1+T_2-T_3.
\end{split}\tag{3.2}
\]
This follows by applying $S_{\rm pl}(b,z_2)=S_{\rm pl}(b,z_1)-\sum_{z_1\le p<z_2}S_{\rm pl}(bp,p)$ twice,
with repeated factors of $N-v$ still allowed.
The final cutoff is $p_3$.
Modulus-one coefficients give $(f,F,f,F)$, and the cutoff factors $(500,-500,500,-1)$
come from $1/u$ and $p_3$.

Let $r_i=\log_Np_i$. Put
$\Delta=\{u\le r_3<r_2<r_1<\beta\},\quad
 g(r)=r_1+2r_2+4r_3,\quad H=\frac{249}{250}.$
Set $PF=\Delta\cap\{g\le H\}$ and $R1=\Delta\cap\{g>H\}$.
The cells carry the inherited global half-open labels
$t_i\le r_i<\rho t_i,\qquad
 \mathcal P(t_i)=\{p\nmid2N:N^{\max(t_i,u)}\le p<N^{\min(\rho t_i,\beta)}\}.$
Clipping at $u$ or $\beta$ preserves each label.
A whole cell is selected when
\[
 h(t)=t_1+2t_2+(4+2\tau)t_3\le K,
 \qquad K=H+2\tau\beta.
 \tag{3.3}
\]
Let $\mathcal F$ be the selected union in $\Delta$.
Its strict complement $\mathcal B$ has $h>K$.
For $g\le H$, $h\le g+2\tau t_3\le K$,
and on $\mathcal F$, $g<\rho(t_1+2t_2+4t_3)\le\rho K$. Hence
\[
 PF\subseteq\mathcal F,\quad\mathcal B\subseteq R1,\quad
 \Delta=PF\ \dot\cup\ (\mathcal F\cap R1)\ \dot\cup\ \mathcal B,
 \quad \mathcal F\cap R1\subseteq\{H<g<\rho K\}.
 \tag{3.4}
\]
Equality belongs to PF or the selected label.
All three pieces retain the negative sign of $T_3$.
The coupled cut $g=H$ is compared with one whole-$K$ product-block estimate.

On a selected prefix block, put
\[
 L_b=N^{\Theta_b-\eta_L-\sum t_i},\qquad
 z=N^u\quad(r=0,1,2),\qquad z=E_3=N^{t_3}\quad(r=3),
 \tag{3.5}
\]
where $\Theta_b$ is the actual analytic cell infimum.
Since $E_3\le p_3$,
$S_{\rm pl}(b,p_3)\le S_{\rm pl}(b,E_3)$, the required upper bound for the negative row.
Its sieve argument exceeds $1.0117$.
The whole-$K$ cells satisfy the admission envelope $h\le1-2\eta_L$.
Here $K<1-2\eta_L,\qquad (1-2\eta_L)-H=19/5000.$
The boundary estimate retains the sharper condition (3.3).
\end{proof}

\subsection{Coefficients and prepared hypotheses}
\label{sec:s2-3}

\begin{lemma}
The four signed rows satisfy H1--H7 of Theorem~\ref{thm:distribution-input}.
Their coefficients are (3.6)--(3.8), with both masks (3.11).
\end{lemma}
\begin{proof}
Fix a complete symmetric prefix/tail block $k$ with its AA37 signs, and put
$C_k(q)=\sum_{d_{AA},(p_j)}\lambda_{d_{AA}}
                   \mathbf1_{\{q=d_{AA}\prod_jp_j\}}.$
The smooth support is
$d_{AA}<L_b^\sigma$, $d_{AA}\mid P_N(L_b^{\sigma^2})$.
It is squarefree, with primes avoiding $2N$.
The empty tail is included once.
Before unequal-coordinate dissection, complete equal-bin prefixes are divided by $M_k=\prod_G|G|!$.
Packing puts $d_{AA}$ wholly in slot 1, and the labelled convolution of slot $s$ is divided by $|J_s|!$. Thus
\[
 \frac{C_k}{M_k}=\frac{C_{\rm slot}}{M_k}
                         (\gamma_1*\gamma_2*\gamma_3),\quad
 C_{\rm slot}=\prod_{s=1}^3|J_s|!,\quad |\gamma_s|\le1.
 \tag{3.6}
\]
The prefix and slot factorials have distinct roles.
The signed remainder and its external scalar are
\[
 R_j=\sum_q(\gamma_1*\gamma_2*\gamma_3)(q)\Delta_j(q;N),
 \qquad c_j=w_kc_{{HB},{\rm box}}\frac{C_{\rm slot}}{M_k}.
 \tag{3.7}
\]
Prime realizations and smooth divisors stay inside the convolution.
Absolute values are taken after the signed $q$-sum.

The local density is $\omega_b(p^a)=p/(p-1)$ for every $a\ge1$, with higher-power constant 3.
Put $U_k=C_k/M_k$.
Let $\Xi_{\rm union}$ be the disjoint forbidden tail/prefix correction and
$\Xi_{\rm allowed}$ the allowed-repeat correction.
Let $W_k$ be the retained squarefree coefficient.
Then $O_k=U_k-\Xi_{\rm union}=W_k+\Xi_{\rm allowed}.$
With $m_A=\sum_v\mathcal A_D(v)$, define
\[
 \Delta_0(q)=\sum_{v\equiv N(q)}\mathcal A_D(v)-\frac{m_A}{\varphi(q)},
 \quad
 \Delta_P(q)=\sum_{v\equiv N(q)}\mathcal A_D(v)
             -\frac{\sum_{(v,q)=1}\mathcal A_D(v)}{\varphi(q)}.
\]
Substitution before HB gives
\[
 R_{AA}=S_{\rm slot}+NC_{sf}
 -\sum_{k,q}w_k\Xi_{\rm union}(q)\Delta_P(q)
 +\sum_{k,q}w_k\Xi_{\rm allowed}(q)(\Delta_0-\Delta_P)(q).
 \tag{3.8}
\]
Every source modulus is eventually below $N^{0.6}$.
Tested primes exceed $N/\log^{30}N$,
so $\Delta_0=\Delta_P$ on the original source.
Its squarefree NC and grouped allowed-repeat difference vanish before HB.
For $X_{Li}=\operatorname{Li}(Y_1)-\operatorname{Li}(Y_0)$, the term
$(m_A-X_{Li})/\varphi(q)$ is kept as the mass-conversion error.
The fixed logarithmic coefficient envelope and
$\sum_{q\le Q}\varphi(q)^{-1}\ll\log Q$ reduce it to PNT.
A sufficiently large fixed logarithmic saving gives $o(U_C)$.

Choose a tested-prime shell $[X,2X)$, with
$N/\log^{31}N\le X\le N$; on it HB has cap $2X^{1/3}$.
In each interior rectangle we choose the largest eligible elementary rough factor,
breaking ties by first index.
Its coefficient is rough M\"obius, rough unit, or normalized rough logarithm, of modulus at most one.
The complement has at most five factors and is bounded by $\tau_5(m)$.
With $\nu=N_{II}$, the full doubled support satisfies
\[
 X^{60\sigma}\le\nu\le X^{1/3+15\sigma},\quad M\nu\asymp X,
 \quad N\le X^{1+\sigma/20},\quad
 |mn-N|\ge\epsilon_A N/2
 \tag{3.9}
\]
This uses the original window $[N^{60\sigma},N^{1/3+14\sigma}]$ and
$\log_N(N/X)\le\sigma/1000$ eventually.
The shift stays $N$.

Lemma~\ref{lem:s2-preparation} supplies the equal-ratio mesh and integer endpoints.
For all four rows, including whole-$K$, it gives
$N_I\ge(2\nu)^{1-2^{-126}}$, $N_I\ge QX^{\sigma/3}$ and $M_I\ge X^{1/2}$.
It also gives H6, H7 and Class I.
The internal packing Lemma~\ref{thm:packing-input} applies at the literal $\Theta_b$,
and since $t_1<\beta<t_{\rm dd}$,
\eqref{p2:eq:row-level-floor} gives $\Theta_b\ge(1+t_1)/2$.
Put $T=t_1+2t_2\le3\beta$.
By (3.3), $t_3\le(K-T)/(4+2\tau)$, so
{\small
\[
\begin{aligned}
\ell_c=\Theta_b-\eta_L-t_1-t_2-t_3
&\ge\frac12-\eta_L-\frac T2-\frac{K-T}{4+2\tau}\\
&\ge\frac12-\eta_L-\frac{3\beta}{2}
              -\frac{K-3\beta}{4+2\tau}
>0.12953>\ell_0=\frac{53}{412}-\frac{\eta_L}{2}.
\end{aligned}
\]
}
Here $1/2-1/(4+2\tau)>0$ and $0<\tau=\sigma^9\le10^{-90}$.
Every selected S2-3 cell therefore has
$L_b^{\sigma^2}\ge N^{\ell_0\sigma^2}$.
Including slot-1 smooth inflation gives
\[
\begin{split}
 Q_1&<\nu X^{-\sigma},\\
 \nu^2Q_2Q_3^2&<X^{1-8\sigma},\\
 A_{\max}^{\theta}\nu Q_1Q_2^5Q_3^2&<X^{2-15\sigma},\\
 (\nu Q_3/Q_2)^{2\theta}\nu Q_1Q_2^5Q_3^2&<X^{2-15\sigma},
 \qquad \theta=7/32,\quad A_{\max}=X^{1+\sigma/20}.
\end{split}\tag{3.10}
\]
The ambient margins are
$(7271/256,52,95221/1280,22407/256)\sigma$.
After scale transfer, dyadic constants and endpoint localization, they exceed
$(28,51,74,87)\sigma$ eventually.
PF and the boundary region retain these margins and $\Theta_b$.

The dispersion argument is the one in
\hyperref[p2:dispersion-proof]{the distribution-input subsection}, with
\cite[Theorem~2.8, (3.17)--(3.20), (3.25)--(3.27), (3.57)--(3.60), (3.70)--(3.73)]{Z1}.
Its S2 parameters are as follows.
\begin{longtable}{@{}p{0.28\textwidth}p{0.67\textwidth}@{}}
\toprule
Step & Data \\
\midrule
First dispersion & $\sum_{m\in I_M}|\beta_m|^2\ll M_I\log^cX$.
For $\alpha^{(r)}=\alpha\mathbf1_{(n,r)=1}$, split at
$\|\alpha^{(r)}\|_2^2=\nu_I\log^{-A_3-1}X$. \\
SW precision & Below: trivial BDH.
Above: $\|\alpha^{(r)}\|_2(2\nu)^{1/2}\gg\nu\log^{-(334+A_3+1)/2}X$.
Request this extra saving before norm-normalized BFI. \\
Frequencies & $\eta_0=\sigma/1000$,
$H_2=QR^2S^2X^{\eta_0}/M_I$, $H_*=40H_2$;
$H_0<|h|\le\min(2H_0,H_*)$. \\
Mellin weight & In \eqref{p2:eq:dispersion-mellin}, set $a=N$.
Take $w=1$ on $[1/2,2]$, supported in $(1/4,4)$.
For $U(t)=m(t)$ and $L_{\rm osc}=(\log X)^{365}X^{\eta_0}$,
$U(t)\ll\min\{1,L_{\rm osc}^2/(1+|t|)^2\}$ and
$\int_{\mathbb R}U(t)\,dt\ll L_{\rm osc}\ll_\sigma X^{3\sigma/40}$.
Here $N/X\le(\log N)^{31}$. \\
Coloring & $d_{\rm disp}=1$, $\ell_0=53/412-1/20000$,
$u_0=N^{\ell_0\sigma^2}$, $\mathcal P_X=[u_0,X^3]\cap\mathbb P$.
Normalized cost: $2^{K_s+K_J}=O_\sigma(1)$. \\
Derivatives & $\epsilon_0=\sigma/2000>0$. \\
\bottomrule
\end{longtable}
Both norm cases give $\ll\nu_I\nu\log^{-A_3}X$.
The second dispersion retains
\[
 (s_1,Nc)=1,\qquad(s_1,q_{\rm disp}n_2)=1.
 \tag{3.11}
\]
Representations are normalized by $U(t)>0$, which is then restored.
At $U(t)=0$ the integrand vanishes.
Both nonmodular Mellin phases of \eqref{p2:eq:dispersion-mellin} are kept,
and coloring precedes aggregation.
Fiber multiplicity is at most $\tau(b)\tau_3(c)$. Hence
\[
 |A_\omega(t)|\ll_\sigma\frac{U(t)}{QR^2S^2},\qquad
 \sum_{b,c,f}|A^\chi_{b,c,f}(t)|^2
 \ll_\sigma U(t)^2X^{o(1)}\frac{\nu^2}{Q^2R^2S^3}.
 \tag{3.12}
\]
The spectral coefficient has the weighted norm
\[
 \sum_{n''\mid N^\infty}(n'')^\theta
 \|\beta^\flat(n'')\|_2^2
 \ll X^{o(1)}(H_2S)^2F,\qquad F\ll\nu/Q.
 \tag{3.13}
\]
For $y=f(h_1s'_1-h_2s_1)\ne0$, specialize to
\[
\begin{gathered}
 (r,s,d,c,n)=(z,p_0,b,c,p_0|y|),\quad
 a_T=\operatorname{sgn}(y)N,\quad(zb,p_0c)=1,\\
 (R_{sp},S_{sp},D_{sp},C_{sp},N_{sp})
 =(Z,\max(p_0,3)/2,B,C,8p_0Y),\\
 Z\in\{S^2,2S^2\},\quad B\ll\nu R,\quad
 C\ll\nu RS,\quad Y\ll H_2SF.
\end{gathered}
\]
with the least prime $p_0\nmid N$.
Then $p_0\ll\log N<z_0(X)$ and $p_0\nmid a_T$, and (3.13) gives the required weighted norm.
Here $p_0\mid c$ remains allowed.

For unit slot 2, its atom $s=1$ has $S=1/2$.
By (3.10), $H_2\ll X^{-9\sigma+\eta_0+o(1)}$,
so $H_*<1$ eventually and the frequency sum is empty.
Every other nonempty spectral scale is at least one.
For general $S$, $H_2/S\le X^{-9\sigma+\eta_0+o(1)}<1$,
hence $1\le H_*\le S$ eventually on nonempty frequency blocks.
Also $F\ge1/2$, as required at the weighted-norm endpoints.
Finally,
$A_{\max}C_{sp}D_{sp}N_{sp}R_{sp}S_{sp}
 \ll X^{\sigma/20+\eta_0+o(1)}\nu^4R^4S^6<X^{10.01}.$
By \cite[(3.89)]{Z1}, the positive-epsilon prefactor is eventually at most $X^{0.19019\sigma}$.
Equations (3.10)--(3.13) give the signed spectral bound.

Put $C_{\rm box}=C_b=334$ and $K_{\rm scale}=31$.
Lemma~\ref{lem:s2-preparation} gives the complete support
\[
 q<N^{e_D},\qquad e_D=153/256-\eta_L/2<0.6,
 \qquad q\le X^{0.6}\quad(N/\log^{31}N\le X\le N).
\]
This support includes repeated-prime moduli and every smooth divisor.
The full hypothesis check is as follows.
\begin{longtable}{@{}p{0.28\textwidth}p{0.67\textwidth}@{}}
\toprule
Condition & S2 data \\
\midrule
H1 & Coefficient (3.6), with primes avoiding $2N$; corrections exactly (3.8). \\
H2 & $d_{AA}$ lies in slot 1; slots 2 and 3 are $u_0$-rough.
Keep repeated primes and empty slots. \\
H3 and LEN & (3.9)--(3.13), the full support above and
Lemma~\ref{lem:s2-preparation}.
Majorant: $m_0=M(1+r_X^r)/2$, $1\le r\le5$,
$|m_0|\le2M\ll X/\nu$.
Ambient length: $N_I\ge(2\nu)^{1-2^{-126}}$. \\
H4 & $a=N$, $(q,N)=1$;
$0<N\le X\log^{31}N\le X^{1+\sigma/20}$. \\
H5 & $N^{60\sigma}\le\nu\le N^{1/3+14\sigma}$ and (3.9),
with Lemma~\ref{lem:s2-preparation}'s endpoints. \\
H6 & $E_{\log}\ll_\sigma N(\log N)^{18-C_b}$;
$E_{\rm bdry}\ll_\sigma N(\log N)^{16-C_b/2}$. \\
H7 & (3.16), using external coefficients only. \\
\bottomrule
\end{longtable}
Theorem~\ref{thm:distribution-input}, or
\cite[Theorem~3.2, (3.13)]{Z1}, therefore applies.
The complete-weight prime-source theorem has a different coefficient family.
\end{proof}

\subsection{Shiu's estimate and the three dispersion terms}
\label{sec:s2-4}

\begin{lemma}
\label{lem:external-classii-sum}
Let $\mathscr F$ be a prepared family of interior Class-II boxes satisfying
the local scale, roughness, support and coprimality hypotheses of
Theorem~\ref{thm:distribution-input} and (3.9)--(3.11), with uniform fixed
source constants. Assume the coefficient and scale bounds (3.16), with
$\mathcal J$ replaced by $\mathscr F$ and $K_0=8042$. Then (3.15) holds
with a fixed exponent $C_{\mathscr F}$ in place of $C_{S2}$, and the
normalized sum satisfies (3.17).
\end{lemma}
\begin{proof}
We first bound the deletion of $f_{\rm gcd}=(m,r_2)>z_0$.
The restoration formula in \cite[\S3.5, (3.45)]{Z1} retains $(m,r_2)=1$
by the exact M\"obius expansion.

Let $MN\asymp x$ and $L=\log x$, with local shift $a$.
Assume fixed $0<\delta<1$, $0<\varepsilon<1/2$, and
\[
 x^{2\delta}\le N\le x^{1/3+\delta/2},\quad
 |I_N|=N_I\asymp NL^{-C_b},\quad
 |\widetilde I_M|=2M_I\asymp ML^{-C_b},
\]
with $I_N\subset[N,2N]$ and $\widetilde I_M\subset[M/4,3M]$.
On the doubled box assume $|mn-a|\ge xL^{-C_b}\ge1$, and finally assume $0<|a|\le x^{1+\varepsilon}$ and $1\le f_{\rm gcd}\le Mx^{-\delta}$.
For fixed nonnegative $B_i$,
\[
 \sum_{n\in I_N}\tau(n)^{B_1}
 \sum_{\substack{m\in\widetilde I_M\\f_{\rm gcd}\mid m}}
 \tau(m)^{B_2}\tau(|mn-a|)^{B_3}
 \ll\frac{M_IN_I}{f_{\rm gcd}}(\tau(f_{\rm gcd})L)^{O_{B_i}(1)}.
 \tag{3.14}
\]
To prove this, write $m=f_{\rm gcd}m'$, so the interval length is $H'=2M_I/f_{\rm gcd}\gg x^{\delta/2}$.
For fixed $n$, put $g=(f_{\rm gcd}n,|a|)$,
$v=f_{\rm gcd}n/g$ and $a'=a/g$.
Then $(v,a')=1$, $|mn-a|=g|vm'-a'|$.
Put $\ell(m')=f_{\rm gcd}nm'-a$; consecutive values differ by
$f_{\rm gcd}n\le(Mx^{-\delta})(2N)\ll x^{1-\delta}$.
Opposite signs would imply
$|\ell(m')|+|\ell(m'+1)|=f_{\rm gcd}n\ge2xL^{-C_b}$, a contradiction, so the sign stays constant.
The values $|vm'-a'|$ lie in one reduced class modulo $v$.
Their interval has length $Y=vH'$ and endpoint $Z\ll x^{1+\varepsilon}$.
Choose fixed $0<b<\min(1/4,\delta/[8(1+\varepsilon)],\delta/8)$.
Since $v\ll x^{1+o(1)}$ and $H'\gg x^{\delta/2}$,
$Y\ge Z^b$ and $v\le Y^{1-b}$ eventually.
Shiu's reduced-progression bound \cite{Shiu1980} gives
$\sum_{m'}\tau(|vm'-a'|)^C
 \ll\frac{Y}{\varphi(v)}L^{2^C-1}\ll H'L^{O_C(1)}.$
For modulus one,
$\sum_{m'}\tau(f_{\rm gcd}m')^C\ll\tau(f_{\rm gcd})^CH'L^{O_C(1)}.$
We use Cauchy--Schwarz and $\tau(g)\le\tau(f_{\rm gcd})\tau(n)$, then the localized divisor moment.
This proves (3.14), including signed shifts.
The modulus condition uses $Y=vH'$ and $Y/v=H'$.

Take $x=X$, $N_{\rm local}=\nu$, $\delta=30\sigma$,
$\varepsilon=\sigma/20$ and $a=N_{\rm ambient}$.
By \cite[(3.6)]{Z1} and (3.9)--(3.10),
$f_{\rm gcd}=(m,r_2)\le r_2\ll M_I X^{-89\sigma}\ll M X^{-30\sigma},$
where $M_I\ll M(\log X)^{-334}$ and $\sigma>0$ is fixed.
Thus (3.14) applies.

In the following divisor sum, $r_2$ is the second two-rough-slot product
of \cite[(3.28)--(3.31)]{Z1}.
Its primes are at least $u_0=N^{\ell_0\sigma^2}$ and $r_2\le X^{0.6}$.
Counting repeats gives
$\Omega(r_2)\le\frac{0.6}{\ell_0\sigma^2},\qquad
 \tau(r_2)\le2^{\Omega(r_2)}\ll_\sigma 1.$
This includes $r_2=1$.
Every nontrivial $f_{\rm gcd}\mid r_2$ satisfies
$f_{\rm gcd}\ge u_0>z_0(X)$ eventually.
For fixed $B\ge0$,
$\sum_{\substack{f_{\rm gcd}\mid r_2\\f_{\rm gcd}>z_0(X)}}\frac{\tau(f_{\rm gcd})^B}{f_{\rm gcd}}
 \le\frac{\tau(r_2)^{B+1}}{u_0}\ll_{\sigma,B}\frac1{u_0}.$
We partition by the exact value $f_{\rm gcd}=(m,r_2)$.
Dropping the remaining gcd condition and summing (3.14) gives
\[
 \sum_{n\in I_\nu}\tau(n)^{B_1}
 \sum_{\substack{m\in\widetilde I_M\\(m,r_2)>1}}
 \tau(m)^{B_2}\tau(|mn-a|)^{B_3}
 \ll_\sigma\frac{M_I\nu_I}{u_0}(\log X)^C,\qquad |I_\nu|=\nu_I.
\]
The divisor weights fix $C$ independently of $r_2$.

The mixed term keeps $qr_1\mid mn_1-a$.
For fixed $(m,n_1)$, the number of admissible pairs $(q,r_1)$ is at most
$\tau(|mn_1-a|)^2$.
Also $\varphi(qr_2)^{-1}\ll(\log X)^C/(Qr_2)$,
$|\lambda_{r_i}|\ll_\sigma1$ and $\sum_{n_2}1\le\nu_I$.
Increasing $B_3$ by two and using $\sum_{r_2\asymp RS}r_2^{-1}\ll1$, we find the positive deletion error
$E_{\rm gcd}\ll_\sigma\frac{M_I\nu_I^2}{Q u_0}(\log X)^C.$
For pure density,
$\sum_{q\asymp Q}\varphi(q)^{-2}\sum_{r_1\asymp RS}\varphi(r_1)^{-1}
\ll(\log X)^C/Q$.
This gives the same bound.
Copy symmetry gives the first-copy deletion bound.
For two congruences, $qr_2\mid mn_2-a$ and $(qr_2,a)=1$
imply $(m,r_2)=1$, so its deleted contribution is empty.
These three bounds include unit rough slots, and both masks (3.11) remain in force.

The first Cauchy--Schwarz inequality is
$|S|^2\ll QM_I(\log X)^C S'$.
Its full divisor-deletion contribution is at most
\[
 (QM_I(\log X)^CE_{\rm gcd})^{1/2}
 \ll_\sigma M_I\nu_I(\log X)^C u_0^{-1/2}
 \ll_\sigma X^{1-\ell_0\sigma^2/2}(\log X)^C.
\]
Since $\ell_0\sigma^2/2=2c_\sigma>c_\sigma$, this term is absorbed by
the power-saving term established below in (3.15).
The exact mixed term retains $(m,r_2)=1$ via \cite[\S3.5, (3.45)]{Z1},
and the divisor sum above gives its full alternative bound.
Restoration and near-diagonal/BDH recombination follow
\cite[(3.34)--(3.47), (3.49)--(3.53)]{Z1}, with both masks and all three dispersion terms kept.

Theorem~\ref{thm:distribution-input} is applied at $A_{\rm work}=A'+4$.
Its local terms are $X_j\log^{-A_{\rm work}}X_j$ and
$X_j^{1-c_\sigma}\log^{C_{Z1}}X_j$, where
$c_\sigma=\frac14(\frac{53}{412}-\frac1{20000})\sigma^2>0$.
The latter includes GR restoration and first-CS deletion.
Here $\ell_0=53/412-1/20000$, $\ell_0\sigma^2/2=2c_\sigma$.
Three internal modulus-slot shells cost at most $\log^3X_j$, so the first term is at most
$X_j\log^{-A'-1}X_j\le X_j\log^{-A'}X_j$.
With shell and source-specific GR logarithms absorbed into $C_{S2}$,
\[
 |R_j|\ll_\sigma X_j\log^{-A'}X_j
       +X_j^{1-c_\sigma}\log^{C_{S2}}X_j.
 \tag{3.15}
\]

For the external family $\mathcal J$, Lemma~\ref{lem:s2-preparation} gives
\[
 \sum_{j\in\mathcal J}|c_j|\le C_\sigma\log^{K_0}N,
 \qquad N/\log^{31}N\le X_j\le N,
 \tag{3.16}
\]
Here $K_0=K_0^{S2}$ is fixed before $A'$.
Put $L=\log N$.
Lemma~\ref{lem:direct-preparation} gives
$K_*=6(C_b+1)+2=2012$.
We use the larger envelope $K_0=24C_b+26$:
$K_0^{S2}=24\cdot334+26=8042,\qquad
 A'=8047,\qquad A_{\rm work}=8051.$
PF and the selected boundary region form one whole-$K$ family.
Their prime realizations and smooth divisors remain internal.

For $\mathscr F$, we replace $\mathcal J$ and $C_{S2}$ in (3.15)--(3.16)
by $\mathscr F$ and any fixed $C_{\mathscr F}$.
With $K_0=8042$, $A'=8047$ and denominator (3.1), we get
\[
 \begin{aligned}
 E_{II}(\mathscr F)&:=\frac1{W_N}\sum_{j\in\mathscr F}|c_j||R_j|\\
 &\ll_\sigma L^{K_0+2-A'}\\
 &\qquad+N^{-c_\sigma}L^{K_0+C_{\mathscr F}+2+31c_\sigma}\\
 &\ll_\sigma L^{-3}+N^{-c_\sigma}L^{8044+C_{\mathscr F}+31c_\sigma}=o(1).
 \end{aligned}
 \tag{3.17}
\]
Here $8044=K_0+2$; the extra 2 comes from $W_N\gg N/L^2$.
By (3.1), $W_N\gg N/L^2$ and
$\tfrac12\log N\le\log X_j\le\log N$.
The bound $X_j\ge N/L^{31}$ gives $31c_\sigma$.
Also $K_0+2-A'=8042+2-8047=-3$,
and for fixed $B$, $e^{-c_\sigma t}t^B\to0$ as $t=\log N\to\infty$.
Here $\sigma>0$ and all source and box parameters are fixed before $N_0$.
For S2 we take $\mathscr F=\mathcal J$, $C_{\mathscr F}=C_{S2}$ and
$E_{II,S2}=E_{II}(\mathcal J)$.
\end{proof}

\subsection{The strict BV complement}
\label{sec:s2-5}

\begin{lemma}
The source on $\mathcal B$ has upper main comparator (3.19)
and remainder $o(U_C)$, at the actual cutoff $p_3$.
\end{lemma}
\begin{proof}
For $r\in\mathcal B$, put
$b=p_1p_2p_3,\quad Q_{BV}=N^{1/2-\kappa}\log^{-B}N,
 \quad\kappa=10^{-6},\quad D_b=Q_{BV}/b,\quad z_b=\sqrt{D_b}.$
On all of $R1$,
\[
 \log_ND_b\ge3/206-\kappa-o(1)>0,\qquad
 2r_3-\log_ND_b\ge63823/309000+\kappa+o(1)>0.
 \tag{3.18}
\]
Write $L_N=\log N$. Since $\mathcal B\subseteq R1$,
$g=r_1+2r_2+4r_3>H=249/250$. Ordering gives $g\le7r_1$, so
$r_1>H/7$ and
\[
 r_1+r_2+2r_3=\frac{g+r_1}{2}>\frac{4H}{7}.
\]
The definition of $D_b$ now gives
\[
\begin{aligned}
 r_3-\log_ND_b
 &=r_1+r_2+2r_3-\frac12+\kappa
       +B\frac{\log\log N}{\log N}\\
 &>\frac{4H}{7}-\frac12+\kappa
       +B\frac{\log\log N}{\log N}\\
 &=\frac{121}{1750}+\kappa
       +B\frac{\log\log N}{\log N}>0.
\end{aligned}
\]
Hence $1<D_b<p_3$ for all sufficiently large $N$, and
$z_b=\sqrt{D_b}<p_3$. Sifting is monotone in its cutoff, so
\[
 S_{\rm pl}(b,p_3)\le S_{\rm pl}(b,z_b).
\]
Apply the $s=2$ upper sieve to the right-hand side.
Its squarefree weight has $|\lambda^+_{b,d}|\le1$,
$d<D_b$ and $d\mid P_N(z_b)$. Thus every prime factor of $d$
is below $p_3$; indeed $d<D_b<p_3$.
For $k=bd$, the three largest distinct prime factors recover $b$ uniquely, also when $d=1$.
The actual aggregate weight
$a_k=\lambda^+_{b,d}\mathbf1_{\{b\text{ belongs to }\mathcal B\}}$
has $\|a\|_\infty\le1$, $\|a\|_2^2\le Q_{BV}$, and the prefix/divisor multiplicity is one.
Here $N\bmod k$ is reduced.
The two-endpoint remainder is at most twice the classical BV maximal sum.
It is $O_A(N\log^{-A}N)=o(U_C)$ for sufficiently large fixed $B$.
The positive upper comparator is enlarged from $\mathcal B$ to $R1$:
\[
 I_{BV}=\int_{R1}\frac{4\,dr_1dr_2dr_3}
 {r_1r_2r_3(1/2-\kappa-r_1-r_2-r_3)}.
 \tag{3.19}
\]
The factor 4 includes Mertens and $F(2)$.
Thus $I_{BV}$ enters directly, with its stated normalization.
The comparator overlaps $\mathcal F\cap R1$.
The source partition (3.4) stays disjoint, with its original boundary conventions.
\end{proof}

\subsection{The signed functional, tail and endpoint}
\label{sec:s2-6}\label{sec:s2-7}

\begin{lemma}
The selected endpoint is (3.22).
With its tail allowance included, the source satisfies (3.23).
\end{lemma}
\begin{proof}
Put $a=2/e^\gamma$. Define
\[
 v(t)=\begin{cases}(1+t)/2,&0\le t\le25/128,\\
 57/64-3t/2,&25/128<t\le1/3,\\0,&\text{otherwise},\end{cases}
 \quad\psi(x,y)=x\mathbf1_{\{x\ge y\}}.
\]
Let $\vartheta(r_1,r_2,r_3)$ be the maximum of
\[
\begin{gathered}
 v(r_1),\ v(r_2),\ v(r_1+r_2),\ v(r_1+r_2+r_3),\\
 \psi(\min(153/256-3r_3/8,1-2r_2),(1+r_1)/2),\\
 \psi(\min(153/256-3r_2/8,1-2r_3),(1+r_1)/2),\\
 \psi(v(r_1+r_3),r_1+2r_2+r_3),\quad
 \psi(v(r_2+r_3),2r_1+r_2+r_3).
\end{gathered}
\]
Put
\[
\begin{split}
 s_0&=19037/64,\quad
 a_1(r)=(\vartheta(r,u,u)-u-r)/u,\\
 a_2(r_1,r_2)&=(\vartheta(r_1,r_2,u)-u-r_1-r_2)/u,\\
 s_3(r)&=(\vartheta(r)-u-r_1-r_2-r_3)/r_3,\\
 P&=\frac au\left[f(s_0)-\int_u^\beta F(a_1(r))\frac{dr}{r}
 +\int_{u\le r_2<r_1<\beta}f(a_2(r_1,r_2))\frac{dr_1dr_2}{r_1r_2}\right],\\
 d\mu&=\frac{dr_1dr_2dr_3}{r_1r_2r_3^2},\quad
 V=\int_\Delta d\mu,\quad V_B=\int_{R1}d\mu,\\
 J_{PF}&=\int_{PF}[F(\max(s_3(r),1))-1]d\mu.
\end{split}
\]
The signed functional is
\[
 T8=P-aV-aJ_{PF}+(aV_B-I_{BV})
   =P-a\int_{PF}F(\max(s_3,1))d\mu-I_{BV}.
 \tag{3.20}
\]
It combines the prefix terms with the third-term comparison.
On $PF$, $\vartheta(r)\ge(1+r_1)/2$ and $g(r)\le H=1-2u$. Thus
$\vartheta(r)-u-r_1-r_2-2r_3\ge\frac{H-g(r)}2\ge0.$
Hence $s_3(r)\ge1$ and $F(\max(s_3,1))=F(s_3)$ on PF, with equality possible on its boundary.
The analytic whole-$K$ argument $1.0117$ uses level loss $\eta_L$.
At depth one and $u\le r<\beta$, the $w$ guard is active and its
branch gives $\vartheta(r,u,u)\ge19101/32000$ throughout the whole
parent cell. Thus the numerator margin is at least
$u-\eta_L=19/10000>0$.

The source/profile numerator margins are
$u-\eta_L-9\beta\tau>0$ on triples and
$u/2-\eta_L-5\beta\tau>0$ on second prefixes.
By continuity and the 2-Lipschitz S2 hull bound, and since $v(r_1)$ dominates at indicator switch-off,
$\Theta_b\ge\vartheta(r)-6\beta\tau$.
Inherited endpoints cost at most another $3\beta\tau$,
and the numerical and analytic losses are $u$ and $\eta_L$.
Missing-coordinate infima retain the floor $153/256-3u/8$
and the relevant $v(r_1+r_2)$ candidate.
Depth-two padding costs at most $u/2+5\beta\tau$.
Increasing $f$ and decreasing $F$ give the signs in (3.20).
PNT/Mertens, harmonic prime measures and equal-bin factorials give the ordered integrals.
Here $\tau$ is fixed throughout.

The boundary region has width
$\delta_K=\rho K-H=\tau(H+2\beta)+2\beta\tau^2.$
Its outer ordered area is at most $\beta^2/2$, and each fiber has width at most $\delta_K/4$.
Also $aF(1)=4$.
Including the dynamic cutoff,
\[
 E_{spill,D}\le\frac{(1+2\tau)\beta^2\delta_K}{8u^4}
 <275000000\tau,\qquad E_{dyn,D}\le2\tau V<9680\tau.
 \tag{3.21}
\]
Thus the cell-selection/dynamic allowance is $275009680\tau$,
while the boundary arithmetic discrepancy is estimated separately.

The directed 192-bit Arb inputs \cite{Johansson2017} are:
\begin{center}
\begin{tabular}{@{}ll@{}}
Prefix & $n=600$ \\
PF upper & $n=1536$ \\
Corner lower & $n=2800$
\end{tabular}
\end{center}
The true sieve functions $f,F$ are used beyond 12.
Put $E(s)=\max(F(s)-1,1-f(s))$.
The finite delay-integral enclosure gives $E(s)<6\cdot10^{-6}$ on $[11,12]$.
The classical linear-sieve asymptotic gives
$F(s)-1=O(e^{-s})$ and $1-f(s)=O(e^{-s})$.
Thus $e^sE(s)$ is bounded for all sufficiently large $s$.
Since $f,F$ are continuous, it is bounded on the remaining compact interval.
Hence $B=\sup_{s\ge11}e^sE(s)<\infty$ before the contraction below.
The same decay gives integrable tails and the vanishing limits used next.
Integrating $(s(F(s)-1))'=f(s-1)-1$ and
$(s(1-f(s)))'=1-F(s-1)$, with their vanishing limits at infinity,
gives the exact identities
\[
 s(F(s)-1)=\int_{s-1}^\infty(1-f(t))dt,\qquad
 s(1-f(s))=\int_{s-1}^\infty(F(t)-1)dt
\]
and hence, for $s\ge12$,
\[
 sE(s)\le\int_{s-1}^\infty E(t)dt\le Be^{1-s}.
\]
Separating $11\le s\le12$ from $s\ge12$ gives
$B\le\max(6\cdot10^{-6}e^{12},eB/12)$.
Since $e/12<1$, this implies
$B\le6\cdot10^{-6}e^{12}<0.976529<1$.
The preceding integral bound therefore gives
$E(s)\le e^{1-s}/s$ for $s\ge12$.

Fix $u\le r_2\le r_1\le\beta$ and write $R=r_1+r_2$, $x=r_3$,
$C=153/256$, $L=57/64$ and $A=(1+r_1)/2$.
The PF fiber has $u\le x<r_2$ and $r_1+2r_2+4x\le H$.
On every affine branch $\phi(x)$ in the hull, its argument is
$s_\phi(x)=c/x+k$, where
$s_\phi=(\phi(x)-u-R-x)/x$.
The following list covers every branch of $v(r_1)$, $v(r_2)$,
$v(R)$, $v(R+x)$, the two unguarded parts of each of the first two
$\psi$ terms, and both branches of each pair guard:
\begingroup\scriptsize
\setlength{\tabcolsep}{1.5pt}
\begin{longtable}{@{}p{0.22\textwidth}p{0.43\textwidth}p{0.26\textwidth}c@{}}
\toprule
Branch $\phi$ & $c$ & Lower bound for $c$ & $k$ \\
\midrule
$v(r_1)$ & $1/2-u-r_1/2-r_2$ & $1/2-u-3\beta/2$ & $-1$ \\
$v(r_2)$ & $1/2-u-r_1-r_2/2$ & $1/2-u-3\beta/2$ & $-1$ \\
$v(R)$, increasing & $1/2-u-R/2$ & $1/2-u-\beta$ & $-1$ \\
$v(R)$, decreasing & $L-u-5R/2$ & $L-u-5\beta$ & $-1$ \\
$v(R+x)$, increasing & $1/2-u-R/2$ & $1/2-u-\beta$ & $-1/2$ \\
$v(R+x)$, decreasing & $L-u-5R/2$ & $L-u-5\beta$ & $-5/2$ \\
$q_1=C-3x/8$ & $C-u-R$ & $C-u-2\beta$ & $-11/8$ \\
$q_1=1-2r_2$ & $1-u-r_1-3r_2$ & $1-u-4\beta$ & $-1$ \\
$q_2=C-3r_2/8$ & $C-u-r_1-11r_2/8$ & $C-u-19\beta/8$ & $-1$ \\
$q_2=1-2x$ & $1-u-R$ & $1-u-2\beta$ & $-3$ \\
$v(r_1+x)$, increasing & $1/2-u-r_1/2-r_2$ & $1/2-u-3\beta/2$ & $-1/2$ \\
$v(r_1+x)$, decreasing & $L-u-5r_1/2-r_2$ & $L-u-7\beta/2$ & $-5/2$ \\
$v(r_2+x)$, increasing & $1/2-u-r_1-r_2/2$ & $1/2-u-3\beta/2$ & $-1/2$ \\
$v(r_2+x)$, decreasing & $L-u-r_1-5r_2/2$ & $L-u-7\beta/2$ & $-5/2$ \\
\bottomrule
\end{longtable}
\endgroup
Here $q_1=\min(C-3x/8,1-2r_2)$ and
$q_2=\min(C-3r_2/8,1-2x)$.
Substitution of $u=1/500$ and $\beta=50/309$ shows that every lower
bound in the table is at least
$c_{\min}=L-u-5\beta=196681/2472000>0$.
The permanent baseline $A=v(r_1)$ is always present.
The guards $q_i\ge A$ switch off only, since each $q_i$ is continuous
and nonincreasing in $x$; at switch-off its value equals $A$.
For the pair guards put
\[
 D_1=v(r_1+x)-(r_1+2r_2+x),\qquad
 D_2=v(r_2+x)-(2r_1+r_2+x).
\]
Since $r_i+x\le2\beta<1/3$, these functions are continuous,
strictly decreasing, with slopes $-1/2$ and $-5/2$ on the two branches.
On the increasing branch $D_i\ge1/2-3\beta=3/206>0$,
so switch-off can occur only on the decreasing branch.
At a switch-off, direct substitution gives
\[
\begin{aligned}
 (r_1+2r_2+x)-A&=-r_1/2+6r_2/5-23/160
                  \le-1507/49440<0,\\
 (2r_1+r_2+x)-A&=7r_1/10-23/160
                  \le-1507/49440<0.
\end{aligned}
\]
Thus removing either guarded branch cannot raise the hull.
All affine crossings and the kink of $v$ are continuous.
The only other drop is $v(R+x)$ at $R+x=1/3$;
its left value is $25/64<A$, so it is inactive there.
Every active branch has $s_\phi'(x)=-c/x^2\le-c_{\min}/x^2$.
The full argument hull $s_3$ is therefore continuous and strictly decreasing on the fiber.
Its finitely many affine pieces have disjoint $s$-images.
Consequently, with no branch-count factor,
\[
\begin{aligned}
 \int_{\{x:s_3(x)>12\}}(F(s_3(x))-1)\frac{dx}{x^2}
 &\le\frac1{c_{\min}}\int_{12}^{\infty}E(s)\,ds\\
 &\le\frac{e^{-11}}{12c_{\min}}.
\end{aligned}
\]
Integrating over the full ordered $(r_1,r_2)$ triangle gives
\[
 \int_{u\le r_2<r_1\le\beta}\frac{dr_1dr_2}{r_1r_2}
 =\frac12\log^2(\beta/u).
\]
The true PF tail above 12 costs less than
\[
 \frac a4\frac{\log^2(\beta/u)}2\frac{e^{-11}}{12c_{min}}
 <4.739184885736468\cdot10^{-5}\quad\text{D units}.
\]
Here $a=2e^{-\gamma}>0$.
The lower endpoint uses $P_{lo}$, upper products $aV$ and
$a(J_{PF})_{hi}$, and the lower corner endpoint.
All prefix arguments exceed 128.
Their repair costs less than $1.154\cdot10^{-54}$ D.
Both repairs fit strictly inside $1/20000$.
Adding the true tail to the finite PF upper bound gives
\[
 \begin{aligned}
 d_{T8}&=\frac{P_{lo}-aV-a(J_{PF})_{hi}+(aV_B-I_{BV})_{lo}}4
                 -\frac1{20000}\\
 &>1.61958135698908352559133404332105729220076714081650.
 \end{aligned}
 \tag{3.22}
\]
This eight-branch bound is in D units.
It includes that allowance.

The source identities and preceding estimates yield
\[
 \frac{S_{2,\rm raw}}{W_N}
 \ge d_{T8}-275009680\tau-E_{AA,S2}-E_{II,S2}-e_{rest,S2}(N),
 \qquad e_{rest,S2}(N)\longrightarrow0.
 \tag{3.23}
\]
The four rows and three third-row pieces form \texttt{S2-T8/4}.
Equation (3.17) includes CRT, colors, zero mode, Poisson, Theorem 2.8
and all three GR/dispersion terms.
At $b=1$, the source-square restoration is empty.
The weighted low-tested-prime and exceptional-source bounds are
$O_u(\log^{-28}N)$ and $O_u(\log^3N/N)$.
Direct proper prime powers have raw bound $X^{1/2+o(1)}$, and their actual scalar sums are retained.
Theorem~\ref{thm:repeated-primes} bounds the forbidden square union by
$O_\sigma(N^{-\ell_0\sigma^2}\log^4N)$.
The NC and allowed-repeat terms vanish in the actual-mass convention of (3.8).
The Li convention retains the weighted source mass conversion.

For H6, the normalized LOG and boundary errors are
$O_\sigma(L^{20-C_b})$ and $O_\sigma(L^{18-C_b/2})$.
Class I is $O_\sigma(L^{-4})$, with $A_I=2018$, and at $C_b=334$ these enter $e_{rest,S2}$.
The strict BV source is $\mathcal B$, with main term in (3.20).
PNT, Mertens, endpoint and mass-conversion errors vanish.
The fixed $275009680\tau$ lies outside $d_{T8}$, while the allowance $1/20000$ lies inside it, in D units.
Finally, \texttt{E\_AA,S2} is the S2 share of the global AA37/G2 error.
Its fixed $O(\sigma)$ part stays separate from $e_{rest,S2}$.
Each contribution enters once.
The joint functional includes the S2 $f$-to-$F$ comparison and R1 corner.
The S1 $f$-to-$F$ and G1 PR-M terms remain separate.
This proves (3.23), which Section~6.3.5 uses in the full error sum.
\end{proof}

\section{The switched rows $S_{11},S_{12},S_{13}$}

This section proves a joint lower bound for the switched rows.

\subsection{The source and its finite signs}

This subsection describes the source and proves the finite-source sign lemma.
Fix the parameters
\[
 \alpha={100\over1149},\quad\beta={50\over309},\quad
 \lambda={93\over400},\quad\mu={25\over128},\quad u={1\over500},
 \qquad \varepsilon_A=\eta_L=10^{-4}.
 \tag{4.1}
\]
Let $0<\sigma\le10^{-10}$, and put
\[
 \tau_g=\tau=\sigma^9,\quad \rho=1+\tau,\quad L=\log N,\quad
 U_C={C(N)N\over L^2},\qquad \mathcal W=4(1-\varepsilon_A)U_C.
 \tag{4.2}
\]
The D-unit denominator is $\mathcal W$.
All sieve primes belong to $\mathcal P(N)=\{p>2:p\nmid N\}$.
Source primes dividing $N$ enter the final error table,
as do repeated source primes in squarefree identities.
Write
\[
 \mathcal A_{\varepsilon_A}
 =\{e=N-p:p\text{ prime},\ p<(1-\varepsilon_A)N\}.
 \tag{4.3}
\]
Let $t_i=\log_Np_i$.
The source forms, each with $\varepsilon_A N<e<N$, are as follows.
\begingroup
\setlength{\tabcolsep}{3pt}
\begin{longtable}{@{}>{\raggedright\arraybackslash}p{0.24\textwidth}>{\raggedright\arraybackslash}p{0.68\textwidth}@{}}
\toprule
Form & Exact representation and constraints \\
\midrule
$D_{11,3}$ & $e=p_1p_2p_3$, $N^\lambda\le p_1<N^{1/3}$, $p_1<p_2$, $p_1p_2^2<N$, $p_3\ge p_2$ \\
$D_{11,4}$ & $e=p_1p_2p_3p_4$, the same outer conditions, $p_2\le p_3\le p_4$ \\
$D_{12}$ & $e=p_1p_2p_3$, $N^\alpha\le p_1<N^{1/3}\le p_2$, $p_1p_2^2<N$, $p_3\ge p_2$ \\
$D_{13}$ & $e=p_1p_2p_3$, $N^\beta\le p_1<N^\lambda\le p_2$, $p_1p_2^2<N$, $p_1p_2p_3^2\ge N$ \\
\bottomrule
\end{longtable}
\endgroup
On $D_{13}$, the two endpoint orders give distinct representations.
The condition $p_3\ge p_2$ is omitted there.
Let $R_{11}=R_{11,3}+R_{11,4}$, $R_{12}$ and $R_{13}$ count representations, and put
\[
 A_G(e)=(2R_{11}(e)+R_{12}(e)+R_{13}(e))
                  \mathbf1_{\varepsilon_A N<e<N}.
 \tag{4.4}
\]
These forms exhaust the sources.
For $D_{11}$, $5\lambda>1$ allows at most two cofactor primes,
and $\varepsilon_A N>N^{2/3}$ excludes cofactor 1.
A composite $D_{12}$ cofactor gives $e\ge N^{\alpha+1}>N$.
Two $D_{13}$ cofactor factors give $e\ge N$, because each factor is at least $\sqrt{N/(p_1p_2)}$.
The linked source parameters are:
\begin{center}
\begin{tabular}{@{}llll@{}}
\toprule
Form & Weight & HB patterns/factors & Cofactor cutoff \\
\midrule
$D_{11,3}$ & $2$ & $27$ / $6\cdot3$ & $p_2$ \\
$D_{11,4}$ & $2$ & $81$ / $24$ & $p_2$ \\
$D_{12}$ & $1$ & $27$ / $6\cdot3$ & $p_2$ \\
$D_{13}$ & $1$ & $27$ / $6\cdot3$ & $\sqrt{N/(p_1p_2)}$ \\
\bottomrule
\end{tabular}
\end{center}
All four forms use rows G1-0 through G1-3.
Their faces are listed in (4.36) and the boundary-family table.
Their multiplicities satisfy (4.6) and the sharper peel bound below.
Every source coordinate is at least $N^\alpha$.
On $D_{11,3}$ and $D_{11,4}$, each remaining factor is at least
$p_2>p_1\ge N^\lambda$.
On $D_{12}$, $p_3\ge p_2\ge N^{1/3}$.
Only on $D_{13}$ do we use its square-root cutoff. There
$t_2<(1-t_1)/2$ and $t_1<\lambda<1/3$, so
\[
 t_3\ge {1-t_1-t_2\over2}>{1-t_1\over4}>
 {1\over6}>\alpha.
 \tag{4.5}
\]
For absolute estimates,
\[
 A_G(e)\le2(3!+4!)+3!+3!=72<100.
 \tag{4.6}
\]
For the peel, $A_G\le3$.
Ordering gives $R_{11},R_{12}\le1$.
If $R_{11}>0$, all factors are at least $N^\lambda$, so $R_{13}=0$.
Otherwise $D_{13}$ chooses one low-band and one high-band factor,
with at most $lh\le2$ choices, where $l+h=3$ and $l,h\ge1$.
Thus $2R_{11}+R_{12}+R_{13}\le3$, before any union.

Let $I_j(e)$ count the admissible divisor tuples,
retaining $P^-(e/d)\ge z$ and every original coprimality condition.
The identity is
\[
 W_{13}=3I_1+I_2-2I_3-I_4-I_5-I_6-I_7
                 +I_8+I_9+I_{10}-2I_{11}-I_{12}-I_{13}.
 \tag{4.7}
\]
With $\tau_S=57/224$, put
\[
 B(r)=
 \begin{cases}
 \beta,&\beta\le r<\mu,\\
 \min\{\beta,\max\{\alpha,57/64-u-(5/2)r-(11/5)\alpha\}\},
             &\mu\le r<\lambda.
 \end{cases}
 \tag{4.8}
\]
\begingroup
\setlength{\tabcolsep}{3pt}
\begin{longtable}{@{}>{\raggedright\arraybackslash}p{0.1\textwidth}>{\raggedright\arraybackslash}p{0.1\textwidth}>{\raggedright\arraybackslash}p{0.51\textwidth}>{\raggedright\arraybackslash}p{0.21\textwidth}@{}}
\toprule
Row & Weight & Divisor variables and range & Cofactor cutoff \\
\midrule
$S_1$ & $3$ & $d=1$ & $N^\alpha$ \\
$S_2$ & $1$ & $d=1$ & $N^\beta$ \\
$S_3$ & $-2$ & $d=q$, $\alpha\le\log_Nq<\mu$ & $N^\alpha$ \\
$S_4$ & $-1$ & $d=q$, $\mu\le\log_Nq<1/4$ & $N^\alpha$ \\
$S_5$ & $-1$ & $d=q$, $1/4\le\log_Nq<\tau_S$ & $N^\alpha$ \\
$S_6$ & $-1$ & $d=q$, $\tau_S\le\log_Nq<1/3$ & $N^\alpha$ \\
$S_7$ & $-1$ & $d=q$, $\mu\le\log_Nq<\lambda$ & $N^\alpha$ \\
$S_8$ & $1$ & $d=q_1q_2$, $\alpha\le\log_Nq_2<\log_Nq_1<\beta$ & $N^\alpha$ \\
$S_9$ & $1$ & $d=q_1q_2$, $\alpha\le\log_Nq_2<\beta\le\log_Nq_1<\mu$ & $N^\alpha$ \\
$S_{10}$ & $1$ & $d=q_1q_2$, $\mu\le r=\log_Nq_1<\lambda$, $\alpha\le\log_Nq_2<B(r)$ & $N^\alpha$ \\
$S_{11}$ & $-2$ & $p_1,p_2$ in the $D_{11}$ outer domain & $p_2$ \\
$S_{12}$ & $-1$ & $p_1,p_2$ in the $D_{12}$ outer domain & $p_2$ \\
$S_{13}$ & $-1$ & $p_1,p_2$ in the $D_{13}$ outer domain & $\sqrt{N/(p_1p_2)}$ \\
\bottomrule
\end{longtable}
\endgroup

\begin{lemma}[Finite-source signs]
Let $N$ be sufficiently large.
Suppose $e\in(\varepsilon_A N,N)$ is squarefree and coprime to $N$.
If $P^-(e)\ge N^\alpha$ and $\Omega(e)\in\{3,4\}$, then
\[
 W_{13}(e)\le0.
 \tag{4.9}
\]
\end{lemma}
\begin{proof}
The exact sum $\sum t_i$ is kept.
Let $k_i$ count primes in
$[\alpha,\beta)$, $[\beta,\lambda)$, $[\lambda,1/3)$ and $[1/3,1)$.
For a second-band exponent $v$, let $m_B(v)$ count first-band primes below $B(v)$.
The first ten rows give
\[
 B_0=3+\mathbf1_{k_1=0}-2k_1-2k_2-k_3+
                    {k_1\choose2}+\sum_{j=1}^{k_2}m_B(v_j),
 \qquad0\le m_B(v_j)\le k_1,
 \tag{4.10}
\]
and hence $W_{13}=B_0-2R_{11}-R_{12}-R_{13}$.
The cases $k_1=2$ and $k_1=1$ give
$B_0\le-k_3$ and $B_0\le1-k_2-k_3$, respectively.
The latter is positive only at $(1,0,0,2)$, of degree 3, where $R_{12}=1$.
Three fourth-band factors with a first-band factor exceed $N$.
For $k_1\ge3$, the lower slab leaves only $(3,0,0,1)$, where $B_0=0$.
It suffices that
\[
 \varepsilon_A N^{\,2/3-3\beta}>1:
 \tag{4.11}
\]
Three first-band factors and a fourth below $1/3$ have exponent
less than $3\beta+1/3<1$.
Four first-band factors have a smaller upper bound.

If $k_1=0,k_2\ge2$, then $B_0\le0$.
If $k_1=0,k_2=1$, then $B_0=2-k_3$.
At degree 3, take the second-band prime as $p_1$ and order the remaining primes as $p_2\le p_3$.
Then $p_1p_2^2\le e<N$, and
\[
 p_1p_2p_3^2=ep_3\ge\varepsilon_A N^{1+\lambda}\ge N
              \quad\text{if }\varepsilon_A N^\lambda\ge1.
 \tag{4.12}
\]
This is a $D_{13}$ representation.
When $k_3=0$, there is also one $D_{12}$ representation, and these cover $k_3=0,1,2$.
At degree 4, $k_3\ge2$, since $\beta+\lambda+2/3>1$.
Finally suppose $k_1=k_2=0$.
Degree 3 gives $k_3\ge1$ and one $D_{11}$ representation.
For $k_3=1$, there is also one $D_{12}$ representation, and they cancel $4-k_3$.
At degree 4, all primes lie in the third band, where $3\lambda+1/3>1$ and $B_0=0$.
\end{proof}

Conditions (4.11)--(4.12) and $\varepsilon_A N>N^{2/3}$ hold eventually,
and the finite radial range is $1+\log\varepsilon_A/L<\sum t_i<1$.
Section~5 gives (5.31)--(5.32) and positive excess at degrees 5 through 11.
Extra $D_{13}$ representations lower $W_{13}$.

Define the weighted sequences
\[
 \mathcal B_j=\{N-e:e\in\mathcal E_j\},\qquad
 \mathcal E_j=\{e:R_j(e)>0,\ \varepsilon_A N<e<N\}.
 \tag{4.13}
\]
Each occurrence has weight $R_j(e)$.
Then
\[
S_j\le\sum_eR_j(e)\mathbf1_{N-e\ {\rm prime}}\le S(\mathcal B_j;\mathcal P(N),N^{1/2})+O(N^{1/2+o(1)}).
\tag{4.14}
\]
Removing source conditions, including $(m,Nd)=1$, enlarges these positive counts.
Their negative coefficients preserve the required lower bound.
The sieve counts primes $N-e\ge N^{1/2}$ with weight $R_j(e)$,
and by (4.6) smaller primes and endpoints cost $O(N^{1/2+o(1)})$.
A source coordinate $q\mid N$ may leave a composite $N-e=qm$, and this excess is $E_{\rm src,exc}$.
Restricting to $p_i\nmid N$ gives equality up to the same error,
so (4.14) has a vanishing D-unit error.

The lower slab is exact:
\[
 p<(1-\varepsilon_A)N
       \ \Longleftrightarrow\ e=N-p>\varepsilon_A N.
 \tag{4.15}
\]
Its complement contributes zero.
The source factors are eventually odd, so $e\ne N$, and the shifted-prime source has $e\le N-3$.
For the fixed target, set
\[
\begin{gathered}
 A=3,\quad C_b=334,\quad N_{\rm cut}=N(1-L^{-333}),\\
 \Sigma_G=(\varepsilon_A N,N_{\rm cut}],\quad
 A_G^{\rm cut}=A_G\mathbf1_{e\le N_{\rm cut}}.
\end{gathered}
 \tag{4.16}
\]
The cut is applied before Buchstab and AA37.
Restoring it once costs
\[
 R0:G1\le100(NL^{-333}+1),\qquad
 E^D_{\rm cut}\le{100(NL^{-333}+1)\over\mathcal W}
                  =O(L^{-331})+O(L^2/N).
 \tag{4.17}
\]
Discrepancies, membership and boundary faces use (4.16).
\subsection{The four sieve rows and their coefficients}

This subsection sets up the four sieve rows and their signed remainder coefficients.
Put $\mathcal B_G^{\rm cut}=2\mathcal B_{11}^{\rm cut}
+\mathcal B_{12}^{\rm cut}+\mathcal B_{13}^{\rm cut}$, in which the source weights occur once.
Three Buchstab steps give
\[
\begin{aligned}
 S(\mathcal B_G^{\rm cut},N^\mu)
 ={}&S(\mathcal B_G^{\rm cut},N^u)
 -\sum_{N^u\le a_1<N^\mu}S((\mathcal B_G^{\rm cut})_{a_1},N^u)\\
 &+\sum_{N^u\le a_2<a_1<N^\mu}
                   S((\mathcal B_G^{\rm cut})_{a_1a_2},N^u)\\
 &-\sum_{N^u\le a_3<a_2<a_1<N^\mu}
                   S((\mathcal B_G^{\rm cut})_{a_1a_2a_3},a_3).
\end{aligned}
 \tag{4.18}
\]
Here $a_i\in\mathcal P(N)$, and the cutoffs exclude their endpoints.
Also $S(\mathcal B_G^{\rm cut},N^{1/2})\le S(\mathcal B_G^{\rm cut},N^\mu)$.
The cells are endpoint-clipped parent intervals $[N^{t_i},N^{\rho t_i})$ with inherited lower labels.
For the full anatomy infimum $\Theta_c$, put
\[
 \ell_c=\Theta_c-\eta_L-\sum_{i=1}^rt_i,\qquad
 D_c=N^{\ell_c},\qquad u_c=D_c^{\sigma^2}.
 \tag{4.19}
\]
AA37 is applied at $D_c$, with relative tail labels.
The smooth divisor satisfies $d<D_c^\sigma$, $d\mid P(u_c;N)$.
All source forms use the following four rows.
\begingroup
\setlength{\tabcolsep}{3pt}
\begin{longtable}{@{}llllp{0.40\textwidth}@{}}
\toprule
Row & Prefix & AA37 & $2A_R$ & Base cutoff and retained cells \\
\midrule
G1-0 & $1$ & upper $F$ & $-1$ & $N^u$; all \\
G1-1 & $a_1$ & lower $f$ & $+1$ & $N^u$; all \\
G1-2 & $a_1>a_2$ & upper $F$ & $-1$ & $N^u$; all \\
G1-3 & $a_1>a_2>a_3$ & lower $f$ & $+1$ & $E=N^{t_3}$ after PR; $\Theta_c-\eta_L-t_1-t_2-(3+2\tau)t_3\ge0$ \\
\bottomrule
\end{longtable}
\endgroup
The first three residuals are at least
\[
 {19101\over32000}-\eta_L,\quad
 {19101\over32000}-\eta_L-\rho\mu,\quad
 {19101\over32000}-\eta_L-2\rho\mu.
 \tag{4.20}
\]
For G1-0--G1-2, the smallest completion coordinate satisfies
$x_3\le u$, so the active $w$ branch is at least $19101/32000$.
If its guard closes, or the parent crosses $\mu$, the $v_1(x_1)$
branch gives the same lower bound. Hence
$\Theta_c\ge19101/32000$.
Dividing by $u$ gives values above 30.
For the fourth row, put
$A_{\rm floor}=53/256-\eta_L-(3/2)\mu\tau$.
The largest anatomy coordinate is $x_1\in[t_1,\rho t_1]$, and the full anatomy contains
$v_1(x)=\min\left\{{1+x\over2},{32-34x\over39},
                         {57\over64}-{3x\over2}\right\},$
also across $\mu$.
The middle branch exceeds the first for $x\le25/107$ and the third for $x\ge175/1568$.
Since $175/1568<25/107$, it is redundant.
Hence
$\Theta_c\ge\min\left\{{1+t_1\over2},
                    {57\over64}-{3\over2}\rho t_1\right\}.$
The increasing branch is evaluated on the left and the decreasing branch on the right.
For $t_2\le t_1\le\mu$, both residual branches decrease.
Thus
\[
 \begin{aligned}
 \Theta_c-t_1-t_2
 &\ge\min\left\{{1\over2}-{3\over2}\mu,
                {57\over64}-\left(1+{3\over2}\rho\right)\mu-\mu\right\}\\
 &=\min\left\{{53\over256},
              {53\over256}-{3\over2}\mu\tau\right\}
 ={53\over256}-{3\over2}\mu\tau.
 \end{aligned}
\]
Subtracting $\eta_L+t_3$ and using the retained-cell condition gives $\ell_c\ge A_{\rm floor}-t_3$ and $\ell_c\ge2t_3$, whence
$\ell_c\ge{2A_{\rm floor}\over3}>{1379\over10000}>\ell_0,
 \qquad \ell_0={53\over412}-{1\over20000}.$
The exact condition is $\ell_c\ge2\rho t_3$.
It gives the lower-sieve range on the whole cell, including equality.
The linked level and orbit data are:
\begin{center}
\begin{tabular}{@{}lllp{0.24\textwidth}@{}}
\toprule
Row & Lower floor for $\ell_c$ & Orbit $M_c$ & Ratio and outside error \\
\midrule
G1-0 & $19101/32000-\eta_L$ & $1$ & $s>30$; AA37 \\
G1-1 & $19101/32000-\eta_L-\rho\mu$ & $1$ & $s>30$; AA37 \\
G1-2 & $19101/32000-\eta_L-2\rho\mu$ & $1$ or $2$ & $s>30$; AA37 \\
G1-3 & $\max(A_{\rm floor}-t_3,2t_3)$ & $1$, $2$ or $6$ & $s\ge2$; AA37 and (4.57) \\
\bottomrule
\end{tabular}
\end{center}
All rows use all four source forms and level (4.19),
with exact signed remainder coefficient (4.26).
The common remaining errors are in the final table.


For the cut source, define
\[
\begin{aligned}
 C(q;N)&=\sum_eA_G^{\rm cut}(e)\mathbf1_{q\mid N-e},&
 X_G&=\sum_eA_G^{\rm cut}(e),\\
 Y_G(q)&=\sum_{(e,q)=1}A_G^{\rm cut}(e),&
 \Delta_0(q)&=C(q;N)-X_G/\varphi(q),\\
 \Delta_P(q)&=C(q;N)-Y_G(q)/\varphi(q).&&
\end{aligned}
 \tag{4.21}
\]
For allowed odd $p\nmid Nb$, take $\omega_b(p^k)=p/(p-1)$ for $k\ge1$,
so that $\omega_b(p^k)/p^k=1/\varphi(p^k)$.
Prefix primes lie above the cutoff and original tail primes lie below it.
Thus $r_{\rm std}((\mathcal B_G)_b,c')=\Delta_0(bc')$, including $c'=1$ and repeated tail primes.
Mertens and $\sum_p(p-1)^{-2}<\infty$ verify the density hypotheses uniformly.

Fix a cell and tail labels, and let $O(q),C_c(q)$ be the ordered masked and unrestricted labelled coefficients.
Put $M_c=\prod_h|G_h|!$.
With the actual masks $K_T,K_P$,
\[
 \Xi_U(q)={1\over M_c}\sum\lambda_d(1-K_TK_P)
                  \mathbf1_{q=d\prod p_j},\qquad
 O(q)={C_c(q)\over M_c}-\Xi_U(q).
 \tag{4.22}
\]
Equal-bin permutations preserve the clipped prime sets, products, labels and summands.
Unequal bins enforce order.
For G1-3, the $a_3$-cutoff difference is removed first, leaving base $E=N^{t_3}$.
The factor $1/M_c$ and the prefix-diagonal correction are kept.

Put $\Phi=M_c^{-1}\sum C_c\Delta_P$,
$NC_{\rm sf}=\sum_{q\ {\rm sf}}O(q)(\Delta_0-\Delta_P)(q)$,
and $\Xi_A(q)=O(q)\mathbf1_{q\ {\rm not\ sf}}$.
Adding and subtracting $\Delta_P$ gives
\[
\begin{gathered}
 R_{AA}=\Phi+NC_{\rm sf}-\sum_q\Xi_U(q)\Delta_P(q)
                   +\sum_q\Xi_A(q)(\Delta_0-\Delta_P)(q),\\
 \Delta_0-\Delta_P=-{X_G-Y_G(q)\over\varphi(q)}.
\end{gathered}
 \tag{4.23}
\]
Allowed repeats are kept with these masks.
Composite sources may retain the comparator difference.
Theorem~\ref{thm:repeated-primes} bounds both squareful terms.
The tested congruence is exactly
\[
 |(\mathcal B_G^{\rm cut})_q|
 =\sum_eA_G^{\rm cut}(e)\mathbf1_{q\mid N-e}
 =\sum_eA_G^{\rm cut}(e)\mathbf1_{e\equiv N\pmod q}.
 \tag{4.24}
\]
Every actual modulus satisfies $(q,2N)=1$ and
\[
 q<N^{e_G},\qquad e_G:=\frac{153}{256}(1+\tau+\sigma)<0.598.
\]
By (4.28), $X>cN$, where $c=\varepsilon_A2^{-24}>0$.
Eventually $N^{e_G}<(cN)^{3/5}<X^{3/5}$.
This verifies the full support condition of Theorem~\ref{thm:distribution-input};
see \cite[Theorem~3.2, (H3)]{Z1}.

Primes below $u_c$ determine $d$ uniquely, and rough positions are at least $u_c$.
Their multiplicities $e_p$ allow at most $m!/\prod_pe_p!$ assignments, and masks remove assignments.
Consequently
\[
 \ell_0={53\over412}-{1\over20000}>0,\qquad
 M_\sigma=\left\lceil{3\over5\ell_0\sigma^2}\right\rceil,\qquad
 |\gamma_q|\le M_\sigma! =: C_\sigma,
 \tag{4.25}
\]
which covers original, unrestricted and grouped coefficients.
Fixed row conventions are included in $C_\sigma$,
and $N_{\rm lab}(\sigma)$ enters aggregate bounds once.
For 1-bounded slot sequences we extract $C_{\rm slot}=\prod_s|J_s|!$.
The D-unit formulas are
\[
 {(2A_R)\nu_c\over M_c\mathcal W}\sum_qC_c(q)\Delta_B(q),
 \qquad
 {(2A_R)\nu_c C_{\rm slot}\over M_c\mathcal W}
                  \sum_q{C_c(q)\over C_{\rm slot}}\Delta_B(q).
 \tag{4.26}
\]
Here $\nu_c=+1$ for upper AA37 and $-1$ for lower AA37.
We form $C_c$ before this sign.
The slot factorial and source weights enter once, and $K_G$ is omitted from the remainder.

\subsection{Preparation of the product source}

This subsection splits the product source into boxes and checks the hypotheses of the distribution input.
Order-3 Heath--Brown \cite{ref6} is applied pointwise to every source coordinate.
Then
\[
 n(\mathbf v)=\prod_{i=1}^sv_i=e,\qquad s\le6k\le24.
 \tag{4.27}
\]
There are $27+81+27+27=162$ patterns per row, giving
648 row/form/pattern instances.\phantomsection\label{p2:g1-external-count}
All source indicators are kept.
Put $\delta=L^{-334}$.
We split each dyadic shell $[2^b,2^{b+1})$, $b\ge1$, into
$J=\lceil\log2/\log(1+\delta)\rceil$ intervals
$[2^b2^{j/J},2^b2^{(j+1)/J})$, $0\le j<J$.
Their common relative width is $\Delta=2^{1/J}-1$.
For $\delta\le1/300$, $\delta/3\le\Delta\le\delta$.
The point 1 is kept as a singleton.
A boundary box crosses a source, window, ordering, cap or cut face.
An interior box has constant source indicators.
Per variable there are $O(L)$ shells and $O(L^{334})$ intervals,
so there are $O(L^{335})$ choices per variable.
The 24-variable fine count is $O(L^{8040})$; two further shells give $K_0=8042$.
The three internal slot-shell sums are counted in $A_{\rm work}-A'=4$.

Choose a nonunit cell, and put
$T:=T_{\rm II}:=N_{II}:=V_{\rm sel},\qquad q:=1+\Delta,\qquad
 C_{\rm sel}=[T,qT),\qquad h_N:=\Delta T.$
Thus $\delta T/3\le h_N\le\delta T$.
The exact interval counts are
\[
 \begin{aligned}
 I_N&:=(T,qT]\subset(T,2T],\\
 H_N&:=\#(C_{\rm sel}\cap\mathbb Z)
      =\lceil qT\rceil-\lceil T\rceil,\\
 N_I&:=\#(I_N\cap\mathbb Z)=\lfloor qT\rfloor-\lfloor T\rfloor.
 \end{aligned}
\]
Both differ from $h_N$ by less than 1.
The ambient count is $N_I$, and roughness stays in its coefficient.
For an integer endpoint $w=2^{b+j/J}$, $w^J=2^{bJ+j}$ gives $J\mid j$, so only $2^b,2^{b+1}$ are integer endpoints.
The factor $\mathbf1_{P^-(n)\ge z_0(N)}$, with $z_0(N)>2$, kills these endpoints.
Hence
$a_{\rm sel}(n)\mathbf1_{C_{\rm sel}}(n)
 =a_{\rm sel}(n)\mathbf1_{I_N}(n)
 \qquad(n\in\mathbb Z).$

\begin{lemma}[Common product boxes]\label{lem:g1-f}
Let $B$ be a retained interior box.
For its lower endpoints $V_i$, put $X=\prod_iV_i$.
Then
\[
 n(\mathbf v)\in\Sigma_G\quad(\mathbf v\in B),\qquad
 \varepsilon_A2^{-24}N<X\le N.
 \tag{4.28}
\]
The largest eligible HB scale in
$[N^{60\sigma},N^{1/3+14\sigma}]$ is chosen, with the first index in a tie.
Write $N_{II}$ for this lower scale and $M$ for the complementary product.
Every box is Class II, and eventually
\[
 MN_{II}\asymp N,\quad
 X^{60\sigma}\le N_{II}\le X^{1/3+15\sigma},\quad
 N\le X^{1+\sigma/20}.
 \tag{4.29}
\]
On its doubled complementary majorant support,
\[
 N-mn\ge NL^{-334}\ge X(\log X)^{-334}>0.
 \tag{4.30}
\]
The thresholds (4.32)--(4.32a) below imply LEN on every box.
\end{lemma}
\begin{proof}
Buchstab, AA37, (4.23) and the G1-3 base $E$ preserve the measure.
Prime-power replacement and logarithm freezing preserve the indicators.
Equation (4.27) gives the support, and the bounds $V_i\le v_i<2V_i$ and $s\le24$ give the scale.

Choose a coordinate with exponent in $[\alpha,1/3]$.
One of its six or fewer factors is at least $N^{\alpha/6}$.
Its lower endpoint lies between half this size and $N^{1/3}$.
Since $\alpha/6>60\sigma$, an eligible choice exists, and the prescribed choice is constant on each box.
Fubini gives
\[
 \sum_{\mathbf v\in B}\prod_ia_i(v_i)H(\prod_iv_i)
       =\sum_{m,n}\beta(m)\alpha_{II}(n)H(mn).
 \tag{4.31}
\]
The selected coefficient is rough Möbius, rough 1, or normalized rough logarithm.
The complementary convolution has at most 23 factors and bound $\tau_{23}$.
Each logarithm stays on its original factor.
Endpoint subtraction is used in the interval and auxiliary-coprimality lemmas of Theorem~\ref{thm:distribution-input}.
Narrow intervals, the relative-norm alternative and both cross-copy masks are kept.
These inputs include more than scalar divisor moments.

Put $\omega=\log_N(N/X)$, and impose
\[
 L\ge\max\{2611/\sigma,300000\}
 \tag{4.32}
\]
\[
 \sigma L\ge2^{130}(\log L+1).
 \tag{4.32a}
\]
For fixed $\sigma>0$ the latter holds eventually, and uniformity as $\sigma\to0$ is unnecessary.
Since $T\ge N^{60\sigma}$,
$60\sigma L\ge60\cdot2^{130}(\log L+1)
              >334\log L+\log6,$
so $h_N\ge T/(3L^{334})>2$.
Thus $T/(6L^{334})\le N_I$ and $N_I\le2\delta T$.
Also
$T^{2^{-126}}\ge\exp(60\sigma L/2^{126})\ge12L^{334},$
since $60\sigma L/2^{126}\ge960(\log L+1)>334\log L+\log12$.
Consequently
$N_I\ge\frac{T}{6L^{334}}
       \ge2T^{1-2^{-126}}\ge(2T)^{1-2^{-126}}.$
With $Y=\lfloor T\rfloor$ and $H=N_I$, the integer intervals coincide.
Then $(2Y)^{1-2^{-126}}\le H\le Y$, using $T\ge2$ and $\delta\le1/300$.
This gives the short-interval input for $k=2^{126}$.
Threshold (4.32) gives
\[
 \omega<{\log(2^{24}/\varepsilon_A)\over L}\le{\sigma\over101},
 \quad {1/3+14\sigma\over1-\omega}\le1/3+15\sigma,\quad
 {2\omega\over1-\omega}<0.02\sigma.
 \tag{4.33}
\]
Also $(1-\omega)(1+\sigma/20)\ge1$, and these prove (4.29) and the packing scale conversion.
The frequency stays $a=N$.

Put $x=L^{-334}$; for $x\le1/300$, $(1+x)^{24}-1\le25x$.
Let $1\le r\le23$ count nonunit complementary cells, with endpoints $V_i,W_i=V_i(1+\Delta)$.
Put $M=\prod V_i$ and $M_+=\prod W_i$.
Unit factors stay singletons.
There is a nonunit complementary factor.
At least three source coordinates have product at least $N^\alpha$.
Only one HB factor was selected.
Thus $1+x/3\le M_+/M\le(1+x)^{23}\le1+25x$.
The ambient length is $M_I=M_+-M$, with $xM/3\le M_I\le25xM$.
The fixed majorant $\psi_I(m)=\phi_0((m-m_0)/M_I)$ in condition C3
of Theorem~\ref{thm:distribution-input} doubles this interval about its midpoint.
Its upper extension is $M_I/2\le12.5xM$.
Since $nM_+\le N_{\rm cut}$ on an interior box,
$mn\le N_{\rm cut}+12.5xN\le N_{\rm cut}+25xN$ on the doubled support.
Hence $N-mn\ge(L-25)NL^{-334}\ge NL^{-334}$ for $L\ge26$.
The function $t/(\log t)^{334}$ increases for $\log t>334$, and as $X\le N$ this proves (4.30).
Threshold (4.32) also gives $NL^{-334}>1$ and $\log X>334$.

Write $X=N^S$, $V_{\rm sel}=N^{e_{\rm sel}}$ and $M=X/V_{\rm sel}$.
The geometry gives
\[
 M_I\ge\tfrac13L^{-334}M,
 \qquad
 S/2-e_{\rm sel}\ge1/6-(15+1/202)\sigma>0
\]
where $S\ge1-\omega$ and $e_{\rm sel}\le1/3+14\sigma$.
Thus $M_I/X^{1/2}\ge\tfrac13L^{-334}N^{1/6-(15+1/202)\sigma}\to\infty$,
and in particular $M_I\ge X^{1/2}$ eventually.
Also $M_I/X^{1/2}\ge\tfrac13L^{-334}N^{0.15}>1$ eventually.
By (4.33),
$(1-\omega)^{-334}
 \le\exp(334\omega/(1-\omega))\le1+4\cdot10^{-10}.$
Since $\log X=(1-\omega)L$,
$L^{-334}\le(\log X)^{-334}\le(1+4\cdot10^{-10})L^{-334}$.
Therefore
\[
\begin{aligned}
 \tfrac17T(\log X)^{-334}&\le N_I\le2T(\log X)^{-334},\\
 \tfrac14M(\log X)^{-334}&\le M_I\le25M(\log X)^{-334}.
\end{aligned}
\]
As $M_+=Mq^r$ and $q^r\to1$ uniformly,
$m_0=M(1+q^r)/2\le2M=2X/T$ eventually.
On $[M,Mq^r]$, $\phi_0((m-m_0)/M_I)=1$, and both Cauchy--Schwarz copies remain in the doubled rectangle after coloring.
Equation (4.30) excludes exact resonance; see \cite[(3.10)]{Z1}.

By the internal Lemma~\ref{thm:packing-input},
\[
\begin{gathered}
 Q\le TN^{-(1+7271/256)\sigma}\le TX^{-(1+7271/256)\sigma},\\
 {7271\over256}-{1\over3}={21557\over768}=28+{53\over768}>28.
\end{gathered}
\]
Thus $QX^{\sigma/3}<TX^{-28\sigma}$.
Threshold (4.32a) gives $14\sigma L>334\log L+\log6$.
Since $\log X\ge L/2$,
$QX^{\sigma/3}<TX^{-28\sigma}
 \le T/(6L^{334})\le N_I.$
The same threshold gives $334\log L+\log3<\sigma L/100$.
Hence
$\log_NM_I\ge{2\over3}
 -\left(14+{1\over101}+{1\over100}\right)\sigma
 >{2\over3}-14.02\sigma.$
Packing also gives $\log_N(QRS)\le153/256+\sigma$.
Since $104.02\sigma<53/768=2/3-153/256$ and $X^{-89\sigma}\ge N^{-89\sigma}$,
$QRS<M_IX^{-89\sigma}$.
Finally, eventual monotonicity gives $\zeta=z_0(N)\ge z_0(X)$.
\end{proof}

Put $L_N=\log N$, $z=z_0(N)$.
The replacement $\mathbf1_{\mathbb P}(x)\mapsto\Lambda_z(x)/\log x$ costs (4.58)--(4.61).
Order-three rough Heath--Brown \cite{ref6} is used on each coordinate shell,
with the inverse Dirichlet series truncated at its upper endpoint.
The pointwise factors are 1 or $z$-rough.
The product inequalities are retained; they are constant on interiors.
Boundary pieces are estimated below.

For a form with $k\in\{3,4\}$ coordinates, let $H_B(\mathbf v)$
denote the HB numerator, including its source-form weight but excluding
the $k$ inverse coordinate logarithms. Then
$|H_B(\mathbf v)|\le2\cdot3^kL_N^k\le162L_N^4$.
Each of these $k$ HB logarithms is normalized by $L_N$, and the frozen coordinate logarithm satisfies $\lambda_{B,j}\ge\alpha L_N-1$.
For $L_N\ge2/\alpha$, the resulting exterior coefficient is at most
$2\cdot3^k(2/\alpha)^k\le162(2/\alpha)^4$ uniformly.
For a selected logarithm,
$\frac{\log n}{L_N}=\frac{\log X}{L_N}\frac{\log n}{\log X},
 \qquad \frac{\log X}{L_N}\le1.$
The selected coefficient is 1-bounded and the complement has bound $\tau_{23}$.
With (4.66) and fixed source-label/slot constants, this gives H7, $K_0=8042$.

The internal Lemma~\ref{thm:packing-input} is applied at $\Theta_c$.
The $d$ of AA37 is placed in slot 1, also when $J_1=\varnothing$.
Empty prime-label slots 2 and 3 carry unit atoms, and those two slots are rough.
The four surplus margins are
\[
 \left({7271\over256},52,{95221\over1280},{22407\over256}\right)\sigma;
 \tag{4.34}
\]
passing to $X\le N$ preserves the first margin.
After subtracting $1/3$ it is $28+53/768>28$, and the other margins exceed the losses in (4.33).
The box discrepancy is
\[
 \Delta_B(q;N)=\sum_{\mathbf v\in B}w(\mathbf v)
 \left[\mathbf1_{\prod v_i\equiv N\ (q)}
       -{\mathbf1_{(\prod v_i,q)=1}\over\varphi(q)}\right].
 \tag{4.35}
\]
For every source form and row, the preparation data are:
\begin{center}
\begin{tabular}{@{}lp{0.68\textwidth}@{}}
\toprule
Object & Common prepared value \\
\midrule
Selected factor & Largest eligible lower scale, first index in a tie; $T$ in (4.29) \\
Selected interval & $I_N=(T,qT]$, $N_I=\lfloor qT\rfloor-\lfloor T\rfloor$ \\
Complement & $\tau_{23}$; $M_I=M(q^r-1)$, $1\le r\le23$ \\
Majorant & $\phi_0((m-m_0)/M_I)$, $m_0\le2M$; doubled support (4.30) \\
External multiplicity & $648$ fixed instances; $O(L^{8042})$ boxes, (4.66) \\
\bottomrule
\end{tabular}
\end{center}
The hypothesis checks are collected below.
\begin{center}
\begin{tabular}{@{}ll@{}}
\toprule
Hypothesis & Verification \\
\midrule
H1, H2 & (4.22)--(4.26), with repeated-prime multiplicities \\
H3, LEN & Full support after (4.24), Lemma~\ref{lem:g1-f}, (4.34)--(4.35) \\
H4 & $a=N$, $K_{\rm scale}=31$, $(q,N)=1$ \\
H5 & The two windows of Lemma~\ref{lem:g1-f} \\
H6, H7 & (4.43), (4.63), (4.66) \\
\bottomrule
\end{tabular}
\end{center}
For H4, eventually $0<N\le X\log^{31}N\le X^{1+\sigma/20}$,
so Theorem~\ref{thm:distribution-input} applies to the plain cut product measure.
The local estimate is \cite[Theorem~3.2, (3.13)]{Z1}.
The spectral frequency is $a_T=\operatorname{sgn}(y)N$.
If $u_1,v_1$ invert $zb$ modulo $p_0c,c$, then $u_1-v_1\equiv0\pmod c$.
Hence $\operatorname{sgn}(y)N|y|(u_1-v_1)/c$ is integral, also for $p_0\mid c$, without any condition $(p_0,c)=1$.

\subsection{Source errors and the Class-II sum}

This subsection bounds the boundary, proper-power, logarithm and comparator errors and the interior Class-II sum.
A source face is
\[
g(\mathbf v)=\sum_ia_i\log v_i=\xi,\qquad
 |a_i|\le2,\quad\sum_i|a_i|\le48.
 \tag{4.36}
\]
Every defining-face coefficient has absolute value at most 2 after
expansion into the HB variables. Since $s\le24$, its coefficient sum
is at most 48.
The face families are:
\begin{center}
\begin{tabular}{@{}p{0.68\textwidth}l@{}}
\toprule
Family & Number \\
\midrule
Two $\Sigma_G$ faces, source windows/orderings, fixed coordinate transports & $\le16$ fixed \\
Coordinate dyadic edges & $\le4$ lattices \\
M\"obius caps $v_i=2X_m^{1/3}$ & $\le12$ lattices \\
Identity caps $n_m=(2X_m^{1/3})^3$ & $\le4$ lattices \\
Product dyadic edge and three slack families & $\le4$ lattices \\
\bottomrule
\end{tabular}
\end{center}
Thus there are at most 24 lattice families, of gap $\gamma=(\log2)/3$.
HB partition edges are refined exactly, and prefix and tail labels are modulus-side conditions.
The Class-II choice is constant on each box.

\begin{lemma}[Lattice thin shells]
Let $N\ge2$, $L=\log N$, $s\ge2$, and $a\in\mathbb Z^s\setminus\{0\}$.
Fix a nonempty $J_0\subseteq\{i:a_i\ne0\}$, and assume $|a_j|\in\{1,2\}$ for $j\in J_0$.
Let $1\le Y_0\le N$ and $0<\eta\le1/4$, and let $\Lambda\subset\mathbb R$ have pairwise distances at least $\gamma>0$.
Put $\kappa=(1-e^{-\gamma/2})^{-1}$.
Then
\[
\begin{aligned}
 &\#\{\mathbf v:\prod v_i\le N,\
 \operatorname{dist}(\sum a_i\log v_i,\Lambda)\le\eta,\
                         \max_{j\in J_0}v_j\ge Y_0\}\\
 &\quad\le s\left[4\eta\kappa N(1+L)^{s-1}
       +(2L/\gamma+2){N\over Y_0}(1+L)^{s-2}\right].
\end{aligned}
 \tag{4.37}
\]
For empty $\Lambda$, the count is zero.
For a singleton one may take $\kappa=1$ and interval-count factor 1.
\end{lemma}
\begin{proof}
Assign each tuple to the least $j\in J_0$ with $v_j\ge Y_0$, and apply \cite[Lemma~B.2, (B.14)]{Z1} on that fibre.
Only $|a_j|\in\{1,2\}$ matters there; the other coefficients enter $c=\sum_{i\ne j}a_i\log v_i$.
Put $Z=N/P$, where $P$ is the complementary product.
After clipping to $[Y_0,Z]$ and merging touching or overlapping intervals, the union has at most $2L/\gamma+2$ components.
Its fibre count is at most $4\eta\kappa Z+2L/\gamma+2$, the count before relaxation in (B.14).
Summing uses
\[
 \sum_{n\le y}\tau_k(n)\le y(1+\log y)^{k-1},\qquad
 \sum_{n\le y}{\tau_k(n)\over n}\le(1+\log y)^k,
 \tag{4.38}
\]
with $k-1$ factors fixed and the product constraint dropped in the harmonic sum.
These give (4.37).
The component count concerns the clipped union, and no restriction relates $\eta$ and $\gamma$.
\end{proof}

The boundary distance is at most $\eta=48L^{-334}$.
For each defining face, take $J_0=\{i:a_i\ne0\}$.
Six or fewer factors of a coordinate at least $N^\alpha$ give $Y_0=N^{\alpha/6}$.
A M\"obius-cap variable is at least $N^{\alpha/3}e^{-\eta}\ge Y_0$.
The same 16 fixed and 24 lattice families suffice.
The worst lattice has $\kappa=(1-2^{-1/6})^{-1}<9.2$, and $2L/\gamma+2\le10L$ for $L\ge2$.
Thus the boundary union $U$ for one form and pattern satisfies
\[
\begin{aligned}
 \#U\le{}&1091174.4\,NL^{-334}(1+L)^{23}\\
 &+24(16+240L)N^{1-\alpha/6}(1+L)^{22}\\
 \le{}&1091175\,NL^{-334}(1+L)^{23}\qquad(L\ge300000).
\end{aligned}
 \tag{4.39}
\]
The leading constant is $24\cdot192(16+24\cdot9.2)=1091174.4$.
The second-line ratio is bounded by
\[
 R(L)=5760e^{-(100/6894)L}L^{334}<3/5\quad(L\ge300000).
 \tag{4.40}
\]
Its logarithmic derivative is negative there.
At 300000, $e>8/3$, $(8/3)^{13}>300000$ and $(8/3)^9>5760$ give $\log R<9-(100/6894)300000+334\cdot13<-3/5$.
Since $e^{3/5}>5/3$, (4.40) follows.
This justifies the rounding at width exponent 334.

Boundary estimates keep the unnormalized HB logarithms.
The full weight $w(\mathbf v)=H_B(\mathbf v)\prod_{j=1}^k(\log p_j)^{-1}$
uses the actual coordinate logarithms; since $\log p_j\ge\alpha L\ge1$,
the numerator bound above gives the deliberately coarser bound below.
Thus for one pattern, including its source-form weight,
\[
 |w(\mathbf v)|\le162L^4,\qquad
 \sum_{q\le N^{3/5}}|\gamma_q|
  \left(\mathbf1_{q\mid N-e}+{1\over\varphi(q)}\right)
 \le C_\sigma\{\tau(N-e)+2(1+L)\}.
 \tag{4.41}
\]
Here $N-e\ge1$.
The reciprocal-totient constant is below 2, since
$q/\varphi(q)=\sum_{d\mid q}\mu^2(d)/\varphi(d)$ and
$\prod_p(1+1/(p(p-1)))<2$.
Put $F_*(e)=\tau(N-e)+2(1+L)$.
The composition bound is $\tau_a\tau_b\le\tau_{ab}$:
one maps the two compositions to a fixed matrix with those margins.
Cauchy--Schwarz and (4.38) give
\[
\begin{aligned}
 \sum_{e<N}\tau_{24}(e)\tau(N-e)^2
 &\le\left(\sum\tau_{576}\right)^{1/2}
       \left(\sum\tau_{16}\right)^{1/2}
 \le N(1+L)^{295},\\
 \sum_{e<N}\tau_{24}(e)F_*(e)^2&\le3N(1+L)^{295}
                              \quad(N\text{ sufficiently large}).
\end{aligned}
 \tag{4.42}
\]
Another Cauchy--Schwarz application on $U$ gives
\[
\begin{aligned}
 |E_{\rm bdry}|&\le162\sqrt{3\cdot1091176}\,C_\sigma\\
 &\qquad\cdot N(1+L)^{163}L^{-167}\\
 &<293108\,C_\sigma\\
 &\qquad\cdot N(1+L)^{163}L^{-167}.
\end{aligned}
 \tag{4.43}
\]
We include 648 instances and the multiplier $1/u=500$, since $u=1/500$;
we also include $N_{\rm lab}(\sigma)$ and $\mathcal W^{-1}$.
Since $C(N)\ge0.66016$,
\[
\begin{aligned}
 E^D_{\rm HB:G1}&\le3.6\cdot10^{10}C_\sigma N_{\rm lab}(\sigma)
                     (1+L)^{163}L^{-165}\\
 &\le3.6\cdot10^{10}C_\sigma N_{\rm lab}(\sigma)
                         e^{163/L}L^{-2}=o_\sigma(1).
\end{aligned}
 \tag{4.44}
\]
This covers cut faces, windows/orderings, dyadic families and caps.
The four parts count row/form/pattern/box occurrences once, including box and dyadic factors.
Hence H6 has $c_R=163$, with D rate $O_\sigma(L^{165-C_b/2})$.
Every fixed $C_b>330$ suffices; $C_b=334$ gives $O_\sigma(L^{-2})$.

The proper-power event concerns one source coordinate,
while other coordinates may contribute further prime factors.
For $n\ge N^\alpha$,
\[
 \mathbf1_{\mathbb P}(n)-{\Lambda(n)\over\log n}
       =-{1\over a}\mathbf1_{\substack{n=p^a\\a\ge2}}.
 \tag{4.58}
\]
Roughness below $N^\alpha$ preserves primes and removes powers.
The same absolute bound is used, telescoping over four or fewer coordinates.
For $Y=N^\alpha$,
\[
 \sum_{\substack{Y\le p^a\le N\\a\ge2}}p^{-a}
 \le2Y^{-1/2}+2(\log_2N)Y^{-2/3}\le3Y^{-1/2}
 \tag{4.59}
\]
Squares are bounded by an integer tail, and for $a\ge3$ we use $O(Y^{-(a-1)/a})$.
Fixing one coordinate at $t$ leaves at most
\[
 {N\over t}(1+\log N)^{\max(r-2,0)}\qquad(1\le r\le4)
\]
completions.
For $r=1$ the sole completion is bounded by $N/t$.
For $r=2$ we count directly, and for $r\ge3$ we use (4.38).
Thus $T_{PP}\le12N^{1-\alpha/2}(1+L)^2$.
By (4.25), the full modulus sum is
$O_\sigma(\tau(N-e)+2(1+L))$ per tuple.
Since $1\le N-e<N$, we use $\tau(N-e)\ll_\delta N^\delta$.
Then
\[
 |E_{PP}|\ll_{\sigma,\delta}N^{1-\alpha/2+\delta}(1+L)^3.
 \tag{4.60}
\]
We take $\delta=\alpha/8$ and absorb logarithms in the spare power.
Dividing by $\mathcal W$ gives
\[
 E^D_{PP:G1}\ll_\sigma N^{-\alpha/4}L^5
      =N^{-25/1149}O_\sigma(L^5)=o_\sigma(1).
 \tag{4.61}
\]
This step precedes HB boxes.
The event $p^2\mid q$ has its separate modulus bound.

For logarithm freezing, each coordinate has at most six HB factors.
If $\lambda_{B,m}$ is its lower logarithm,
\[
 \left|{1\over\log p_m}-{1\over\lambda_{B,m}}\right|
 \le {6L^{-334}\over(\alpha L-1)^2}.
 \tag{4.62}
\]
Up to four inverse logarithms change by at most $c_\alpha L^{-336}$.
This is multiplied by the unnormalized numerator bound $162L^4$ above.
The frozen inverse logarithms contribute to (4.62),
while the normalized interior coefficient $162(2/\alpha)^4$ excludes them.
Summing exact differences over disjoint boxes,
with $\sum\tau_{24}\le N(1+L)^{23}$, (4.42) and Cauchy--Schwarz, gives
\[
 |E_{LOG}|\ll_\sigma N(1+L)^{161}L^{-334},\qquad
 E^D_{LOG}\ll_\sigma L^{-171}\ll_\sigma L^{166-C_b}=L^{-168}.
 \tag{4.63}
\]
Hence H6 holds with $C_{\rm box}=C_b=334$ and $c_{\log}=164$.

For the comparator error, source primes are at least $N^\alpha$ before HB.
Thus
\[
 X_G-Y_G(q)
 \le100\sum_{\substack{p\mid q\\p\ge N^\alpha}}{N\over p}
 \le100\omega(q)N^{1-\alpha}.
 \tag{4.64}
\]
Also $\sum_q|O(q)|\omega(q)/\varphi(q)\ll_\sigma L^3$, by (4.25) and reciprocal-totient summation.
Equations (4.23)--(4.24) give
\[
 E^D_{NC:G1}\ll_\sigma N^{-\alpha}L^5=o_\sigma(1).
 \tag{4.65}
\]
The two squareful terms in (4.23) contribute
$O_\sigma(N^{-\ell_0\sigma^2}L^4)$ once, by Theorem~\ref{thm:repeated-primes}.
Class II includes GR cross-copy restoration.

The external precisions are
\[
\begin{gathered}
 A_{\rm raw}=3,\quad K_0=24C_b+26=8042,\\
 A'=A_{\rm raw}+K_0+2=8047,\quad
 A_{\rm work}=A'+4=8051.
\end{gathered}
 \tag{4.66}
\]
These use the \hyperref[p2:g1-external-count]{box count above}.
Three internal slot-shell sums use three of the four extra logarithms.
Choose
\[
\begin{aligned}
 A_E&\ge16770+c_*,\qquad C_1=A_E+2c_*+3,\\
 C_M&=5C_1+A_E+334+c_*+3,\qquad A_3=A_E+336,\\
 A''&=A_{BFI}+(334+A_3+1)/2.
\end{aligned}
 \tag{4.67}
\]
Here $16770=2A_{\rm work}+2C_b$.
Fix $c_*\ge\max\{c_{\rm pre}+2,c_1,c_{\rm diag},c_1(B_0),0\}$
and above all remaining preceding logarithmic losses, with
$c_{\rm pre},c_1,c_{\rm diag}$ as in (1.29).
Then (4.67) implies every lower bound in (1.29), in particular
$A_E\ge16772+c_{\rm pre}$.
Precisions are fixed before the threshold.
Theorem~\ref{thm:distribution-input} gives
\[
 |S_B|\ll X(\log X)^{-8047}
      +X^{1-c_\sigma}(\log X)^C,\qquad
 c_\sigma={\ell_0\sigma^2\over4}>0.
 \tag{4.68}
\]
This includes local restoration and GR at precision $A_{\rm work}$.
The prime-source condition $C_{\rm box}\ge\lceil2A+32\rceil$ uses
the global target $A$ of \cite[Theorem~1.2]{Z1}.
The local target $A'$ does not enter this box-width requirement.
The product-source H6 exponents are those above.
Lemma~\ref{lem:external-classii-sum} applies with $C_{\mathscr F}=C$:
by (4.2) and (3.1), $\mathcal W=W_N$, and (4.28) gives $X\ge N/L^{31}$ eventually.
Set $C_{\rm pow}=C+31c_\sigma$, which is a fixed finite exponent.
Hence
\[
 E^D_{II:G1}\ll_\sigma L^{-3}
 +N^{-c_\sigma}L^{8042+C_{\rm pow}+2}=o_\sigma(1).
 \tag{4.69}
\]
Scale conversion preserves $c_\sigma$, with no $N^{-\sigma}$ saving.
The stronger (4.28) is not needed here.
Boundary, comparator, proper-power and logarithm sums have no $L^{8042}$ factor.
The raw boundary rate is $O_\sigma(NL^{-4})$, with D rate $O_\sigma(L^{-2})$.
All these errors vanish at the fixed target.

\subsection{The dynamic-cutoff peel}

This subsection bounds the error of the dynamic-cutoff peel.
For a retained G1-3 cell, put $b=a_1a_2a_3$ and $E=N^{t_3}$.
The finite union bound is
\[
 S((\mathcal B_G)_b,a_3)\ge S((\mathcal B_G)_b,E)
       -\sum_{\substack{E\le r<a_3\\r\ {\rm prime}}}
                             S((\mathcal B_G)_{br},E).
 \tag{4.45}
\]
Here $a_1>a_2>a_3$, $r<a_3$, $r\nmid N$ and $(br,N)=1$.
The last two follow from $\mathcal P(N)$ and the coordinate conditions.
The negative upper-chain sign permits enlargement of this nonnegative difference.
Only the $E$-base enters HB.

\begin{lemma}[Uniform peel bound]
Let $0<\tau<8\cdot10^{-6}$ and $\rho=1+\tau$.
Suppose $(m,N)=1$ and $P^-(m)\ge z$, with
\[
 \log_Nm\le153/256+2\tau,\quad
 N^{u/\rho^3}\le z\le N^{\rho\mu},\quad w=\min(z,N^\alpha).
 \tag{4.46}
\]
Eventually, uniformly in $m$,
\[
 S((\mathcal B_G)_m,z)\le3\cdot502002\,{U_0\over m},
 \qquad U_0=2U_C.
 \tag{4.47}
\]
\end{lemma}
\begin{proof}
Write $N-e=mk$.
By the representation bound above, $A_G^{\rm cut}\le3$.
Put $\mathcal E_h=\{e:A_G^{\rm cut}(e)\ge h\}$ for $h=1,2,3$.
Then $A_G^{\rm cut}=\sum_{h=1}^3\mathbf1_{\mathcal E_h}$.
Each layer has the same roughness and two-form bound.
The layers are summed once, including coefficient 2 on both $R_{11}$ forms.

Enlarge the range to $1\le k\le\lfloor(N-1)/m\rfloor$, of length $K\le N/m$;
both $k$ and $N-mk$ are $w$-rough.
If $p<w$, $p\mid N$ and $p\mid k$, then $p\mid e$, which source roughness excludes.
Thus the all-primes two-form majorant applies.
For squarefree $d\mid P(w)$, CRT gives
\[
\begin{gathered}
 \nu_N(p)=1\ (p\mid N),\quad \nu_N(p)=2\ (p\nmid N),
 \quad \nu_N(2)=1,\\
 A_d=K{\nu_N(d)\over d}+R_d,\quad |R_d|\le\nu_N(d).
\end{gathered}
 \tag{4.48}
\]
For $p<w$, $p\nmid m$ follows from $P^-(m)\ge z\ge w$.
For $p\ge w$, the same rules are extended only in the Euler product.
A prime endpoint costs $g(w)=O(1/w)$ in the auxiliary sum.
Selberg least common multiples stay on $P(w)$.
Put $f(p)=\nu_N(p)/p$ and $g(p)=f(p)/(1-f(p))$.
They are extended squarefree multiplicatively, and we put
$G_N(w)=\sum_{d\le w,d\mid P(w)}g(d)$.
Finite-poset diagonalization gives
\[
 \sum_{d,e}\lambda_d\lambda_ef([d,e])
 =\sum_r{y_r^2\over g(r)},\quad
 y_r=\sum_{r\mid d}\lambda_df(d),\quad
 1=\lambda_1=\sum_r\mu(r)y_r.
 \tag{4.49}
\]
Cauchy--Schwarz gives minimum $1/G_N(w)$ at $y_r=\mu(r)g(r)/G_N(w)$.
Finite inversion gives
\[
 \lambda_d=\mu(d)\prod_{p\mid d}(1-f(p))^{-1}
     {\sum_{r\le w/d,(r,d)=1}g(r)\over G_N(w)},\\
 |\lambda_d|\le\prod_{p\mid d}(1-f(p))^{-1}
              \ll(d/\varphi(d))^2.
 \tag{4.50}
\]
The Dirichlet series factors as $\zeta(1+s)^2H_N(s)$,
where $H_N(s)=\sum h_N(r)r^{-s}$.
Its nonconstant local coefficients are
$g(p)-2/p$, $p^{-2}-2g(p)/p$ and $g(p)/p^2$.
For $p\nmid2N$, absolute logarithmic moments through order 2 converge.
For $p\mid N$, the absolute local mass is $O(1/p)$.
We split at $Y=\log(3N)$.
Mertens gives $\sum_{p\mid N,p\le Y}1/p\ll\log\log Y+1$.
Above $Y$ we use $\sum_{p\mid N}\log p\le\log N$.
The $j$-th logarithmic moment is $O((\log Y)^{j-1})$, for $j=1,2$.
Expanding absolute local products gives a fixed $B$ with
\[
 \sum_r|h_N(r)|(1+\log r)^2\ll(\log\log(3N))^B,\qquad
 H_N(0)={1\over2C(N)}.
 \tag{4.51}
\]
At zero, 2 contributes $1/2$.
A generic odd prime contributes $(1-1/(p-1)^2)^{-1}$.
For $p\mid N$, this is multiplied by $(p-2)/(p-1)$.
Convolution with $\sum_{n\le y}\tau(n)/n=\frac12\log^2y+O(\log(2y))$
and the second logarithmic moment for the tail $r>w$ give, uniformly,
\[
 G_N(w)={\log^2w\over4C(N)}
       \left(1+O\left({(\log\log(3N))^B\over\log w}\right)\right).
 \tag{4.52}
\]
Thus the dimension-two convention is
$G_N(w)\prod_{p<w}(1-\nu_N(p)/p)=e^{-2\gamma}/2+o(1)$.
Equations (4.48), (4.50) give
\[
 S([1,K],w)\le {K\over G_N(w)}+O(w^2\log^5w);
 \tag{4.53}
\]
the remainder is at most
$(\sum_{d\le w}|\lambda_d|\nu_N(d))^2
\ll w^2\log^2w(\log\log(3w))^4$.
Uniformly in (4.46),
\[
 \log_Nw\ge2/1001,\qquad
 w^2m\le N^{253/256+(153/64)\tau}<N^{1-0.0116}.
 \tag{4.54}
\]
These use $(1+8\cdot10^{-6})^3<1001/1000$, $u/\rho^3\ge2/1001$ and $\alpha>2/1001$,
together with $w\le z\le N^{\rho\mu}$ for the uniform upper saving.
For one layer, the main coefficient in $U_0/m$ units is eventually at most
$2.001/(2/1001)^2=501251.00025<502002$.
The strict slack absorbs the remainder in (4.53).
This proves (4.47) on the $N$-scale, including short $K$.
\end{proof}

For $m=br$ in (4.45), the retained-cell condition gives
\[
 \begin{aligned}
 \log_Nm&\le\rho(t_1+t_2+2t_3)\\
 &\le\rho[\Theta_c-\eta_L-(2\rho-1)t_3]\\
 &\le\rho(153/256-\eta_L)<153/256+2\tau.
 \end{aligned}
 \tag{4.55}
\]
On G1-1, G1-2 and G1-3, $\Theta_c\le\overline H=153/256$.
The two-branch peak is $(1+\mu)/2=153/256$.
Masked branches are zero or below it, and extra $w$-branches are zero or at most $(153-96c)/256\le153/256$.
The inherited label has $N^{t_3}\le a_3<N^{\rho t_3}$,
even for a clipped parent interval.
Since $N^u\le a_3<N^\mu$,
$t_3>u/\rho\ge u/\rho^3$ and $t_3<\mu\le\rho\mu$.
Hence $E\in[N^{u/\rho^3},N^{\rho\mu}]$.
All prime factors of $m$ are at least $E$, and $(m,N)=1$.
We apply (4.47) with $z=E$; $w$ may be smaller.
For fixed positive $\tau$, Mertens gives eventually
\[
 M_h=\sum_{p\ {\rm in\ parent\ bin}\ h}{1\over p}\le2\tau,
 \qquad\sum_hM_h\le4.59,
 \tag{4.56}
\]
Here $\log(\mu\rho/u)<4.59$.
The pair $a_3,r$ belongs to one parent bin. Thus
\[
\sum(a_1a_2a_3r)^{-1}\le4.59^2\sum_hM_h^2
\le2\cdot4.59^3\tau.
\]
The raw difference is below $3\cdot97100000\tau U_0$.
Its adverse D bound is
\[
 E^D_{\rm PR:peel}\le {3\cdot97100000\over2(1-\varepsilon_A)}\tau
 ={485500000000\over3333}\tau<145664566.4567\tau.
 \tag{4.57}
\]
The source multiplier enters once, without 500, label counts or tail factorials.
The slack includes the Selberg remainder.

\subsection{Canonical profiles and the joint lower bound}

This subsection compares the rows with canonical profiles and proves the joint lower bound for the four rows.
Put
$v(s)=v_1(s)-u,\qquad \psi(a,b)=a\mathbf1_{a\ge b},$
with the \hyperlink{p2:def:six-branch-saved}{six-branch definition}
of Section~\ref{p2:sec:2.1}, where $v_1$ is the unbuffered function of \eqref{p2:eq:row-level-function}.
Thus
\[
 \Theta_*(t_1,t_2,t_3)=\Theta_3^{\rm saved}(t_1,t_2,t_3).
 \tag{4.70}
\]
For $u\le t_3\le t_2\le t_1\le\mu$, put
\[
\begin{aligned}
v_1(s)&=\min\left\{\frac{1+s}{2},\frac{32-34s}{39},
                     \frac{57}{64}-\frac32s\right\},\\
v(s)&=v_1(s)-u,\qquad \psi(x,y)=x\mathbf1_{x\ge y}.
\end{aligned}
\]
The six branches are
\[
\begin{aligned}
\Theta_*(a,b,c)=\max\{&v(a),v(b),v(a+b),v(a+b+c),\\
 &\psi(v(a+c),a+2b+c),\psi(v(b+c),2a+b+c)\}.
\end{aligned}
\]
The $u$ subtraction precedes each $\psi$ guard.
For the actual coordinates, set
\[
 a_3={\Theta_*-t_1-t_2-t_3\over t_3},\qquad
 J=\int_u^\mu{dt_1\over t_1}\int_u^{t_1}{dt_2\over t_2}
                  \int_u^{t_2}{1-f(a_3)\over t_3^2}\,dt_3.
 \tag{4.71}
\]
Here $f$ is extended by zero below 2.
Above every finite table endpoint its true values are used.

The function $v_1$ is $3/2$-Lipschitz, and each exponent changes by at most $\mu\tau$ in a parent cell.
An active $\psi$-branch has unshifted gap $u-(9/2)\mu\tau-4\mu\tau>0$ throughout.
After $\eta_L$, its source margin is
\[
 g=(u-\eta_L)-(9/2)\mu\tau>0.001899>2\mu\tau.
 \tag{4.72}
\]
The other branches and maxima preserve this comparison across switches, and so do the two full-anatomy branches.
If a parent crosses $\mu$, then $t_1\le\mu\le\rho t_1$, so $|x_1-\mu|\le\mu\tau$ for its largest actual exponent.
Since $v_1(\mu)=153/256$,
$v_1(x_1)-\eta_L\ge153/256-(3/2)\mu\tau-\eta_L$.
Every canonical branch is at most $153/256-u$, so the margin is at least $u-\eta_L-(3/2)\mu\tau\ge g>0$.
Parents below $\mu$ use the activation argument above, and clipping shrinks the infimum domain.

If $f(a_3)>0$, the numerator exceeds $2t_3$,
and (4.72), for actual exponents then inherited labels, gives $\ell_c\ge2\rho t_3$.
Thus positive canonical lower terms lie in retained whole cells, and elsewhere their contribution is zero.
On retained cells we use Mertens and the monotonicity of $sf(s)$, keeping the numerator fixed.
The $E$-base dominates the actual-cutoff lower main term, so canonical restriction loses zero.
The arithmetic difference (4.57) remains, with possibly nonzero mass.

Let $H_{\log}=\log(\mu/u)$.
By (4.20), the first three residuals satisfy
\[
 0\le R\le500(F(30)-1)+500H_{\log}(1-f(30))
                         +250H_{\log}^2(F(30)-1)=:R_{\rm up}.
 \tag{4.73}
\]
Three elementary integrations give
\[
 500(1-H_{\log}+H_{\log}^2/2)
 -\int_u^\mu{dt_1\over t_1}\int_u^{t_1}{dt_2\over t_2}
                            \int_u^{t_2}{dt_3\over t_3^2}
 ={1\over\mu}
 \tag{4.74}
\]
which cancels the constant-one parts of $+F,-f,+F,-f$.
The actual upper main chain is therefore at most
\[
 G_1={1\over\mu}+R+J,\qquad0\le R\le R_{\rm up},
 \tag{4.75}
\]
with the preceding errors.

Let $X_j=\sum_eR_j(e)\mathbf1_{\varepsilon_A N<e<N}$,
and $X_j^{\rm cut}=\sum_eR_j(e)\mathbf1_{\varepsilon_A N<e\le N_{\rm cut}}$.
Then $X_G=2X_{11}^{\rm cut}+X_{12}^{\rm cut}+X_{13}^{\rm cut}$, and (4.17) bounds the mass difference.
PNT and partial summation give
\[
 X_j=\big((1-\varepsilon_A)I_j+o(1)\big){N\over L},
 \tag{4.76}
\]
where
\[
\begin{aligned}
 I_{11}={}&\int_\lambda^{1/3}
       {\log((1-2t)/t)\over t(1-t)}\,dt\\
 &+\int_\lambda^{1/4}{dt\over t}
       \int_t^{(1-t)/3}{\log((1-t-2v)/v)\over v(1-t-v)}\,dv,\\
 I_{12}={}&\int_\alpha^{1/3}{\log(2-3t)\over t(1-t)}\,dt,\\
 I_{13}={}&\int_\beta^\lambda
       {\log((1-t-\lambda)/\lambda)\over t(1-t)}\,dt.
\end{aligned}
 \tag{4.77}
\]
The last-prime density on the radial face is $1/(1-t-v)$.
Integrating the companion variable gives $I_{12},I_{13}$.
For $I_{11}$ the Buchstab argument is at most $214/93$, and its value is $v/(1-t-v)$ plus
$v\log((1-t-2v)/v)/(1-t-v)$ for $v\le(1-t)/3$.
This gives both integrals and the positive correction.
Set $e=yN$, $\varepsilon_A<y<1$.
On compact exponent boxes, the last logarithm has radial-face error
$O_{\varepsilon_A}(1/L)$, and $\int_{\varepsilon_A}^1dy=1-\varepsilon_A$.
PNT is applied first and fixed partitions are then refined,
and the positive exponent lower bound gives domination.
Monomial-face strips, including $D_{13}$'s endpoint-swap strip, have vanishing main mass.
Representation conditions are retained before the limit.
The upper cut costs (4.17).
This proves (4.76) at finite $\varepsilon_A$.

With $V_N(z)=(1+o(1))2e^{-\gamma}C(N)/\log z$, (4.75)--(4.76) give
\[
 {2S_{11}+S_{12}+S_{13}\over4(1-\varepsilon_A)U_C}
 \le K_GG_1+E^D_{G1},\qquad
 K_G={2I_{11}+I_{12}+I_{13}\over2e^\gamma}.
 \tag{4.78}
\]
The factor $1-\varepsilon_A$ cancels once here.
It remains in (4.57) and in the finite target of Section~1.
At $a=9/5$, that target is $8\log(50009/39991)$, with zero-epsilon limit $8\log(5/4)$.

The directed bounds are
\[
\begin{gathered}
 R_{\rm up}<4.4\cdot10^{-6},\qquad G_1\le\overline G,\\
 \overline G=6.0884403625974206828946573628908904\ldots<6.088440363.
\end{gathered}
 \tag{4.79}
\]
Appendix~A.5 gives the delay-function, quadrature and mass enclosures.
Its endpoint $K_{\rm hi}$ is \eqref{eq:g1-khi}.
Here $\overline G$ is the 4096-cell outward upper-corner sum in Appendix~A.5.
It is a conservative bound, not the continuous integral value.
Let $r_0,r_1,r_2,J_{\rm hi}$ denote the directed upper parts of $\overline G$.
Put
\[
\begin{aligned}
 T_0&=500+r_0,&T_1&=-500H_{\log}+r_1,\\
 T_2&=250H_{\log}^2+r_2,&
 T_3&=-(500-500H_{\log}+250H_{\log}^2-1/\mu)+J_{\rm hi}.
\end{aligned}
 \tag{4.83}
\]
Exactly, $\sum T_r=1/\mu+r_0+r_1+r_2+J_{\rm hi}=\overline G$.
Their signs are $+,-,+,-$, so $-K_GT_r$ has signs $-,+,-,+$.
We therefore use upper $K_G$ for rows 0,2 and lower $K_G$ for rows 1,3.
The four directed row enclosures are in Appendix~A.5,
and the four terms are cancelled before rounding.
The joint reference value is
\[
\begin{aligned}
 L_{G1}&=
 -{243950428265194656254206921939947762868871361874409552095097
 \over100433627766186892221372630771322662657637687111424552206336}\\
 &=-2.42897158741611974640561365334839795714\ldots .
\end{aligned}
 \tag{4.84}
\]
Put
$\delta_{G1}=4.98\cdot10^{-54},\qquad
 E_{G1}=L_{G1}-\delta_{G1},$
and let $R_{G1}$ sum the exact four row lower endpoints.
Appendix~A.5 proves $L_{G1}-R_{G1}<\delta_{G1}$, so $E_{G1}$ is used.
The main base contains the reference value $L_{G1}$ once,
and the adverse sum contains $\delta_{G1}$ once.
The row enclosures include $K_G$.
The cap for $G_1-2$ is $4.088440363$, which gives the weaker value $-K_{\rm hi}(2+4.088440363)=-K_{\rm hi}6.088440363$,
with about $1.61\cdot10^{-10}$ D loss.
Equations (4.13)--(4.78) transfer these enclosures to the arithmetic rows.

Successive differences are taken in the order
\[
 \begin{gathered}
 \text{source restoration, source conditions, PR-M, AA37,}\\
 \text{squareful/NC, proper powers, LOG, boundary, interior Class II}.
\end{gathered}
\]
AA37 includes fixed-parameter and asymptotic terms.
Only the $E$-base enters G1-3 Heath--Brown \cite{ref6}.
After monotonicity the exact first-contribution partition is retained,
and the favourable fixed-cutoff and G1 cutoff differences are discarded.
Each difference or part occurs once below.
\begingroup
\setlength{\tabcolsep}{3pt}
\begin{longtable}{@{}>{\raggedright\arraybackslash}p{0.24\textwidth}>{\raggedright\arraybackslash}p{0.68\textwidth}@{}}
\toprule
Contribution & Adverse D bound at fixed $\sigma,\tau$ \\
\midrule
Upper source restoration & (4.17) \\
Lower slab & zero: outside the source by (4.15) \\
G1 source-tag exceptions & $O_\sigma(N^{-\alpha/2+o(1)})$; the DIRECT $O_\sigma(L^3/N)$ term is separate, as in (1.23) \\
G1-3 peel & $145664566.4567\tau$, once \\
Matched canonical restriction loss & zero, by (4.72) \\
AA37 & its common $C_E\sigma$ share once, plus $o_\sigma(1)$ \\
Squareful terms & Theorem~\ref{thm:repeated-primes}: $O_\sigma(N^{-\ell_0\sigma^2}L^4)$ \\
NC & (4.65) \\
Tested proper powers & (4.61) \\
LOG & (4.63) \\
HB boundary & (4.44), one term with four parts \\
Interior Class II, including GR & (4.69) \\
Prime mass/PNT and Mertens & $o_\sigma(1)$, from (4.76)--(4.78) \\
\bottomrule
\end{longtable}
\endgroup
For $p\mid N$, there are $O((N/p)L^C)$ complementary products.
There are $O(L)$ such source primes, all at least $N^\alpha$.
A positive fraction of $\alpha$ absorbs divisor/logarithmic losses.
Labels and HB boxes preserve these sums.

Theorem~\ref{thm:distribution-input} and Theorem~\ref{thm:repeated-primes} in Appendix~B.1,
together with the directed endpoints above, give
\[
 \sum_{r=0}^3 Q_{G1-r,0}(N)
 \ge L_{G1}-\delta_{G1}-145664566.4567\tau-E^D_{AA:G1}(\sigma)
                    -E^D_{{\rm remaining}:G1}(N;\sigma),
 \tag{4.85}
\]
where the last term sums the vanishing errors in the table.
The AA37 term $C_E\sigma$ and the PR term $O(\tau)$ stay fixed.
Analytic constants and the residual margin are fixed first.
Then $\sigma$ is chosen, followed by $\tau=\sigma^9$, labels, widths and precisions.
Finally $N_0$ is chosen for all length and spectral conditions.
These include BFI/SW, Poisson and the W-prime bound at width 334.
The threshold depends on fixed $\sigma>0$.
Thus (4.85) holds for all four arithmetic rows, with $E_{G1}=L_{G1}-\delta_{G1}$.

\section{Direct BV rows, the O1 switched term, and weak G4}

This section treats the rows R1, S5-top and S6 by prime BV,
and the switched term O1 and weak G4 by one convolution lemma.

\subsection{Sources and classical inputs}

Fix even $N$ and put $L=\log N$.
Fix $\varepsilon=\varepsilon_A=10^{-4}$ before $N\to\infty$.
Here $q$ denotes a modulus, $\lambda_d$ sieve weights
and $\ell$ a sieve prime.
Put $\rho=1+\tau$, and write $h$ for an interval width.
The exponent $K_0$ measures the fixed logarithmic loss.
We use
\begingroup\small
\[
 X=(1-\varepsilon)N,\quad
 \mathcal A_\varepsilon=\{N-p_0:p_0<X,\ p_0\text{ prime}\},\quad
 \mathcal P_N=\{\ell>2:\ell\text{ prime},\ \ell\nmid N\},
\]

\[
 P_N(z)=\prod_{\ell<z,\,\ell\in\mathcal P_N}\ell,\qquad
 C(N)=\prod_{\ell>2}\left(1-\frac1{(\ell-1)^2}\right)
       \prod_{\ell\mid N,\,\ell>2}\frac{\ell-1}{\ell-2},\qquad
 U_C=\frac{C(N)N}{L^2}.
\tag{5.1}
\]

\endgroup
Here $C(N)\ge C_2>0$, and the raw coefficients satisfy $\Sigma=4D$.
For finitely supported $f$, put
\[
 \Delta_f(q,a)=\sum_{e\equiv a\pmod q}f(e)
       -\frac1{\varphi(q)}\sum_{(e,q)=1}f(e),\qquad (a,q)=1.
\tag{5.2}
\]

Both sums run over the full support.
Remove prime powers from the $\psi$-estimate \cite[(2.20)]{Z1}
and apply partial summation. Increasing $B_0(A)$ absorbs the conversion
losses. Thus, for each $A>0$, there is $B_0(A)$ such that
\[
 \sum_{q\le T^{1/2}(\log T)^{-B_0(A)}}
 \max_{(a,q)=1}\left|\pi(T;q,a)-\frac{\operatorname{Li}(T)}{\varphi(q)}\right|
 \ll_A T(\log T)^{-A}.
\tag{5.3}
\]

It is applied with residue $N\bmod q$.
The sieve input is \cite[Theorem~2.5, (2.25)]{Z1}, under these hypotheses:

\begin{center}
\begin{tabular}{@{}ll@{}}
\toprule
Sequence & Finite, with nonnegative multiplicities; $X_{\mathcal A}>0$ \\
Density & $\omega$ multiplicative; $\omega(\ell)/\ell=1/(\ell-1)$ on $\mathcal P_N$ \\
Outside $\mathcal P_N$ & $\omega(\ell)=0$ \\
Parameters & $y\ge z\ge2$; (5.5) with fixed $K\ge2$ \\
\bottomrule
\end{tabular}
\end{center}
Then
\[
 S(\mathcal A,z)\le X_{\mathcal A}V_N(z)
 \{F(s)+e^{\sqrt K}Q_1(s)(\log y)^{-1/3}\}
 +\sum_{d<y,\,d\mid P_N(z)}|r_d|,
 \qquad s=\frac{\log y}{\log z}.
\tag{5.4}
\]

The density obeys $0<\omega(\ell)<\ell$, since $\ell>2$.
Also $F(s)=2e^\gamma/s$ on $[1,3]$, including $s=1$.
Uniformly in $N$,
\[
 \prod_{w\le\ell<z,\,\ell\in\mathcal P_N}
 (1-1/(\ell-1))^{-1}
 \le\frac{\log z}{\log w}\left(1+\frac K{\log w}\right),\qquad 2\le w<z.
\tag{5.5}
\]

By Mertens' theorem \cite{Mertens1874}, the logarithm of the product is
$\log(\log z/\log w)+O(1/\log w)$;
here one uses $1/(\ell-1)=1/\ell+O(\ell^{-2})$ and partial summation.
Removing $\ell\mid N$ decreases the product,
and a larger fixed $K$ covers bounded $w$.
This proves (5.5).
For $z>2$,
\[
 V_N(z)=2\prod_{\ell<z}(1-1/\ell)
 \prod_{2<\ell<z}\left(1-\frac1{(\ell-1)^2}\right)
 \prod_{2<\ell<z,\,\ell\mid N}\frac{\ell-1}{\ell-2}
\tag{5.6}
\]

Hence, uniformly for $z\ge N^c$, fixed $c>0$,
\[
 V_N(z)=\frac{2e^{-\gamma}C(N)}{\log z}(1+o(1)).
\tag{5.7}
\]

The tails are $1+O(1/z)$ and $\exp(O(\omega(N)/z))$,
so $2e^{-\gamma}$ occurs once.
The primitive-character large sieve \cite[(2.19)]{Z1} is
\[
 \sum_{q\le Q}\sum_{\chi\ (q)}^*
 \left|\sum_{h\le H_{\rm LS}}c_h\chi(h)\right|^2
 \ll (Q^2+H_{\rm LS})\|c\|_2^2.
\tag{5.8}
\]

Here $m\sim M$ means $M<m\le2M$.
We use the complete-support convention \cite[(2.15)--(2.17)]{Z1},
power-log Siegel--Walfisz \cite[(2.21)]{Z1},
and PNT with the de la Vallée Poussin error.
Through Siegel--Walfisz, the threshold is ineffective.
The convolution estimate is proved in \S{}5.4.
% Sieve provenance: Iwaniec, Rosser's sieve, Acta Arith. 36 (1980),
% Theorem 1, (1.4); the precise imported form is [Z1, Theorem 2.5].

\subsection{One prime-BV calculation for three rows}

One application of prime BV covers the three negative rows. Put
\[
 u=\frac1{500},\quad \alpha=\frac{100}{1149},\quad
 \beta=\frac{50}{309},\quad t_* =\frac{7109}{28000},\quad
 \tau_S=\frac{57}{224},\quad H=1-2u=\frac{249}{250}.
\]

For a permitted squarefree prefix $b$, define
\begingroup\small
\[
 \mathcal A_b=\{m:bm=N-p_0,\ p_0<X\text{ prime}\},\qquad
 S_b(z)=\#\{m\in\mathcal A_b:(m,P_N(z))=1\}.
\tag{5.9}
\]

\endgroup
Dropping the condition $(m,Nb)=1$ gives an upper bound.
The three negative rows are
\[
 \begin{aligned}
 R_{\rm BV}&=\sum_{b=p_1p_2p_3,\ (r_1,r_2,r_3)\in R_1}S_b(p_3),\\
 T_5&=\sum_{b=p\in\mathcal P_N,\ N^{t_*}\le p<N^{\tau_S}}S_b(N^\alpha),\\
 T_6&=\sum_{b=p\in\mathcal P_N,\ N^{\tau_S}\le p<N^{1/3}}S_b(N^\alpha),
 \end{aligned}
\tag{5.10}
\]

where $r_i=\log p_i/L$ and $r_3<r_2<r_1$.
Put
\[
 \Delta_3=\{u\le r_3<r_2<r_1<\beta\},\qquad
 R_1=\Delta_3\cap\{r_1+2r_2+4r_3>H\}.
\tag{5.11}
\]

All prefix primes belong to $\mathcal P_N$.
The radial condition is $\varepsilon N<bm<N$;
the product $bd$ has no radial cutoff.
For $\kappa\in\{0,10^{-6}\}$, put
\[
 Q=N^{1/2-\kappa}L^{-B},\quad D_b=Q/b,\quad
 \delta_N=B\log L/L,\quad a_\kappa(r)=1/2-\kappa-\sum r_i.
\tag{5.12}
\]

For R1 put $z_b=\sqrt{D_b}$; for T5 and T6 put $z_b=N^\alpha$.
In all three rows, let $s_b=\log D_b/\log z_b$.
The parameter checks are as follows.
\begin{longtable}{@{}p{0.25\textwidth}p{0.70\textwidth}@{}}
\toprule
Row & Uniform positive margins before the vanishing $\delta_N$ loss \\
\midrule
R1 & $a_\kappa\ge3/206-\kappa$; $2r_3-a_\kappa\ge63823/309000+\kappa$; $r_3\ge13147/103000>u$ \\
T5 & $a_\kappa\ge55/224-\kappa$; $t_* -\alpha=5368241/32172000>0$; $a_0/\alpha\in[12639/4480,7917759/2800000]\subset(2,3)$ \\
T6 & $a_\kappa\ge1/6-\kappa$; $\tau_S>\alpha$; eventually $(a_\kappa-\delta_N)/\alpha\in(1,3)$ \\
\bottomrule
\end{longtable}

\begingroup\sloppy
For R1, the linear forms are evaluated at its closure vertices:
$(\beta,\beta,\beta)$,
$(H/7,H/7,H/7)$,
$(\beta,64441/463500,64441/463500)$,
$(\beta,\beta,13147/103000)$.
Thus $z_b<p_3$ and $S_b(p_3)\le S_b(z_b)$, and we apply (5.4).
For $d<D_b$, its primes lie below those of $b$.
Hence $(b,d)=1$, $(bd,2N)=1$, and
\par\endgroup
\[
 |(\mathcal A_b)_d|=\pi_{<}(X;bd,N),\quad
 r_{b,d}=\pi_{<}(X;bd,N)-\frac{\operatorname{Li}(X)}{\varphi(bd)}.
\tag{5.13}
\]

The source mass is $\operatorname{Li}(X)/\varphi(b)$.
For R1, $b$ consists of the three largest primes of $q=bd$.
For T5 or T6, $b$ is the largest prime.
So each row counts each modulus at most once, including $d=1$,
although different rows may share moduli.

For a half-open source cell $C$, put
$\gamma_C(q)=\sum_{b\in C,\,bd=q}\lambda_b^+(d)$.
Then $|\gamma_C(q)|\le1$, $q<Q$, and
$\|\gamma_C\|_1,\|\gamma_C\|_2^2\le Q$.
The absolute remainder obeys the same multiplicity bound.
Choose $A>2$, then $B$. By (5.3),
\[
 \sum_b\sum_{d<D_b}|r_{b,d}|
 \ll_{A,\varepsilon}NL^{-A}+Q=o(U_C).
\tag{5.14}
\]

The strict endpoint costs $Q$.
For $\kappa>0$, one may eventually take $B=0$.
By (5.7) and the formula for $F$ above,
\[
 V_N(z_b)F(s_b)=(1+o(1))\frac{4C(N)}{L(a_\kappa-\delta_N)}.
\tag{5.15}
\]

For R1, $s_b=2$.
PNT and partial summation give
$\sum_{p\in\mathcal P_N}(p-1)^{-1}\delta_{\log p/L}\to dt/t$.
The convergence is uniform away from zero.
Removing $p\mid N$ costs $O(\omega(N)N^{-u})$,
and replacing $1/\varphi(b)$ by $1/b$ costs a relative $O(N^{-u})$.
The domains are bounded Jordan polytopes.
Since $a_\kappa$ stays positive, the kernels converge uniformly.
Thus
\[
 \begin{aligned}
 R_{\rm BV}&\le4(1-\varepsilon)U_C J_{R1}(\kappa)+o(U_C),&
 J_{R1}(\kappa)&=\iiint_{R_1}\frac{dr_1dr_2dr_3}{r_1r_2r_3a_\kappa(r)},\\
 T_j&\le4(1-\varepsilon)U_C J_j(\kappa)+o(U_C),&
 J_j(\kappa)&=\int_{C_j}\frac{dt}{t(1/2-\kappa-t)},
 \end{aligned}
\tag{5.16}
\]

where $C_5=[t_*,\tau_S)$ and $C_6=[\tau_S,1/3)$.
The R1 contribution is $-J_{R1}$ in $D$,
and its raw comparator is $4J_{R1}$.

\subsection{Comparators, endpoints and finite losses}

This subsection evaluates the comparators and collects the finite losses.
Put $\mathrm{IBV}=4J_{R1}$,
$A=1/2-\kappa-r_1-r_2$, and $\ell=(H-r_1-2r_2)/4$.
Then
\[
 \int_\ell^{r_2}\frac{4\,dr_3}{r_1r_2r_3(A-r_3)}
 =\frac4{r_1r_2A}\log\frac{r_2(A-\ell)}{\ell(A-r_2)},
\tag{5.17}
\]

and the outer limits are $H/7\le r_1\le\beta$,
$(H-r_1)/6\le r_2\le r_1$.
For $c_\gamma=2e^{-\gamma}$, the S2 replacement is
\[
 S2_{\rm canonical}
 =\mathrm{prefix}-c_\gamma V_{\rm full}-c_\gamma J_{PF}
       +(c_\gamma V_B-\mathrm{IBV}),
\tag{5.18}
\]

where $V_B$ integrates $1/(r_1r_2r_3^2)$ over R1.
The cutoff in $J_{PF}$ is $g=r_1+2r_2+4r_3\le H$.
Thus R1 is replaced once by $\mathrm{IBV}$,
with coefficient $-1$ in (5.18).

For lower labels $t_i$, put $q_g=1+\tau_g$ and
\[
 \ell_g=t_1+2t_2+(4+2\tau_g)t_3,\quad K=H+2\tau_g\beta,
 \quad F_{\rm lab}=\{\ell_g\le K\},\quad B_{\rm lab}=\{\ell_g>K\}.
\tag{5.19}
\]

The half-open cells partition the source, with
$PF\subset F_{\rm lab}$ and $B_{\rm lab}\subset R1$.
The region between them satisfies $H<g\le q_gK$.
Here $\tau_g=\tau=\sigma^9$; the estimates hold for $\tau_g\le\tau$.
The PF-side cost in $D$ is
$(2.75\cdot10^8+9680)\tau_g=275009680\tau$ (Appendix~A.4).
The PF estimate and the full S2 row are proved in \S{}3.
At $g=H$, $(p_1p_2^2p_3^4)^{250}=N^{249}$ forces $p_1\mid N$,
so such a point is absent from the permitted source.
Parent-grid faces follow the half-open convention.
For S5 and S6,
\[
 4J_5(0)=8\left[\log\frac{t}{1/2-t}\right]_{t_*}^{\tau_S},\qquad
 4J_6(0)=8\log\frac{110}{57}.
\tag{5.20}
\]

The disjoint domains are $[1/4,t_*)$, $[t_*,\tau_S)$,
and $[\tau_S,1/3)$.
The first is treated directly, the other two by BV.
A prime equal to $N^{t_*}$, $N^{\tau_S}$ or $N^{1/3}$ divides $N$,
so permitted primes avoid these rational endpoints.
The S5-top raw value is $4J_5(0)=0.0365816642384326$.
Its row coefficient is $-1$, so its $D$ entry is
\[
 -J_5(0)=-\frac{0.0365816642384326\ldots}{4}
        =-0.00914541605960815\ldots,
\tag{5.20a}
\]

before the level loss.
It is obtained by dividing (5.15) times $\operatorname{Li}(X)/\varphi(p)$
by $4(1-\varepsilon)U_C$; see Appendix~A.6.
The sources are specified in Sections~2--4,
and the errors are summed in Section~6.3.
For finite $N$, the denominator is $a_\kappa-\delta_N$.
If $\delta_N\le a_{\min}/2$, then
\[
 0\le\frac1{a_\kappa-\delta_N}-\frac1{a_\kappa}
 \le\frac1{a_\kappa}\frac{\delta_N}{a_{\min}-\delta_N}.
\tag{5.21}
\]

Let $E_{N6}(N)\ge0$ collect the adverse vanishing errors:
sieve $O(L^{-1/3})$, (5.14), endpoints $O(Q)$,
and Euler/PNT conversion $o(1)$.
By (5.21), it tends to zero uniformly.
Its logarithmic part obeys $E_{\log}\le18.03\delta_N$ in $D$.
By (5.21), (5.17), (A.13) and (A.14), the R1, S5-top and S6 contributions to $E_{\log}/\delta_N$ are respectively $<2.169459$, $<0.074494$ and $<15.778471$, using row minima $3/206-10^{-6}$, $55/224$ and $1/6-10^{-6}$; hence their sum is $18.022424<18.03$.
We use $\kappa=10^{-6}$ for R1 and S6,
and $\kappa=0$ for S5-top, whose positive-$\kappa$ increment is zero.
For S6,
\[
 E_{S6}=J_6(10^{-6})-J_6(0)
 \le 0.000006484294227956901
\tag{5.22}
\]

in $D$ units.
S5 keeps its vanishing logarithmic loss.
Appendix~A.6 gives the outward enclosures.

\subsection{A common complete-support convolution lemma}

Both applications below need a mean-value estimate of
Bombieri--Vinogradov type \cite{Bombieri1965,Vinogradov1965}
for a convolution with complete supports. The form used here is stated
and proved in this subsection.

\begin{lemma}[Bilinear mean value for complete supports]
\label{lem:complete-support-bv}
Put $MN_{\rm box}=x$, $\mathcal L=\log x$.
Fix $\eta,B_{\rm div}>0$ and assume
$x^\eta\le M\le x^{1-\eta}$ and
$x^\eta\le N_{\rm box}\le x^{1-\eta}$.
The supports are $(M,2M]$ and $(N_{\rm box},2N_{\rm box}]$.
For every $A_2>0$, assume
\[
 \left|\sum_{n\equiv l(k),(n,d)=1}b_n
 -\frac1{\varphi(k)}\sum_{(n,dk)=1}b_n\right|
 \ll_{A_2}\|b\|_2N_{\rm box}^{1/2}\tau(d)^{B_{\rm div}}
       (\log(2N_{\rm box}))^{-A_2},
\tag{5.23}
\]

uniformly in $x,d,k,l$ and the coefficient family,
for $d,k\ge1$ and $(l,k)=1$.
Then, for every $A>0$, there is $B_1$ such that
\[
 \sum_{q\le Q}\max_{(r,q)=1}|\Delta_{a*b}(q,r)|
 \ll\|a\|_2\|b\|_2x^{1/2}\mathcal L^{-A},\qquad
 Q\le x^{1/2}\mathcal L^{-B_1}.
\tag{5.24}
\]

Here $\Delta$ is as in (5.2).
The constant $B_1$ depends on $A,\eta,B_{\rm div}$
and the hypothesis constants.
\end{lemma}

\begin{proof}
Put $A(\chi)=\sum_m a_m\chi(m)$ and
$B(\chi)=\sum_n b_n\chi(n)$.
Orthogonality gives
\[
 \Delta_{a*b}(q,r)=\frac1{\varphi(q)}
 \sum_{\chi\ne\chi_0}\overline{\chi(r)}A(\chi)B(\chi).
\tag{5.25}
\]

Write $q=ef$, with primitive nonprincipal $\psi\pmod f$ inducing $\chi$.
Since $\varphi(ef)\ge\varphi(e)\varphi(f)$, the left side of (5.24) is at most
\[
 \sum_{e\le Q}\frac1{\varphi(e)}
 \sum_{1<f\le Q/e}\frac1{\varphi(f)}
 \sum_{\psi\ (f)}^*|A_e(\psi)B_e(\psi)|,
\tag{5.26}
\]

where $A_e(\psi)=\sum_{(m,e)=1}a_m\psi(m)$,
and $B_e$ is defined likewise.
Split the conductors at $F=\mathcal L^C$.
For $f\le F$, subtract the principal average on each reduced class.
By (5.23),
\[
 |B_e(\psi)|\ll\varphi(f)\|b\|_2N_{\rm box}^{1/2}
 \tau(e)^{B_{\rm div}}(\log(2N_{\rm box}))^{-A_2}.
\]

Cauchy--Schwarz gives
$\sum_\psi^*|A_e(\psi)|\ll\varphi(f)M^{1/2}\|a\|_2$.
Thus small conductors cost at most
$F^2\tau(e)^{B_{\rm div}}\|a\|_2\|b\|_2x^{1/2}(\log(2N_{\rm box}))^{-A_2}$.
For $V<f\le2V$, use (5.8) and Cauchy--Schwarz:
\begingroup\small
\[
 \sum_{V<f\le2V}\frac1{\varphi(f)}\sum_\psi^*|A_eB_e|
 \ll (\log\log x)
 \left(V+\sqrt M+\sqrt{N_{\rm box}}+\frac{\sqrt x}{V}\right)\|a\|_2\|b\|_2.
\tag{5.27}
\]

\endgroup
There are $O(\mathcal L)$ conductor blocks,
and $\sum_{e\le Q}1/\varphi(e)\ll\mathcal L$.
The divisor-weighted sum obeys
$\sum_{e\le Q}\tau(e)^{B_{\rm div}}/\varphi(e)\allowbreak
\ll\mathcal L^{2^{B_{\rm div}}+1}$.
After division by $\|a\|_2\|b\|_2x^{1/2}$, the bound is
\begingroup\small
\[
 \ll \eta^{-A_2}\mathcal L^{2C+2^{B_{\rm div}}+1-A_2}
 +\mathcal L^{3-B_1}+\mathcal L^3x^{-\eta/2}+\mathcal L^{3-C}.
\tag{5.28}
\]

\endgroup
Take $C=A+3$, $B_1\ge A+3$,
$A_2\ge A+2C+2^{B_{\rm div}}+1$.
This proves (5.24) on the full supports.
\end{proof}

The next lemma verifies that hypothesis for primes in short intervals.

\begin{lemma}[Prime interval coefficients]
Fix $H_{\rm int},c,C>0$.
Let $N^c\le Y\le N^C$ and $N_{\rm box}=Y/2$.
If $J\subset[Y/2,Y]$ has length $|J|\asymp YL^{-H_{\rm int}}$,
then $b=\mathbf1_{\mathbb P\cap J}$ satisfies (5.23)
with $B_{\rm div}=1$, uniformly in $J$.
\end{lemma}

\begin{proof}
Endpoint PNT gives
\[
\|b\|_2^2\asymp YL^{-H_{\rm int}-1}.
\]

Fix $A_2$.
Put $s=A_2+(H_{\rm int}+1)/2$ and $R=YL^{-s}\to\infty$.
The target bound is then comparable to $R\tau(d)$.
For $k\le L^{C_0}$, use endpoint power-log Siegel--Walfisz.
After removing the primes dividing $d$, it gives
\[
|\Delta|\ll_{B,C_0}YL^{-B}+\omega(k)+\tau(d).
\]

Take $B>s+3$; then the first term is $O(R)$.
Also $\omega(k)\le\log k/\log2=O_{C_0}(\log L)=o(R)$.
Since $\tau(d)\ge1$, the bound is $O(R\tau(d))$, uniformly in $d$.
For $k>L^{C_0}$, positivity gives
\[
|\Delta|\le |J|/k+1+\pi(J)/\varphi(k).
\]

The function $\log\log(3t)/t$ eventually decreases.
Since $1/\varphi(k)\ll\log\log(3k)/k$,
\[
\sup_{k>L^{C_0}}1/\varphi(k)\ll L^{-C_0}\log\log(3L^{C_0}).
\]

Using $|J|\ll YL^{-H_{\rm int}}$ and
$\pi(J)\ll YL^{-H_{\rm int}-1}$ in the displayed bound above, we get
\[
\begin{aligned}
 |J|/k+1+\pi(J)/\varphi(k)
 &\ll YL^{-H_{\rm int}-C_0}+1\\
 &\quad+YL^{-H_{\rm int}-1-C_0}\log\log(3L^{C_0})
 \ll R.
\end{aligned}
\]
Here $R=YL^{-s}\to\infty$, and $C_0>s+3$.
Hence (5.23) holds with $B_{\rm div}=1$.
\end{proof}

In (5.26), one takes $d=e$, $k=f$, $ef\le Q$,
using (5.23) for $f\le F$ and (5.8) otherwise.
The two applications use these parameters:
\begin{center}
\begin{tabular}{@{}llll@{}}
\toprule
Source & $H_{\rm int}$ & $x_b$ & $\eta$ \\
\midrule
O1 & $2$ & $[\varepsilon N/8,2N]$ & $0.08$ \\
G4 & $H_g=A+4$ & $\asymp_\varepsilon N$ & $1/4$ \\
\bottomrule
\end{tabular}
\end{center}
The G4 partition has equal widths, including its last interval.

\subsection{O1: signs, full convolution and mass}

This subsection proves the upper bound for the switched excess.
We use the finite source and the thirteen-row weight of \S{}1.
Put
\[
 S_{\rm sw}=\sum_{n\in\mathcal A_\varepsilon,\ \Omega(n)\ge3,\ P^-(n)\ge N^\alpha}
                   \max(W_{13}(n),0).
\tag{5.29}
\]

Let $N_j(n)$ count the permitted source pairs, $j=11,12,13$,
with their exact indicators from \S{}4.
Let $k_i$ denote the four bin counts defined below.
Set $m_B(v)=\#\{u_i:u_i<B_{O1}(v)\}$.
For squarefree $(n,N)=1$, the pointwise identity is
\begingroup\small
\[
\begin{aligned}
 W_{13}(n)&=R(n)-2N_{11}(n)-N_{12}(n)-N_{13}(n),\\
 R(n)&=3+\mathbf1_{k_1=0}-2k_1-2k_2-k_3\\
      &\quad+\binom{k_1}{2}+\sum_jm_B(v_j).
\end{aligned}
\tag{5.29a}
\]

\endgroup
For $k\ge5$, all three negative source counts vanish:
\begin{center}
\begin{tabular}{@{}lp{0.70\textwidth}@{}}
\toprule
Row & Obstruction \\
\midrule
S11 & Five factors $\ge N^\lambda$; $5\lambda>1$. \\
S12 & One factor $\ge N^\alpha$ and four $\ge N^{1/3}$. \\
S13 & Three remaining factors would have product $\ge(N/p_ap_b)^{3/2}>N/p_ap_b$, contradicting $e<N$. \\
\bottomrule
\end{tabular}
\end{center}
Thus $W_{13}=R$.
For $k=3,4$, put $t_1<\cdots<t_k$, $s=\sum t_i$.
Take $N$ such that $0<1-s<10^{-3}$.
The following conditions guarantee pairs.
\begin{longtable}{@{}p{0.25\textwidth}p{0.70\textwidth}@{}}
\toprule
Condition & Guaranteed count \\
\midrule
$G11:\ k_1=k_2=0,\ k_3\ge1$ & $N_{11}\ge1$ \\
$G12:\ k_4=k-1,\ k_1+k_2+k_3=1$ & $N_{12}\ge1$ \\
$G13:\ k=3,\ k_1=0,\ k_2=1$ & $N_{13}\ge1$ \\
\bottomrule
\end{longtable}

In each case the pair consists of the two smallest prime factors.
Its upper pair constraint follows from
$1-t_1>s-t_1=t_2+\cdots+t_k>2t_2$.
All remaining factors exceed the second prime.
For G13, $t_3>1-s$ gives the square-root cutoff.
Extra S13 pairs, including radial endpoints, carry a negative coefficient.
Put
\begingroup\small
\[
 U(k)=3+\mathbf1_{k_1=0}-2k_1-2k_2-k_3
          +\binom{k_1}{2}+k_1k_2,\qquad
 U^*(k)=U(k)-2\mathbf1_{G11}-\mathbf1_{G12}-\mathbf1_{G13}.
\tag{5.29b}
\]

\endgroup
Since $m_B(v)\le k_1$, $W_{13}(n)\le U^*(k)$.
Apply (4.9)--(4.12) on $\varepsilon N<n<N$,
and discard vectors with $\sum k_i\inf I_i\ge1$
or $\sum k_i\sup I_i\le1-10^{-3}$.
The remaining ten vectors for $k=3$ have $U^*\le0$,
and so do the nineteen for $k=4$.
Hence their positive weight vanishes.

For the remaining cells put $\mu=25/128$, $\lambda=93/400$, and
\[
 C_0=\frac{57}{64}-u-\frac{11}{5}\alpha=\frac{6408241}{9192000},
 \quad v_\beta=\frac{2(C_0-\beta)}5,
\]

\[
 B_{O1}(v)=\begin{cases}\beta,&\beta\le v\le v_\beta,\\
 C_0-5v/2,&v_\beta<v<\lambda.
 \end{cases}
\tag{5.30}
\]

Here $2(C_0-\alpha)/5>\lambda$, so the lower clamp is inactive.
For $5\le k\le11$, the coordinate bins are
$I_1=[\alpha,\beta)$, $I_2=[\beta,\lambda)$,
$I_3=[\lambda,1/3)$, $I_4=[1/3,1]$.
A cell is specified by the following data:
\begin{center}
\begin{tabular}{@{}lp{0.70\textwidth}@{}}
\toprule
Counts & $k_1+\cdots+k_4=k$ \\
Order & $u_1<\cdots<u_{k_1}<v_1<\cdots<v_{k_2}<w_1<\cdots<w_{k_3}<z_1<\cdots<z_{k_4}$ \\
Ranks & $k_1\ge m_1\ge\cdots\ge m_{k_2}\ge0$ \\
Flat branch & $v_j\le v_\beta$, $m_j=k_1$ \\
Slope branch & $v_j>v_\beta$, $u_{m_j}<C_0-5v_j/2\le u_{m_j+1}$ \\
\bottomrule
\end{tabular}
\end{center}
The last line uses infinite outer sentinels.
Each nonempty cell is convex,
and each of its inequalities involves at most two coordinates.
Its weight is
\[
 w_C=3+\mathbf1_{k_1=0}-2k_1-2k_2-k_3+\binom{k_1}{2}+\sum_jm_j.
\tag{5.31}
\]

Let $\mathcal C$ be the set of cells with $w_C\ge1$;
their weights are at most 36.
They partition the ordered prime vectors, with the stated face conventions.
Define
\begingroup\small
\[
\begin{aligned}
 C_N&=C\cap\{1-\log(1/\varepsilon)/L<\sum t_i<1\},\\
 V_C(s)&=\int_{C\cap\{\sum t_i=s\}}\frac{dt_1\cdots dt_{k-1}}{t_1\cdots t_k},\\
 I_{HL}&=\sum_Cw_CV_C(1).
\end{aligned}
\tag{5.32}
\]

\endgroup
The radial interval is exactly $\varepsilon N<e<N$.
Put $g(e)=\sum_Cw_Ca_C(e)$,
where $a_C$ selects permitted squarefree factorizations in $C_N$.
Squareful or non-coprime rough exceptions number $O(N^{1-\alpha}L)$.
Each has at most eleven prime factors.
So each one-prime row is at most 11,
and each two-prime row is at most $\binom{11}{2}$.
Thus $|W_{13}|\le455$, and
\[
 S_{\rm sw}\le\sum_e g(e)\mathbf1_{\mathbb P}(N-e)+O(N^{1-\alpha}L).
\tag{5.33}
\]

\emph{Full-support majorant.}
Put $\Delta_0=1+L^{-2}$, $h=\log\Delta_0/L\le L^{-3}$,
and $J_j=[N^\alpha\Delta_0^{j-1},N^\alpha\Delta_0^j)$.
Let $T_j$ be its exponent interval.
A box $\mathbf j=(j_1,\ldots,j_k)$ is outer for $C$ if
$\overline{T_{j_1}\times\cdots\times T_{j_k}}\cap\overline{C_N}\ne\varnothing$.
Let $\mathbf1_{\mathbf j}(e)$ count labelled prime factorizations
with factor $i$ in $J_{j_i}$.
Put
\[
 c(\mathbf j)=\sum_{C:\,\mathbf j\text{ outer for }C}w_C,\qquad
 f^+(e)=\sum_{k=5}^{11}\sum_{\mathbf j}c(\mathbf j)\mathbf1_{\mathbf j}(e).
\tag{5.34}
\]

Then $g\le f^+$.
For $J=J_j$, put
\[
 b_J=\mathbf1_{\mathbb P\cap J},\qquad
 a_J(m)=\sum_{k=5}^{11}\sum_{\mathbf j:\,j_1=j}c(\mathbf j)
  (\mathbf1_{\mathbb P\cap J_{j_2}}*\cdots*
   \mathbf1_{\mathbb P\cap J_{j_k}})(m).
\tag{5.35}
\]

Thus $f^+=\sum_Ja_J*b_J$.
For fixed $m$, $k=\Omega(m)+1$,
and each labelled factorization fixes one interval tuple.
There are at most $(k-1)!$ such factorizations, so
\[
 0\le a_J(m)\le K_{\rm box}:=\max_{5\le k\le11}(k-1)!\sum_{C\in\mathcal C_k}w_C,
\tag{5.36}
\]

uniformly in $N,J$.
Every original cell constraint remains, and outer boxes satisfy
$\sum t_i\in[1-\log(1/\varepsilon)/L-kh,1+kh]$,
$t_1\in[\alpha,1/5+h]$.
For $J=[Y/\Delta_0,Y)$, split $a_J$ dyadically at $M$;
then, eventually,
\[
 x_b=M(Y/2)\in[\varepsilon N/8,2N],\quad
 \eta=0.08\le\frac{\log(Y/2)}{\log x_b}\le1-\eta,
\tag{5.37}
\]

and there are $O_\varepsilon(L^3)$ blocks.
Their norms satisfy
$\|a_{J,M}\|_2\ll K_{\rm box}\sqrt M$,
$\|b_J\|_2\ll\sqrt YL^{-3/2}$.
The interval lemma gives (5.23) with $H_{\rm int}=2$.
Use saving 6 and a single level $D=N^{1/2}L^{-B}$,
$B\ge B_1(6)+1$.
Each block then costs $O(NL^{-15/2})$, giving
\[
 \sum_{d\le D}\max_{(r,d)=1}|\Delta_{f^+}(d,r)|
 \ll_\varepsilon NL^{-9/2}.
\tag{5.38}
\]

Each cofactor has at least four rough prime factors,
and the sieve divisors remain outside the convolution.
All block scales are comparable to $N$.

\emph{Sieve and mass.}
Sieve $\{N-e\}$ with multiplicity $f^+(e)$ at $z=\sqrt D$,
and put $X^+=\sum f^+$.
For $d\mid P_N(z)$,
\[
 r_d=\Delta_{f^+}(d,N)-\frac1{\varphi(d)}\sum_{(e,d)>1}f^+(e).
\tag{5.39}
\]

Every $e$ has at most eleven prime factors, all $\ge N^\alpha$.
Put
\[
 R_{\rm nc}:=\sum_{\substack{d<D\\d\mid P_N(z)}}
 \frac1{\varphi(d)}\sum_{(e,d)>1}f^+(e).
\]

For squarefree $d=\ell r$ with prime $\ell\mid(e,d)$,
$(\ell,r)=1$ and $\varphi(d)=(\ell-1)\varphi(r)$.
Hence
\[
 \begin{aligned}
 R_{\rm nc}
 &\le\sum_e f^+(e)\sum_{\ell\mid e}\frac1{\ell-1}
       \sum_{r<D/\ell}\frac{\mu^2(r)}{\varphi(r)}\\
 &\le22X^+N^{-\alpha}\sum_{r<D}\frac{\mu^2(r)}{\varphi(r)}
 \ll X^+N^{-\alpha}L.
 \end{aligned}
\tag{5.40}
\]

Here $1/(\ell-1)\le2N^{-\alpha}$ eventually.
By (5.42), the normalized cost is $O(N^{-\alpha}L^2)$;
it is included in $E_{S_{\rm sw}}$.
By (5.33), source exceptions cost $O(N^{1-\alpha}L)$ before normalization.
Since $U_C\gg N/L^2$, their normalized cost is $O(N^{-\alpha}L^3)$.
Primes $N-e<z$ in $g$ cost at most $36z$.
Outer-box multiplicities remain nonnegative for nonpositive $N-e$.
Apply (5.4) to the finite multiset.
Each box has mass
\[
\prod_i\pi(J_{j_i}).
\]

Uniform PNT approximates it by
\[
\int_{T_{\mathbf j}}N^{\sum t_i}\prod_i dt_i/t_i
\]

with exponentially small relative error.
The excess of the outer boxes lies within $\sqrt k h$ of the facets.
Radial facets cost $O(hN)$.
Every other facet has at most two nonzero normal coordinates,
and its normal is not parallel to $(1,\ldots,1)$.
Each radial slice has excess $O(h)$, costing $O(hN/L)$.
Thus the total weighted excess is $O(NL^{-3})$.

The slice mass $V_C(s)$ is Lipschitz near 1.
All coordinates exceed $\alpha$, so denominators and derivatives are bounded.
Changing $s$ shifts each two-coordinate inequality by $O(|1-s|)$.
If a constraint uses the eliminated coordinate, its gradient has at least $k-2\ge3$ identical nonzero entries.
Thus the slice change is $O(|1-s|)$.
Radial integration gives
\[
 \int_{C_N}N^{\sum t_i}\prod_i\frac{dt_i}{t_i}
 =\frac NL\left[(1-\varepsilon)V_C(1)
 +\int_0^{\log(1/\varepsilon)}e^{-v}
       \{V_C(1-v/L)-V_C(1)\}\,dv\right].
\tag{5.41}
\]

Hence
\[
 X^+=\frac NL\{(1-\varepsilon)I_{HL}+O_\varepsilon(L^{-1})\}.
\tag{5.42}
\]

Thus (5.40) is $O(N^{1-\alpha})=o(U_C)$.
For $s=2$, $F(2)=e^\gamma$ and
$\log z=L/4-(B/2)\log L$.
By (5.7),
\[
 V_N(z)F(2)=\frac{8C(N)}L\{1+O(\log L/L)+o(1)\}.
\]

Equations (5.33), (5.38)--(5.42) and the sieve give
\[
 \boxed{\ S_{\rm sw}\le[8(1-\varepsilon)I_{HL}+o(1)]U_C\ }.
\tag{5.43}
\]

The errors are $O(U_CL^{-1/3})$, $O(N^{1/4})$,
$O_\varepsilon(NL^{-9/2})$ and $O(N^{1-\alpha}L)$,
including the radial and box errors.
Each is $o(U_C)$.
Appendix~A.6 encloses (5.32).
The low-$\Omega$ cells are handled by (4.9)--(4.12) and (5.29a)--(5.29b).

\subsection{Weak G4: interval partition and injection}

This subsection bounds the remaining balanced semiprimes.
Fix $a$ with $3/2<a<2$ and
$\frac{a-1}{a}-\varepsilon>\frac13$. Put
\[
 b_{G4}=\frac{a-1}{a}-\varepsilon>\frac13.
\]

For $a=9/5$, this is the cutoff in (1.3).
Count ordered prime pairs by
\[
 g_N(e)=\#\{(m,p'):m,p'\text{ prime},\ N^{b_{G4}}\le m\le\sqrt N,
                      \ \varepsilon N\le mp'<N,\ mp'=e\}.
\tag{5.44}
\]

Unique factorization gives $g_N\le2$;
both orders and all coprimality types are included.
For $A>2$, put $H_g=A+4$ and
\[
 R=\left\lceil\frac{(1/2-b_{G4})L}{\log(1+L^{-H_g})}\right\rceil,
 \quad \xi=N^{(1/2-b_{G4})/R},\quad u_j=N^{b_{G4}}\xi^j.
\tag{5.45}
\]

Partition into the intervals $[u_{j-1},u_j)$, with the last upper endpoint closed.
Each relative width is comparable to $L^{-H_g}$,
and there are $O(L^{H_g+1})$ intervals.
For an interval with upper endpoint $Y$, take $b_J=\mathbf1_{\mathbb P\cap J}$ and
\[
 a_J^+(p')=\mathbf1_{p'\text{ prime},\ \varepsilon N/Y\le p'<N\xi/Y},\qquad
 a_J^-(p')=\mathbf1_{p'\text{ prime},\ \varepsilon N\xi/Y\le p'<N/Y}.
\tag{5.46}
\]

Then $g^-\le g_N\le g^+$, with $g^\pm=\sum_Jb_J*a_J^\pm$.
For $h=g^+-g_N$, $0\le h\le2$.
Its support lies in
$[\varepsilon N/\xi,\varepsilon N)\cup[N,N\xi)$,
a set of length $W\ll NL^{-H_g}$.
PNT or Brun--Titchmarsh gives $\sum h\ll NL^{-H_g-1}$.
Thus
\[
 \sum_{q\le Q}\max_{(r,q)=1}|\Delta_h(q,r)|
 \ll\sum_{q\le Q}\left(2(W/q+2)+(\sum h)/\varphi(q)\right)
 \ll NL^{-H_g+1}+Q.
\tag{5.47}
\]

Write $\Delta_{g_N}=\Delta_{g^+}-\Delta_h$, and split $a_J^+$ into $O_\varepsilon(1)$ dyadic pieces $(M_i,2M_i]$.
Here $x_b=M_iY/2\asymp_\varepsilon N$ and
$\log(Y/2)/\log x_b\in[b_{G4}-o(1),1/2+o(1)]$.
Thus $\eta=1/4$ is admissible.
The norms are $\|a_{J,i}\|_2\ll\sqrt{M_i}$ and
$\|b_J\|_2\ll\sqrt YL^{-(H_g+1)/2}$.
The interval lemma gives (5.23) for every exact interval.
With $A_{\rm block}=A+H_g/2+2$ in (5.24),
each block costs $O(NL^{-A_{\rm block}-(H_g+1)/2})$,
and their sum is $O(NL^{-A-3/2})$.
A single level $Q=N^{1/2}L^{-B}$ works throughout, with
$B\ge B_1(A_{\rm block},1/4,1)+1$.
Together with (5.47),
\[
 \sum_{q\le Q}\max_{(r,q)=1}|\Delta_{g_N}(q,r)|\ll_{A,a,\varepsilon}NL^{-A}.
\tag{5.48}
\]

For example, one may take $A=4$, $H_g=8$, $A_{\rm block}=10$.
Apply this to
\[
 G_3(b_{G4})=\sum_{m\in\mathcal P_N,\ N^{b_{G4}}\le m\le\sqrt N}
            S(\mathcal A_m;\mathcal P_N,m),
\tag{5.49}
\]

with the conditions $P^-(c)\ge m$, $(c,Nm)=1$.
A counted term satisfies $N-p_0=mc$, $p_0<X$.
Eventually,
\[
 1<\varepsilon\sqrt N<c<N/m\le N^{1-b_{G4}}<N^{2b_{G4}}\le m^2.
\]

A composite $c$ would satisfy $c\ge m^2$, so $c$ is prime.
The designated $m$ gives an injection into (5.44),
with multiplicity one in all ordering and coprimality cases.
The closed lower endpoint covers $e>\varepsilon N$,
and the square and endpoint conventions cost $O(\sqrt N)$.

Sieve with multiplicity $g_N(e)$ at $D_s=Q$, $z=\sqrt Q$, $s=2$.
Primes $p_0<z$ cost at most $2z$. Put $X_g=\sum g_N$.
For $d\mid P_N(z)$, $N$ is reduced and
\[
 |\mathcal B_d|=X_g/\varphi(d)+r_d,\qquad
 r_d=\Delta_{g_N}(d,N)-\frac1{\varphi(d)}\sum_{(e,d)>1}g_N(e).
\tag{5.50}
\]

Eventually both prime factors exceed
$\min(N^{b_{G4}},\varepsilon\sqrt N)>N^{1/3}>z$,
so the non-coprime term vanishes.
For comparison, at $y=z=Q$ one has $s=1$, $F(1)=2e^\gamma$
and $\log Q=L/2-B\log L$.
Thus (5.7) gives
\[
 V_N(Q)F(1)=\frac{4C(N)}{\log Q}(1+o(1))
 =\frac{8C(N)}L(1+o(1)).
\]
Since $Q_1(1)$ is bounded and $K$ is fixed, the correction in (5.4) is
\[
 V_N(Q)e^{\sqrt K}Q_1(1)(\log Q)^{-1/3}
 =O_K(C(N)L^{-4/3})=o(C(N)/L).
\]
Hence the sieve factor has the same leading constant.
The G4 bound below uses $z=\sqrt Q$ as above.
By (5.48) and (5.4), the remainder is $o(U_C)$.
Uniform PNT gives
\[
 X_g=\sum_m\{\pi(N/m)-\pi(\varepsilon N/m)\}+o(N/L)
 =\frac{(1-\varepsilon)N}{L}
     \int_{b_{G4}}^{1/2}\frac{dt}{t(1-t)}+o(N/L).
\tag{5.51}
\]

Endpoints cost $O(\sqrt N)$.
PNT is applied to the inner count, then partial summation to the outer sum.
Since $\log(N/m)=(1-t)L$, the integral is
$\log((1-b_{G4})/b_{G4})$.
Also $V_N(\sqrt Q)F(2)=(8+o(1))C(N)/L$.
Hence
\[
 \boxed{\ \sum_{p_0\text{ prime}}g_N(N-p_0)\le
 \left[8(1-\varepsilon)\log\frac{1-b_{G4}}{b_{G4}}+o(1)\right]U_C\ }.
\tag{5.52}
\]

The injection gives
$G_3(b_{G4})\le\sum_{p_0\text{ prime}}g_N(N-p_0)+O(\sqrt N)$,
which proves weak G4.
Sections~2--4 give G1 and the nineteen OL estimates.

\subsection{One signed allocation}

Theorem~\ref{thm:signed-assembly} and \S{}6.3 give (1.21),
with $\Sigma_W=\Sigma$ and $r_\Sigma(N)=o(1)$.

For general $a$, the pointwise bound and (5.52) give
\[
 D_{(1,a)}(N)\ge
 \left[(1-\varepsilon)\left\{D_{\rm cert}
           -8\log\frac{1-b_{G4}}{b_{G4}}\right\}-o(1)\right]U_C.
\tag{5.54}
\]

\begingroup\sloppy
Cancel the common positive factor $1-\varepsilon$.
For $a=9/5$, $\varepsilon=1/10000$, $b_{G4}=39991/90000$,
the criterion is then (1.22):
$D_{\rm cert}>8\log(50009/39991)=1.788388483459608\ldots$.
The raw bad-count coefficient is
$8(9999/10000)\log(50009/39991)$, that is, $1.788209644611262\ldots$.
Its limit is $8\log(5/4)=1.785148410513678\ldots$.
The fixed-$\varepsilon$ construction here uses (1.22).
For fixed $a\in(3/2,2)$, put $\rho_0=(a-1)/a$ and
$r_c=1/(1+\exp(D_{\rm cert}/8))$.
The same source argument works if
\par\endgroup
\[
 \varepsilon=10^{-4},\qquad
 \rho_0-10^{-4}>\max\{1/3,r_c\}.
\tag{5.56}
\]

Then $b_{G4}=\rho_0-10^{-4}>1/3$ and
$D_{\rm cert}>8\log((1-b_{G4})/b_{G4})$.
With $\varepsilon$ kept fixed, require $N>\varepsilon^{-(a-1)/(a\varepsilon)}$
and the analytic threshold.
The signed allocation is as follows.
\begin{longtable}{@{}p{0.25\textwidth}p{0.70\textwidth}@{}}
\toprule
Term & Position and multiplicity \\
\midrule
Full R1 comparator $4J_{R1}$ & One replacement in (5.18). Appendix~A.4 includes the R1-$\kappa$ comparison in S2. The extra allowance in (6.6) is conservative. \\
PF/BV selected cells & The PF-side loss of (5.19), once. \\
S5 BV top & One replacement on $[t_*,\tau_S)$ in the split row. \\
S6 level increment & $E_{S6}$ of (5.22), once. \\
Finite BV errors & $E_{N6}(N)\ge0$, including the logarithmic loss; $E_{N6}(N)\to0$. \\
AA37 perturbation & $C_E\sigma W_{\rm AA}/4\le g_{\rm AA}$ by (1.27), once in (6.6). Cutoff and selection costs are in (6.4). \\
O1 main term & $-S_{\rm sw}/[4(1-\varepsilon)U_C]\ge-2I_{HL}-o(1)$, inside $D$. Its enclosure enters once. \\
O1 errors & $E_{S_{\rm sw}}=o(1)$ from (5.43): box, radial, exceptional, small-prime, non-coprime and BV. The non-coprime cost is $O(N^{-\alpha}L^2)$; finite-source exceptions cost $O(N^{-\alpha}L^3)$. \\
Weak G4 & $G_3(b_{G4})$ is subtracted once outside $D$ in (5.54). \\
\bottomrule
\end{longtable}

Sections~2--4 and~6.3 prove the source and profile bounds.
Each fixed cost enters once,
and direct low-prime removal costs $O_u(L^{-28})$ after normalization.

\section{Proof of Theorem~\ref{thm:goldbach}}\label{sec:proof-theorem-1}

\subsection{One signed assembly}

Put $L=\log N$, $\rho=1+\tau$ and
$W_{\rm norm}=4(1-\varepsilon_A)U_C$.
This scalar differs from the pointwise weight.
The signed source is
\[
\begin{split}
W_{\rm norm}D_{\rm src}(N)={}&3S_1+S_2-2S_3-S_4-S_5-S_6-S_7\\
&+S_8+S_9+S_{10}-2S_{11}-S_{12}-S_{13}-S_{\rm sw}.
\end{split}\tag{6.1}
\]

\begin{theorem}[Signed assembly]\label{thm:signed-assembly}
Fix one actual construction with
$u=1/500$, $\alpha=100/1149$, $\beta=50/309$,
$0<\sigma\le10^{-10}$, $\tau=\sigma^9$ and
$\ell_0=53/412-1/20000>0$.
Assume four conditions.
\begin{enumerate}
\item The nineteen row chains and outer chains cover (6.1).
They retain exact source maps, multiplicities, half-open domains,
selectors and restrictions.
\item The source, AA37, comparator, detection, boundary, Class-I,
Class-II, measure and profile estimates below hold.
They retain actual domains, signed modulus sums and norms.
Uniformity holds after the displayed family count.
\item Directed endpoints bound their specified profiles.
S1, S2-T8/4 and G1 enter jointly, once each.
The S1-3 prefactor comparison is separate from its dynamic cutoff.
\item Each fixed loss enters once under the paired tail convention.
The AA37 perturbation fits $g_{\rm AA}$.
The explicit selector, cutoff and peel bounds remain.
Parameters follow the order below.
\end{enumerate}
There are finite signed terms $a_j(N)$ and slacks $s_k(N)\ge0$ with
\[
\begin{gathered}
D_{\rm src}(N)=B_{\rm main}+\sum_j a_j(N)+\sum_k s_k(N),\\
D_{\rm src}(N)\ge B_{\rm main}-F_{\rm pers}-E_{\rm van}(N),
\qquad E_{\rm van}(N)\longrightarrow0.
\end{gathered}\tag{6.2}
\]
Hence the fixed-$\epsilon_A$ threshold holds eventually when
\[
D_{\rm lower}\le B_{\rm main}-F_{\rm pers},\qquad
D_{\rm lower}>T.\tag{6.3}
\]
\end{theorem}

\begin{proof}
The signed row coefficient is included in $Q_{r,k}$.
All quantities below are in D units.
\begin{longtable}{@{}p{0.13\textwidth}p{0.82\textwidth}@{}}
\toprule
Stage & Quantity \\\midrule
$Q_{r,0}$ & Actual signed master source. \\
$Q_{r,1}$ & Declared row source after its exact identity. \\
$Q_{r,2}$ & Low/upper cuts, source tags, squarefree/coprime restoration. \\
$Q_{r,3}$ & Exact selector, dynamic-cutoff and peel partition, with main-comparison error terms. \\
$Q_{r,4}$ & Actual-level AA37 main term, fixed/asymptotic perturbations and remainder. \\
$Q_{r,5}$ & Exact comparator decomposition. \\
$Q_{r,6}$ & $\Lambda/\log$ detection, proper powers and LOG freezing. \\
$Q_{r,7}$ & Boundary, II, I-1 and I-2 partition. \\
$Q_{r,8}$ & Continuum mass after prime-measure/Euler conversion. \\
$Q_{r,9}$ & Sign-correct specified profile after domination. \\
$Q_{r,10}$ & Directed endpoint with the paired tail convention. \\
\bottomrule
\end{longtable}
Finite telescoping yields
\[
Q_{r,0}=Q_{r,10}+\sum_{k=1}^{10}(Q_{r,k-1}-Q_{r,k}).\tag{6.4}
\]
For S2 we use (3.3)--(3.4):
\[
\begin{gathered}
\mathcal F=\Delta\cap\{h\le K\},\quad
\mathcal B=\Delta\cap\{h>K\},\quad R1=\Delta\cap\{g>249/250\},\\
\Delta=PF\ \dot\cup\ (\mathcal F\cap R1)\ \dot\cup\ \mathcal B.
\end{gathered}
\]
Here $h$ is evaluated at inherited labels, and $K$ differs from $K_0$.
S10 Ret/Drop, S5 low/top and S4/S7 A/B retain their half-open endpoints.
The modulus masks are $1-K_T$, $K_T(1-K_P)$ and
$K_TK_P\mathbf1_{\rm rough\ square}$.
They give, before the signed row coefficient,
\[
R_{\rm AA}=S_{\rm slot}+NC_{\rm sf}
-\sum\Xi_{\rm union}\Delta_P
+\sum\Xi_{\rm allowed}(\Delta_0-\Delta_P).\tag{6.4a}
\]
The first two masks split $\Xi_{\rm union}$ disjointly,
and the third gives $\Xi_{\rm allowed}$.
NC uses the remaining squarefree moduli.
A failure of both conditions belongs to the tail.

For $O\in\{\mathrm{S1},\mathrm{S2\text{-}T8/4},\mathrm{G1}\}$,
put $Q_{O,9}=\sum_{r\in O}Q_{r,9}$ and $Q_{O,10}=E_O$.
Then
\[
\begin{split}
\sum_{r\in O}Q_{r,0}={}&Q_{O,10}+(Q_{O,9}-Q_{O,10})\\
&+\sum_{r\in O}\sum_{k=1}^{9}(Q_{r,k-1}-Q_{r,k}),
\qquad Q_{O,9}-Q_{O,10}\ge0.
\end{split}\tag{6.4d}
\]
Signed differences are kept and favourable slacks are discarded.
Each tuple is assigned to its first failed condition.
The low direct primes, the upper G1 cut, the excluded factors and the squarefree/coprime restoration are removed.
The source is then retained.
The unique D1 Lemma~\ref{lem:source-replacement} and the G1
cut-measure comparison prove these changes.

A box meeting a source face is a boundary box.
Otherwise the largest eligible lower-endpoint exponent is chosen for II.
Ties are broken by the fixed variable index, and then I-1 and I-2 are used.
The DIRECT preparation lemmas verify (SH1)--(SH5)
\cite[\S3.11, (3.98)--(3.101)]{Z1}.
Every G1 coordinate exceeds $N^\alpha$,
and its two Class-I families are empty by the product-source proof.
The tables bound every adverse difference,
and the fixed-parameter argument gives $E_{\rm van}\to0$.
This proves (6.2).
\end{proof}

\subsection{Sources, local inputs and directed endpoints}

The map fixes each sign and domain.
Positive terms use lower endpoints, and negative terms use upper endpoints.
\begin{longtable}{@{}p{0.24\textwidth}p{0.71\textwidth}@{}}
\toprule
Contribution & Source and endpoint direction \\\midrule
$3S_1$ & (2.19): coefficients $(3,-3,3,-3)$, orbits $1,2,6$; one joint $\tfrac34\mathcal S_1$ from (A.4). \\
$S_2$ & Four terms; $PF$, $\mathcal F\cap R1$, $\mathcal B$ in (3.4). Whole-$K$ uses II; strict complement uses BV-R1. One lower $d_{T8}$ from (A.7). \\
$-2S_3,S_8,S_9$ & Absolute single-prefix profiles with level reduced by $u=1/500$. \\
$-S_4,-S_7$ & A/B partitions within each row; upper endpoints. \\
$-S_5$ & S5-low on $[1/4,t_*)$ and disjoint BV-top; upper endpoints. \\
$-S_6$ & Separate BV-S6 upper estimate. \\
$S_{10}$ & Ret $\dot\cup$ Drop, with its directed lower estimate. \\
$-2S_{11}-S_{12}-S_{13}$ & G1 chain (4.18), separate cutoff peel; source weights $(2,1,1)$ occur inside $A_G^{\rm cut}$. Child signs $(-,+,-,+)$ give upper/lower/upper/lower profiles; one joint $E_{\rm G1}$. \\
$-S_{\rm sw}$ & O1 (1.9), (1.12), upper $2H_{L,\rm upper}$. \\
\bottomrule
\end{longtable}
The master coefficients are
$(3,1,-2,-1,-1,-1,-1,1,1,1,-2,-1,-1,-1)$.

D1/G1 replace the sources before Buchstab/HB, with their signed corrections.
DIRECT uses Lemmas~\ref{lem:direct-preparation} and~\ref{lem:source-replacement}.
G1 uses (4.14), the monotonicity $N^{1/2}\to N^\mu$ and its cut-measure comparison.
Sections~2--4 verify the actual source and local hypotheses of
\cite[Theorem~3.2, (3.13)]{Z1} and the packing of \cite[Theorem~4.1]{Z1}.
The complete-weight prime-source theorem has a different coefficient family.
For AA37 we apply \cite[Theorem~2.6, (2.30)--(2.33b)]{Z1} with
\[
D=L_c,\quad z=z_k,\quad R=1,\quad
\epsilon_s=\sigma,\quad\delta_s=\tau=\sigma^9.
\]
Section~1 supplies the density and sieve mass.
Sections~2--4 supply the levels and cutoffs.
The three terms of this theorem are $c_kX_kV_N(z_k)M_k(s_k)$,
$-|c_k|X_kV_N(z_k)E_{s,k}$ and $a_kR_{{\rm AA},k}$.
Outer contributions and adverse source differences complete (6.1).

The nineteen applications are indexed by $k=(r,b,\mathcal C)$,
that is, by row, parent-prefix cell and signed remainder.
The remainder $R_{AA,k}$ is taken before the exterior coefficient $a_k$.
It includes the source coefficient and the AA37 sign once,
and it excludes the orbit $M_k$.
For G1, $a_k=(2A_R)\nu_k$, with $\nu_k=+1$ for upper remainders
and $-1$ for lower remainders.
The quantities $C_k,NC_{{\rm sf},k},\Xi_{U,k},\Xi_{A,k}$ are formed before this sign.
Put $\gamma_k=C_k/C_{{\rm slot},k}$.
By \cite[(2.32)--(2.33)]{Z1}, (3.8) and (4.22)--(4.26),
\[
\begin{split}
\frac1{W_{\rm norm}}\sum_k a_kR_{AA,k}
={}&\frac1{W_{\rm norm}}\sum_k\frac{a_kC_{{\rm slot},k}}{M_k}
                  \sum_q\gamma_k(q)\Delta_{P,k}(q)\\
&+\frac1{W_{\rm norm}}\sum_k a_kNC_{{\rm sf},k}
-\frac1{W_{\rm norm}}\sum_{k,q}a_k\Xi_{U,k}(q)\Delta_{P,k}(q)\\
&+\frac1{W_{\rm norm}}\sum_{k,q}a_k\Xi_{A,k}(q)
                    (\Delta_{0,k}-\Delta_{P,k})(q).
\end{split}\tag{6.4b}
\]
This identity is finite at fixed $N$.
Here $r(\mathcal A_k^{\rm pl},d)=\Delta_0(b_kd)$ by (2.17c) and (4.21),
with the AA37 remainder convention of \cite[(2.32)--(2.33)]{Z1}.
For S2 we use the per-block terms underlying (3.8).
The DIRECT comparator differences vanish, while G1 keeps their bounds.
HB is applied to the first term. Its Class-II factor is
\[
\frac{a_kh_{k,B}}{M_kW_{\rm norm}}S^{\rm PSDT}_{k,B}.\tag{6.4c}
\]
Here $h_{k,B}$ comes from the HB/log/box expansion of the same source.
Division by $C_{\rm slot}$ leaves this factor once.
The other terms enter the boundary, proper-power, LOG, I-1 or I-2 classes.
Each block keeps its tested measure, grid, coefficients, cutoff and normalization.

The hypotheses are checked as follows.
\begin{longtable}{@{}p{0.20\textwidth}p{0.75\textwidth}@{}}
\toprule
Input & Actual-source check \\\midrule
C1--C5, LEN & H1--H5 of \cite[Theorem~3.2]{Z1}; Sections~2--4 check coefficients, packing, shift, norms and scales. $a=N$, $N/L^{31}\le X\le N$, $q\le X^{0.6}$, $C_{\rm box}=C_b=334$, $K_{\rm scale}=31$. Direct preparation gives tested measure and integer intervals; G1-LOCAL uses the plain cut measure, fixed majorant and doubled support. \\
H6 & Source-wide freezing and boundary bounds, separately; Section~4.4 includes proper powers. \\
H7 & Both direct families have exponent $2012$; G1 has envelope $K_0=8042$. The finite sum preserves it, including representation multiplicities \cite[(3.91)]{Z1}. \\
Sources and masks & Sections~2.1--2.5, 3.1--3.2, 4.1--4.2; D1/G1 replacement; (6.4a)--(6.4b). \\
Profiles and cutoffs & Sections~3.5, 4.6; (2.22), (2.38)--(2.40), (3.21), (4.45)--(4.57). Squarefree/II: (4.64)--(4.65), (4.68)--(4.69). BV: Sections~3.5, 5.2--5.5. \\
Joint endpoints & Appendix A and the next two comparisons; signs from (6.1). \\
Fixed losses & (1.26)--(1.29), (6.6), parameter order below. \\
\bottomrule
\end{longtable}
The repeated-prime bound is
$232\Lambda_\sigma C_\sigma N^{-\ell_0\sigma^2}(1+L)^4$,
proved in (B.1)--(B.7).
Thus the four assembly hypotheses hold for these families.

For G1 set
\[
\begin{gathered}
E_{\rm G1}=L_{\rm G1}-\delta_{\rm G1},\qquad
\delta_{\rm G1}=4.98\cdot10^{-54},\\
L_{\rm G1}-R_{\rm G1}
=\frac{500103}{\scriptstyle100433627766186892221372630771322662657637687111424552206336},\\
L_{\rm G1}-R_{\rm G1}<4.98\cdot10^{-54},\qquad E_{\rm G1}<R_{\rm G1}.
\end{gathered}
\]
Here $R_{\rm G1}$ sums the four exact row lower endpoints.
We insert $L_{\rm G1}$ from (4.84) in $B_{\rm main}$ and
$\delta_{\rm G1}$ in $F_{\rm pers}$, once each.
Their difference is $E_{\rm G1}=L_{\rm G1}-\delta_{\rm G1}$ in (4.85).
The signed allocation underlying the row endpoints is (A.9).
The row bounds prove (6.4d).
S1 and S2 enter only through their joint endpoints.
Sections~3.1--3.5 assign all four S2 terms, the three pieces, the selector and the tail.

The S1-3 coefficient is $-3$ in (2.19).
Its single $p_3\to E_3$ replacement is (2.20).
The single comparison with $\Phi$ is (2.21)--(2.22).
Thus $E_{1,\rm pref}$ occurs once, at $Q_{\rm S1\text{-}3,9}$, with adverse value $0$.
Its cutoff cost lies inside $62500\tau$, and S2-3 has its own $E_3$.
The correction $E_{1,fF}$ changes S1-1 from $f$ to $F$ instead.

\subsection{The complete error sum}

Every entry below is a nonnegative bound in D units or a proved zero.
The finite logarithmic exponents $C$ are chosen after the earlier constants.
They may differ by row.
\begingroup\small
\setlength{\tabcolsep}{3pt}
\renewcommand{\arraystretch}{1.15}
\begin{longtable}{@{}p{0.17\textwidth}p{0.4\textwidth}p{0.34\textwidth}@{}}
contribution & Scope and exact error term & Bound \\ \hline
\endfirsthead
contribution & Scope and exact error term & Bound \\ \hline
\endhead
Source identities & The 19 signed identities, (6.4b) & $0$ \\
Source restoration & Squarefree and $(n,N)=1$ restoration on 15 direct sources, before AA37 & $O_u(N^{-u/(1+\tau)}L_N^3)$ \\
Direct lower cut & Original $p\le N/L^{30}$ tail, counted before boxes & $O_u(L^{-28})$ \\
Switched upper cut & Actual $e>N(1-L^{1-C_b})$ strip & $O_\sigma(L^{3-C_b}+L^2/N)$ \\
Source endpoint & The original $p<(1-\epsilon_A)N$ source face & $0$; straddling boxes belong to boundary \\
Exceptional source factors & Direct $p\mid N$ or switched tagged factor dividing $N$, at first removal & $O_\sigma(L^3/N+N^{-\alpha/2+o(1)})$ \\
S2 selector & Extra selected cells and dynamic comparison & $275009680\tau$ \\
S10 selector & Dropped whole cells & $7.160\tau$ \\
S1 dynamic cutoff & S1-3 cutoff and dynamic comparisons; prefactor excluded & $62500\tau$ \\
G1 cutoff peel & G1-3 four-prefix peel, never HB-expanded & $145664566.4567\tau$ \\
G1 profile boundary & Matched-profile modulus-label strips; peel separate & $0$ adverse \\
$g_{\rm BV}$ & S5-top $E_{5,\kappa}=4\cdot10^{-8}$ plus an extra allowance $g_{R1}=5.71\cdot10^{-6}$; this allowance is not the corner bound: $d_{T8}$ already includes the R1 corner at $\kappa=10^{-6}$ & $0.00000575$, once \\
$\delta_{\rm G1}$ & Joint-versus-row G1 comparison above & $4.98\cdot10^{-54}$, once \\
$g_\tau$ & Replaces the four selector/cutoff/peel bounds above & $2\cdot10^{-24}$, once \\
Fixed sieve perturbation & Integrated absolute AA37 error on all rows & $g_{\rm AA}=0.00002853328$, once \\
Asymptotic sieve perturbation & Remaining AA37 perturbation & $o_\sigma(1)$; rate unspecified \\
Squareful tail, repeated prefix, allowed rough repetition & Disjoint masks of (6.4a): first forbidden tail equality; first prefix repetition after allowed tail; retained rough repetition. Tests: $\Delta_P$, $\Delta_P$, $\Delta_0-\Delta_P$. & Theorem~\ref{thm:repeated-primes}: $O_\sigma(N^{-\ell_0\sigma^2}L^4)$ separately and together. Third term $0$ on DIRECT. \\
Squarefree normalization & Retained squarefree full moduli only & Direct $0$ on the operative tested-prime source; G1 $O_\sigma(N^{-\alpha}L^5)$ \\
Comparator density & $\omega_b(p^k)=p/(p-1)$ for every $k$; discrepancy in the preceding two rows & $0$ additional \\
Direct prime powers & Tested-side proper prime powers & $O_\sigma(N^{-1/2+o(1)})$ \\
Switched prime powers & At most four tested coordinates, each at least $N^\alpha$ & $O_\sigma(N^{-25/1149+o(1)}L^5)$ \\
Logarithm freezing & Exact $\log^{-1}$-freezing difference over all disjoint boxes & G1 $O_\sigma(L^{166-C_b})$; direct $O_\sigma(L^{20-C_b})$ \\
Logarithmic identity & Freezing is selected; Laplace truncation is unused & $0$ \\
Switched boundary & All four switched rows, one grouped contribution & $3.6\times10^{10}C_\sigma N_{\rm lab}(\sigma)(1+L)^{163}L^{2-C_b/2}$ \\
Direct boundary & All 15 direct rows, one grouped contribution & $3.94\times10^6C_\sigma N_{\rm lab}(\sigma)(1+L)^{16}L^{2-C_b/2}$ \\
Class I-1, Class I-2 & First-failure direct classes; count $K_I$ before choosing $A_i$ & $O_\sigma(L^{K_I+2-A_i})$ \\
Class II & All actual blocks with absolute exterior coefficients & Both terms of (6.5), once \\
Prime-measure and Euler-product conversion & PNT, Li, radial/product prime measure, sieve product and large $p\mid N$ factors & $o_\sigma(1)$ each; \S{}3.5, \S{}4.6 (4.76)--(4.78), direct-BV terms in \S{}5. O1 (5.33)--(5.43) enters $E_{\rm Ssw}$. \\
Profile comparison & Signed domination of the specified profile at the actual level & $0$ adverse by the row comparisons in \S\S{} 2--4; favourable slack discarded \\
$E_{1,\rm pref}$ & S1-3 prefactor, (2.20)--(2.21), at $Q_{\rm S1-3,9}$; dynamic cutoff excluded & $0$ adverse, once \\
$E_{1,fF}$ & S1-1 correction; the S2 correction is already inside its joint endpoint & One error $10^{-390}$ \\
S2 infinite tail & Joint S2 A=1 tail beyond the finite prefix & One error $1/20000$, already included when $d_{T8}$ is used \\
Directed enclosure & Directed numerical enclosure of the specified profile & $0$ additional \\
R1 BV remainder, S5 BV remainder, S6 BV remainder & Three actual BV/sieve bridge errors & $o(1)$, including $O(L^{-1/3})$, $O_A(L^{2-A})$ and the recorded $O(\log L/L)$ measure conversion \\
S6 level increment & Selected S6 $\kappa=10^{-6}$ increment over the nominal S6 baseline & One error $0.000006484294227956901$ \\
Switched-excess error & Full O1 (5.33)--(5.43): source exceptions, outer boxes, prime/radial mass, noncoprime correction, small primes, sieve error, level/Euler conversion and product BV & $E_{\rm Ssw}=o(1)$, once; product BV costs $O(L^{-5/2})$. \\
\end{longtable}
\endgroup

We take $H_L=I_{HL}$, the coordinate measure and $t_k$ of (5.30)--(5.32).
Set $W=W_{\rm norm}$ and put
\[
E_{\rm Ssw}(N)=\max\{0,S_{\rm sw}(N)/W-2I_{HL}\}.
\]
Then $-S_{\rm sw}/W\ge-2I_{HL}-E_{\rm Ssw}$, and
$I_{HL}\le H_{L,\rm upper}$ gives the sole main term $-2H_{L,\rm upper}$.
For fixed geometry and BV exponents, (5.43) gives
\[
\begin{split}
E_{\rm Ssw}\ll_{\varepsilon_A}{}&
L^{-1/3}+\frac{\log L}{L}+L^{-1}+L^{-2}
+\eta_P(N)+\eta_E(N)\\
&+L^{-5/2}+N^{-3/4}L^2+N^{-\alpha}L^3=o(1).
\end{split}
\]
Here $\eta_P,\eta_E\to0$ are the prime-box PNT/Euler errors of Section~5.5.
The radial and outer-box masses contribute $L^{-1},L^{-2}$.
By (5.40),
\[
R_{\rm nc}\le22X^+N^{-\alpha}\sum_{r<D}\frac{\mu^2(r)}{\varphi(r)}
\ll X^+N^{-\alpha}L.
\]
Since $X^+\ll N$ and $U_C\gg N/L^2$, its contribution is $O(N^{-\alpha}L^3)=o(1)$.
It enters $E_{\rm Ssw}$ once.
The $N^{-\alpha}L^3$ envelope covers (5.33) and (5.40), including source exceptions,
and $N^{-3/4}L^2$ covers the $36z$ endpoint.
The direct-BV conversions and the outer errors below are kept separate.

For Class II put $c_\sigma=\ell_0\sigma^2/4>0$ and
$C=C(\sigma,C_b,A',31,\mathcal C_{\rm src})$.
By \cite[Theorem~3.2, (3.13)]{Z1},
\[
|S_B|\ll_{\sigma,31,C_b,A',\mathcal C_{\rm src}}
XL^{-A'}+X^{1-c_\sigma}L^C.
\]
We sum the exterior coefficients $\le L^{K_0}$, use $X\ge N/L^{31}$ and divide by $W_{\rm norm}$:
\[
E_{\rm II}(N)\ll_\sigma L^{K_0+2-A'}
+N^{-c_\sigma}L^{K_0+C_{\rm pow}+2}.\tag{6.5}
\]
Here $C_{\rm pow}=C+31c_\sigma$ absorbs the fixed factor from
$X\ge N/L^{31}$.
Both terms are kept.
The cited proof contains GR, colouring, cross-copy coprimality, CRT,
zero mode, moments, Poisson, W-prime norm, Theorem T and endpoints
\cite[Theorem~3.2, \S\S3.5--3.9]{Z1}.
The actual-source H7 sum includes them once
\cite[\S3.10, (3.91)--(3.92)]{Z1}.
The outer errors are $O(N^{-\alpha}L^3)$ for rough exceptions,
$O(N^{-1/2}L^2)$ for squares/source factors dividing $N$,
and $o(1)$ for G4.
They enter (6.9) once, separately from O1.

\subsection{Selected endpoints and persistent costs}

We use the adverse endpoints of Appendix A.2--A.6. Terminating decimals are rational.
\begingroup\small
\setlength{\tabcolsep}{3pt}
\renewcommand{\arraystretch}{1.15}
\begin{longtable}{@{}p{0.22\textwidth}p{0.30\textwidth}p{0.40\textwidth}@{}}
Input & Certified value or enclosure & Use \\ \hline
\endfirsthead
Input & Certified value or enclosure & Use \\ \hline
\endhead
\multicolumn{3}{@{}p{\textwidth}@{}}{S7 at $\lambda=93/400$: canonical raw value $2.299260427561925683718\ldots$ in $[2.299209653294856,2.299402520476652]$; use the upper endpoint.} \\
\multicolumn{3}{@{}p{\textwidth}@{}}{S1 in D units: $3S_1/4\ge\underline D_1$, the lower endpoint of the S1 ball in Appendix A.2.} \\
\multicolumn{3}{@{}p{\textwidth}@{}}{Selected S2 endpoint: $d_{T8}>\underline d_2$.} \\
\multicolumn{3}{@{}p{\textwidth}@{}}{Independent S2 enclosure: $T8/4\in[1.62063469556829,1.62071322785677]$.} \\
\multicolumn{3}{@{}p{\textwidth}@{}}{S2 refinement: PF $0.000944946137955914$, corner $0.000058392441251611$, tail $0.00005$. The sum differs from the endpoint improvement by less than $2\cdot10^{-15}$.} \\
Switched integral & $H_{L,\rm upper}=0.0778506$ & $-2H_{L,\rm upper}=-0.1557012$. \\
Signed reference sum & $B_{\rm main}$ & $=1.83403742109459478063$. \\
Persistent errors & $F_{\rm pers}$ & $<0.000040767574227956903$; decomposition (6.6). \\
\end{longtable}
\endgroup

The two lower endpoints used above are
{\small
\[
\begin{aligned}
\underline D_1&=9.6765387489412441290823874931869409094969802156943270964
                 -6.00\cdot10^{-56},\\
\underline d_2&=1.61958135698908352559133404332105729220076714081650.
\end{aligned}
\]
}

The printed reference sum is the following thirteen-term sum of the
adverse endpoints in Appendix~A:
\[
\begin{aligned}
B_{\rm print}={}&\underline D_1+\underline d_2-\frac{U_{3,A}}2-\frac{U_4}4
-\frac{U_{5,\rm low}}4-\frac{U_7}4\\
&+\frac{L_{8,A}}4+\frac{L_{9,A}}4+\frac{L_{10}}4+L_{G1}
-\overline J_5-\overline J_6-2H_{L,\rm upper}.
\end{aligned}
\]
Here $\underline D_1$ is the lower endpoint of the S1 ball in Appendix~A.2,
and $\underline d_2$ is the selected S2 lower endpoint printed above.
The four $U$ endpoints and four $L$ endpoints are the adverse directed
rational endpoints printed in (A.7), (A.10), and Appendices~A.2--A.6.
The upper endpoints for the two remaining integrals are
{\small
\[
\begin{aligned}
\overline J_5&=\frac{0.03658166423843258896399704402764682632822246022922592990
                    +8.84\cdot10^{-57}}4,\\
\overline J_6&=1.3148581959157321533513247276782285310758770468886237619585.
\end{aligned}
\]
}
They come from (A.14) and (A.13), respectively.
The switched-integral endpoint is $H_{L,\rm upper}=0.0778506$.
Every summand is a printed rational endpoint; the sum uses no unrounded
analytic integral value.
Exact rational addition gives
$B_{\rm print}>1.83403742109641331345957914542194044932304169727869681073064$.
We take the smaller rational $B_{\rm main}=1.83403742109459478063$.
Thus $B_{\rm main}$ is obtained from the printed endpoint data.

We use the net-tail value $d_{T8}$ of (A.7).
Put $\delta_{T8}=1/20000$; this tail is already inside $d_{T8}$.
The persistent sum is
\[
\begin{split}
F_{\rm net}={}&0.000006484294227956901+0.00000575+0.00002853328\\
&+4.98\cdot10^{-54}+2\cdot10^{-24}+10^{-390}
<0.000040767574227956903.
\end{split}\tag{6.6}
\]
The order is the S6 increment, $g_{\rm BV}$, $g_{\rm AA}$,
the G1 reconciliation, $g_\tau$, $E_{1,fF}$.
By (1.27), $C_E\sigma W_{\rm AA}/4\le g_{\rm AA}$.
The exact BV decomposition is
\[
g_{\rm BV}=0.00000575
=\underbrace{0.00000004}_{E_{5,\kappa}}
+\underbrace{0.00000571}_{g_{R1}}.
\]
Here $E_{5,\kappa}=4\cdot10^{-8}$ is optional and conservative.
For comparison, Appendix~A.6 bounds the hypothetical S5-top increment
at $\kappa=10^{-6}$; the selected S5-top uses zero shift.
The selected S2 corner already uses $\kappa=10^{-6}$ and its increment
by (3.19) and (3.22).
The extra $g_{R1}\ge0$ is conservative. Both terms enter (6.6) once.

The four selector/cutoff/peel bounds sum to
\[
\begin{split}
\bigl(275009680+7.160+62500\\
\qquad+145664566.4567\bigr)\tau
&=420736753.6167\tau\\
&\le4.207367536167\times10^{-82}<2\times10^{-24}.
\end{split}
\]
The deliberately looser bound $4.201\times10^{12}\tau$ also covers this sum.
The allowance $2\times10^{-24}$ enters (6.6) once.
The independent S2 enclosure includes every tail:
\[
T8/4>d_{T8}.\tag{6.6a}
\]
Its lower bound improves $d_{T8}$ by about $0.001053338579$.
The input table gives the three improvements and the rounding error.
Sections~3.1--3.5 and~6.2 prove the actual-source arithmetic,
non-II, cutoff, AA37 and joint comparisons.

The equivalent pre-tail convention transfers only $\delta_{T8}$:
\[
\begin{aligned}
S2_{\rm pre}&=S2_{\rm net}+\delta_{T8},\\
B_{\rm pre}&=B_{\rm net}+\delta_{T8},\\
F_{\rm pre}&=F_{\rm net}+\delta_{T8}<0.000090767574227956903.
\end{aligned}\tag{6.6b}
\]
Thus $B_{\rm pre}-F_{\rm pre}=B_{\rm net}-F_{\rm net}$ exactly.
Here
\[
S2_{\rm pre}>1.61963135698908352559133404332105729220076714081650.
\]
The net value is smaller by $0.00005$,
and all other endpoints and error terms agree.
We use the net pair and subtract each error term once.
Then
\[
\begin{aligned}
B_{\rm main}-F_{\rm pers}&=B_{\rm net}-F_{\rm net}\\
&>1.83403742109459478063-0.000040767574227956903\\
&=1.833996653520366823727\\
&>1.83399665352036682372\\
&>1.78838848345960807361670991333>T.
\end{aligned}
\]
The strong outward margin exceeds $0.04560817006075875010$ D.
We retain the smaller safety margin
$D_{\rm lower}-T-g_{\rm AA}>0.04557963678075875010$.
This spends an extra $g_{\rm AA}=0.00002853328$ in the final vanishing-error
threshold; it is a reserve and is not charged a second time in (6.6).
The computation uses the full-precision (A.7), $H_{L,\rm upper}$ and the single G1 correction.
The displayed endpoint sum is $1.83403742109641\ldots$;
its excess over the conservative base is $1.82\cdot10^{-12}$.
Shorter intervals are outward summaries.
By (6.2) and (6.4b),
\[
\begin{gathered}
D_{\rm lower}=1.83399665352036682372,\qquad
D_{\rm src}(N)\ge D_{\rm lower}-E_{\rm van}(N),\\
E_{\rm van}(N)\longrightarrow0.
\end{gathered}\tag{6.3a}
\]

\subsection{Fixed parameters, then the limit}

First fix the source geometry and the classical constants.
Theorem~\ref{thm:sieve-mass} gives $W_{\rm AA}\le W_*=10^{14}$,
uniformly in tail-label counts.
The constant $C_E$ is absolute \cite[Theorem~2.6, (2.31)]{Z1}.
Choose
\[
0<\sigma\le10^{-10},\qquad
\sigma\le\frac{4g_{\rm AA}}{C_EW_*}\quad\text{if }C_E>0.
\]
Since $\eta_L/2500=4\cdot10^{-8}$ and AA37 has depth $R=1$
\cite[Theorem~2.6]{Z1}, the omitted restrictions
$\sigma\le\eta_L/2500$ and $\sigma<3^{-R}$ are automatic.
When $C_E=0$, only $0<\sigma\le10^{-10}$ is required.
Then $\tau=\sigma^9$ fixes the labels, factorial bounds and positive power savings.
Appendix A holds for $0<\sigma\le10^{-10}$ and $0\le\tau\le10^{-90}$.

By (4.20), with $A_0=53/256-\eta_L-\frac32\mu\tau$,
\[
\ell_c\ge\max(A_0-t_3,2t_3)\ge2A_0/3>1379/10000>\ell_0.
\]
Fix $A_{\rm raw}=3$, $C_b=334$ and then
\[
K_0=24C_b+26=8042,\quad A'=A_{\rm raw}+K_0+2=8047,\quad
A_{\rm work}=A'+4=8051.\tag{6.7}
\]
The G1 boundary exponent is $165-C_b/2$.
The condition $C_b>330$ gives decay, and the choice $334$ gives $O_\sigma(L^{-2})$.
The three slot-shell sums cost at most three logarithms, which leaves $A'$.
After counting the full Class-I family, we take $A_i\ge K_I+4$.
Choose the exceptional-cutoff and moment precisions by (4.67), including
$c_*\ge c_{\rm pre}+2$, so that all bounds in (1.29) hold.
The SW/BFI precisions are chosen next.
The width $334$ makes the source-wide boundary and freezing errors vanish.
The selected local precision gives the Class-II logarithmic term
$O_\sigma(L^{-3})$.
The resulting rates are
\begin{center}
\begin{tabular}{@{}lcc@{}}
Error & G1 & DIRECT \\\hline
Boundary & $O_\sigma(e^{163/L}L^{-2})$ & $O_\sigma(L^{-149})$ \\
Logarithm & $O_\sigma(L^{-168})$ & $O_\sigma(L^{-314})$
\end{tabular}
\end{center}
The upper G1 cut is $O_\sigma(L^{-331}+L^2/N)$;
Class I is $O_\sigma(L^{-2})$; the logarithmic term of (6.5) is $O_\sigma(L^{-3})$.
For fixed $v>0$ and finite $B$,
\[
\log(N^{-v}L^B)=-v\log N+B\log\log N\longrightarrow-\infty.
\]
Theorems~\ref{thm:repeated-primes} and~\ref{thm:sieve-mass}, with (6.5), give vanishing row errors.
The PNT/Mertens, BV/sieve and asymptotic AA37 errors are $o_\sigma(1)$.
The finite identities and the complete error table cover the sum.
Finally choose $N_0$. The choice includes all local theorem thresholds and requires $X\ge N/L^{31}$ and $\log X\ge L/2$. The endpoint slack conditions are
\[
\begin{aligned}
202(31\log L+1)&\le\sigma L,
&202\log(2^{24}/\varepsilon_A)&\le\sigma L,\\
202\cdot30\log2&\le\sigma L,
&\sigma L&\ge2^{130}(\log L+1).
\end{aligned}
\]
The last is the short-interval threshold (4.32a).
We also impose the GR roughness condition
$1/(\log\log X)^3\le\ell_0\sigma^2$ and all fixed short-interval moment and Brun--Bonferroni thresholds. The choice also includes (1.14) and the final slack in (1.23).
These are eventual conditions at fixed positive $\sigma$. For fixed $\kappa>0$, $C>0$, $d>0$ and real $B$, the quantity $CN^{-\kappa}L^B$ is at most $d$ when
\[
 L\ge\max\{1,(2\max(B,0)/\kappa)^2,
                         (2/\kappa)\log^+(C/d)\}.
\]
This follows from $\log L\le\sqrt L$. It shows directly that the
power errors in (1.30) are bounded by the last threshold choice.

For $N\ge N_0$ we also require
\[
E_{\rm van}(N)<\tfrac12(0.04557963678075875010).
\]
The weaker allowance is $D_{\rm lower}-T>0.04557963678075875010$.
The last threshold choice may be ineffective.

\subsection{From positive weight to a representation}

The results of Sections~3--5 hold for each fixed $0<\epsilon_A\le10^{-4}$.
Choose $\epsilon_A=10^{-4}$, $a=9/5$ and
\[
b=\frac{a-1}{a}-\epsilon_A=\frac{39991}{90000},\qquad
\frac{1-b}{b}=\frac{50009}{39991}.
\]
The exact threshold satisfies
\[
\begin{gathered}
1.78838848345960807361670991332<T=8\log(50009/39991),\\
T<1.78838848345960807361670991333,\\
D_{\rm lower}-T>0.04557963678075875010.
\end{gathered}
\]
We fix $\sigma$ before the limit.
By the definition of $\Sigma_W$ in (1.9), (6.1) gives
$\Sigma_W=W_{\rm norm}D_{\rm src}(N)$.
Theorem~\ref{thm:signed-assembly}, (6.2) and (6.3a) give
\[
\Sigma_W\ge
[4(1-\epsilon_A)D_{\rm lower}-o(1)]U_C.\tag{6.8}
\]
The pointwise sieve inequality subtracts the rough exceptional-set error.
Then the count is at least $[(1-\epsilon_A)D_{\rm lower}-o(1)]U_C$.
The fixed-$\epsilon_A$ G4 bound costs
$[8(1-\epsilon_A)\log((1-b)/b)+o(1)]U_C$.
Since $1-\epsilon_A>0$, $D_{\rm lower}>T$, and $U_C\to\infty$,
the integer count is eventually positive.

For a surviving semiprime $n=rq>\epsilon_A N$, with distinct primes
$r<q$ and $r<N^b$, one has $q>\epsilon_A N^{1-b}$ and
\[
(1-b)(a-1)-b=a\epsilon_A>0.
\]
Thus $r<q^{a-1}=q^{4/5}$ for
$N>\epsilon_A^{-(a-1)/(a\epsilon_A)}$.
For a prime survivor, $r=1$.
The positive integer count gives (1.1) in the quantitative form
\[
\operatorname{Good}(N)\ge
\big[(1-\epsilon_A)(D_{\rm lower}-T)-o(1)\big]U_C,
\qquad T=8\log(50009/39991).\tag{6.9}
\]
This proves Theorem~\ref{thm:goldbach}. $\square$

\appendix
{
\def\UrlBreaks{\do a\do b\do c\do d\do e\do f\do g\do h\do i\do j\do k\do l\do m\do n\do o\do p\do q\do r\do s\do t\do u\do v\do w\do x\do y\do z\do A\do B\do C\do D\do E\do F\do G\do H\do I\do J\do K\do L\do M\do N\do O\do P\do Q\do R\do S\do T\do U\do V\do W\do X\do Y\do Z\do 0\do 1\do 2\do 3\do 4\do 5\do 6\do 7\do 8\do 9\do /\do _\do .\do -}
\newcommand{\CertificateInterval}[2]{\left[\begin{gathered}{}#1,\\{}#2\end{gathered}\right]}
\section{Numerical certificates}

The source comparisons hold by Sections~2--5.
Section~6 uses each selected inequality once,
with the signs of (6.1).
The tables give full rational certificates and independent checks.
Their directed interval calculations use Arb's midpoint-radius arithmetic
\cite{Johansson2017}; rational endpoint checks are given explicitly.

\subsection{Exact interval and normalization conventions}

Each terminating decimal is an exact rational,
and each endpoint includes its displayed radius.
The extra digits fix the direction of each bound.
Set
\[
 B(c,r)=[c-r,c+r].                                      \tag{A.1}
\]
Shorter intervals are rounded outward.
The numerical-certificate supplement archive \path{paper-ii-public-certificates-v11.tar.gz} has SHA-256
\path{cb9d7b96ce1323c9118d845f4dd701a0d830e8c610090deec88bd8792648e100}.
Each printed file SHA-256 is the digest of that file in the archive.
An endpoint certificate bounds its endpoint,
and a functional enclosure bounds the integral.
Paths are relative to the extracted numerical-certificate supplement.

The common parameters are
\[
\begin{aligned}
 u&=\frac1{500},\quad \alpha=\frac{100}{1149},\quad
 \beta=\frac{50}{309},\\
 \mu&=\frac{25}{128},\quad \lambda=\frac{93}{400},\quad
 \tau_S=\frac{57}{224},\\
 t_*&=\frac{7109}{28000}.
\end{aligned}\tag{A.2}
\]
The parameter-family certificates require
$0<\sigma\le10^{-10}$ and $\tau=\sigma^9\le10^{-90}$.
They apply to the same source and scalar profile.
In Section~4 the G1 source width is $C_b=334$.

The two numerical outputs and their component manifests are:

{
\normalsize
\renewcommand{\arraystretch}{1.14}
\begin{longtable}{@{}p{\textwidth}@{}}
\toprule
\textbf{\path{endpoint-recert/out/compare_hp.json}, S1, row and G1} \\
\emph{SHA-256}: \path{cff4d6dc0168e66a29a915473c3391c383c168fb81ac71ecec87e90016331bf8} \\
\emph{Manifest SHA-256}: \path{8760a5a78fff44ca739cd12660d3a86765dd2f3d6872eae111d933f354f54730} \\
\midrule
\textbf{\path{s2-joint/outputs/ASSEMBLY.json}, S2} \\
\emph{SHA-256}: \path{21083883403f3903fff3ce781e086a26f931dd809e445578cfcba546295e9e5d} \\
\emph{Manifest SHA-256}: \path{d26488e736e5798544e758338acb299d417b74a309568f5ce20f738e399cf96d} \\
\bottomrule
\end{longtable}
}

Certificate paths are relative to the supplement root.
The root \path{SHA256SUMS} lists each retained filename and digest,
including the selected S2 components.
The true $F/f$ profiles are computed by \path{endpoint-recert/code/tm.py}, with SHA-256
\path{d34b11f02353473f714d00dbb02d73f9a91be5275d9589fc0db1de3aa851ec16}.
The exact comparisons are in \path{endpoint-recert/out/compare_hp.json}, with SHA-256
\path{cff4d6dc0168e66a29a915473c3391c383c168fb81ac71ecec87e90016331bf8}.
The bounds come from Taylor models and directed quadrature.

Throughout, $U_C=C(N)N/\log^2N$, with $C(N)$ as in \S~1.
The signed master coefficients, in the order
$S_1,\ldots,S_{13},S_{\rm sw}$, are
\[
 (3,1,-2,-1,-1,-1,-1,1,1,1,-2,-1,-1,-1).               \tag{A.3}
\]
Thus S1 contributes $3S_1/4$ and S3 contributes $-S_3/2$.
The multiplier is $-1/4$ for S4/S5/S6/S7
and $1/4$ for S8/S9/S10.
The joint S2 and G1 values are already in D units.
The G1 source multiplicities $(2,1,1)$ sit inside its source weight,
so they add no DIRECT-row coefficients.
The switched term contributes $-S_{\rm sw}/\mathcal W$.
Up to its stated vanishing error, this is bounded below by $-2H_L$.

Each selected endpoint has an independent check,
and each row and joint contribution enters once.

\subsection{S1: joint enclosure}

For the joint profile in Section~2,
\[
 12.9020516<S_1<12.9023390,\qquad
 \frac34S_1>9.6765387.                                 \tag{A.4}
\]
It covers S1-0 through S1-3 and includes the first $-F$ subtraction.
The subtraction uses $u$ and $s_0=19037/64$.
The integral of Section~2.3 satisfies
\[ I_*<2161.074150710003. \]
Both bounds refer to the joint profile of Section~2.3.

The selected enclosure is computed by
\path{s1-tail/certify_s1_corrected_tail.py},
with SHA-256
\path{da9f572cbb28ca06a00e050785eefd9cd4f499269a627deaaa9a593c4ab210a3};
its output
\path{s1-tail/corrected-certificate.json}
has SHA-256
\path{ebf20d49c73cbf7ae11314e3a1e7e82005a75c42a6c2fe91547698f877011f4f}.
The infinite factorial tail is
\[
 \sum_{m\ge40}\frac{(m+2)^2}{(m+1)m!}
 \le\frac{42^2}{41\,40!}\frac{39}{38},                 \tag{A.5}
\]
since the ratio of consecutive summands is at most $1/39$.
The endpoint (A.4) is used.

{
\normalsize
\renewcommand{\arraystretch}{1.14}
\begin{longtable}{@{}p{\textwidth}@{}}
\toprule
\emph{Exact outer intervals containing the computed bounds} \\
\midrule
\textbf{S1 lower bound} \\
$B(12.9020516652549921721098499909159212126626402875924361285,4.64\cdot10^{-56})$ \\
\midrule
\textbf{S1 upper bound} \\
$B(12.9023389323670212025588110758944404798760246915400198042,3.39\cdot10^{-56})$ \\
\midrule
\textbf{D lower bound from S1} \\
$B(9.6765387489412441290823874931869409094969802156943270964,6.00\cdot10^{-56})$ \\
\midrule
\textbf{Infinite-tail majorant $J_{\rm upper}$} \\
$B(0.000255821764915303175726784912190615770902363950678941419476,3.34\cdot10^{-61})$ \\
\midrule
\textbf{Prefix deficit} \\
$[0,7.792\cdot10^{-383}]$ \\
\bottomrule
\end{longtable}
}

The prefix deficit is subtracted before computing the lower endpoints in (A.4).
Thus (A.4) includes the deficit once.
The source-to-profile comparison is in \S~2.3.

The independent calculation is \path{endpoint-recert/code/s1_cert.py}, with SHA-256
\path{718e495fb18ea71ef0538c73ac03f96a0cac77f2274ffbe5c1a093b76605a508}.
It re-derives the cancellation and tail bound from the delay equations.
Its output is \path{endpoint-recert/out/s1_cert.json}, with SHA-256
\path{f16972e2b73258532182357454193b5002b553dc465b829f9af9b2c0680d4137}.
It gives the sharper bounds
\[
 S_1\in[12.9020760,12.9023390],\qquad
 \frac34S_1>9.6765570.                                 \tag{A.5a}
\]
The computed lower-bound ball is
{\normalsize
\[
B(12.9020760472276414637285015429505331626940241680712780695,
4.38\cdot10^{-56})\mathpunct{,}
\]
and its $3/4$ multiple is
\[
B(9.67655703542073109779637615721289987202051812605345855210,
5.80\cdot10^{-56})\mathpunct{.}
\]
}
These bounds imply (A.4) and its tail estimate.
They also give $I_*<2161.074150710003$.
The selected endpoint is (A.4).

\subsection{The seven separate profile endpoints}

The seven directed endpoints below come from three certificate groups.
Option-A gives the absolute S3-A/S8-A/S9-A values.
Direct-row gives S4, S5-low and S7,
and Retained-row gives S10-Ret.
Each endpoint has an independent comparison.
For S4, S5-low, S7 and S10 the final sum uses the adverse endpoint of the
first printed interval; the following ball is an independent check.
For the three Option-A rows it uses the displayed ball endpoint in the
adverse direction.
S7 is taken at $\lambda=93/400$.
S10 uses the source-trimmed upper inner boundary.

{
\normalsize
\renewcommand{\arraystretch}{1.14}
\begin{longtable}{@{}p{\textwidth}@{}}
\toprule
\emph{Row; coefficient giving its D contribution; certificate group.} \\
\emph{Raw-profile outer interval; next line: independent endpoint comparison.} \\
\midrule
\textbf{S3-A (Option A)}\quad $-1/2$\quad Option-A \\
$B(10.45234847108890216860085651838875393044983358511004,5.69\cdot10^{-51})$ \\
Upper endpoint verified in the archived comparison record \\
\midrule
\textbf{S4}\quad $-1/4$\quad Direct-row \\
$[3.336569865049636,3.336975348289866]$ \\
$B(3.33672680275562584206999257242,4.25\cdot10^{-30})$ \\
\midrule
\textbf{S5-low}\quad $-1/4$\quad Direct-row, direct cells $[1/4,t_*)$ \\
$[0.238854512468062,0.238888433191047]$ \\
$B(0.238875420270830663069285634601,4.07\cdot10^{-31})$ \\
\midrule
\textbf{S7}\quad $-1/4$\quad Direct-row, $\lambda=93/400$ \\
$[2.299209653294856,2.299402520476652]$ \\
$B(2.29926042756192568371844970138,1.46\cdot10^{-30})$ \\
\midrule
\textbf{S8-A (Option A)}\quad $+1/4$\quad Option-A \\
$B(2.36595898835331722398869549955202092760264561327680,6.59\cdot10^{-51})$ \\
Lower endpoint verified in the archived comparison record \\
\midrule
\textbf{S9-A (Option A)}\quad $+1/4$\quad Option-A \\
$B(1.341531658234294936348198541015408702983637197761145323037489950,4.21\cdot10^{-64})$ \\
Lower endpoint verified in the archived comparison record \\
\midrule
\textbf{S10-Ret}\quad $+1/4$\quad Retained-row \\
$[0.858847455777940,0.867486588163666]$ \\
$B(0.862377086132397387439154957615,6.98\cdot10^{-32})$ \\
\bottomrule
\end{longtable}
}

{
\normalsize
\renewcommand{\arraystretch}{1.14}
\begin{longtable}{@{}p{\textwidth}@{}}
\toprule
\textbf{Direct-row} \\
\emph{Source code and SHA-256}: \path{row-enclosures/cert_s4_s7.py}: \path{d98866b45f1ddb8f3f93411bfe09cb11c62569cc80fa775a77a3eab1ad4384aa} \\
\emph{Output and SHA-256}: \path{row-enclosures/result.json}: \path{0eb5fa52ed180a071633e50f0d099f5c4e1f9bb001add71765fe1c2f033b8f7a}; S4, S5 direct cells, S7 \\
\midrule
\textbf{Option-A} \\
\emph{Source code and SHA-256}: \path{option-a/arb_enclose_optionA_independent.py}: \path{9c1880609692d30be9f282a218280c291dbb782a7d7cfab35732e9203687c36b} \\
\emph{Output and SHA-256}: \path{option-a/arb_enclose_optionA_out.json}: \path{5a8786a27da1a52550eb0c077ae11c48072aff6d1765cebbad25b24e2a556b27} \\
\midrule
\textbf{Retained-row} \\
\emph{Source code and SHA-256}: \path{s10/arb_s10_exact.py}: \path{74a130143e5c9e40ac530dd688526e63acde93a95f5e8271b0f07984afc5a98f} \\
\emph{Output and SHA-256}: \path{s10/s10-s10.json}: \path{f6fba094d222f362341bc63ba0388cbf6a7556accdd14d5b28b882d7e0626cf6} \\
\bottomrule
\end{longtable}
}

Direct-row uses 192-bit Arb with $\tau\in[0,10^{-90}]$.
S5-low sums two direct cells on $[1/4,t_*)$, with outward rounding.
The BV-top cell is disjoint from them. The independent
calculation
\path{endpoint-recert/code/rows_cert.py} has SHA-256
\path{8b2ef756bbdd5530f614a66b8c5eb5dddd17922555d979515b1ca85876064bc5}.
It produces all seven results in \path{endpoint-recert/out/rows_cert_hp.json}, with SHA-256
\path{456d03d3ea9954bb71fedc8e20ea35cd463454bb9a0919b3ace9f64e09f54464}.
The exact rational comparisons are in \path{endpoint-recert/out/compare_hp.json} of Appendix A.1.
They include the selected S8/S9 directions.
S3/S4/S5/S7 include the adverse $\tau=10^{-90}$ effects,
and S8/S9/S10 use the actual-coordinate profiles.
Their fixed-$\tau$ comparisons are proved in Sections~2.4--2.5,
especially (2.40).

The independent S10 retained-strip bound is $4.82\tau$ in D units.
It lies within the single separate $7.160\tau$ bound.

The absolute Option-A rows use one source,
with the subtraction at $u=1/500$.
These rows give the three selected absolute endpoints.

The independent comparison gives
{\normalsize
\[
\begin{gathered}
 S_{3,\mathrm{OptA}}\in\\
B(10.45234847108890216860085651838875393044983358511004,
   5.69\cdot10^{-51}),
\end{gathered}
\]
}
and the matched change is
{\normalsize
\[
\begin{gathered}
 \Delta D_{\mathrm{OptA}}\in\\
B(-0.0110499568842910017666362322996150028835154631047,
   9.53\cdot10^{-50}).
\end{gathered}\tag{A.6}
\]
}
Equation (A.6) is a comparison value only.
The signed sum uses the three absolute endpoints.

\subsection{S2: one four-row endpoint}

Let $P_-$, $J_+$ and $C_-$ be directed endpoints.
They bound the prefix below, PF above and the BV corner below.
On $\Delta=\{u\le r_3<r_2<r_1\le\beta\}$, set
\[
 V_{\rm full}=V
 =\int_\Delta\frac{dr_1\,dr_2\,dr_3}{r_1r_2r_3^2}
 =\frac{\log^2(\beta/u)-2\log(\beta/u)+2}{2u}-\frac1\beta.
\]
An outward evaluation gives
{\small
\[
\begin{aligned}
 V_{\rm full}\in[&
3122.423962424988515509469383070552351626538598751986375091552742,\\
&3122.423962424988515509469383070552351626538598751986375091552743].
\end{aligned}
\]
}
The endpoints in (A.7) are
\[
\begin{aligned}
P_-&=3513.1687300568893335013223065223679611140065161989243,\\
J_+&=0.359440819140453278302418897934329259372468992946659201284,\\
C_-&=-0.0574908755554524530180814901226778275854309897822443970870.
\end{aligned}
\]
Use four times the printed interval for $1/(2e^\gamma)$ below to bound
$2/e^\gamma$. Define $d_{T8}$ by
\[
 d_{T8}=\frac{P_--(2/e^\gamma)V_{\rm full}
                    -(2/e^\gamma)J_++C_-}{4}
              -\frac1{20000}.                         \tag{A.7}
\]
The directed assembly gives
\[
 d_{T8}>
 1.61958135698908352559133404332105729220076714081650.    \tag{A.8}
\]
Let $k_-,k_+$ be the printed endpoints for $k=1/(2e^\gamma)$ below,
and let $V_-,V_+$ be the endpoints above. Exact rational interval
arithmetic in (A.7) gives
\[
\begin{aligned}
 d_{T8}\in\bigl[&P_-/4-k_+(V_++J_+)+C_-/4-1/20000,\\
                 &P_-/4-k_-(V_-+J_+)+C_-/4-1/20000\bigr].
\end{aligned}
\]
This is the selected joint lower value in D units.
It covers S2-0 through S2-3(K/PF)
and includes the tail $1/20000$ once.
The term $275009680\tau$ is separate,
and so are AA37 and the other finite-$N$ errors.
The grid is the global one, $\tau_g=\tau=\sigma^9$.
The selector and PR costs satisfy
\[
\begin{aligned}
 &(275009680+7.160+62500+145664566.4567)\tau\\
 &=420736753.6167\tau\le4.207367536167\cdot10^{-82}.
\end{aligned}
\]
The BV corner is $c_\gamma V_B-\mathrm{IBV}$ from (5.18). IBV is
the R1 kernel (5.17) at $\kappa=10^{-6}$ in \S~5.3. The independent
nested Arb source is \path{s2-joint/code/corner.py}, with SHA-256
\path{7cbc5d276b54b52d657fc52622db1a3542c5079d72a39e3be3d12eec1ddb21f2}.
Its line 11 sets $\kappa=10^{-6}$. It gives
\[
 c_\gamma V_B-\mathrm{IBV}\in
 B(-0.057257305790446009369486009082468758,4.67\cdot10^{-37}).
\]
The output is \path{s2-joint/outputs/corner.json}, with SHA-256
\path{360841b6008cb08be571d23177887df14be286ebde7ba3f3fe2c7571f31be9d8}.
This value exceeds $C_-$ by more than $2.3356\cdot10^{-4}$, so $C_-$
is a lower endpoint at $\kappa=10^{-6}$. Increasing $\kappa$
decreases $1/2-\kappa-r_1-r_2-r_3$ and lowers the corner, so its
$\kappa=0$ value is larger.
For the R1 level increment, the S2 corner uses $\kappa=10^{-6}$,
so $d_{T8}$ includes this increment.
Section~\ref{sec:proof-theorem-1} keeps an additional conservative R1 allowance.

\paragraph{Corner increment bound.} R1 is the tetrahedron
$\{r_3\le r_2\le r_1\le\beta,\ r_1+2r_2+4r_3\ge H\}$. Its vertices are
$(\beta,\beta,\beta)$, $(H/7,H/7,H/7)$, $(\beta,x,x)$, and
$(\beta,\beta,y)$, where $x=(H-\beta)/6$ and
$y=(H-3\beta)/4=13147/103000$. On R1, $r_i\ge y$ and
$a_0:=1/2-r_1-r_2-r_3\ge3/206$. In D units,
\[
\begin{aligned}
 J_{R1}(\kappa)-J_{R1}(0)&=\int_{R1}\frac{\kappa\,dr}{r_1r_2r_3\,a_0(a_0-\kappa)}\\
&\le\frac{\mathrm{vol}(R1)\,\kappa}{y^3\,\tfrac3{206}(\tfrac3{206}-\kappa)}\\
 &=5.74468\ldots\cdot10^{-6}<5.745\cdot10^{-6}
\end{aligned}
\]
For $\kappa=10^{-6}$,
$\mathrm{vol}(R1)=2.53345\ldots\cdot10^{-6}$. Its exact rational value
is the absolute determinant of three edge vectors at a vertex divided by $6$.
This bounds the actual R1 increment, which is already included in $d_{T8}$.
The extra allowance $g_{R1}=5.71\cdot10^{-6}$ in Section~6 is a separate
conservative subtraction; it is not used as a bound for this increment.

{
\normalsize
\renewcommand{\arraystretch}{1.14}
\begin{longtable}{@{}p{\textwidth}@{}}
\toprule
\emph{Exact outward enclosures} \\
\midrule
\textbf{Prefix, n=600} \\
$P\in\CertificateInterval{3513.1687300568893335013223065223679611140065161989243}{3513.1687300568893335013223065223679611140065161989248}$ \\
\emph{Source code SHA-256}: \path{s2-selected/sources/numeric/split_prefix_exact.py}: \path{492ef149b0460f59732102ddeb252fa8fb30fce03df2fd953376ad4fb5417ada} \\
\emph{Output SHA-256}: \path{s2-selected/sources/numeric/prefix-u-n600-corrected.json}: \path{01c72d03c23ff51eb903e299ecb63a90496485c0e6549e05c3818077f2e88e08} \\
\midrule
\textbf{PF outer, n=1536} \\*
$J\in\CertificateInterval{0.352576588183432872639904522863463727436871987468524955851}{0.359440819140453278302418897934329259372468992946659201284}$ \\
\emph{Source code SHA-256}: \path{s2-selected/sources/numeric/exact_outer_s2.py}: \path{1cf453f317366d9c0b191aac5457b518fae48b54d7321da80bf577d1aa5f69f0}; \path{s2-selected/sources/numeric/arb_pf_replay.py}: \path{e99d81397544d0814b49c8263d61a4a9389dc6688488f4ee0359cc1fc76bfcb2} \\
\emph{Output SHA-256}: \path{s2-selected/sources/numeric/outer-n1536-truepsi.json}: \path{431d567fe714e412da6ba528b02cab0e251cbb616aa1fd8d9c8a11efca67d323} \\
\midrule
\textbf{BV corner, n=2800} \\
$C\in\CertificateInterval{-0.0574908755554524530180814901226778275854309897822443970870}{-0.0570237847462003399142149129742403275854309897822443970868}$ \\
\emph{Source code SHA-256}: \path{s2-selected/sources/numeric/corner_arb.py}: \path{7d5586c43d7951cecd86fdb85fca277efc91e7cf9b85fbd75ada5fb01726d861} \\
\emph{Output SHA-256}: \path{s2-selected/sources/numeric/corner-arb.json}: \path{2bc56db0d72a56afbd88aac7111afc2f2eaa06aa65a4c509a96107765415afc3} \\
\midrule
\textbf{Tail above 12, in D} \\
$[0,0.000047391848857364673265889344422276912]$ \\
\emph{Source code SHA-256}: \path{s2-selected/certify_tail_seed.py}: \path{42565c0188066a45c1f166a49a8fb64c95d6bd43874047bade6af9f856428fef} \\
\emph{Output SHA-256}: \path{s2-selected/TAIL-SEED.json}: \path{0363a93ec43884b70ca0b54e06e71ea8a64210c64910e4e4a3bc3a366df04c53} \\
\midrule
\textbf{Rowwise prefix-repair allowance, in D} \\
$[0,1.154\cdot10^{-54}]$ \\
\emph{Justification}: \\
The preceding tail majorant, within the single $1/20000$ allowance \\
\midrule
\textbf{Joint lower-value assembly} \\
\emph{Outward enclosure assembled from the printed inputs in (A.7)}: \\
$\CertificateInterval{1.6195813569890835255913340433210572922007671408165057072212979}{1.6195813569890835255913340433210572922007671408165057072212983}$ \\
\emph{Full-precision assembler source SHA-256}: \path{s2-selected/verify_interface.py}: \path{fae4b0cb684502ed7d1d75537cb5456868ba83094eb9157f14e30cf58574de71} \\
\emph{Output SHA-256}: \path{s2-selected/INTERFACE-VERIFICATION.json}: \path{839dbd7f3d75d39f52aca3920325d63dd33dcc4346479f0871720d93194fc469} \\
This output uses the retained full-precision prefix endpoint, not the shortened $P_-$ in (A.7); the printed-input enclosure above is evaluated directly from (A.7). \\
\bottomrule
\end{longtable}
}

\paragraph{Independent quadrature.} The calculation recorded in
\path{s2-joint/outputs/ASSEMBLY.json},
SHA-256
\path{21083883403f3903fff3ce781e086a26f931dd809e445578cfcba546295e9e5d},
independently evaluates the PF and corner integrals in the \S~3 functional and gives
\[
\begin{aligned}
 \frac{T_8}{4}&\in
 [1.62063469556829,1.62071322785677]\quad\text{D units},\\
 1.62063469556829-d_{T8}&>0.001053338579.
\end{aligned}\tag{A.8a}
\]
Interval (A.8a) covers the continuous functional and the infinite tail.
The selected $d_{T8}$ uses coarser adverse quadrature endpoints
and also includes the $1/20000$ allowance.
The three improvements are
\[
\begin{aligned}
 \frac{2e^{-\gamma}}4(J_{+,\rm old}-J_{+,\rm new})
   &\simeq0.000944946137955914,\\
 \frac{C_{-,\rm new}-C_{-,\rm old}}4
   &\simeq0.000058392441251611,\\
 \text{tail allowance}&=0.000050000000000000.
\end{aligned}
\]
Here $J_{+,\rm old}$ and $J_{+,\rm new}$ are the printed upper PF
endpoints. Similarly, $C_{-,\rm old}$ and $C_{-,\rm new}$ are the
printed lower corner endpoints. For the new ball we take its outward lower
endpoint. Because of endpoint display rounding, the sum differs from
$1.62063469556829-d_{T8}$ by less than $2\cdot10^{-15}$.
The difference exceeds $0.001053338579$,
so the continuous lower endpoint exceeds $d_{T8}+1/20000$.
The independent prefix-only assembly gives $P_{\rm def}<1.1\cdot10^{-51}$.
The coarser joint prefix-and-tail bound is $10^{-49}$.
The rowwise repair allowance $1.154\cdot10^{-54}$ in (A.7) covers its
own tail-majorant correction and is already inside the single $1/20000$
allowance used in $d_{T8}$. These independent deficits are not added again.
Section~\ref{sec:proof-theorem-1} bounds the finite-$N$ corrections,
including the selector and AA37 terms.
The selected value is (A.8).

The component sources are:

{
\normalsize
\renewcommand{\arraystretch}{1.14}
\begin{longtable}{@{}p{\textwidth}@{}}
\toprule
\textbf{True F/f, 320-bit Arb Taylor models} \\
\emph{Source code SHA-256}: \path{s2-joint/code/ff_models.py}: \path{7a6b137417573ba0e717ffa4444e7fb6671cb1b325ff826cac26630591bd4624} \\
\emph{Output SHA-256/result}: \path{s2-joint/outputs/ff_report.json}: \path{cfb7d637cb623590478e3069b14a2fa2b85aeab1eb296985523b1ccf60f56f01}; $E(s)\le1.302\cdot10^{-55}$ for all $s\ge33$ \\
\midrule
\textbf{PF directed quadrature, 24,000,005 leaves} \\
\emph{Source code SHA-256}: \path{s2-joint/code/jpf.c}: \path{e5a4680d2a6ab4bed93bf8628dae78acabc79bd66bbfcbc6a526ca03cf564973} \\
\emph{Output SHA-256/result}: $J_{PF}\in[0.355795040748992,0.356074784134854]$; four shard output hashes below \\
\midrule
\textbf{BV corner, rigorous nested Arb integration} \\
\emph{Source code SHA-256}: \path{s2-joint/code/corner.py}: \path{7cbc5d276b54b52d657fc52622db1a3542c5079d72a39e3be3d12eec1ddb21f2} \\
\emph{Output SHA-256/result}: \path{s2-joint/outputs/corner.json}: \path{360841b6008cb08be571d23177887df14be286ebde7ba3f3fe2c7571f31be9d8} \\
\midrule
\textbf{Joint directed assembly, including prefix deficit and full tail} \\
\emph{Source code SHA-256}: \path{s2-joint/code/assemble.py}: \path{4b521ae19ef7742e3bcbb4eb28a0f48c2a161195c5ab0e9b4f8456ba6ddf5c45} \\
\emph{Output SHA-256/result}: \path{s2-joint/outputs/ASSEMBLY.json}: \path{21083883403f3903fff3ce781e086a26f931dd809e445578cfcba546295e9e5d}; (A.8a) \\
\bottomrule
\end{longtable}
}

All paths in this table are relative to the supplement root.
The hashes for \path{s2-joint/outputs/jpf_shard0.json} through \path{s2-joint/outputs/jpf_shard3.json} are,
in that order,
\path{14b67605807cfcef9ec96b134fe50b9c8e4b4319c3f33a5ebf99a8eab0cf3c67},
\path{d163c9c8d8d90dbfa5603b49f379867df0912a87e088f0e71924300500bf0fee},
\path{e6f4077dcc2c05317eb2a1b26e4bba519dc7b98cb1960038a2a4885a8bb1f1de},
and
\path{623eebd522d2b99122cf4395cfe760b0d10cc89ec31ac5ca8a70009e8c729b6d}.
The bundle manifest SHA-256 is
\path{d26488e736e5798544e758338acb299d417b74a309568f5ce20f738e399cf96d},
and the manifest contains 16 entries.

The independent assembly gives $P_{\rm def}<1.1\cdot10^{-51}$.
It also bounds the
no-growth tail after $s=33$, and checks the classical A=1 tail argument.
Its PF and corner enclosures are independent of the primary quadrature;
the prefix term $P$ and full volume $V$ are shared inputs.
That section proves the source comparison and the Class-II bounds.
It also bounds the selector, the AA37/G2 terms and the other errors.
The joint four-row endpoint enters once.

\subsection{G1: signed allocation and radial multiplier}

For the exact allocation in \S~4,
\[
 K_G=\frac{2I_{11}+I_{12}+I_{13}}{2e^\gamma},\qquad
 P_r:=-K_GT_r,\qquad \sum_{r=0}^3T_r=\overline G.
                                                               \tag{A.9}
\]
Here $P_r$ is a pure profile target. The corresponding arithmetic
contribution is $Q_{G1-r,0}(N)$; its lower comparison, including the
fixed and vanishing errors, is (4.85). The allocation identity (A.9)
does not assert equality with that arithmetic contribution.
The multiplicities $(2,1,1)$ enter $K_G$ once. The source code is
\path{g1-endpoints/certify_signed_endpoints.py},
with SHA-256 \path{e0eb449e97eedc1ea5aa8ab40de847a1baec1e8316271bf873ead904d845a864}.
Its output \path{g1-endpoints/outputs/result.json}
has SHA-256
\path{0a940b64e5f762fbfc1439709a3092de9749ecfdf448ef467ff52248e1794c69}.
A separate exact rational calculation verifies the retained source intervals.
Its output is
\path{g1-rational/exact-rational-check.json}, with SHA-256
\path{7d7048c800af14d06963bd8fcf4f17ff45517e41fa9fd682bc583ddfb88dd4b6}.
The outward decimal enlargements of these intervals and their product are displayed below.

{
\normalsize
\renewcommand{\arraystretch}{1.14}
\begin{longtable}{@{}p{\textwidth}@{}}
\toprule
\emph{Exact rational outward intervals} \\
\midrule
\textbf{$I_{11}$, including four-prime correction} \\
$\CertificateInterval{0.2208790666923091612026153538231208031879235648368125703408}{0.2208790666923091612026153538231208031879235648368125703424}$ \\
\midrule
\textbf{$I_{12}$} \\
$\CertificateInterval{0.5740619734452483407647685136781950850886922314288474504985}{0.5740619734452483407647685136781950850886922314288474504987}$ \\
\midrule
\textbf{$I_{13}$} \\
$\CertificateInterval{0.4052907329856379464660009714315655397835085202678768274511}{0.4052907329856379464660009714315655397835085202678768274513}$ \\
\midrule
\textbf{$K_G$} \\
$\CertificateInterval{0.3989480791070578466390443009830933753748442947478110112211}{0.3989480791070578466390443009830933753748442947478110112223}$ \\
\midrule
\textbf{$\overline G$} \\
$\CertificateInterval{6.0884403625974206828946573628908904359431953974002641181857}{6.0884403625974206828946573628908904359431953974002641181858}$ \\
\midrule
\textbf{Joint $-K_G\overline G$, in D} \\
$\CertificateInterval{-2.4289715874161197464056136533483979571384269074906618794340}{-2.4289715874161197464056136533483979571384269074906618794266}$ \\
\bottomrule
\end{longtable}
}

These intervals certify G1-0 through G1-3.
All four targets are in D units.
The exact row endpoints below use the narrower exported dyadic $K_G$
and $T_r$ intervals from the retained source.

{
\normalsize
\renewcommand{\arraystretch}{1.14}
\begin{longtable}{@{}p{\textwidth}@{}}
\toprule
\textbf{G1-0}\quad K endpoint for lower bound: upper \\
\emph{Exact lower endpoint}: $\frac{-19564356883289720593180343123839226092697830325553783422353}{98079714615416886934934209737619787751599303819750539264}$ \\
\emph{Exact upper endpoint}: $\frac{-626059420265271058981770979962855234966330570417721069513945}{3138550867693340381917894711603833208051177722232017256448}$ \\
\midrule
\textbf{G1-1}\quad K endpoint for lower bound: lower \\
\emph{Exact lower endpoint}: $\frac{179266388678009569180357614458786189164679660590048624716243}{196159429230833773869868419475239575503198607639501078528}$ \\
\emph{Exact upper endpoint}: $\frac{717065554712038276721430457835144756658718642360194498866525}{784637716923335095479473677900958302012794430558004314112}$ \\
\midrule
\textbf{G1-2}\quad K endpoint for lower bound: upper \\
\emph{Exact lower endpoint}: $\frac{-821300653301490295585240971954986923239672986567876292003797}{392318858461667547739736838950479151006397215279002157056}$ \\
\emph{Exact upper endpoint}: $\frac{-821300653301490295585240971954986923239672986567876292002013}{392318858461667547739736838950479151006397215279002157056}$ \\
\midrule
\textbf{G1-3}\quad K endpoint for lower bound: lower \\
\emph{Exact lower endpoint}: $\frac{270036186059109561485627059871971764166199228966336131972949}{196159429230833773869868419475239575503198607639501078528}$ \\
\emph{Exact upper endpoint}: $\frac{135018093029554780742813529935985882083099614483168065986769}{98079714615416886934934209737619787751599303819750539264}$ \\
\bottomrule
\end{longtable}
}

The exact joint reference value is given below.
{\normalsize
\[
 L_{G1}=-\frac{243950428265194656254206921939947762868871361874409552095097}
 {100433627766186892221372630771322662657637687111424552206336}.
                                                               \tag{A.10}
\]
}
Let $R_{G1}$ be the sum of the four exact row lower endpoints in the
table. The exact difference is
{\normalsize
\[
\begin{gathered}
 L_{G1}-R_{G1}\\
 =\frac{500103}
 {100433627766186892221372630771322662657637687111424552206336}\\
 <4.98\cdot10^{-54}.
\end{gathered}
\]
}
Put
\[
 \delta_{G1}=4.98\cdot10^{-54},\qquad
 E_{G1}=L_{G1}-\delta_{G1}.
\]
The printed decimal intervals are rounded outward separately from the
narrower source intervals used for the exact row certificates.
For an exact check of the endpoint arithmetic, enclose $1/(2e^\gamma)$ by
{\small
\[
\CertificateInterval{0.2807297417834425849120716073954403933828551934625765840770795380225439}{0.2807297417834425849120716073954403933828551934625765840770795380225440}.
\]
}
This rational enclosure comes from the Arb constant evaluation of
the retained source, at 256-bit precision.
We multiply its lower endpoint by $2I_{11,\rm lo}+I_{12,\rm lo}+I_{13,\rm lo}$,
and its upper endpoint by $2I_{11,\rm hi}+I_{12,\rm hi}+I_{13,\rm hi}$.
Exact rational multiplication and outward rounding to $10^{-58}$ then give
the displayed $K_G$ interval.
Its upper endpoint exceeds the upper product by less than
$9.17\cdot10^{-59}$.
The joint interval also contains the displayed interval product.
Put $K_{\rm up}=K_{\rm hi}$ from \eqref{eq:g1-khi}.
Exact multiplication gives
\[
 0<L_{G1}+K_{\rm up}\overline G_{\rm hi}
 <9.36\cdot10^{-58}<\delta_{G1}.
\]
Thus $E_{G1}<-K_{\rm up}\overline G_{\rm hi}$. The exact difference
also gives $E_{G1}<R_{G1}$.
The rational row endpoints keep their independent, sharper certificates.
The selected endpoint is the joint one, $E_{G1}$.
Section~\ref{sec:proof-theorem-1} includes $L_{G1}$ once,
and its error sum includes $\delta_{G1}$ once.
The four row values certify this comparison.
The final sum uses this joint lower endpoint.
The primitive (A.10c) is an earlier alternative.
No completed validated complex quadrature check of its analytic
neighborhoods and branches is used for the selected endpoint.

A separate implementation computes the four-prime correction in two ways
and checks the same $K$ directions. Its source \path{endpoint-recert/code/g1_K.py}
has SHA-256
\path{fd749f9a795f6876c9852c516789a9bbf937ec667c76a06f5662fed480aab52e};
its output \path{endpoint-recert/out/g1_K.json} has SHA-256
\path{e3ea5971a526c195ea33e7eb6b7b0c756d77fa1e0560d3b83eedf1785536cb07}.
The output confirms both $K$ directions. Its residual source \path{endpoint-recert/code/g1_J.py}
has SHA-256
\path{e7e6f75367c3af5d2e97a0c653d866328df962caef8791305203a6ee18be5715};
the output \path{endpoint-recert/out/j_final.json} has SHA-256
\path{7ad65327cda9adf324a5325b664ce3788d6119ea40f31c0c007eb2416f42069c}.
It proves the exact-rational residual enclosure, rounded outward:
\[
 J\in[0.964498,0.967120],\qquad J<0.968435989360.                                \tag{A.10a}
\]
The residual-profile $\tau$-effect is at most
$5\cdot10^4\tau\le5\cdot10^{-86}$. The selected residual slack absorbs it.
The comparisons give all four T allocations and lower targets,
with signs $(-,+,-,+)$.
The $K$ directions are
$(\mathrm{upper},\mathrm{lower},\mathrm{upper},\mathrm{lower})$,
with the weights $(2,1,1)$ entering once. Section~4 proves the source application for
$C_b=334$ and its error bounds.

\paragraph{The 4096-cell upper-corner calculation.}
This paragraph gives the 4096-cell trapezoidal calculation of $\overline G$ used in (4.79).
The residual enclosure (A.10a) comes from the separate Taylor-model source and output identified above.
The true delay functions are

\[
 F(s)={2e^\gamma\over s}\ (1\le s\le3),\quad
 f(s)=0\ (s\le2),\quad
 f(s)={2e^\gamma\log(s-1)\over s}\ (2\le s\le3),
\]

and thereafter $(sF)'=f(s-1)$, $(sf)'=F(s-1)$.
On open integer intervals, their second derivatives are bounded by 12.
The quadrature uses trapezoidal panels of width $2^{-16}$ aligned at integers,
with integral remainder at most $h^3$.
Arb carries all propagated radii, and outward rational rounding uses the
$2^{-128}$ lattice.
Monotonicity encloses $\Phi(a)=\int_1^a(1-f(s))\,ds$. Beyond 30, its tail is bounded by
$(a-30)(1-f(30))$; no constant table extrapolation is used.

For fixed rational $t_1,t_2$, the $t_3$ range is split at each kink,
activation switch and pairwise affine crossing in (4.70).
On each segment $a_3=c/t_3+k$, and

\[
 \int_p^q{1-f(c/t+k)\over t^2}\,dt
     ={\Phi(c/p+k)-\Phi(c/q+k)\over c}\quad(c\ne0).
 \tag{A.10b}\label{eq:g1-phi-primitive}
\]

The constant-argument case is direct. Each segment integral lies
between 0 and $1/p-1/q$. The residual
$\Theta_*-t_1-t_2-t_3$ is nonincreasing in $t_1,t_2$.
Its affine derivatives are negative, and the $\psi$ active sets are down-sets.
Hence the inner integral is nondecreasing in both outer variables.
The outer integration uses $t=ue^x$, outward-enclosed logarithmic cells, and upper
and lower corner sums. Diagonal triangles have half the squared cell width.
On a grid of 4096 cells, the outward-rounded sums are accumulated as exact
integers. This proves the directed enclosure (4.79).

For an independent evaluation of the four-prime correction, put
$z=v/(1-t)$. A primitive of
$\log((1-2z)/z)/(z(1-z))$ is

\[
\begin{aligned}
 P(z)={}&-\operatorname{Li}_2(2z)-\tfrac12\log^2z
 -\tfrac12\log^2(1-z)-\log2\log(1-z)\\
 &-\operatorname{Li}_2(1/[2(1-z)])
       +\log z\log(1-z)+\operatorname{Li}_2(z).
\end{aligned}
 \tag{A.10c}\label{eq:g1-dilog-primitive}
\]

The formula is checked by differentiation. The analytic neighborhoods
and the real branches would require a validated complex-quadrature check
before this primitive could serve as an endpoint certificate.
Nested integration of the original double integral gives an
independent enclosure. At 200-bit precision, the signed enclosures are:

{
\normalsize
\setlength{\tabcolsep}{3pt}
\begin{longtable}{@{}p{\textwidth}@{}}
\toprule
\textbf{$I_{11}$} \\
\emph{Abbreviated centre}: $0.22087906669230916120261535382312080318$ \\
\emph{Width less than}: $1.476\cdot10^{-57}$ \\
\midrule
\textbf{$I_{12}$} \\
\emph{Abbreviated centre}: $0.57406197344524834076476851367819508508$ \\
\emph{Width less than}: $6.348\cdot10^{-59}$ \\
\midrule
\textbf{$I_{13}$} \\
\emph{Abbreviated centre}: $0.40529073298563794646600097143156553978$ \\
\emph{Width less than}: $3.672\cdot10^{-59}$ \\
\midrule
\textbf{$K_G$} \\
\emph{Abbreviated centre}: $0.39894807910705784663904430098309337537$ \\
\emph{Width less than}: $8.588\cdot10^{-58}$ \\

\bottomrule
\end{longtable}
}


These centres are not directed endpoints. The source-derived widths above
refer to the narrower retained intervals. For the printed decimal
integral intervals, the exact outward upper endpoint is

\[
 K_{\rm hi}=
 0.3989480791070578466390443009830933753748442947478110112223.
 \tag{A.10d}\label{eq:g1-khi}
\]


{
\normalsize
\setlength{\tabcolsep}{3pt}
\begin{longtable}{@{}p{\textwidth}@{}}
\toprule
\textbf{G1-0} \\
\emph{Abbreviated centre}: $-199.47403966257515881115583710023774090$ \\
\emph{$K_G$ direction}: upper \\
\midrule
\textbf{G1-1} \\
\emph{Abbreviated centre}: $+913.88106797076246277430172018777105602$ \\
\emph{$K_G$ direction}: lower \\
\midrule
\textbf{G1-2} \\
\emph{Abbreviated centre}: $-2093.45188381184444127509736016826295322$ \\
\emph{$K_G$ direction}: upper \\
\midrule
\textbf{G1-3} \\
\emph{Abbreviated centre}: $+1376.61588391624101756554586342738124014$ \\
\emph{$K_G$ direction}: lower \\

\bottomrule
\end{longtable}
}



The narrower source-derived two-sided joint width is below $5.248\cdot10^{-57}$.
The exact gap is the fraction displayed after (A.10).

\subsection{Switched mass, BV endpoints and finite threshold}

Write $H_L:=I_{HL}$ as in (5.32), equivalently (1.12), with the parameters
and cell weight (5.30)--(5.31). The measure is the coordinate measure $dt_1\cdots dt_{k-1}$,
where $t_k=1-\sum_{i<k}t_i$.
It carries no factorial or Euclidean surface-measure factor.
Then

\[
 0\le H_L<0.0778506,\qquad -2H_L>-0.1557012.            \tag{A.11}
\]
Both implementations cover every positive-weight stratum $5\le k\le11$,
and the coverage is exact rational.
Their upper endpoints lie in
\[
 [0.0778505820459738841964296666351,
   0.0778505820459738841964296666353].                   \tag{A.12}
\]
The selected source is
\path{switched-mass/original/scripts/refine.py}, with SHA-256
\path{83ff0a33baf95be8400a69c466965baae130b3cd788acd2d83b83fd2cc45398b}.
The separate calculation encloses the same switched mass.
The independent source
\path{switched-mass/independent/independent_replay.py}
has SHA-256
\path{c275b1904b506cf82315d4e27f44f1949487e67044923501fd0628738a7744ad}.
Their common compact output
\path{switched-mass/replay_summary.json}
has SHA-256
\path{f0dac63f31707d0ecfe9ffd08fdfe21d4e5b2718b19a83af713ba518d2077e06}.
Exact rational coverage and overlap identities give the scalar input for
the switched term of \S~5.

For the S6 BV row, put
\[
 J(\kappa)=\int_{\tau_S}^{1/3}\frac{dt}{t(1/2-\kappa-t)}.
\]
Exact logarithms and independent Arb arithmetic give the enclosures.
{\normalsize
\[
\begin{gathered}
 J(0)\in\\\CertificateInterval{1.3148581959157321533513247276782285310758770468886237619584}{1.3148581959157321533513247276782285310758770468886237619585},\\
 J(10^{-6})\in\\\CertificateInterval{1.3148646802099601102520747017315401230927649569798117109688}{1.3148646802099601102520747017315401230927649569798117109689},\\
 -[J(10^{-6})-J(0)]\in\\
 \CertificateInterval{-6.48429422795690074997405331159202\cdot10^{-6}}{
 -6.48429422795690074997405331159201\cdot10^{-6}}.
 \end{gathered}\tag{A.13}
\]
}
Here $S_6=4J$, so the nominal D contribution is $-J(0)$.
The increment in (A.13) enters once at level $1/2-10^{-6}$ of
\S~5. The source
\path{row-enclosures/cert_s4_s7.py}
has SHA-256
\path{d98866b45f1ddb8f3f93411bfe09cb11c62569cc80fa775a77a3eab1ad4384aa};
its \path{row-enclosures/result.json} has SHA-256
\path{0eb5fa52ed180a071633e50f0d099f5c4e1f9bb001add71765fe1c2f033b8f7a}.
The independent 256-bit calculation is
\path{bv-increment/checks/charge_check.py},
with SHA-256
\path{1a7be201efe55e801c2ea7c2dc0807ac543d323838165f259b982a064e460ac6},
and its output \path{bv-increment/checks/charge_check.out}
has SHA-256
\path{e99ef910bb85c1bd2119314c37a0deddcf1bc11f43c1d4fe50e8aaf956748015}.
It verifies the exact-log enclosures in (A.13).

\paragraph{S5 BV-top normalization.} The moved-cell value is
enclosed by
\[
 B(0.03658166423843258896399704402764682632822246022922592990,
   8.84\cdot10^{-57}).                                 \tag{A.14}
\]
This raw row is $4J_5(0)$, with master coefficient $-1$ by
(5.15)--(5.20a). Its D contribution is $-J_5(0)$, with exact interval
{\normalsize
\[
\begin{gathered}
 -J_5(0)\in\\
-\frac14 B\left(\begin{gathered}0.03658166423843258896399704402764682632822246022922592990,\\
             8.84\cdot10^{-57}\end{gathered}\right),
\end{gathered}\tag{A.14a}
\]
}
which is about $-0.00914541605960815$ D. Section~5 proves the
normalization. S5-low uses its direct-cell endpoint, and S5 BV-top uses
the moved-cell exact-log value (A.14).
Their supports are disjoint at $t_*$,
and each contribution enters once.

For comparison, if S5-top used $\kappa=10^{-6}$, its negative coefficient
would require an upper bound for $J_5(\kappa)$. Since $t\ge t_*$ and
$1/2-t\ge1/2-\tau_S$,
\[
\begin{aligned}
 0\le J_5(\kappa)-J_5(0)
 &=\int_{t_*}^{\tau_S}
   \frac{\kappa\,dt}{t(1/2-t)(1/2-t-\kappa)}\\
 &\le\frac{(\tau_S-t_*)\kappa}
 {t_*(1/2-\tau_S)(1/2-\tau_S-\kappa)}\\
 &=\frac{25088}{672019919285}<3.734\cdot10^{-8}<4\cdot10^{-8}.
\end{aligned}
\]
Thus $-J_5(\kappa)\ge-J_5(0)-E_{5,\kappa}$,
with $E_{5,\kappa}=4\cdot10^{-8}$ in D units.
Section~\ref{sec:proof-theorem-1} subtracts this error once.
The separate conservative BV allowance is
$g_{\rm BV}=5.75\cdot10^{-6}=E_{5,\kappa}+g_{R1}
=4\cdot10^{-8}+5.71\cdot10^{-6}$.
The R1 increment itself is already inside the selected S2 endpoint.
Equations (A.14)--(A.14a) give the nominal $\kappa=0$ value.

For an independent exact-rational S5-top certificate, the logarithmic
ratio in (5.20) is
$R=392787/390995$, so $z=(R-1)/(R+1)=896/391891$. Define
\[
 P_5=16\sum_{k=0}^{15}\frac{z^{2k+1}}{2k+1},\qquad
 R_5=\frac{16z^{33}}{33(1-z^2)}.
\]
Then $P_5<4J_5(0)<P_5+R_5$ by the positive atanh series.
Exact rational comparisons give $R_5<3.447\cdot10^{-88}$ and
{\normalsize
\[
\begin{gathered}
 [P_5,P_5+R_5]\subset\\
B(0.03658166423843258896399704402764682632822246022922592990,
   8.84\cdot10^{-57}).
\end{gathered}
\]
}
This proves the directed $-1/4$ normalization. The printed rational
series is the certificate. Section~6.3 proves its source application
and arithmetic error bounds.

The finite-epsilon target is
\[
\begin{aligned}
 T&=8\log\frac{50009}{39991},\\
 1.78838848345960807361670991332&<T
 <1.78838848345960807361670991333.
\end{aligned}\tag{A.15}
\]
It is proved by exact rational arithmetic. For
$z=5009/45000$, set
\[
 P=16\sum_{k=0}^{23}\frac{z^{2k+1}}{2k+1},\qquad
 R=\frac{16z^{49}}{49(1-z^2)}.
\]
Then $P<T<P+R$, and exact rational comparisons give (A.15).
This series and its tail bound certify (A.15),
and its upper endpoint enters the final strict comparison.

\subsection{Selected numerical endpoints}

The table below lists the selected values.
The full endpoints and proofs are in A.2--A.6.
\begin{longtable}{@{}p{9em}p{\dimexpr\textwidth-9em-4\tabcolsep\relax}@{}}
\toprule
Contribution & Selected value and direction \\
\midrule
S1 & Joint lower endpoint (A.4), with multiplier $3/4$ \\
S2 & Joint lower value (A.8), including the $1/20000$ tail \\
S3/S8/S9 & Absolute Option-A endpoints; signs $-,+,+$ \\
S4/S5-low/S7 & Direct-row upper endpoints; multiplier $-1/4$ \\
S10-Ret & Retained-row lower endpoint; multiplier $1/4$ \\
G1 & $E_{G1}=L_{G1}-\delta_{G1}$, once \\
S5-top & Raw-to-D factor $-1/4$; separate error $E_{5,\kappa}$ \\
S6 & Nominal value $-J(0)$ and increment (A.13) \\
Switched term & $-2H_L>-0.1557012$ by (A.11) \\
Threshold & Upper endpoint (A.15) \\
\bottomrule
\end{longtable}
Sections~2--5 prove the source applications,
and Section~6.3 combines the terms with the signs of (6.1).
}

\section{Auxiliary estimates}

\subsection{Repeated-prime terms}

This subsection shows that terms with a repeated large prime are negligible.
Fix $0<\sigma\le10^{-10}$ and put
\[
 \tau=\sigma^9,\qquad
 \ell_0=53/412-\eta_L/2=264897/2060000,\qquad
 L_N=\log N,\quad u_0=N^{\ell_0\sigma^2}.
\]
Only the normalization $\mathcal W=4(1-\varepsilon_A)U_C$ of (4.2) is used.
The signed sum is that of Theorem~\ref{thm:signed-assembly} and Section~6.3.

\subsubsection*{Blocks, masks and signs}

A block $c$ fixes a row, a literal parent prefix cell
and, if present, its whole-cell selector.
It also fixes an AA37 sign $\nu$,
and finally $R_1^\nu$ or $\mathbf D\in\mathcal D^\nu(L_c)$.
The labelled prefix intervals are clipped sets $[N^{t_i},N^{\rho t_i})$.
Put $L_c=N^{\ell_c}$, where $\ell_c\ge\ell_0$, and $u_c=L_c^{\sigma^2}$.
Let $m_c$ count all labelled rough positions, prefix and tail together.
Put $M_c=\prod_h|G_h|!$ and write $C_c(q)$ for the unrestricted coefficient.
The nineteen rows and their coefficients are
\[
\begin{array}{c|c|c}
 \text{rows}&c_r&T^{(r)}\\ \hline
 \mathrm{S1\hbox{-}0\hbox{--}3}&3,-3,3,-3&A_D\\
 \mathrm{S2\hbox{-}0\hbox{--}3}&\pm1&A_D\\
 \mathrm{S3\hbox{-}A}&-2&A_D\\
 \mathrm{S4,S5\hbox{-}low,S7}&-1&A_D\\
 \mathrm{S8\hbox{-}A,S9\hbox{-}A,S10\hbox{-}Ret}&+1&A_D\\
 \mathrm{G1\hbox{-}0\hbox{--}3}&2A_R=\pm1&A_G^{\rm cut}
\end{array}
\]
There are fifteen DIRECT rows and four G1 rows (\S\S~2--3, \S~4.2).
On DIRECT rows, $c_r$ includes the O1 coefficient and the Buchstab sign.
The G1 weights $2,1,1$ are part of $A_G^{\rm cut}$ in (4.16).

The masks are those of C1 on DIRECT rows and (4.22)--(4.23) on G1.
Define
\begingroup\small
\[
\begin{aligned}
 \Xi_{U,c}(q)&=\frac1{M_c}\sum_{\rm full}\lambda_d(1-K_TK_P)
              \mathbf1_{q=d\prod_jp_j},\\
 \Xi_{A,c}(q)&=\frac1{M_c}\sum_{\rm full}\lambda_dK_TK_P
              \mathbf1_{q=d\prod_jp_j}\mathbf1_{q\ \rm not\ sf}.
\end{aligned}
 \tag{B.1}
\]
\endgroup
The full sum runs over $d$ and every labelled prime tuple.
Call $c$ active if some realization has $\lambda_d\ne0$.
Otherwise both coefficients vanish, including for empty intervals.
On DIRECT rows, $\Xi_{U,c}=\Xi_T/M_c+\Xi_P$, because
$1-K_TK_P=(1-K_T)+K_T(1-K_P)$ is a disjoint decomposition.
For S1-2, S1-3, S2-2, S2-3, S8-A, G1-2 and G1-3,
orbit invariance gives $O_c=C_c/M_c-\Xi_{U,c}$.
For the other rows the same identity holds, since $M_c=1$.
Thus $O_c\mathbf1_{q\ \rm not\ sf}=\Xi_{A,c}$ in every row.

Define $\Delta_0^{(r)},\Delta_P^{(r)}$ by (4.21).
By C1 and (4.23), (4.26),
\[
 \mathcal E_{\rm SQ}
 =\frac1{\mathcal W}\sum_c c_{r(c)}\nu_c
  \Big[-\sum_q\Xi_{U,c}(q)\Delta_P^{(r)}(q)
       +\sum_q\Xi_{A,c}(q)(\Delta_0^{(r)}-\Delta_P^{(r)})(q)\Big].
 \tag{B.2}
\]
Here $\nu_c=\pm1$ is the AA37 remainder sign.
Both terms use the original tested measure, before HB expansion.
The first contains the forbidden-tail and forbidden-prefix parts,
the second the allowed repeats.
Only the C1 $\gamma$-term (Class II) uses HB expansion,
dyadic shells, short boxes and $K_0$ of (4.66).
The terms $NC_{\rm sf}$ and $\sum W_c(\Delta_0-\Delta_P)$ are treated separately in (4.65).
Both vanish on DIRECT rows.

\subsubsection*{Coefficient and test bounds}

The first lemma records the support, roughness and finiteness facts used below.

\textbf{Lemma B.1.} For sufficiently large $N$, depending on $\sigma$:
\begin{enumerate}
\item Every nonzero-weight full realization has $q<Q:=N^{3/5}$ and $(q,2N)=1$.
\item $u_c\ge u_0\ge2$, and every labelled rough coordinate is at least $u_c$.
\item On active blocks,
$m_c\le M_\sigma:=\lceil3/(5\ell_0\sigma^2)\rceil$.
Hence $m_c!/M_c\le M_\sigma!=C_\sigma$, the constant of (4.25).
\item Each row has finitely many blocks $N_{{\rm lab},r}(\sigma)$, independently of $N$.
\end{enumerate}

\emph{Proof.} The final cubic constraint and decreasing labels give $\prod_iD_i<L_c$.
For a single coordinate, use $D_1<L_c^{1/2}$; the empty product is $1$.
Let $T_c$ be the sum of prefix lower exponents.
The literal intervals and $d<L_c^\sigma$ give
\[
 q<N^{\rho T_c}L_c^{\rho+\sigma}
 \le N^{(\rho+\sigma)(T_c+\ell_c)}
 \le N^{(153/256)(1+\tau+\sigma)}<N^{3/5}.
\]
Here $T_c+\ell_c=\Theta_c-\eta_L\le153/256$; use (4.20) for G1.
This covers $R_1^\nu$ before both masks, with empty tail.
All primes lie in $\mathcal P(N)=\{p>2:p\nmid N\}$,
and $d\mid P(u_c;N)$, so $(q,2N)=1$.

Next, $u_c=N^{\ell_c\sigma^2}\ge N^{\ell_0\sigma^2}=u_0$.
Tail intervals start at $D_i\ge L_c^{\sigma^2}$.
Prefix coordinates satisfy $p\ge N^{u/\rho}>u_c$ by \S~2 and (4.20).
For any active block, choose a nonzero-weight realization.
Then $u_0^{m_c}\le q<N^{3/5}$, so $m_c\ell_0\sigma^2<3/5$.
The argument does not use $C_c\ne0$.

Each prefix coordinate has at most
$G_\sigma:=1+\lceil\log(\rho/u)/\log\rho\rceil$ grid labels on $[u/\rho,1]$.
The endpoints
$u,\alpha,\beta,\mu,\lambda,1/4,t_*,\tau_S,1/3,t_{\rm dd},v_\beta$
split any such grid interval into at most twelve literal pieces.
Clipping can only reduce this count. Set $B_\sigma:=12G_\sigma$.
There are at most three coordinates and two signs.
Tail vectors number $<e^{8\sigma^{-3}}$, besides $R_1$.
Thus one explicit bound is
\[
 N_{{\rm lab},r}(\sigma)\le2P_\sigma T_\sigma,\quad
 P_\sigma=1+B_\sigma+B_\sigma^2+B_\sigma^3,\quad
 T_\sigma=1+\lceil e^{8\sigma^{-3}}\rceil,
\]
with $B_\sigma=12G_\sigma$ allowed in this upper bound. $\square$

Put $\mathcal H_0=\{q<Q:p^2\mid q\text{ for some prime }p\ge u_0\}$.
Set $B_c=m_c!/M_c$ on active blocks, and $B_c=0$ otherwise.
The next lemma shows that both coefficients are supported on integers
divisible by the square of a large prime.

\textbf{Lemma B.2.} For every block and integer $q$,
\[
 |\Xi_{U,c}(q)|+|\Xi_{A,c}(q)|
 \le B_c\mathbf1_{\mathcal H_0}(q)
 \le C_\sigma\mathbf1_{\mathcal H_0}(q).
 \tag{B.3}
\]

\emph{Proof.} Inactive blocks contribute nothing. Write $q=q_<q_\ge$ at the cutoff $u_c$.
The support $d\mid P(u_c;N)$ and roughness force $d=q_<$ uniquely.
The multiset $q_\ge=\prod_pp^{e_p}$ fills the $m_c$ positions in at most
$m_c!/\prod_pe_p!\le m_c!$ ways.
The two masks in (B.1) are disjoint indicators,
so $|\lambda_d|\le1$ gives
\[
 M_c(|\Xi_{U,c}(q)|+|\Xi_{A,c}(q)|)
 \le\#\{\text{full tuples}:q=d\textstyle\prod_jp_j\}\le m_c!.
\]
A forbidden tuple has $p_i=p_j\ge u_c$.
For an allowed nonsquarefree tuple, $d$ is squarefree and coprime to $q_\ge$.
Thus in either case $p^2\mid q$ for some $p\ge u_0$.
Each multiplicity enters once, through $e_p!$, and no graph factor occurs. $\square$

The tail argument is \cite[Lemma~4.8]{Z1},
with prime-source cost \cite[Appendix~B, (B.5)--(B.6)]{Z1}.
Lemma B.2 also covers arbitrary Goldbach prefixes and G1 masks.
The next bounds use their actual tested measures;
for the test functions a trivial bound suffices.

\textbf{Lemma B.3.} Put $A_{\max}=100$.
For every row and $q<Q$,
\[
 |\Delta_P^{(r)}(q)|,\ |(\Delta_0^{(r)}-\Delta_P^{(r)})(q)|
 \le A_{\max}\Big(\frac Nq+1+\frac N{\varphi(q)}\Big)
 \le A_{\max}\Big(\frac{N(2+L_N)}q+1\Big).
 \tag{B.4}
\]
The second discrepancy vanishes on DIRECT rows.

\emph{Proof.} The tested measures are supported in $[1,N]$,
since $v<N$ on DIRECT rows and $e\le N_{\rm cut}<N$ on G1.
For G1, the tested variable is the product $e$; the condition
$e\equiv N\pmod q$ forces $(e,q)=1$, while $Y_G(q)$ retains that
coprimality restriction.
By (4.6), $A_D\le1$ and $A_G^{\rm cut}\le A_G\le72<100$.
Thus $C(q)\le A_{\max}(N/q+1)$ and $0\le Y(q)\le X\le A_{\max}N$.
Use $|\Delta_P|\le C+Y/\varphi(q)$ and
$|\Delta_0-\Delta_P|=(X-Y(q))/\varphi(q)$.
For the distinct primes $p_1<\cdots<p_k$ of $q$,
\[
 \frac q{\varphi(q)}\le\prod_i\frac{i+1}i=k+1
 \le1+\log_2q\le1+\frac{(3/5)L_N}{\log2}\le1+L_N.
\]
On DIRECT rows, $v>N/L_N^{30}>N^{3/5}>q$ is prime,
hence $Y_A(q)=X_A$. $\square$

The exceptional set of integers with a repeated large prime is sparse.

\textbf{Lemma B.4.} For $u_0\ge2$,
\[
 \sum_{q\in\mathcal H_0}\frac1q\le\frac{2(1+L_N)}{u_0},\qquad
 \#\mathcal H_0\le\frac{2Q}{u_0}.
 \tag{B.5}
\]

\emph{Proof.} Write $q=np^2$ and overcount over $p\ge u_0$.
The harmonic sum is then at most $(1+\log Q)\sum_{k\ge\lceil u_0\rceil}k^{-2}$.
The count is at most $Q\sum_{k\ge\lceil u_0\rceil}k^{-2}$.
For $K=\lceil u_0\rceil$, this last sum is
$\le1/(K-1)\le1/(u_0-1)\le2/u_0$. $\square$

For either test function, (B.4)--(B.5) give
\[
 \sum_{q\in\mathcal H_0}|T^{(r)}\text{-test}(q)|
 \le A_{\max}\frac{2N(2+L_N)(1+L_N)+2Q}{u_0}
 \le\frac{6A_{\max}N(1+L_N)^2}{u_0}.
 \tag{B.6}
\]
Here $Q\le N$ and $2(2+L)(1+L)+2\le6(1+L)^2$,
equivalently $4L^2+6L\ge0$.
These lemmas combine into the main estimate of this subsection.

\begin{theorem}[Repeated-prime estimate]\label{thm:repeated-primes}
Fix $0<\sigma\le10^{-10}$ and put
$\Lambda_\sigma=\sum_r|c_r|N_{{\rm lab},r}(\sigma)$ over the nineteen rows.
For sufficiently large even $N$,
\begingroup\small
\[
\begin{aligned}
 \frac1{\mathcal W}\sum_c|c_{r(c)}|\sum_q\Big(&|\Xi_{U,c}(q)\Delta_P^{(r)}(q)|\\
 &+|\Xi_{A,c}(q)(\Delta_0^{(r)}-\Delta_P^{(r)})(q)|\Big)\\
 &\le232\Lambda_\sigma C_\sigma N^{-\ell_0\sigma^2}(1+L_N)^4.
\end{aligned}
 \tag{B.7}
\]
\endgroup
This bounds $|\mathcal E_{\rm SQ}|$,
separately and jointly for the three parts in (B.2).
Their size is $O_\sigma(N^{-\ell_0\sigma^2}L_N^4)$,
and also $O_\sigma(N^{-\ell_0\sigma^2}L_N^{K_{\rm SQ}+4})$ for fixed $K_{\rm SQ}\ge0$.
The allowed-repeat part vanishes on DIRECT rows.
\end{theorem}

\emph{Proof.} Disjointness in (B.3) bounds the combined integrand by
$C_\sigma\mathbf1_{\mathcal H_0}(q)A_{\max}\allowbreak(N(2+L_N)/q+1)$.
By (B.6), each block contributes at most $6A_{\max}C_\sigma N(1+L_N)^2/u_0$.
Also
\[
 \mathcal W=4(1-\varepsilon_A)C(N)N/L_N^2
 \ge4(1-\varepsilon_A)C_2N/L_N^2\ge2.59N/L_N^2.
\]
Here $C(N)\ge C_2>13/20$, since
\[
 C_2\ge\prod_{\substack{3\le p\le2000\\p\ {\rm prime}}}
 \left(1-\frac1{(p-1)^2}\right)
 \left(1-\frac1{2000}\right)>\frac{13}{20}.
\]
The finite product is compared in rational arithmetic.
For the tail, $\prod(1-a_j)\ge1-\sum a_j$ and
$\sum_{m\ge2001}m^{-2}\le\sum_{m\ge2001}(m(m-1))^{-1}=1/2000$.
The weighted block count is $\Lambda_\sigma$.
Finally use $600/2.59<232$ and $L_N^2\le(1+L_N)^2$. $\square$

Fix $\sigma$, hence $\Lambda_\sigma,C_\sigma$, and take $N$ last
(\S~\ref{p2:sec:fixed-parameters}).
Then $\mathcal E_{\rm SQ}=o_\sigma(1)$.
Above $N_0(\sigma)$ the saving is $\ell_0\sigma^2/2$, and the proof gives $K_{\rm SQ}=0$.
Prefix multiplicities enter Lemma B.2, while row counts enter only $\Lambda_\sigma$.

\subsection{Sieve remainder mass}

This subsection proves (1.26), with the normalization (1.27).
The AA37 input is \cite[Theorem~2.6, (2.30)--(2.33b)]{Z1},
taken from \cite[Theorem~4]{ref5} with $R=1$.
For a finite nonnegative sequence, write
$|\mathcal A_d|=X\omega(d)/d+r(\mathcal A,d)$, with $X>0$.
The density is multiplicative and nonnegative on prime powers.
We require $0<\omega(p)<p$, $K\ge2$ and $L_\omega>1$.
The product and prime-power hypotheses are \cite[(2.24), (2.27)]{Z1}, for $z>w\ge2$.
Substitute
\[
 \begin{gathered}
 D=L_c,\quad\epsilon_s=\sigma<1/3,\quad\delta_s=\sigma^9=\tau,\quad s=\frac{\log D}{\log z},\\
 \mathcal G(D)=\{D^{\sigma^2(1+\tau)^j}:j\ge0\}.
 \end{gathered}
\]
The supports in \cite[(2.32)--(2.33)]{Z1}, including $R_1^\nu$, are
\[
 d<D^\sigma,\quad d\mid P(D^{\sigma^2}),\quad
 P(y)=\prod_{p<y,\,p\in\mathcal P}p,\quad
 D_i\le p_i<\min(D_i^{1+\tau},z).
\]
Positions with equal labels are distinct,
and the nonincreasing labels obey \cite[(2.29)]{Z1}.
Thus $D_1<D^{1/2}$ and $|\mathcal D^\nu(D)|<e^{8\sigma^{-3}}$.
Also $|\phi_d^\nu|,|\lambda_d^\nu(\mathbf D)|\le1$, independently of $\mathcal A,X,\omega$.
The upper and lower ranges are $2\le z\le D$ and $2\le z\le D^{1/2}$.
By \cite[(2.31)]{Z1},
\[
 E_s\le C_E\sigma+C_E\sigma^{-8}e^{K+L_\omega}(\log L_c)^{-1/3},
\]
with absolute $C_E$.

For a source block $\mathcal A_b$, take $X_{r,b}=X_r/\varphi(b)$ and
$\omega(p^k)/p^k=1/\varphi(p^k)$ on permitted primes (\S~4.3, after (4.21)).
The prime-power sum is $\sum_{w\le p<z,\,p\in\mathcal P}(p-1)^{-2}$.
Mertens gives the product bound after the primes dividing $Nb$ are omitted.
The source checks give the density identity, the prefix exclusion and the cutoffs.
Here $X_r=X_A=\sum_vA_D(v)$ on DIRECT rows,
and $X_r=X_G=\sum_eA_G^{\rm cut}(e)$ on G1.
Each error is $X_{r,b}V(z_{r,b})E_s$,
with no factor $F$ or $f$.
Hence
\[
\begin{aligned}
 E^D_{\rm AA}&\le\frac{C_E\sigma}4W_{\rm AA}(N)
 +\frac{C_E\sigma^{-8}e^{K+L_\omega}}{4(\ell_0L_N)^{1/3}}W_{\rm AA}(N),\\
 W_{\rm AA}(N)&:=
 \frac{\sum_r|c_r|\sum_bX_{r,b}V(z_{r,b})}{(1-\varepsilon_A)U_C}.
\end{aligned}
 \tag{B.8}
\]
The sum runs over the nineteen rows and their actual prefixes $b$.
Each prefix lies in one parent cell, and the empty prefix is $b=1$.
The labels $\mathbf D$, $R_1/R^\nu$, $e^{8\sigma^{-3}}$, $M_\sigma!$,
HB patterns and boxes occur only in the remainder,
so they add no factor to $W_{\rm AA}$.
For DIRECT rows, $\sum|c_r|=24$: eleven rows in \S~2 give 20, and S2 gives 4.
The subsequent estimate uses the convenient coarser ceiling $45$.
The four G1 coefficients have absolute value 1.
The ceiling $45$ is a harmless enlargement of the exact DIRECT coefficient sum $24$.
The BV pieces S5-top, S6 and BV-R1 are treated in \S~5.

\subsubsection*{Three uniform bounds}

Three uniform estimates bound this mass. Put $h=1/1000$.

\textbf{Lemma B.8 (exponent floor).} Every Buchstab prefix prime and AA37 cutoff is at least $N^h$.
Every G1 source prime is at least $N^\alpha>N^h$.

\emph{Proof.} The smallest fixed cutoff is $N^u$, with $u=1/500$.
The others are E3 on S1-3, $E=N^{t_3}$ on G1-3,
and PF/K on S2-3; all are at least $N^{u/\rho^3}$, and $u/(1+\tau)^3>h$.
Other prefix ranges start at $N^\alpha$ or higher; see (4.4)--(4.5).
AA37 tail primes are not covered by this prefix assertion. $\square$

\textbf{Lemma B.9 (density product).} Uniformly in $N$, rows, cells and $\sigma$,
for $z\ge N^h$ and an absolute threshold $N\ge N_0$,
\[
 V(z)=\prod_{\substack{2<p<z\\p\nmid N}}\left(1-\frac1{p-1}\right)
 \le\frac{4C(N)}{\log z}\le\frac{4000C(N)}{L_N}.
 \tag{B.9}
\]

\emph{Proof.} Prefix primes are $\ge z$,
so $\mathcal P(Nb)$ and $\mathcal P(N)$ agree below $z$.
Since $2\mid N$ and $1-1/(p-1)=(1-1/p)(1-1/(p-1)^2)$,
\[
 V(z)=2\prod_{p<z}\left(1-\frac1p\right)c_2(z)
 \prod_{\substack{2<p<z\\p\mid N}}\frac{p-1}{p-2},\qquad
 c_2(z)=\prod_{2<p<z}\left(1-\frac1{(p-1)^2}\right).
\]
By (1.2), the last product is at most $C(N)/C_2$.
Also $c_2(z)/C_2\to1$ and Mertens gives
$\prod_{p<z}(1-1/p)\sim e^{-\gamma}/\log z$.
Use $2e^{-\gamma}<1.13<4$ and $\log z\ge hL_N$. $\square$

\textbf{Lemma B.10 (source masses).} For an absolute threshold $N\ge N_0$, at every legal cut,
\[
 X_A\le\frac{2N}{L_N},\qquad X_G\le\frac{2560000N}{L_N}.
 \tag{B.10}
\]

\emph{Proof.} Rosser--Schoenfeld \cite{RosserSchoenfeld1962} gives
$X_A\le\pi(N)<1.25506N/\log N$.
For $k=3,4$, count ordered prime $k$-tuples, with repetitions,
whose entries are $\ge N^h$ and whose product is $\le N$.
Fix the first $k-1$ entries, with product $Q'$.
If a last entry exists, then $N/Q'\ge N^h$,
and its choices number at most
$\pi(N/Q')\le2(N/Q')/\log(N/Q')\le2N/(hQ'L_N)$.
Dropping the other constraints gives
\[
 T_k(N)\le\frac2h H_N^{k-1}\frac N{L_N},\qquad
 H_N:=\sum_{N^h\le p\le N}\frac1{p-1}\le8.
\]
Also $\sum_{N^h\le p\le N}1/p\to\log(1/h)=\log1000<7$
and $\sum_{m\ge N^h}1/(m(m-1))\to0$.
The tuple lengths in $R_{11,3},R_{11,4},R_{12},R_{13}$ are $3,4,3,3$.
By Lemma B.8, their entries are $\ge N^\alpha$ and $e<N$.
Using the weights $2,2,1,1$ of (4.4),
\[
 X_G\le\sum_eA_G(e)\le4T_3+2T_4
 \le(4\cdot2000\cdot8^2+2\cdot2000\cdot8^3)N/L_N
 =2560000N/L_N.
\]
The bound uses no representation ceiling and no factorial. $\square$

\textbf{Lemma B.11 (prefix harmonic sum).} For each fixed row,
the prefix depth may vary from $0$ to $3$ across its cells. Hence
\[
 \sum_b\frac1{\varphi(b)}\le\sum_{j=0}^3H_N^j
 \le1+8+64+512=585,
 \tag{B.11}
\]
where the sum covers all actual prefixes in all cells.

\emph{Proof.} For a prefix of exact depth $j$,
use $\varphi(b)\ge\prod_i(p_i-1)$ and drop order, distinctness,
selectors and ranges. Its total is at most $H_N^j\le8^j$.
Sum the four possible depths, including $j=0$.
The parent cells partition the prefixes into disjoint half-open pieces,
including A/B children, Option-A pieces and retained whole cells.
Thus a prefix is counted once and no cell-count factor occurs. $\square$

The three lemmas combine into the uniform mass bound.

\begin{theorem}[Sieve remainder mass]\label{thm:sieve-mass}
There is $N_0$, depending only on $h$ and the fixed source conventions,
such that for every $0<\sigma\le10^{-10}$ and even $N\ge N_0$,
the bound (1.26) holds, independently of the tail labels and cells.
\end{theorem}

\emph{Proof.} By (B.9)--(B.11),
\[
\begin{aligned}
 \sum_r|c_r|\sum_bX_{r,b}V(z_{r,b})
 &\le\frac{4000C(N)}{L_N}\,585
 \left(45\frac{2N}{L_N}+4\frac{2560000N}{L_N}\right)\\
 &=4000\cdot585\cdot10240090\frac{C(N)N}{L_N^2}.
\end{aligned}
\]
Dividing by $(1-\varepsilon_A)U_C=(1-\varepsilon_A)C(N)N/L_N^2$ gives
\[
 W_{\rm AA}(N)\le\frac{4000\cdot585\cdot10240090}{1-1/10000}
 =\frac{239618106000000000}{9999}<10^{14}=W_*.
\]
The last comparison is
$239618106000000000<9999\cdot10^{14}=999900000000000000$. $\square$

\textbf{Corollary (persistent and vanishing AA37 errors).}
Fix $C_E$ first.
For $C_E>0$, choose rational $\sigma$ satisfying (1.27):
\[
 \sigma\le\frac{4g_{\rm AA}}{C_EW_*},\qquad
 g_{\rm AA}=0.00002853328.
\]
Then $C_E\sigma W_{\rm AA}/4\le C_E\sigma W_*/4\le g_{\rm AA}$.
This bounds the actual AA37 perturbation on all nineteen rows by one
$g_{\rm AA}$ charge. Any further subtraction of the same numerical
allowance in the final error total is an additional conservative reserve.
For $C_E=0$ the condition is empty.
The selector and cutoff errors retain their stated $\tau$ bounds.
The second term of (B.8) is at most
\[
 \frac{W_*C_E\sigma^{-8}e^{K+L_\omega}}{4(\ell_0L_N)^{1/3}}
 =o_\sigma(1).
\]
Here $K,L_\omega$ are absolute constants for the formal density,
and $\log L_c\ge\ell_0L_N$.
The bound (1.26) is applied before $\sigma$ is chosen in (1.27).
Pointwise finiteness alone would allow unbounded $\sigma W(\sigma)$;
the uniform bound removes that dependence.

\subsection{Source and interval references}

The source identity is Lemma~\ref{lem:source-replacement},
with its prefix and sifting data kept fixed.
Section~6 bounds both difference weights,
including the G1 representation multiplicities.

Hypothesis (5.23) follows from the Prime interval coefficients lemma in Section~5.
Its endpoint power-log Siegel--Walfisz input is \cite[(2.21)]{Z1}.
Endpoint subtraction requires
\[
 J\subset[Y/2,Y],\qquad |J|\asymp YL^{-H},\qquad
 N^c\le Y\le N^C,
\]
with fixed $H,c,C$ and $c,C>0$.
Then $\#\{n\in J:n\equiv l\pmod k\}\le|J|/k+1$.
The O1 and G4 supports have $H=2$ and $H=H_g$ (\S\S~5.5--5.6).
Thus the interval hypothesis holds, also with the auxiliary divisor $d$ of (5.23).

\section{Finite source checks}\label{app:finite-source-checks}

Here we check source identities and cutoffs on selected instances.
The main text proves these estimates uniformly in $N$.

The script uses integer arithmetic and \texttt{fractions.Fraction},
with a smallest-prime-factor sieve through $400000$.
For each seed
\[
  1,\qquad 2,\qquad 20261022,
\]
the script runs with second argument $150$ and Python \texttt{random.Random(seed)}.
It draws $150$ even integers
$N=2\cdot\texttt{randrange}(5000,60000)$. Each draw selects
$z\in\{5,7,11,13,17\}$ and $V_0\in\{200,500,1000,2000\}$.
For each of the three prefix families below, it draws
$D\in\{200,1000,5000\}$ separately.
It tests the primes $v$ with
$100v<99N$. The plain sample also requires $v>V_0$.
There are three prefix families:
$b=1$; the first at most six primes $p\in[z,4z)$, $p\nmid N$;
and the first at most six ordered pairs $(p_1,p_2)$.
The pairs satisfy $z\le p_2<p_1<6z$ and $p_i\nmid N$.
Thus the tested prefix condition is $p_i\ge z$ and $(p_i,N)=1$.

For the three-step Buchstab identity, each seed gives $30$ further even integers
$N=2\cdot\texttt{randrange}(3000,15000)$. Here $z_0=5$,
$z_1\in\{17,23,29\}$ and $V_0=300$.
The plain primes satisfy $V_0<v$ and $100v<99N$.
The switched checks use $30$ further even integers
$N=2\cdot\texttt{randrange}(20000,100000)$.
Set $y_1\in\{5,7,11\}$, $z_{\rm sw}\in\{7,11,13\}$,
and interval endpoints $N/100,99N/100$.
The ordered triples satisfy
$y_1\le p_1<p_2\le p_3$, $p_1^3<N$, $p_1p_2^2<N$ and
$N/100<p_1p_2p_3<N$, with the coordinate bounds $p_i^3<8N$.
The tagged triples also have $p_i\nmid N$
and $p_3\ne p_2$. In these three batches the rational cutoffs are compared
by cross-multiplying integers.

\begingroup
\small
\setlength{\tabcolsep}{3pt}
\renewcommand{\arraystretch}{1.16}
\begin{longtable}{@{}p{0.105\textwidth} p{0.465\textwidth} p{0.365\textwidth}@{}}
\hline
\textbf{Check} & \textbf{Finite assertion} & \textbf{Recorded result} \\
\hline\endhead
D1 & The literal cofactor row equals $S(\mathcal A_{r,b},\mathcal P_N,z)$ on each selected direct fibre. & 1,350 direct-case batches; all passed \\
D2 & The first-true partition of the ambient/plain source differences is exhaustive. & 0 escapes among 3,988,123 ambient tuples; CUT-LOW 96,876, EXC 272, SRC-RESTORE 121,017 \\
D3 & The sifted source identity is the signed sum of the plain source and its two set differences. & Exact on every tested fibre \\
D4 & Escapes of the earlier source-difference predicate are removed by the sieve. & 1,295,622 escapes; each has a cofactor prime $<z$ not dividing $N$; 0 exceptions \\
D5 & Legendre's identity uses the remainder of the plain tested sequence. & Exact in every tested fibre \\
D6 & The Rosser $\beta=2$ lower and upper inequalities \cite{ref4} hold, with the stated signed-remainder decompositions. & Exact in every tested fibre \\
D7 & For plain sampled primes $v>V_0$, $\Delta_0(q)=\Delta_P(q)$ whenever $q\le V_0$. & Exact for every such modulus and prime \\
D8 & The three-step Buchstab identity and its cutoff monotonicity hold on the plain sequence. & 90/90 cases \\
G1 & The tagged-ambient/plain-cut weight identity, tag partition, switching comparison, and plain-sequence Legendre identity hold. & 90/90 cases; 0 tag escapes; 12,337 coordinate-divides-$N$ tags and 9,701 repeated-prime tags; all switching discrepancies occur at $N-e\le\sqrt N$ or $N-e=1$ \\
\hline
\end{longtable}
\endgroup

Each direct-row run has 450 batches.
Together they contain 3,988,123 ambient tuples,
and D2 and D4 report the combined counts.
In D4, an ``escape'' is a tuple missed by the earlier difference predicates.
Its cofactor has a sieving prime below $z$, coprime to $N$.
The Buchstab and G1 counts also combine the three seeds.
G1 counts coordinate-divisibility tags and repeated-prime tags separately,
using a first-true partition: a coordinate-divisibility tag takes precedence
over a repeated-prime tag when both predicates hold.

Each seed also gives five small examples.
One further illustration has about 260 digits.
The large examples use fixed-base Miller--Rabin screens.
Their primality is uncertified.
They illustrate the cutoffs only;
neither the proof nor the numerical bounds use them.

The SHA-256 identifiers are:
\begin{itemize}[leftmargin=2em]
\item Checker \texttt{lemma\_s\_ambient\_check.py}:
  \texttt{81d74731\allowbreak{}cf82e911\allowbreak{}ea509e3d\allowbreak{}cead1fcf\allowbreak{}885e4d3c\allowbreak{}a06e3d66\allowbreak{}b53fc749\allowbreak{}717138fc}.
\item Seed $1$ output:
  \texttt{997cb3e3\allowbreak{}2195d3a3\allowbreak{}1e61557c\allowbreak{}a87474e8\allowbreak{}d0d785e3\allowbreak{}feb396cb\allowbreak{}d95164c9\allowbreak{}e3b35fe2}.
\item Seed $2$ output:
  \texttt{3139f095\allowbreak{}2534a707\allowbreak{}2f1a0938\allowbreak{}8305f426\allowbreak{}092cdfab\allowbreak{}07eab568\allowbreak{}92e8992c\allowbreak{}f14ff09a}.
\item Seed $20261022$ output:
  \texttt{a6f69703\allowbreak{}a169e18f\allowbreak{}a05cd697\allowbreak{}399a6713\allowbreak{}2b1fb3cd\allowbreak{}8383c5f4\allowbreak{}d10d24c8\allowbreak{}4d00d04d}.
\item Environment record (Python 3.12.3, Linux x86\_64):
  \texttt{7f0e02eb\allowbreak{}4aff3c83\allowbreak{}e7e8cacd\allowbreak{}9a831669\allowbreak{}782f676f\allowbreak{}fc5dfec0\allowbreak{}47ce26ca\allowbreak{}200e5271}.
\end{itemize}
These files are in the numerical-certificate archive;
see the Data Availability statement.
Its README gives the verification and reproduction commands.

\section*{Use of AI tools}
The author devised the method of proof.
Codex (OpenAI) and Claude Code (Anthropic) are the AI tools used.
They assisted with derivation checks, code, typesetting, and review.
The author checked every step of the proof, every table and every reference.
The author takes full responsibility for the content of the paper.

\section*{Statements and declarations}
\textbf{Competing interests.} The author declares no competing interests.

\textbf{Funding.} No external funding was received.

\textbf{Data availability.} The numerical-certificate supplement is provided as
\path{paper-ii-public-certificates-v11.tar.gz}. It contains the source
scripts, full-precision outputs, finite-check run records, verification
logs, and environment records. The archive SHA-256 is
\texttt{cb9d7b96\allowbreak{}ce1323c9\allowbreak{}118d845f\allowbreak{}4dd701a0\allowbreak{}d830e8c6\allowbreak{}10090dee\allowbreak{}c88bd879\allowbreak{}2648e100}.
After extraction, the root \path{SHA256SUMS} has SHA-256
\texttt{18eb23d9\allowbreak{}6be1c096\allowbreak{}bd2d32a3\allowbreak{}6a34d7e5\allowbreak{}d9da978e\allowbreak{}5494e7ae\allowbreak{}02851a7b\allowbreak{}74adb182}.
The archive README gives verification and reproduction commands.
No DOI has been assigned.
\IfFileExists{references.tex}{\begin{thebibliography}{99}

\bibitem{Chen1966} J. R. Chen,
On the representation of a large even integer as the sum of a prime and
the product of at most two primes (in Chinese),
\emph{Kexue Tongbao} \textbf{17} (1966), 385--386.

\bibitem{Chen1973} J. R. Chen,
On the representation of a large even integer as the sum of a prime and
the product of at most two primes,
\emph{Scientia Sinica} \textbf{16} (1973), 157--176.

\bibitem{Wu2004} J. Wu,
Chen's double sieve, Goldbach's conjecture and the twin prime problem,
\emph{Acta Arithmetica} \textbf{114} (2004), no.~3, 215--273.

\bibitem{Wu2008} J. Wu,
Chen's double sieve, Goldbach's conjecture and the twin prime problem, 2,
\emph{Acta Arithmetica} \textbf{131} (2008), no.~4, 367--387.
DOI: \href{https://doi.org/10.4064/aa131-4-5}{10.4064/aa131-4-5}.

\bibitem{Lichtman} J. D. Lichtman,
Primes in arithmetic progressions to large moduli, and Goldbach beyond
the square-root barrier, with an appendix with S. Drappeau,
\href{https://arxiv.org/abs/2309.08522v1}{arXiv:2309.08522v1}, 2023.

\bibitem{LiLiu2026} J. Li and J. Liu,
Theorem $(1+1.9)$ on the Goldbach Conjecture,
\href{https://arxiv.org/abs/2606.05224v2}{arXiv:2606.05224v2}, 2026.

\bibitem[Z1]{Z1} R. Zhang,
Shifted primes in arithmetic progressions to large moduli with linear sieve remainder weights,
preprint, Zenodo, 2026.
DOI: \href{https://doi.org/10.5281/zenodo.23615525}{10.5281/zenodo.23615525}.

\bibitem{ref1} J.-M. Deshouillers and H. Iwaniec,
Kloosterman sums and Fourier coefficients of cusp forms,
\emph{Inventiones Mathematicae} \textbf{70} (1982), 219--288.
DOI: \href{https://doi.org/10.1007/BF01390728}{10.1007/BF01390728}.

\bibitem{ref2} E. Bombieri, J. B. Friedlander and H. Iwaniec,
Primes in arithmetic progressions to large moduli,
\emph{Acta Mathematica} \textbf{156} (1986), 203--251.
DOI: \href{https://doi.org/10.1007/BF02399204}{10.1007/BF02399204}.

\bibitem{ref3} S. Drappeau,
Sums of Kloosterman sums in arithmetic progressions, and the error term
in the dispersion method,
\emph{Proceedings of the London Mathematical Society} \textbf{114}
(2017), no.~4, 684--732.
DOI: \href{https://doi.org/10.1112/plms.12022}{10.1112/plms.12022};
\href{https://arxiv.org/abs/1504.05549v4}{arXiv:1504.05549v4}.

\bibitem{ref4} H. Iwaniec,
Rosser's sieve,
\emph{Acta Arithmetica} \textbf{36} (1980), 171--202.

\bibitem{ref5} H. Iwaniec,
A new form of the error term in the linear sieve,
\emph{Acta Arithmetica} \textbf{37} (1980), 307--320.
DOI: \href{https://doi.org/10.4064/aa-37-1-307-320}{10.4064/aa-37-1-307-320}.

\bibitem{ref6} D. R. Heath-Brown,
Prime numbers in short intervals and a generalized Vaughan identity,
\emph{Canadian Journal of Mathematics} \textbf{34} (1982), 1365--1377.
DOI: \href{https://doi.org/10.4153/CJM-1982-095-9}{10.4153/CJM-1982-095-9}.

\bibitem{ref7} H. H. Kim,
Functoriality for the exterior square of $GL_4$ and the symmetric fourth
of $GL_2$,
\emph{Journal of the American Mathematical Society} \textbf{16}
(2003), 139--183, with Appendix~1 by D. Ramakrishnan and Appendix~2
by H. H. Kim and P. Sarnak.
DOI: \href{https://doi.org/10.1090/S0894-0347-02-00410-1}{10.1090/S0894-0347-02-00410-1}.

\bibitem{ref11} H. Iwaniec and E. Kowalski,
\emph{Analytic Number Theory}, American Mathematical Society Colloquium
Publications \textbf{53}, American Mathematical Society, Providence, RI, 2004.

\bibitem{Shiu1980} P. Shiu,
A Brun--Titchmarsh theorem for multiplicative functions,
\emph{Journal f\"ur die reine und angewandte Mathematik} \textbf{313}
(1980), 161--170.

\bibitem{Bombieri1965} E. Bombieri,
On the large sieve,
\emph{Mathematika} \textbf{12} (1965), 201--225.

\bibitem{Vinogradov1965} A. I. Vinogradov,
The density hypothesis for Dirichlet $L$-series (in Russian),
\emph{Izvestiya Akademii Nauk SSSR. Seriya Matematicheskaya}
\textbf{29} (1965), 903--934.

\bibitem{Siegel1935} C. L. Siegel,
\"Uber die Classenzahl quadratischer Zahlk\"orper,
\emph{Acta Arithmetica} \textbf{1} (1935), 83--86.

\bibitem{Walfisz1936} A. Walfisz,
Zur additiven Zahlentheorie.~II,
\emph{Mathematische Zeitschrift} \textbf{40} (1936), 592--607.
DOI: \href{https://doi.org/10.1007/BF01218882}{10.1007/BF01218882}.

\bibitem{Mertens1874} F. Mertens,
Ein Beitrag zur analytischen Zahlentheorie,
\emph{Journal f\"ur die reine und angewandte Mathematik} \textbf{78}
(1874), 46--62.

\bibitem{Buchstab1937} A. A. Buchstab,
Asymptotische Absch\"atzungen einer allgemeinen zahlentheoretischen Funktion,
\emph{Matematicheskii Sbornik (N.S.)} \textbf{2(44)} (1937),
no.~6, 1239--1246.

\bibitem{RosserSchoenfeld1962} J. B. Rosser and L. Schoenfeld,
Approximate formulas for some functions of prime numbers,
\emph{Illinois Journal of Mathematics} \textbf{6} (1962), 64--94.
DOI: \href{https://doi.org/10.1215/ijm/1255631807}{10.1215/ijm/1255631807}.

\bibitem{Johansson2017} F. Johansson,
Arb: Efficient arbitrary-precision midpoint-radius interval arithmetic,
\emph{IEEE Transactions on Computers} \textbf{66} (2017), no.~8,
1281--1292.
DOI: \href{https://doi.org/10.1109/TC.2017.2690633}{10.1109/TC.2017.2690633}.

\end{thebibliography}
}{}
\end{document}
